% !TEX encoding = UTF-8 Unicode
\documentclass[a4paper,9pt]{amsart}
\usepackage[english]{babel}
\usepackage[utf8]{inputenc}
\usepackage[T1]{fontenc}

%For Japanese kanji and kana: but use XeLaTeX !
%\usepackage{xeCJK}
%\setCJKmainfont{MS Mincho} % for \rmfamily
%\setCJKsansfont{MS Gothic} % for \sffamily

\usepackage{csquotes}
\usepackage[style=numeric,%alphabetic,%
	useprefix,%
	giveninits=true,%
	hyperref,%
	doi=false,%
	url=false,%
	isbn=false,%
	backend=bibtex,%
	maxbibnames=99%
	]{biblatex}
\bibliography{./BIB}

%%%%%%%%%%%%%%%%%%%%%%%%%%%%%%%%%%%%%%%%%%%%%%%%%%%%%%%%%%%%%%%
% PACCHETTI
%%%%%%%%%%%%%%%%%%%%%%%%%%%%%%%%%%%%%%%%%%%%%%%%%%%%%%%%%%%%%%%
\usepackage{amssymb}
\usepackage{mathrsfs}
\usepackage{hyperref}
\usepackage[all]{xy} % For \xymatrix
\usepackage[usenames,dvipsnames]{xcolor}
\hypersetup{colorlinks,%
citecolor=Black,%
filecolor=Black,%
linkcolor=Black,%
urlcolor=Black}
%\hypersetup{colorlinks,%
%citecolor=OliveGreen,%
%filecolor=magenta,%
%linkcolor=RoyalBlue,%
%urlcolor=cyan}
%\usepackage{mparhack}	% Per correggere le note a margine, che a volte vengono dalla parte sbagliata.
%\usepackage{todonotes} % Per le note a margine in stile Thomas Z.
%\usepackage{comment}	% Per i commenti lunghi: \begin{comment} ... \end{comment}
\usepackage{marvosym}	% Per il fulmine: \Lightning
\usepackage{enumitem}	% Per personalizzare gli elenchi
%\usepackage{booktabs}	% Per i comandi \toprule \midrule \bottomrule , che fanno righe orrizontali nelle tabelle
%\usepackage{array}		% Per formattare meglio le celle nelle tabelle.
\usepackage{mathtools}	% Per \mathclap
\mathtoolsset{showonlyrefs=true} % Per mostrare solo i numeri che servono. With Mathtools: See https://tex.stackexchange.com/questions/4728/how-do-i-number-equations-only-if-they-are-referred-to-in-the-text
%\usepackage[pdftex]{graphicx}	% Per il frontespizio...
%\usepackage{makeidx}	% Per fare l'indice analitico. Già in amsart, amsbook and amsproc. 
\makeindex
\allowdisplaybreaks % per avere equazioni su più pagine
%\usepackage{tikz}
%
%\usepackage{framed} % for the environment ``framed''. For more complicated things, use \usepackage{mdframed}.
%\usepackage{refcheck} % Mostra i label non usati ``Unused label...'' nel log file. See https://tex.stackexchange.com/questions/209782/get-list-of-unused-labels
\usepackage[normalem]{ulem} % for strikeout: \sout{striketouted} (preferable over 'soul')
%\usepackage{soul} % for strikeout: \st{striketouted} (see 'ulem' package, which is preferable)

%%%%%%%%%%%%%%%%%%%%%%%%%%%%%%%%%%%%%%%%%%%%%%%%%%%%%%%%%%%%%%%
% pgfplots FOR PLOTTING FUNCTIONS: 
%\documentclass{article}
\usepackage{pgfplots}
%
%\begin{document}
%\begin{tikzpicture}[>=stealth]
%    \begin{axis}[
%        xmin=-4,xmax=4,
%        ymin=-2,ymax=2,
%        axis x line=middle,
%        axis y line=middle,
%        axis line style=<->,
%        xlabel={$x$},
%        ylabel={$y$},
%        ]
%        \addplot[no marks,blue,<->] expression[domain=-pi:pi,samples=100]{sin(deg(2*x))+1/2} 
%                    node[pos=0.65,anchor=south west]{$y=\sin(2x)+\frac{1}{2}$}; 
%    \end{axis}
%\end{tikzpicture}
%\end{document}
\usepgfplotslibrary{external} % To make pictures only once.
\tikzexternalize % activate!
\tikzsetexternalprefix{figures/}
% From the documentation: If a figure shall be remade, one can simply delete all or selected graphics files and regenerate them. Alternatively, one can use the command \tikzset{external/force remake} somewhere in the document to remake every following picture automatically.
% -> there is a way to change the names of the files generated.
%%%%%%%%%%%%%%%%%%%%%%%%%%%%%%%%%%%%%%%%%%%%%%%%%%%%%%%%%%%%%%%

%%%%%%%%%%%%%%%%%%%%%%%%%%%%%%%%%%%%%%%%%%%%%%%%%%%%%%%%%%%%%%%
% COMANDI
%%%%%%%%%%%%%%%%%%%%%%%%%%%%%%%%%%%%%%%%%%%%%%%%%%%%%%%%%%%%%%%
\newcommand{\scr}[1]{\mathscr{#1}}
\newcommand{\frk}[1]{\mathfrak{#1}}
\newcommand{\bb}[1]{\mathbb{#1}}
\newcommand{\cal}[1]{\mathcal{#1}}
%
\newcommand{\N}{\mathbb{N}}	% Numeri naturali
\newcommand{\Z}{\mathbb{Z}}	% Numeri interi
\newcommand{\Q}{\mathbb{Q}}	% Numeri razionali
\newcommand{\R}{\mathbb{R}}	% Numeri reali
\newcommand{\C}{\mathbb{C}}	% Numeri complessi
\newcommand{\K}{\mathbb{K}}	% Campo di numeri
\newcommand{\Id}{\mathrm{Id}}	% Mappa identità
%\newcommand{\Span}{\mathrm{span}}	% Span
%\newcommand{\Ker}{\mathrm{Ker}}	% Ker
\newcommand{\dd}{\mathrm{d}}	% 'd' di derivata
\newcommand{\did}{\,\mathrm{d}}	% 'd' in integrali
\newcommand{\diff}{\mathrm{d}}	% 'd' di derivata
\newcommand{\de}{\partial}		% Derivata parziale
\newcommand{\Diff}{\mathrm{D}}
\renewcommand{\div}{\operatorname{div}}	% div di Divergenza.
\newcommand{\into}{\hookrightarrow}		% (->  %)
%\newcommand{\dto}{\dashrightarrow}
\newcommand{\ol}[1]{\overline{#1}}	% overline breve
%\DeclareMathOperator{\arctanh}{arctanh}
\newcommand{\IF}{\Leftarrow}	% <=
\newcommand{\THEN}{\Rightarrow}	% =>
\newcommand{\IFF}{\Leftrightarrow}	% <=>
\newcommand{\ddx}{\framebox[\width]{$\Rightarrow$} }
\newcommand{\ssx}{\framebox[\width]{$\Leftarrow$} }
\newcommand{\ci}{\framebox[\width]{$\subset$} }
\newcommand{\ic}{\framebox[\width]{$\supset$} }
\newcommand{\lebox}{\framebox[\width]{$\le$} }
\newcommand{\gebox}{\framebox[\width]{$\ge$} }
\newcommand{\real}{\frk{Re}}
\newcommand{\imag}{\frk{Im}}
%\renewcommand{\Re}{\operatorname{Re}}
%\renewcommand{\Im}{\operatorname{Im}}
\newcommand{\im}{\operatorname{im}}
% \hel is the left-lower corner for restriction of measures.
\newcommand{\hel}
{
\hskip2.5pt{\vrule height6pt width.5pt depth0pt}
\hskip-.2pt\vbox{\hrule height.5pt width6pt depth0pt}
\,
}
\newcommand{\Lin}{\mathtt{Lin}}
\newcommand{\End}{\mathtt{End}}
\newcommand{\Orth}{\mathtt{O}} % Orthogonal group
\newcommand{\GL}{\mathtt{GL}}
\newcommand{\J}{\mathtt{J}}
\newcommand{\ad}{\mathrm{ad}}
\newcommand{\Sym}{\mathtt{Sym}}
\newcommand{\jet}{\mathfrak{j}}
\newcommand{\Jet}{\mathtt{J}}
\newcommand{\HorDer}{\mathtt{HD}}
\newcommand{\bbar}[1]{\overline{\overline{#1}}}
\newcommand{\Der}{\mathtt{Der}}
\newcommand{\Aut}{\mathtt{Aut}}
\newcommand{\Poly}{\scr P}
\newcommand{\UniEnvAlg}{\scr U}
\newcommand{\Tensor}{\scr T}
\newcommand{\HH}{\bb H}
\newcommand{\Ad}{\operatorname{Ad}}
\newcommand{\w}{\mathbf{w}}
\newcommand{\Haus}{\scr H}
\newcommand{\dist}{\mathrm{dist}}
\newcommand{\spt}{\mathrm{spt}}
\newcommand{\loc}{\mathrm{loc}}
\newcommand{\laplacian}{\triangle}
\newcommand{\grad}{\nabla}
\newcommand{\trace}{\mathrm{trace}}
\newcommand{\singspt}{\operatorname{singSpt}}
\newcommand{\interior}{\operatorname{int}}
\newcommand{\llangle}{\langle\!\langle}
\newcommand{\rrangle}{\rangle\!\rangle}
\newcommand{\abs}[1]{\left\lvert #1 \right\rvert}
\newcommand{\norm}[1]{\left\lVert #1 \right\rVert}
\newcommand{\nnorm}[1]{\left\lvert\!\left\lvert\!\left\lvert #1 \right\rvert\!\right\rvert\!\right\rvert}

\newcommand{\TODO}{{\bf\color{Red}TODO}}
\newcommand{\angryface}{}

\newcommand{\lcontr}{\lrcorner}
\newcommand{\rcontr}{\raisebox{\depth}{\scalebox{1}[-1]{$\lrcorner$}}} % serve graphicx

% Per \one, la mappa che vale 1
\usepackage{dsfont}
\newcommand{\one}{{\mathds 1\!}}

% Mean integral
\def\Xint#1{\mathchoice
      {\XXint\displaystyle\textstyle{#1}}%
      {\XXint\textstyle\scriptstyle{#1}}%
      {\XXint\scriptstyle\scriptscriptstyle{#1}}%
      {\XXint\scriptscriptstyle\scriptscriptstyle{#1}}%
      \!\int}
   \def\XXint#1#2#3{{\setbox0=\hbox{$#1{#2#3}{\int}$}
        \vcenter{\hbox{$#2#3$}}\kern-.5\wd0}}
%   \def\ddashint{\Xint=}
%   \def\dashint{\Xint-}
   \def\fint{\Xint-}

%%%%%%%%%%%%%%%%%%%%%%%%%%%%%%%%%%%%%%%%%%%%%%%%%%%%%%%%%%%%%%%
% TEOREMI
%%%%%%%%%%%%%%%%%%%%%%%%%%%%%%%%%%%%%%%%%%%%%%%%%%%%%%%%%%%%%%%
\theoremstyle{plain}
\newtheorem{proposition}{Proposition}[section]
\newtheorem{theorem}[proposition]{Theorem}
\newtheorem{lemma}[proposition]{Lemma}
\newtheorem{corollary}[proposition]{Corollary}
\newtheorem{thm}{Theorem}[section]
\renewcommand{\thethm}{\Alph{thm}}
%\newtheorem{question}{Question}

\theoremstyle{definition}
\newtheorem{definition}[proposition]{Definition}
\newtheorem{remark}[proposition]{Remark}
%\declaretheorem[name=Remark,sibling=proposition,qed={\lower-0.3ex\hbox{$\blacklozenge$}}]{remark} % but with \usepackage{thmtools}
\newtheorem{exercise}[proposition]{Exercise}%[section]
%\AtEndEnvironment{exercise}{\hfill$\blacksquare$\endexercise}
%\AtEndEnvironment{exercise}{\hfill$\clubsuit$\endexercise}
\AtEndEnvironment{exercise}{\hfill$\Diamond$\endexercise}



\theoremstyle{remark}



%%%%%%%%%%%%%%%%%%%%%%%%%%%%%%%%%%%%%%%%%%%%%%%%%%%%%%%%%%%%%%%
% Subsectioning for notes:
\renewcommand{\thesubsection}{{\bf\S\arabic{section}.\arabic{subsection}}}


%%%%%%%%%%%%%%%%%%%%%%%%%%%%%%%%%%%%%%%%%%%%%%%%%%%%%%%%%%%%%%%

\title[Introduction to PDE]{Introduction to\\Partial Differential Equations\\--Lecture Notes--}
\author[Nicolussi Golo]{Sebastiano Nicolussi Golo}
%\date{\today}
\date{\today. \IfFileExists{./.gittex}{\input{./.gittex}}{}}

\begin{document}
\maketitle

\setcounter{tocdepth}{2}
\phantomsection
\addcontentsline{toc}{section}{Contents}
\tableofcontents

%%%%%%%%%%%%%%%%%%%%%%%%%%%%%%%%%%%%%%%%%%%%%%%%%%%%%%%%%%%%%%%
%%%%%%%%%%%%%%%%%%%%%%%%%%%%%%%%%%%%%%%%%%%%%%%%%%%%%%%%%%%%%%%
\section*{Introduction}


%%%%%%%%%%%%%%%%%%%%%%%%%%%%%%%%%%%%%%%%%%%%%%%%%%%%%%%%%%%%%%%
\subsection{Bibliography}
I will closely follow \cite{MR2597943}.
Other references are:
\begin{itemize}
\item \fullcite{MR2597943}
\item \fullcite{MR1357411}
\item \fullcite{zbMATH01211561}
\end{itemize}


%%%%%%%%%%%%%%%%%%%%%%%%%%%%%%%%%%%%%%%%%%%%%%%%%%%%%%%%%%%%%%%
\newpage
\setcounter{part}{-1}
\part{Preliminaries}

In this first preliminary part, we will touch superficially several notions that we will refer to along the lecture notes.

%%%%%%%%%%%%%%%%%%%%%%%%%%%%%%%%%%%%%%%%%%%%%%%%%%%%%%%%%%%%%%%
\section{Pairings, dot and scalar products, etc...}

%\subsection{}
There is a flourishing art of making products of stuff in mathematics and physics.
But notation lacks.
So, here are my choices.

%%%%%%%%%%%%%%%%%%%%%%%%%%%%%%%%%%%%%%%%%%%%%%%%%%%%%%%%%%%%%%%
\subsection{Dot Product}\label{subs6889e81a}
If $a,b\in\C^n$, then
\begin{equation}
	a\cdot b = \sum_{j=1}^n a_jb_j .
\end{equation}

%%%%%%%%%%%%%%%%%%%%%%%%%%%%%%%%%%%%%%%%%%%%%%%%%%%%%%%%%%%%%%%
\subsection{Scalar or inner product}\label{subs6889e7fa}
If $a,b\in\C^n$, then
\begin{equation}
	\langle a,b \rangle = \sum_{j=1}^n a_j\bar b_j = a\cdot \bar b .
\end{equation}
More generally, if $\mu$ is a measure and $a,b\in L^2(\mu;\C^n)$ are complex-valued functions
\begin{equation}
	\langle a,b \rangle = \langle a,b \rangle_\mu = \int a(x)\cdot \bar b(x) \did\mu(x) .
\end{equation}

These are inner products, with the properties ($\lambda\in\C$)
\begin{align}
	\langle a,b \rangle = \overline{ \langle b,a \rangle } , &&
	\langle a,\lambda b \rangle = \bar \lambda \langle a,b \rangle, \\
	\langle a,a \rangle \in[0,+\infty) , &&
	\langle \lambda a,b \rangle = \lambda \langle a,b \rangle .
\end{align}

%%%%%%%%%%%%%%%%%%%%%%%%%%%%%%%%%%%%%%%%%%%%%%%%%%%%%%%%%%%%%%%
\subsection{Pairing}\label{subs6889e83c}
If $V$ is a vector space (over some field $\bb K$, e.g., $\bb K=\C$) with algebraic dual $V^*$ (that is, the space of all linear maps $V\to\bb K$), then,
for all $a\in V^*$ and $b\in V$,
\begin{equation}
	\langle a|b \rangle = a[b] \in \bb K \qquad\text{ (which is the evaluation of $a$ in $b$) .}
\end{equation}
Sometimes, we can make it more precise
\begin{equation}
	{}_{V^*}\langle a|b \rangle_V = a[b] .
\end{equation}
In general, if $b\in V$ and $a\in\operatorname{Lin}(V;W)$ for some linear space $W$,
\begin{equation}
	\langle a|b \rangle = a[b] \in W .
\end{equation}
With this in mind, it is clear that
\begin{equation}
	\langle a|b \rangle = \langle b|a \rangle .
\end{equation}
For instance,
${}_{V^*}\langle a|b \rangle_V = {}_{V}\langle b|a \rangle_{V^*}$.

Pairings have the following properties ($\lambda\in\bb K$):
\begin{align}
	\langle a|b \rangle = \langle b|a \rangle , \\
	\langle \lambda a|b \rangle = \lambda \langle a|b \rangle , \\
	\langle a|\lambda b \rangle = \lambda \langle a|b \rangle .
\end{align}

%%%%%%%%%%%%%%%%%%%%%%%%%%%%%%%%%%%%%%%%%%%%%%%%%%%%%%%%%%%%%%%
\subsection{Example in $L^2(\mu)$}
To put everything together, we see that, if $a,b\in L^2(\mu)$, define $B,\bar B\in L^2(\mu)^*$ as 
\begin{align}
	B[\phi] &= \int b\cdot\phi\did \mu,
	\qquad\forall \phi\in L^2(\mu)  ,\\
	\bar B[\phi] &= \langle \phi,b \rangle,
	\qquad\forall \phi\in L^2(\mu) .
\end{align}
Then
\begin{equation}
	\langle a,b \rangle 
	= \int a\cdot \bar b \did\mu 
	= \langle \bar B|a \rangle
	= \overline{\langle B|\bar a \rangle}
\end{equation}
and
\begin{equation}
	\langle B|a \rangle
	= \langle a|\bar b \rangle 
	= \langle b|\bar a \rangle .
\end{equation}

%%%%%%%%%%%%%%%%%%%%%%%%%%%%%%%%%%%%%%%%%%%%%%%%%%%%%%%%%%%%%%%
\section{Derivatives}\label{sec6889e87f}

Let $n,m\in\N$.
We denote by $e_1,\dots,e_n$ the standard basis of $\R^n$.

Let $\Omega\subset\R^n$ be an open set and $f:\Omega\to\C^m$ and $x\in\Omega$.
The derivative of $f$ at $x$, if it exists, is the $\R$-linear map $\Diff f(x):\R^n\to\C^m$ such that
\begin{equation}
	\lim_{y\to x} \frac{ \norm{ f(y) - f(x) - \Diff f(x)[y-x] } }{ \norm{y-x} } = 0 .
\end{equation}

The partial derivative of $f$ are
\begin{equation}
	\frac{\de f}{\de x_j}(x) 
	\coloneqq \Diff f(x)[e_j] 
	= \lim_{h\to0} \frac{ f(x+h e_j) - f(x) }{ h } \in \C^m .
\end{equation}

The gradient of $f$ is (if it exists)
\begin{equation}
	\grad f
	\coloneqq \left( \frac{\de f}{\de x_1} , \dots, \frac{\de f}{\de x_n} \right) .
\end{equation}
The gradient is mostly used for scalar functions, that is, when $m=1$.

If $\alpha\in\N^n$ is a multiindex, we set $\abs{\alpha} \coloneqq \sum_{j=1}^n\alpha_j$ and we denote by $\de^\alpha f$ or $\Diff^\alpha f$ the partial derivative of $f$ given by 
\begin{equation}
	\Diff^\alpha f = \left(\frac{\de }{\de x_1}\right)^{\alpha_1} \cdots \left(\frac{\de }{\de x_n} \right)^{\alpha_n} f .
\end{equation}
If a function $f(x,y)$ depends on several variables, we write $\Diff_xf$ or $\de_xf$ to denote the derivatives of $f$ in the directions of $x$.

The space of functions $\Omega\to\C^m$ that are continuous, differentiable and with continuous derivatives up to order $k\in\N$ is denote by $C^k(\Omega;\C^m)$.
If $m=1$, then we just write $C^k(\Omega)$.
We write $C^k(\bar\Omega;\C^m)$ for the space of functions that are smooth of class $C^k$ on a neighborhood of $\bar\Omega$.

If the domain of a function $u$ is described as a product of open sets $X\subset\R^m$ and $Y\subset\R^n$, 
so $u:X\times Y\to\C$, we may require different regularity in the two entries.
So, we say that $u\in C^{a;b}(X\times Y)$ for some $a,b\in\N$ if, for every $y\in Y$, $u(\cdot,y)\in C^a(X)$, and for every $x\in X$, $u(x,\cdot)\in C^b(Y)$.

Notice that $C^{a,b}(X)$ (with comma in place of semicolumn) is usually reserved for functions of class $C^a$ whose $a$-th order derivative is $b$-Hölder. 
In this case, we would have $a\in\N$ and $b\in[0,1]$.

%%%%%%%%%%%%%%%%%%%%%%%%%%%%%%%%%%%%%%%%%%%%%%%%%%%%%%%%%%%%%%%
\section{Measures}

%%%%%%%%%%%%%%%%%%%%%%%%%%%%%%%%%%%%%%%%%%%%%%%%%%%%%%%%%%%%%%%
\subsection{$L^p$ spaces}\label{subs6889e917}
Let $(X,\scr M,\mu)$ be a measure space, that is, $X$ is a set endowed with a $\sigma$-algebra $\scr M\subset 2^X$ and $\mu:\scr M\to[0,+\infty]$ is a measure.
In fact, we will mostly have $X=\Omega\subset\R^n$ open subset of $\R^n$,
$\scr M$ the Borel $\sigma$-algebra and $\mu$ the Lebesgue measure.

For $u:X\to\C$ measurable and $p\in[1,+\infty)$, define
\begin{equation}
	\|u\|_{L^p} 
	\coloneqq \|u\|_{L^p(X)}
	\coloneqq \left( \int_X \abs{u(x)}^p \did\mu(x) \right)^{1/p} .
\end{equation}
The space $L^p(X)$ or $L^p(\mu)$ is the Banach space of all equivalence classes of measurable functions $u:X\to\C$ with $\|u\|_{L^p}<\infty$, where two functions are identified if they are equal $\mu$-almost everywhere.

Sometimes, when we write ``$u\in L^p(X)$'', we mean that $u$ is a specific function, not an equivalence class.


%%%%%%%%%%%%%%%%%%%%%%%%%%%%%%%%%%%%%%%%%%%%%%%%%%%%%%%%%%%%%%%
\subsection{Hölder and Minkowski inequalities}

%Anyway, the following inequalities have really nothing to do with $\R^n$.
Let $(X,\scr M,\mu)$ be a measure space.
Let $u,v:X\to\C$ and $u_j:X\to\C$, $j\in\{1,\dots,k\}$, be measurable functions.

%{\bf Hölder inequality:}
%If $u,v:X\to\C$ are measurable functions and $p,q\in[1,\infty)$ are such that $\frac1p+\frac1q=1$, then
%\begin{equation}
%	\int_X u(x)v(x) \did\mu(x) 
%	\le \left( \int_X u(x)^p \did\mu(X) \right)^{1/p}
%		\left( \int_X v(x)^q \did\mu(X) \right)^{1/q} .
%\end{equation}

\begin{equation}\label{Hölder}
\begin{array}{c}
	\forall p,p'\in[1,+\infty),
	\text{ with }\frac{1}{p} + \frac{1}{p'} = 1, \\
%	\quad\THEN\quad
	\int_X u(x)v(x) \did\mu(x) \le \|u\|_{L^p(X)} \|v\|_{L^{p'}(X)} .
\end{array}
	\tag{Hölder}
\end{equation}

\begin{equation}\label{General Hölder}
\begin{array}{c}
	\forall \{p_j\}_{j=1}^k\subset[1,+\infty),
	\text{ with }\sum_{j=1}^k \frac1{p_j} = 1, \\
%	\quad\THEN\quad
	\int_X \prod_{j=1}^k u_j(x) \did\mu(x) \le \prod_{j=1}^k \|u_j\|_{L^{p_j}(X)} .
\end{array}
	\tag{General Hölder}
\end{equation}

\begin{equation}\label{Minkowski}
	\forall p\in[1,\infty]\qquad
	\|u+v\|_{L^p} \le \|u\|_{L^p(X)} + \|v\|_{L^p(X)} .
	\tag{Minkowski}
\end{equation}




%%%%%%%%%%%%%%%%%%%%%%%%%%%%%%%%%%%%%%%%%%%%%%%%%%%%%%%%%%%%%%%
\subsection{Approximation with continuous functions}
\begin{theorem}\label{thm680f2023}
	Let $\Omega\subset\R^n$ open and $p\in[1,\infty)$.
	The space $C_c(\Omega)$ of continuous functions with compact support in $\Omega$ is dense in $L^p(\Omega)$.
\end{theorem}
\begin{proof}
	\cite[Proposition (7.9)]{MR767633}
\end{proof}


%%%%%%%%%%%%%%%%%%%%%%%%%%%%%%%%%%%%%%%%%%%%%%%%%%%%%%%%%%%%%%%
\subsection{Dominated Convergence Theorem}
\begin{theorem}[LDCT, Lebesgue Dominated Convergence Theorem {\cite[Theorem 2.24]{MR767633}}]\label{thm6798ec42}
	% These two lines are for making the link to LDCT. It seems I need to compile three times before it works.
	\makeatletter\def\@currentlabel{\text{LDCT}}\makeatother
	\label{LDCT}
	Let $(X,\mu)$ be a measure space, $T$ a topological space with $\hat t\in T$, and $f:T\times X\to\R$ a function such that there exists $g\in L^1(\mu)$ with $|f(t,x)|\le g(x)$ for all $(t,x)\in T\times X$.
	Suppose that, for $\mu$-a.e.~$x\in X$, the limit $\lim_{t\to \hat t} f(t,x)$ exists.
	Then
	\[
	\lim_{t\to \hat t} \int_X f(t,x) \did \mu(x)
	= \int_X (\lim_{t\to \hat t} f(t,x)) \did \mu(x) .
	\]
\end{theorem}





%%%%%%%%%%%%%%%%%%%%%%%%%%%%%%%%%%%%%%%%%%%%%%%%%%%%%%%%%%%%%%%
\subsection{Functions defined by integrals}


%\subsection{Integral with parameter}
\begin{theorem}[Integral with parameter]\label{thm679f29de}
	Let $X\subset\R^m$ be an open set and $(Y,\mu)$ a measure space.
	Let $K:X\times Y\to\C$ be a measurable function and set 
	\begin{equation}
		F(x) \coloneqq \int_Y K(x,y) \did\mu(y) ,
	\end{equation}
	for all $x\in X$ such that $y\mapsto K(x,y)$ is $\mu$-integrable,
	i.e., such that $\int_Y \abs{K(x,y)} \did\mu(y)<\infty$.

	\begin{enumerate}[label={\bf (\ref{thm679f29de}.\arabic*)}]
	\item\label{thm679f29de_1}
	If  
	\begin{equation}\label{eq679f27fe}
		\exists g\in L^1(Y)\quad \forall (x,y)\in X\times Y
		\qquad \abs{K(x,y)} \le g(y) 
	\end{equation}
	then $F$ is defined for all $x\in X$ and $F\in L^\infty(X)$.
	In particular, $\|F\|_{L^\infty(X)} \le \|g\|_{L^1(Y)}$.
	
	\item\label{thm679f29de_2}
	If~\eqref{eq679f27fe} holds, and if
	\begin{equation}\label{eq679f28e5}
		\forall y\in Y
		\qquad
		K(\cdot,y)\in C^0(X),
	\end{equation}
	then $F\in C^0(X)$.
	
	\item\label{thm679f29de_3}
	If~\eqref{eq679f27fe} and~\eqref{eq679f28e5} hold, and if 
	\begin{equation}\label{eq679f2b96}
		\forall y\in Y
		\ :\   K(\cdot,y)\in C^1(X),
		\text{ and }\exists g_1\in L^1(Y)\forall (x,y)\in X\times Y 
		\ :\ |D_xK(x,y)| \le g_1(y) ,
	\end{equation}
	then $F\in C^1(X)$ and, for every $x\in X$,
	\begin{equation}\label{eq679f2bf4}
		D_x F(x) = \int_Y D_xK(x,y) \did\mu(y) .
	\end{equation}
	\end{enumerate}
\end{theorem}
\begin{proof}
	Proof of~\ref{thm679f29de_1}.
	The condition~\eqref{eq679f27fe}  implies directly that $K(x,\cdot)\in L^1(Y)$ and that, for every $x\in X$,
	\[
	|\int_Y K(x,y) \did\mu(y)| 
	\le \int_Y |K(x,y)| \did\mu(y) 
	\le \|g\|_{L^1(Y)} .
	\]
	
	Proof of~\ref{thm679f29de_2}.
	Let $\{x_j\}_{j\in\N}\subset X$ be a sequence converging to $x_\infty\in X$.
	Condition~\eqref{eq679f27fe} allows us to apply the Lebesgue Dominated Convergence Theorem~\ref{thm6798ec42},
	\begin{align*}
	F(x_\infty)
	&= \int_Y K(x_\infty,y) \did\mu(y) \\
	[\text{by }K(\cdot,y)\in C^0(X)]
	&= \int_Y  \lim_{j\to\infty} K(x_j,y) \did\mu(y)  \\
	[\text{by \ref{LDCT} and~\eqref{eq679f27fe}}]
	&= \lim_{j\to\infty}\int_Y K(x_j,y) \did\mu(y)  
	= \lim_{j\to\infty} F(x_j).
	\end{align*}	
	This shows that $F\in C^0(X)$.
	
	Proof of~\ref{thm679f29de_3}.	
	Fix $\hat x\in X$ and $i\in\{1,\dots,m\}$ and $\hat h>0$ such that $B(\hat x,\hat h)\Subset X$.
	For $h\in(-\hat h,\hat h)$, we have
	\[
	\frac{F(\hat x + he_i) - F(\hat x)}{h} 
		= \int_Y \frac{ K(\hat x + he_i,y) - K(\hat x,y) }{h} \did\mu(y) .
	\]
	Notice that for every $h\in(-\hat h,\hat h)$ and every $y\in Y$ there exists $h'\in (-h,h)$ such that 
	\[
	\frac{ K(\hat x + he_i,y) - K(\hat x,y) }{h}
	= \frac{\de K}{\de x_i} (\hat x + h'e_i , y) .
	\] 
	Hence, by~\eqref{eq679f2b96},
	\begin{equation}\label{lem68889cd0}
	\left|\frac{ K(\hat x + he_i,y) - K(\hat x,y) }{h}\right| 
	\le \left|	\frac{\de K}{\de x_i} (\hat x + h'e_i , y) \right| 
	\le g_1(y) .
	\end{equation}
	Therefore, we can apply the Lebesgue Dominated Convergence Theorem~\ref{thm6798ec42},
	to obtain
	\begin{align*}
		\int_Y \frac{\de K}{\de x_i}(\hat x,y) \did y 
		&= \int_Y \lim_{h\to0}\frac{ K(\hat x + he_i,y) - K(\hat x,y) }{h} \did\mu(y) \\ 
		[\text{by \ref{LDCT} and~\eqref{lem68889cd0}}]
		&= \lim_{h\to0}\int_Y \frac{ K(\hat x + he_i,y) - K(\hat x,y) }{h} \did\mu(y) \\
		&= \lim_{h\to0}\frac{F(\hat x + he_i) - F(\hat x)}{h} .
	\end{align*}
	This shows that $F$ is differentiable at every point, with identity~\eqref{eq679f2bf4}.
	By~\ref{thm679f29de_2}, $D_xF$ is continuous.
\end{proof}

\begin{corollary}\label{cor67b2f278}
	Let $X\subset\R^m$ be an open set and $(Y,\mu)$ a measure space.
	Let $K:X\times Y\to\C$ be a measurable function and set 
	\begin{equation}
		F(x) = \int_Y K(x,y) \did\mu(y) ,
	\end{equation}
	for all $x\in X$ such that $y\mapsto K(x,y)$ is integrable.
	
	Suppose that
	\begin{equation}\label{eq67b2f2bb}
	\begin{array}{c}
		\forall y\in Y
		\ :\   K(\cdot,y)\in C^\infty(X),
		\text{ and }\\
	\forall\alpha\in\N^m\;\exists g_\alpha\in L^1(Y)\;\forall (x,y)\in X\times Y 
		\ :\ |D^\alpha_xK(x,y)| \le g_1(y) .
	\end{array}
	\end{equation}
	
	Then $F\in C^\infty(X)$.
\end{corollary}

%%%%%%%%%%%%%%%%%%%%%%%%%%%%%%%%%%%%%%%%%%%%%%%%%%%%%%%%%%%%%%%
The following Proposition~\ref{prop6798f690} is a direct consequence of Theorem~\ref{thm679f29de}
and the Hölder inequality.

\begin{proposition}\label{prop6798f690}
	Let $X\subset\R^m$ and $Y\subset\R^n$ be open sets, $k\in\N$ and $p\in[1,\infty]$.
	Let $K:X\times Y\to\C$ be a measurable function such that:
	\begin{enumerate}
	\item	for every $y\in Y$, the function $x\mapsto K(x,y)$ belongs to $C^k(Y)$;
	\item	for every $E\Subset X$ and $\alpha\in\N^n$ with $|\alpha|\le k$, there exists $g_{\alpha,E}\in L^p(Y)$ such that, for every $x\in E$ and every $y\in Y$,
	\[
	\left|\frac{\de^{|\alpha|} K}{\de x^\alpha}(x,y)\right| \le g_{\alpha,E} (y) .
	\]
	\end{enumerate}
	
	For $q\in[1,\infty]$ with $1/p+1/q=1$ and $f\in L^q(Y)$, define $T_Kf$ as
	\[
	T_Kf (x) = \int_Y K(x,y) f(y) \did y .
	\]
	Then $T_Kf\in C^k(X)$ and, for every $\alpha\in\N^m$ with $|\alpha|\le k$,
	\begin{equation}\label{eq6798ef4d}
	\frac{\de^{\abs{\alpha}}}{\de x^\alpha} T_Kf = T_{\frac{\de^{\abs{\alpha}}K}{\de x^\alpha}} f .
	\end{equation}
\end{proposition}

%{\color{Blue} TODO: write this proof using Theorem~\ref{thm679f29de} .
%\begin{proof}
%	First, notice that, for every $x\in X$, we can estimate using the Hölder inequality
%	\[
%	\int_Y |K(x,y) f(y)| \did y 
%	\le \int_Y  g_{0,E} f(y)\did y 
%	\overset{\text{Hölder}}\le \|g_{0,E}\|_{L^q} \|f\|_{L^p} ,
%	\]
%	where $E$ is any compact set in $X$ with $x\in E$.
%	Therefore $T_Kf (x)$ is well defined and $T_Kf$ is locally bounded.
%	
%	If $\{x_j\}_{j\in\N}\subset X$ is a sequence converging to $x_\infty\in X$, then 
%	\[
%	|K(x_j,y) f(y)| \le g_{0,E}(y) f(y)
%	\qquad
%	\forall y\in Y,
%	\]
%	where $E=\{x_j\}_{j\in\N}\cup\{x_\infty\}$ is a compact subset of $X$
%	and $y\mapsto g_{0,E}(y) f(y)$ is an element of $L^1(Y)$.
%	By the Dominated Convergence Theorem~\ref{thm6798ec42},
%	\begin{align*}
%	\lim_{j\to\infty} T_Kf(x_j)
%	&= \lim_{j\to\infty}\int_Y K(x_j,y) f(y) \did y  \\
%	[\text{by \ref{LDCT}}] 
%	&= \int_Y  \lim_{j\to\infty} K(x_j,y) f(y) \did y  \\
%	[\text{by }K(\cdot,y)\in C^0(X)] 
%	&= \int_Y K(x_\infty,y) f(y) \did y \\
%	&= T_Kf (x_\infty).
%	\end{align*}
%	This shows that $Tf\in C^0(X)$.
%	
%%	Let $\{\phi_\ell\}_{\ell\in\N}$ be a partition of the unity on $X$, i.e., $\phi_\ell\in C^\infty_c(X)$ and $0\le\phi_\ell\le1$ for each $\ell$, and $\sum_{\ell\in\N}\phi_\ell = 1$ on $X$.
%	
%	It is clear that if~\eqref{eq6798ef4d} holds for $k=1$, then we can iterate it to obtain the statement for every $k$.
%	
%	Fix $\hat x\in X$ and $i\in\{1,\dots,m\}$ and $\hat h>0$ such that $B(\hat x,\hat h)\Subset X$.
%	For $h\in(-\hat h,\hat h)$, % with $|he_i|<{\rm dist}(\hat x,\de X)$,
%	we have
%%	\begin{align*}
%%		&\hspace{-1cm}\left| \frac{T_Kf(\hat x + he_i) - T_Kf(\hat x)}{h} - T_{\frac{\de K}{\de x^i}}f(\hat x) \right| \\
%%		&\le \int_Y \left| \frac{ K(\hat x + he_i,y) - K(\hat x,y) }{h} - \frac{\de K}{\de x^i}(\hat x,y) \right| |f(y)| \did y \\
%%		&= \sum_\ell \int_Y \left| \frac{ K(\hat x + he_i,y) - K(\hat x,y) }{h} - \frac{\de K}{\de x^i}(\hat x,y) \right| |\phi_\ell(y)f(y)| \did y . 
%%	\end{align*}
%%	
%	\[
%	\frac{T_Kf(\hat x + he_i) - T_Kf(\hat x)}{h} 
%		= \int_Y \frac{ K(\hat x + he_i,y) - K(\hat x,y) }{h} f(y) \did y .
%	\]
%%		&= \sum_\ell \int_Y \frac{ K(\hat x + he_i,y) - K(\hat x,y) }{h} \phi_\ell(y)f(y) \did y . 
%%	\end{align*}
%	Notice that for every $h\in(-\hat h,\hat h)$ and every $y\in Y$ there exists $h'(h,y)\in (-h,h)$ such that 
%	\[
%	\frac{ K(\hat x + he_i,y) - K(\hat x,y) }{h}
%	= \frac{\de K}{\de x_i} (\hat x + h'e_i , y) .
%	\] 
%	Hence, 
%	\[
%	\left|\frac{ K(\hat x + he_i,y) - K(\hat x,y) }{h}\right| 
%	\le \left|	\frac{\de K}{\de x_i} (\hat x + h'e_i , y) \right| 
%%	\sup\left\{ \left| \frac{\de K}{\de x_i}(x,y)\right| : x\in \bar B(\hat x,\hat h) \right\} 
%	\le g_{i,\bar B(\hat x,\hat h)}(y) .
%	\]
%	Therefore, we can apply the Dominated Convergence Theorem~\ref{thm6798ec42},
%	to obtain
%	\begin{align*}
%		\lim_{h\to0}\frac{T_Kf(\hat x + he_i) - T_Kf(\hat x)}{h} 
%		&= \lim_{h\to0}\int_Y \frac{ K(\hat x + he_i,y) - K(\hat x,y) }{h} f(y) \did y \\
%		[\text{by \ref{LDCT}}] 
%		&= \int_Y \lim_{h\to0}\frac{ K(\hat x + he_i,y) - K(\hat x,y) }{h} f(y) \did y \\
%		&= \int_Y \frac{\de K}{\de x_i}(\hat x,y) f(y) \did y .
%	\end{align*}
%	
%\end{proof}
%}

%%%%%%%%%%%%%%%%%%%%%%%%%%%%%%%%%%%%%%%%%%%%%%%%%%%%%%%%%%%%%%%
\subsection{Fundamental Theorem of Calculus of Variations}
\index{Fundamental Theorem of Calculus of Variations}
\index{theorem!Fundamental -- of Calculus of Variations}

\begin{theorem}[Fundamental Theorem of Calculus of Variations]\label{thm67c2eb77}
	If $u\in L^1_\loc(\R^n)$ is such that
	\begin{equation}\label{eq67c2ebaf}
		\forall\phi\in C^\infty_c(U) ,
		\qquad
		\int_U \phi\cdot u \did x = 0 ,
	\end{equation}
	then $u=0$ almost everywhere in $\R^n$.
\end{theorem}

\begin{exercise}
	Prove the Fundamental Theorem of Calculus of Variations~\ref{thm67c2eb77}.
\end{exercise}

%%%%%%%%%%%%%%%%%%%%%%%%%%%%%%%%%%%%%%%%%%%%%%%%%%%%%%%%%%%%%%%
\section{Divergence theorem}

Given an open set $U\subset\R^n$,
the divergence of a vector field $v \in C^1(U;\C^n)$ is 
\[
\div(v) = \sum_{j=1}^n \frac{\de v_j}{\de x_j} \in C(U) .
\]

\begin{theorem}[Divergence Theorem {\cite[Lemma 10.1]{MR0447753}}]\label{teo67a7a31e}
	Let $U\subset\R^n$ be an open set with $C^1$ boundary $\de U$.
	We denote by $\nu:\de U\to\bb S^{n-1}$ the outward normal.
	Let $v\in C^1(\bar U;\C^n)$ be a vector field.
	Then
	\begin{equation}\label{eq67a7a3f9}
		\int_U \div(v) \did x = \int_{\de U} v\cdot \nu \did S(x) .
	\end{equation}
\end{theorem}

Recall that the integral in $\did S$ is the surface integral over $\de U$.
One can think of $\did S$ as the Hausdorff measure of dimension $n-1$.

\begin{remark}
	You might have seen the divergence Theorem~\ref{teo67a7a31e} for real vector fields.
	In the case of a complex vector field, we have
	$v = v_1 + i v_2$ with $v_1,v_2\in  C^1(\bar U;\R^n)$.
	Then
	\begin{align*}
		\int_U \div(v) \did x
		&= \int_U \div(v_1) \did x + i \int_U \div(v_2) \did x \\
		&= \int_{\de U} v_1\cdot \nu \did S(x) + i \int_{\de U} v_2\cdot \nu \did S(x) \\
		&= \int_{\de U} (v_1+iv_2)\cdot \nu \did S(x) .
	\end{align*}
	Notice that $\nu$ has necessarily real components and thus $\bar \nu = \nu$.
	In other words, we could also write $v\cdot \nu = v \cdot \bar \nu = \langle v,\nu \rangle$, as the hermitian product of complex vectors.
\end{remark}

\begin{remark}
	Formula~\eqref{eq67a7a3f9} is usually paired with the following formula for the divergence:
	if $f\in C^1(U)$ and $v\in C^1( U;\C^n)$, then
	\begin{equation}\label{eq67a7a587}
		\div(fv) = \grad f\cdot v + f\div(v) .
	\end{equation}
	Indeed,
	\begin{align*}
		\div(fv) 
		&= \sum_{j=1}^n \frac{\de (fv_j)}{\de x_j} \\
		&= \sum_{j=1}^n \left( \frac{\de f}{\de x_j} v_j + f \frac{\de v_j}{\de x_j} \right) \\
		&= \grad f\cdot v + f\div(v) .
	\end{align*}
	It follows that
	\[
	\int_U (\grad f\cdot v + f\div(v)) \did x = \int_{\de U} f v\cdot \nu \did S(x) .
	\]
\end{remark}


%%%%%%%%%%%%%%%%%%%%%%%%%%%%%%%%%%%%%%%%%%%%%%%%%%%%%%%%%%%%%%%
\section{Coarea formulas}

%%%%%%%%%%%%%%%%%%%%%%%%%%%%%%%%%%%%%%%%%%%%%%%%%%%%%%%%%%%%%%%
\subsection{Intro to surface measures}
If $\Sigma\subset\R^n$ is an $m$-dimensional immersed $C^1$-submanifold,
we denote by $S^m$ the $m$-dimensional surface measure on $\Sigma$.
We can describe $S^m$ as the Hausdorff $m$-dimensional measure, or, being the submanifold smooth, as an integral of $m$-differential forms on $\Sigma$.
We can also obtain several formulas for its explicit use.
However, we will use a list of properties of these measures, and we do not need further details.

We need the symmetries of $S^m$: it is translation invariant, rotation invariant and scales properly under dilations. 
More precisely, if $\Sigma\subset\R^m$ is an immersed submanifold and if $O\in\Orth(n)$ is an orthogonal matrix, $v\in\R^n$ and $r>0$, then
\begin{equation}\label{eq67a718d1}
	\int_{r(O\Sigma + v)} u(x) \dd S^m(x) 
	= r^m \int_{\Sigma} u(r(Ox+v)) \dd S^m(x) ,
	\qquad\forall u\in C^0(\R^n) .
\end{equation}

%%%%%%%%%%%%%%%%%%%%%%%%%%%%%%%%%%%%%%%%%%%%%%%%%%%%%%%%%%%%%%%
\subsection{Coarea formula}
\begin{theorem}[Coarea Formula:{\cite[Theorem 2.93\& Remark 2.94]{MR1857292}}]\label{teo67a715db}
	Let $\Omega\subset\R^n$ be an open set and $F:\Omega\to\R^k$ a $C^1$-submersion, that is, a $C^1$-smooth map with surjective differential at each point.
	As a consequence, we have that $F(\Omega)$ is open in $\R^k$ and that, for every $y\in F(\Omega)$, the set $F^{-1}(y)\subset\Omega$ is an immersed submanifold of dimension $n-k$.
	
	Then, for every $u\in L^1(\Omega)$ with compact support,
	\begin{equation}\label{eq67a715e0}
		\int_\Omega u(x) J(DF(x)) \did x 
		= \int_{F(\Omega)} \int_{F^{-1}(y)} u(x) \dd S^{n-k}(x) \dd y ,
	\end{equation}
	where
	\begin{equation}\label{eq67a71bad}
		J(DF(x)) = \sqrt{ \det(DF(x)\times DF(x)^T) }
		= \sqrt{\sum_{B\in\{k\times k\text{ minors of }DF(x)\}} \det(B)^2 }  .
	\end{equation}
\end{theorem}

%If $k=1$, that is, if $F:\R^n\to\R$ is of class $C^1$, then
%\[
%J(DF(x)) = |\grad F(x)| .
%\]
\begin{exercise}
	Compute $J(DF)$ as in~\eqref{eq67a71bad} when $k=1$ and when $k=n-1$.
\end{exercise}

%%%%%%%%%%%%%%%%%%%%%%%%%%%%%%%%%%%%%%%%%%%%%%%%%%%%%%%%%%%%%%%
\subsection{Consequences of the coarea formula: spherical integrals}

\begin{proposition}
	If $\Omega\subset\R^n$ is open, then 
	\begin{equation}\label{eq67a71e03}
		\int_{\Omega} u(x) \did x 
		= \int_0^\infty \int_{\de B(0,r)\cap \Omega} u(x) \did S^{n-1}(x) \did r 
		\qquad \forall u\in C^0(\Omega).
	\end{equation}
\end{proposition}
\begin{proof}
	Consider the function $F:\R^n\setminus\{0\}\to\R$, $F(x) = |x|$.
	Then $J(DF(x)) = |\grad F(x)| = 1$ for all $x\in \R^n\setminus\{0\}$.
	So, Coarea Formula~\eqref{eq67a715e0} gives immediately the identity~\eqref{eq67a71e03}.
%	The second identity in~\eqref{eq67a71e03} is a consequence of~\eqref{eq67a718d1}.
\end{proof}



%%%%%%%%%%%%%%%%%%%%%%%%%%%%%%%%%%%%%%%%%%%%%%%%%%%%%%%%%%%%%%%
\subsection{Consequences of the coarea formula: on the sphere}
From Theorem~\ref{teo67a715db}, we can deduce seemingly more general results.
For instance, a coarea formula on the sphere:

\begin{proposition}\label{prop67a71b34}
	Let $\bb S^{n-1}$ the unit sphere in $\R^n$ centered at $0$.
	If $f\in C^1(\R^n;\R)$, then, for every $u\in C^0(\bb S^{n-1})$,
	\begin{equation}
		\int_{\bb S^{n-1}} u(x) |\grad f(x)-(\grad f(x)\cdot x)x| \did S^{n-1}(x) 
		= \int_\R \int_{\bb S^{n-1}\cap \{f=z\}} u(x) \did S^{n-2}(x) \did z .
	\end{equation}
	For example, if $f(x) = x_n$, then, for every $u\in C^0(\bb S^{n-1})$,
	\begin{equation}\label{eq67a72eb2}
		\int_{\bb S^{n-1}} u(x) \sqrt{1-x_n^2} \did S^{n-1}(x) 
		= \int_{-1}^1 \int_{\bb S^{n-1}\cap \{f=z\}} u(x) \did S^{n-2}(x) \did z .
	\end{equation}
	Moreover, we have 
	\begin{equation}\label{eq67a72eb7}
		\int_{\bb S^{n-1}} u(x) \did S^{n-1}(x) 
		= \int_{-1}^1 \int_{\bb S^{n-1}\cap \{x_3=z\}} u(x) \did S^{n-2}(x) \frac1{\sqrt{1-z^2}}\did z ,
	\end{equation}
	and
	\begin{equation}\label{eq67a735b3}
		\int_{\de B(0,r)} u(x) \did S^{n-1}(x) 
		= \int_{-r}^r \int_{\de B(0,r)\cap \{x_3=z\}} u(x) \did S^{n-2}(x) \frac1{\sqrt{r^2-z^2}}\did z .
	\end{equation}
\end{proposition}
\begin{proof}
	Fix $u\in C^0(\bb S^{n-1})$ and $\epsilon>0$.
	Define $\Omega_\epsilon = B(0,1+\epsilon)\setminus B(0,1)$ and 
	$\tilde u\in C^0(\Omega_\epsilon)$ by $\tilde u(x) = u(x/|x|)$.
%	Notice that $J(Df(x)) = |\grad f(x)|$.
%	\begin{align*}
%		\int_{\Omega_\epsilon} \tilde u(x) |\grad f(x)| \did x 
%		&\overset{\eqref{eq67a71e03}}= \int_1^{1+\epsilon} \int_{\de B(0,r)} u(x/r) |\grad f(x)| \dd S^{n-1}(x) \\
%		&\overset{\eqref{eq67a718d1}}= \int_1^{1+\epsilon} r^{n-1} \int_{\bb S^{n-1}} u(x) |\grad f(rx)| \dd S^{n-1}(x) \dd r\\
%	\end{align*}
	
	Define $F:\R^n\setminus\{0\}\to \R^2$ by $F(x) = (|x|,f(x))$.
	Notice that
	\begin{align*}
		DF(x) &= \begin{pmatrix} \cdots & x/|x| & \cdots \\ \cdots & \grad f(x) & \cdots \end{pmatrix} , \\
		DF(x)DF(x)^T &= \begin{pmatrix}1&\frac{x}{|x|}\cdot\grad f(x) \\ \frac{x}{|x|}\cdot\grad f(x) & \grad f(x) \cdot \grad f(x)\end{pmatrix} , \\
		\det( DF(x)DF(x)^T ) &= |\grad f(x)|^2 - \left(\frac{x}{|x|}\cdot\grad f(x)\right)^2 \\
			&= \left| \grad f(x) - \left(\frac{x}{|x|}\cdot\grad f(x) \right) \frac{x}{|x|} \right|^2 .
	\end{align*}
	On the one hand, using the coarea formula for $F$, we have
	\begin{align*}
	\int_{\Omega_\epsilon} \tilde u(x) J(DF(x)) \did x 
	&\overset{\eqref{eq67a715e0}}= \int_{F(\Omega_\epsilon)} \int_{\{F(x)=y\}} \tilde u(x) \did S^{n-2}(x) \did y \\
	&= \int_1^{1+\epsilon} \int_{\R} \int_{\{|x| = r, f(x) = z\}} u(x/r) \did S^{n-2}(x) \did z \did r .
	\end{align*}
	On the other hand, using the coarea formula for spheres, that is~\eqref{eq67a71e03},
	\begin{align*}
	\int_{\Omega_\epsilon} \tilde u(x) J(DF(x)) \did x 
	&\overset{\eqref{eq67a71e03}}= \int_1^{1+\epsilon} \int_{\de B(0,r)} \tilde u(x) J(DF(x)) \did S^{n-1}(x) \did r \\
	&\overset{\eqref{eq67a718d1}}= \int_1^{1+\epsilon} r^{n-1} \int_{\de B(0,1)}  u(x) J(DF(rx)) \did S^{n-1}(x) \did r .
	\end{align*}
	Therefore,
	\begin{align*}
	&\hspace{-1cm}\int_{\bb S^{n-1}} u(x) |\grad f(x)-(\grad f(x)\cdot x)x| \did S^{n-1}(x) \\
	&= \lim_{\epsilon\to0} \fint_1^{1+\epsilon} r^{n-1} \int_{\de B(0,1)}  u(x) J(DF(rx)) \did S^{n-1}(x) \did r \\
	&= \lim_{\epsilon\to0} \fint_1^{1+\epsilon} \int_{\R} \int_{\{|x| = r, f(x) = z\}} u(x/r) \did S^{n-2}(x) \did z \did r  \\
	&= \int_\R \int_{\bb S^{n-1}\cap \{f=z\}} u(x) \did S^{n-2}(x) \did z .
	\end{align*}
	
	Formula~\eqref{eq67a72eb2} follows from a direct computation.
	
	The subsequent formula~\eqref{eq67a72eb7} is instead the result of a limit.
	For $u\in C^0(\bb S^{n-1})$, set 
	\[
	u_\epsilon(x) := \frac{\min\{ 1 , (1-x_n^2)/\epsilon \}}{\sqrt{1-x_3^2}} \cdot u(x) .
	\]
%	Then $u_\epsilon(x) \to \frac{u(x)}{\sqrt{1-x_3^2}}$ for all $x$ with $x_n\neq0$.
%	Use then the Dominated Convergence Theorem~\ref{thm6798ec42}.
%	\footnote{This part of the proof needs some further care...} 
	Then we plug $u_\epsilon$ into~\eqref{eq67a72eb2}.
	On the one hand, we have
	\begin{align*}
	&\int_{\bb S^{n-1}} u_\epsilon(x) \sqrt{1-x_n^2} \did S^{n-1}(x) \\
	&= \int_{\bb S^{n-1}\cap \{1-x_n^2\ge\epsilon\} } u(x) \did S^{n-1}(x) 
		+ \int_{\bb S^{n-1}\cap \{1-x_n^2<\epsilon\} } u(x) \frac{1-x_n^2}{\epsilon} \did S^{n-1}(x) .
	\end{align*}
	If we take the limit $\epsilon\to0$ we obtain
	\[
	\lim_{\epsilon\to 0} \int_{\bb S^{n-1}\cap \{1-x_n^2\ge\epsilon\} } u(x) \did S^{n-1}(x)
	= \int_{\bb S^{n-1}} u(x) \did S^{n-1}(x)
	\]
	and
	\[
	\lim_{\epsilon\to 0} \int_{\bb S^{n-1}\cap \{1-x_n^2<\epsilon\} } u(x) \frac{1-x_n^2}{\epsilon} \did S^{n-1}(x)
	= 0 .
	\]
	On the other hand, we have
	\begin{align*}
		&\hspace{-1cm}\int_{-1}^1 \int_{\bb S^{n-1}\cap \{x_3=z\}} u_\epsilon(x) \did S^{n-2}(x) \did z \\
		&= \int_{-\sqrt{1-\epsilon}}^{\sqrt{1-\epsilon}} \int_{\bb S^{n-1}\cap \{x_3=z\}} \frac{u(x)}{\sqrt{1-x_n^2}} \did S^{n-2}(x) \did z\\
		&\qquad+ \int_{[-1,1]\setminus[-\sqrt{1-\epsilon},\sqrt{1-\epsilon}]} \int_{\bb S^{n-1}\cap \{x_3=z\}} u(x) \frac{\sqrt{1-z^2}}{\epsilon} \did S^{n-2}(x) \did z .
	\end{align*}
%		\begin{equation}\label{eq67a72eb2}
%		\int_{\bb S^{n-1}} u(x) \sqrt{1-x_n^2} \did S^{n-1}(x) 
%		= \int_{-1}^1 \int_{\bb S^{n-1}\cap \{f=z\}} u(x) \did S^{n-2}(x) \did z .
%	\end{equation}
	Taking the limit $\epsilon\to0$, we obtain the right-hand side of~\eqref{eq67a72eb7}.
	
	The last formula~\eqref{eq67a735b3} follows from the previous~\eqref{eq67a72eb7}:
	take $\tilde u(x):= u(rx)$. Then
	\begin{align*}
	\int_{r\de B(0,r)} u(x) \did S^{n-1}(x)
	&= r^{n-1}  \int_{\de B(0,1)} u(rx) \did S^{n-1}(x) \\
	&\overset{\eqref{eq67a72eb7}}= r^{n-1} \int_{-1}^1 \int_{\bb S^{n-1}\cap \{x_3=z\}} u(rx) \did S^{n-2}(x) \frac1{\sqrt{1-z^2}}\did z \\
	&\overset{\bar z=rz}= r^{n-1} \int_{-r}^r \int_{\bb S^{n-1}\cap \{x_3=\bar z/r\}} u(rx) \did S^{n-2}(x) \frac1{\sqrt{1-(\bar z/r)^2}} \frac{\did \bar z}{r} \\
	&\overset{\bar x = rx}= \int_{-r}^r \int_{r\bb S^{n-1}\cap \{\bar x_3=\bar z\}} u(\bar x) \did S^{n-2}(\bar x) \frac{r}{\sqrt{r^2-\bar z^2}} \frac{\did \bar z}{r} \\
	&= \int_{-r}^r \int_{r\bb S^{n-1}\cap \{\bar x_3=\bar z\}} u(\bar x) \did S^{n-2}(\bar x) \frac{1}{\sqrt{r^2-\bar z^2}} \did \bar z .
	\end{align*}
\end{proof}

%%%%%%%%%%%%%%%%%%%%%%%%%%%%%%%%%%%%%%%%%%%%%%%%%%%%%%%%%%%%%%%
%%%%%%%%%%%%%%%%%%%%%%%%%%%%%%%%%%%%%%%%%%%%%%%%%%%%%%%%%%%%%%%
\subsection{Change of variables}

\begin{theorem}[Change of variables]\label{thm679f367e}
	Let $X,Y\subset\R^n$ be open sets and $\phi:X\to Y$ a diffeomorphism.
	Then, for every $u\in L^1(Y)$,
	\begin{equation}\label{eq679f3376}
		\int_X u(\phi(x)) |J\phi(x)| \did x = \int_Y u(y) \did y ,
	\end{equation}
	where $J\phi(x) = \det(D\phi(x))$.
\end{theorem}

Theorem~\ref{eq679f3376} holds for the {\it non-oriented} integral. 
When dealing with oriented integrals, we need to take care of the sign of the jacobian $J\phi$.
For example, line integrals are usually oriented:
If $a<b$ and $f\in L^1([a,b])$,
then $\int_a^bf(s)\did s$ is a {\it oriented} integral, while $\int_{[a,b]}f(s)\did s$ is {\it not oriented}.
This distinction becomes clear by the identities
\begin{equation}\label{eq679f35fd}
	\int_{[a,b]} f(s) \did s = \int_a^b f(s) \did s =  - \int_b^a f(s) \did s .
\end{equation}
The root of the distinction is the following.
A non-oriented integral is the integral of a function $f$ over a measure space $(X,\mu)$:
$\int_X f \did\mu$. 
A oriented integral is the integral of a differential form:
for example, $\int_{\R^2} f(x,y) \dd x\wedge\dd y = - \int_{\R^2} f(x,y) \dd y\wedge\dd x$.
In the case of integrals of differential forms, the change of variables states:
$\int \phi^* \omega = \int \omega$
and $\phi^*(\dd x_1\wedge\dots\wedge\dd x_n) = J\phi\cdot (\dd x_1\wedge\dots\wedge\dd x_n)$.

On $\R^n$ with $n>1$, we usually think of ``$\dd x$'' as $\dd\cal L^n(x)$, where $\cal L^n$ is the Lebesgue measure, and not as the volume form $\dd x_1\wedge\dots\wedge\dd x_n$.
On $\R$, instead, we usually think of ``$\dd x$'' as a 1-form.

Forgetting the details, the punchline is that, on line integrals, we need to keep in mind the indetities~\eqref{eq679f35fd}.





%%%%%%%%%%%%%%%%%%%%%%%%%%%%%%%%%%%%%%%%%%%%%%%%%%%%%%%%%%%%%%%
%%%%%%%%%%%%%%%%%%%%%%%%%%%%%%%%%%%%%%%%%%%%%%%%%%%%%%%%%%%%%%%

%\newpage
\section{Spherical averages (or means)}\label{sec67bac9ff}

Let $X\subset\R^m$ and $Y\subset\R^n$ be open sets.
For $u\in L^1_\loc(X\times Y)$, $x\in X$, $y\in Y$ and $r>0$ with $B(y,r)\subset Y$, define
\begin{equation}\label{eq67bacbb5}
\begin{aligned}
	\phi_u(x,y;r) &:= \fint_{B(y,r)} u(x,z) \did z 
		= \fint_{B_Y(0,1)} u(x,y+rz) \did z , \\
	\psi_u(x,y;r) &:= \fint_{\de B(y,r)} u(x,z) \did S(z)
		= \fint_{\de B(0,1)} u(x,y+rz) \did S(z) .
\end{aligned}
\end{equation}

If $u$ is continuous, it is clear that 
\begin{equation}\label{eq67bad1c9}
	u(x,y) = \lim_{r\to 0}\phi_u(x,y;r) = \lim_{r\to 0}\psi_u(x,y;r) .
\end{equation}

\begin{lemma}\label{lem67cd4b64}
	If $u\in C^{a;b}(X\times Y)$, with $a,b\in\N$, then, for every $\epsilon>0$,
	$\phi_u,\psi_u \in C^{a;b;\lfloor b/2 \rfloor}(X\times Y_\epsilon\times(0,\epsilon))$,
	where $Y_\epsilon = \{y\in Y: d(y,\de Y)>\epsilon\}$.
	
	Moreover, for all $\alpha\in\N^m$ and $\beta\in\N^n$, with $|\alpha|\le a$ and $|\beta|\le b$,
	\begin{equation}\label{eq67a36bcb}
		D_x^\alpha D_y^\beta \phi_u 
		= \phi_{D_x^\alpha D_y^\beta u}
		\qquad\text{and}\qquad
		D_x^\alpha D_y^\beta \psi_u 
		= \psi_{D_x^\alpha D_y^\beta u} .
	\end{equation}
	and
	\begin{equation}\label{eq67bacb9a}
	\begin{aligned}
		\de_r\phi_u(x,y;r) &= \frac{n}{r} (\psi_u(x,y;r)- \phi_u(x,y;r) ) , \\
		\de_r\psi_u(x,y;r) &= \frac{r}{n} \phi_{\laplacian_yu}(x,y;r) 
			\overset{\eqref{eq67a36bcb}}= \frac{r}{n} \laplacian_y\phi_{u}(x,y;r) .
	\end{aligned}
	\end{equation}
\end{lemma}
\begin{proof}
	Notice that both definitions of $\phi_u$ and $\psi_u$ falls into the framework of Theorem~\ref{thm679f29de}, which then implies~\eqref{eq67a36bcb}.
	Moreover,
	\begin{align*}
		\phi_u(x,y;r) 
		&= \frac{1}{\alpha_nr^n} \int_{B(y,r)} u(x,z) \did z \\
		&= \frac{1}{\alpha_nr^n} \int_0^r \int_{\de B_Y(y,s)} u(x,z) \did S(z) \did s \\
%		&= \frac{n}{r^n} \int_0^r s^{n-1} \fint_{\de B_Y(y,s)} u(x,z) \did S(z) \did s \\
		&= nr^{-n} \int_0^r s^{n-1} \psi_u(x,y;s) \did s .
	\end{align*} 
	Therefore,
	\begin{align*}
		\de_r\phi_u(x,y;r)
		&= -n^2 r^{-n-1} \int_0^r s^{n-1} \psi(x,y;s) \did s
			+ nr^{-1} \psi_u(x,y;r) \\
		&= \frac{n}{r} (- \phi_u(x,y;r) + \psi_u(x,y;r)) .
	\end{align*}
	Finally,
	\begin{align*}
		\de_r\psi_u(x,y;r)
		&= \frac1{n\alpha_n}  \de_r \int_{\de B(0,1)} u(x,y+rz) \did S(z) \\
		[\text{by Theorem~\ref{thm679f29de}}]
		&= \frac1{n\alpha_n} \int_{\de B(0,1)} (D_yu)(x,y+rz) \cdot z \did S(z) \\
		[\text{by Divergence Theorem and~\eqref{eq67a371f8}}]
		&= \frac{r}{n\alpha_n} \int_{B(0,1)} (\laplacian_yu)(x,y+rz) \did z \\
		&= \frac{r}{n} \phi_{\laplacian_yu}(x,y;r) ,
	\end{align*}
	where
	\begin{equation}\label{eq67a371f8}
	\div_z((D_yu)(x,y+rz)) 
	= r(\div_y(D_yu))(x,y+rz)
	= r(\laplacian_yu)(x,y+rz) .
	\end{equation}
\end{proof}

%%%%%%%%%%%%%%%%%%%%%%%%%%%%%%%%%%%%%%%%%%%%%%%%%%%%%%%%%%%%%%%
\section{Change of coordinates}

\subsection{Differential operator}\label{subs67d7dce6}
Let $\Omega\subset\R^n$ be open.
A \emph{linear differential operator with smooth coefficients} on $\Omega$ is a map $P:C^\infty(\Omega) \to C^\infty(\Omega)$ of the form
\begin{equation}
	P\phi = \sum_{\alpha\in\N^n} P_\alpha \Diff^\alpha \phi ,
\end{equation}
for some $P_\alpha\in C^\infty(\Omega)$, all zero but for finitely many indices.

%.
%We usually assume that $P$ is linear and local, that is, if $\spt(P\phi) \subset\spt(\phi)$.
%This ensures that $P$ is 

%%%%%%%%%%%%%%%%%%%%%%%%%%%%%%%%%%%%%%%%%%%%%%%%%%%%%%%%%%%%%%%
\subsection{An abstract overview}\label{subs67d7dc9c}
Let $\Omega_1,\Omega_2\subset\R^n$ be open sets and $\Phi:\Omega_1\to\Omega_2$ a diffeomorphism.
We have pull-backs and push-forwards of both functions and differential operators:
\begin{itemize}
\item	Pull-back of functions: $\Phi^*:C^\infty(\Omega_2)\to C^\infty(\Omega_1)$, $\Phi^*\psi = \psi\circ\Phi$ for $\psi\in C^\infty(\Omega_2)$;
\item	Push-forward of functions: $\Phi_*:C^\infty(\Omega_1)\to C^\infty(\Omega_2)$, $\Phi_*\phi = \phi\circ\Phi^{-1}$ for $\psi\in C^\infty(\Omega_2)$;
\item	Push-forward of differential operators: if $P_1:C^\infty(\Omega_1)\to C^\infty(\Omega_1)$ is a differential operator, then $\Phi_*P_1 = \Phi_*\circ P_1 \circ \Phi^*$;
\item	Pull-back of differential operators: if $P_2:C^\infty(\Omega_2)\to C^\infty(\Omega_2)$ is a differential operator, then $\Phi^*P_2 = \Phi^*\circ P_2 \circ \Phi_*$.
\end{itemize}
The following diagrams might help:
\begin{equation}\label{eq67d7e093}
	\xymatrix{
	\Omega_1 \ar[r]^{\Phi} & \Omega_2 \\
	C^\infty(\Omega_1)\ar[d]_{P_1} & C^\infty(\Omega_2) \ar[l]_{\Phi^*}\ar[d]^{\Phi_*P_1} \\
	C^\infty(\Omega_1)\ar[r]^{\Phi_*} & C^\infty(\Omega_2) \\
	}
	\qquad
	\xymatrix{
	\Omega_1 \ar[r]^{\Phi} & \Omega_2 \\
	C^\infty(\Omega_1)\ar[d]_{\Phi^*P_2}\ar[r]^{\Phi_*} & C^\infty(\Omega_2) \ar[d]^{P_2} \\
	C^\infty(\Omega_1) & C^\infty(\Omega_2)\ar[l]_{\Phi^*} \\
	}
\end{equation}

Notice that, for all functions $\phi\in C^\infty(\Omega_1)$, $\psi\in C^\infty(\Omega_2)$, 
and all differential operators $P_1:C^\infty(\Omega_1)\to C^\infty(\Omega_1)$ and $P_2:C^\infty(\Omega_2)\to C^\infty(\Omega_2)$,
\begin{align*}
	\Phi_*\phi &= (\Phi^{-1})^*\phi = (\Phi^*)^{-1}\phi ;\\
	(\Phi_*P_1)\psi & = (P_1(\phi\circ \Phi))\circ\Phi^{-1} ; \\
	(\Phi^*P_2)\phi &=  (P_2(\phi\circ\Phi^{-1}))\circ \Phi ; \\
	\Phi_* \Phi^* &= \Id,\text{ and }\Phi^*\Phi_* = \Id .
\end{align*}

Notice that, if $P_j,Q_j$ are  differential operators, then
\begin{equation}\label{eq67d7eb47}
	\Phi^*(P_2\circ Q_2) = (\Phi^*P_2)\circ (\Phi^*Q_2),
	\text{ and }
	\Phi_*(P_1\circ Q_1) = (\Phi_*P_1)\circ(\Phi_*Q_1) .
\end{equation}

\begin{exercise}
	A differential operator of order zero on $\Omega$ is of the form $P\phi = f\cdot\phi$ for some $f\in C^\infty(\Omega)$.
	Compute $\Phi_*P$ and $\Phi^*P$ for differential operators $P$ of order zero.
	Is it coherent with pull-back and push-forward of functions?
\end{exercise}

\begin{exercise}\label{ex67d7e500}
	Let $\Phi:\Omega_1\to\Omega_2$ be a diffeomorphism between open subsets of $\R^n$.
	For $j\in\{1,\dots,n\}$, compute $\Phi_*\de_j$ and $\Phi^*\de_j$.
\end{exercise}

\begin{exercise}[Harder]
	Consider $\Phi:\R^n\to\R^n$, $\Phi(x) = Ax$ for $A\in\GL(\R^n)$, that is $A$ is a $n\times n$ invertible matrix.
	For $\alpha\in\N^n$, compute $\Phi^*\Diff^\alpha$ and $\Phi_*\Diff^\alpha$.
\end{exercise}



%%%%%%%%%%%%%%%%%%%%%%%%%%%%%%%%%%%%%%%%%%%%%%%%%%%%%%%%%%%%%%%
\subsection{An example: Laplacian in polar coordinates}
Consider 
\begin{equation}
	\Phi : (0,+\infty) \times (-\pi,\pi) \to \R^2 \setminus (\{0\}\times[-\infty,0]),
	\qquad
	\Phi(r,\theta) = (r\cos(\theta),r\sin(\theta)) .
\end{equation}
This function $\Phi$ is a \emph{polar parametrization} of the plane:
\emph{polar coordinates} are in fact the inverse function of $\Phi$.
In other words, we can define functions 
\begin{equation}\label{eq67d7eca2}
	r,\theta: \R^2 \setminus (\{0\}\times[-\infty,0])\to\R ,
	\text{ such that }\Phi(r(x,y),\theta(x,y)) = (x,y) .
\end{equation}

Anyway, with the function $\Phi$ above, we want to compute $\Phi^*\laplacian$.
To do this, using Exercise~\ref{ex67d7e500}, we first compute
\begin{equation}
	\Diff\Phi(r,\theta) 
	= \begin{pmatrix} \cos\theta & - r\sin\theta \\ \sin\theta & r\cos\theta \end{pmatrix} ,
	\qquad
	\Diff\Phi(r,\theta)^{-1}
	= \frac1r \begin{pmatrix} r \cos\theta & r\sin\theta \\ - \sin\theta & \cos\theta\end{pmatrix} .
\end{equation}
Therefore,
\begin{align*}
 	(\Phi^*\de_x)(r,\theta) 
	&= \Diff\Phi(r,\theta)^{-1} \left(\de_x|_{\Phi(r,\theta)}\right)
	= \cos\theta \de_r - \frac{\sin\theta}{r} \de_\theta \\
	(\Phi^*\de_y)(r,\theta)
	&= \Diff\Phi(r,\theta)^{-1} \left(\de_y|_{\Phi(r,\theta)}\right)
	= \sin\theta\de_r + \frac{\cos\theta}{\de r} .
\end{align*}
Using now~\eqref{eq67d7eb47}, we obtain
\begin{equation}\label{eq67d7ec01}
	\Phi^*\laplacian  
	= (\Phi^*\de_x)^2 + (\Phi^*\de_y)^2
	= [...]
	= \de_r^2 + \frac{1}{r^2} \de_\theta^2 + \frac1r \de_r .
\end{equation}
This formula represents ``the laplacian in polar coordinates''.

%%%%%%%%%%%%%%%%%%%%%%%%%%%%%%%%%%%%%%%%%%%%%%%%%%%%%%%%%%%%%%%
\subsection{Another interpretation of the previous change of variables}
We have another interpretation of the computation we made in~\eqref{eq67d7ec01}.
Define the vector fields
\begin{align*}
	\vec v_r(x,y) &= \cos\theta \de_x + \sin\theta\de_ y 
		= \frac{x\de_x-y\de_y}{\sqrt{x^2+y^2}} , \\
	\vec v_\theta(x,y) &= r\sin\theta \de_r + r\cos\theta \de_y 
		= -y\de_x + x\de_y ,
\end{align*}
where we see $\theta$ and $r$ as the functions defined in~\eqref{eq67d7eca2}.
These are the push-forward vector fields $\Phi_*\de_r$ and $\Phi_*\de_\theta$, respectively.
As such, one usually simply writes $\de_r$ and $\de_\theta$ for $\vec v_r$ and $\vec v_\theta$.
We keep the distinction here for purely educational purposes.

At this point, we can write $\de_x$ and $\de_y$ in terms of $\vec v_r$ and $\vec v_\theta$:
\begin{align*}
 	\de_x = \cos\theta \vec v_r - \frac{\sin\theta}{r} \vec v_\theta ,\\
	\de_y = \sin\theta \vec v_r + \frac{\cos\theta}{r} \vec v_\theta .
\end{align*}

We can now perform the computation~\eqref{eq67d7ec01} again as
\begin{equation}
	\laplacian
	= \de_x^2 + \de_y^2
	= (\cos\theta \vec v_r - \frac{\sin\theta}{r} \vec v_\theta)^2	
		+ (\sin\theta \vec v_r + \frac{\cos\theta}{r} \vec v_\theta)^2
	= [...]
	= \vec v_r^2 + \frac{1}{r^2} \vec v_\theta^2 + \frac1r \vec v_r ,
\end{equation}
where, we recall, $r,\theta$ are the functions defined in~\eqref{eq67d7eca2},
and ``$\vec v_r = \de_r$'', and ``$\vec v_\theta = \de_\theta$''.

%%%%%%%%%%%%%%%%%%%%%%%%%%%%%%%%%%%%%%%%%%%%%%%%%%%%%%%%%%%%%%%
\subsection{Exercise: Laplacian in polar coordinates in arbitrary dimension}
\begin{exercise}
	Consider the polar coordinates in $\R^n$ given by the function
	$\Phi:(0,+\infty)\times \R^{n-1}\to\R^n$,
	\begin{equation}
		\Phi(r,\theta_1,\dots,\theta_{n-1})
		= r ( \cos\theta_1, \sin\theta_1\cos\theta_2, \sin\theta_1\sin\theta_2\cos\theta_3, \dots, \sin\theta_1\cdots \sin\theta_{n-1} ) .
	\end{equation}
	\begin{enumerate}
	\item	Determine domain and image of $\Phi$ so that it becomes a diffeomorphism;
	\item	Compute the laplacian in polar coordinates in $\R^n$.
	\end{enumerate}
\end{exercise}

%%%%%%%%%%%%%%%%%%%%%%%%%%%%%%%%%%%%%%%%%%%%%%%%%%%%%%%%%%%%%%%
\section{Mollifiers}
\index{mollifier}
\index{standard family of mollifiers}
\index{family of mollifiers}

Let $\rho\in C^\infty_c(\R^n)$ be such that $\spt(\rho)=B(0,1)$, $0\le \rho\le 1$, $\rho(-x)=\rho(x)$ for all $x$, and $\int\rho\did x = 1$.
Define $\rho_\epsilon(x) = \rho(x/\epsilon)/\epsilon^n$.
We call the family $\{\rho_\epsilon\}_{\epsilon>0}$ an \emph{approximation of the identity on $\R^n$}, or \emph{a family of mollifiers}.

For example, one can take
\begin{equation}\label{eq688c7260}
	\rho(x) = \begin{cases}
	k \exp\left(\frac{1}{|x|^2-1}\right) & \text{ if }|x|<1 \\
	0 & \text{ otherwise} ,
	\end{cases}
\end{equation}
where $k$ normalizes the integral.
The family of mollifiers given by the function $\rho$ in~\eqref{eq688c7260} is called \emph{standard family of mollifiers}.

\begin{exercise}
	Show that the function 
	\begin{equation}
		\rho_0(x) = \begin{cases}
		\exp\left(\frac{1}{|x|^2-1}\right) & \text{ if }|x|<1 \\
		0 & \text{ otherwise} ,
		\end{cases}
	\end{equation}
	is $C^\infty$-smooth on $\R^n$ and compute $\int_\R^n \rho_0(x)\did x$.
	Show also that $\rho_0$ is not analytic.
	Does it exist a family of analytic mollifiers?	
%	Finally, why have I put the square in the definition of $\rho$? Do you think we could do without?
%	
%	{\it Hint:} I have put the square to help you. The square is itself a hint.
\end{exercise}

\TODO
PROPRIETA' DEI MOLLIFICATORI.

\begin{proposition}\label{prop688c72e1}
	$u*\rho_\epsilon\in C^\infty$ \TODO
\end{proposition}



%%%%%%%%%%%%%%%%%%%%%%%%%%%%%%%%%%%%%%%%%%%%%%%%%%%%%%%%%%%%%%%
%%%%%%%%%%%%%%%%%%%%%%%%%%%%%%%%%%%%%%%%%%%%%%%%%%%%%%%%%%%%%%%
%%%%%%%%%%%%%%%%%%%%%%%%%%%%%%%%%%%%%%%%%%%%%%%%%%%%%%%%%%%%%%%
%%%%%%%%%%%%%%%%%%%%%%%%%%%%%%%%%%%%%%%%%%%%%%%%%%%%%%%%%%%%%%%
\newpage
\part{Four classical PDE}

\section{Transport Equation}
\index{transport equation}
\index{equation!transport --}

This is \cite[Section 2.1]{MR2597943}.
The \emph{transport equation} is the initial value problem
\begin{equation}\label{eq67a79b9b}
	\begin{cases}
	\de_t u + b\cdot \grad u = f & \text{ on }\R^n\times I , \\
	u = g & \text{ on }\R^n\times\{t_0\} ,
	\end{cases}
\end{equation}
where $I\subset\R$ is an interval, $t_0\in I$, $b\in\R^n$, $f:\R^n\times I\to\C$ and $g:\R^n\to\C$.
The function $u$ is intended as a function in two variables, $u=u(x,t)$, where $x\in\R^n$ and $t\in I$.
The derivative $\de_tu$ is the derivative along the second variable, $t$, 
while the gradient $\grad u$ is the derivative of $u$ in the variable $x$.

If $f=0$, then we call~\eqref{eq67a79b9b} the \emph{homogeneous transport equation}.
If $f\neq0$, then~\eqref{eq67a79b9b} is the \emph{nonhomogeneous transport equation}.

A solution to~\eqref{eq67a79b9b} is easily found.

\begin{theorem}\label{thm67d7b17c}
	Let $I\subset \R$ be an open interval, $t_0\in \bar I$, $b\in\R^n$, $f\in C^0(\R^n\times \bar I)$, and $g\in C^1(\R^n)$.
	Then the function $u\in C^1(\R^n\times \bar I)$ defined by
	\begin{equation}\label{eq67a79d37}
		u(x,t) = g(x-(t-t_0)b) + \int_{t_0}^t f(x+(r-t)b,r) \did r
	\end{equation}
	is the unique solution to~\eqref{eq67a79b9b}.
\end{theorem}
\begin{proof}
	To show that $u$ is a solution to~\eqref{eq67a79b9b}, we just compute the derivatives:
	\begin{align*}
	\de_t u(x,t) &= - \grad g(x-(t-t_0)b)\cdot b + f(x,t) - \int_{t_0}^t \grad f(x+(r-t)b,r)\cdot b \did r ,\\
	\grad u(x,t) &= \grad g(x-(t-t_0)b) + \int_{t_0}^t \grad f(x+(r-t)b,r) \did r .
	\end{align*}
	It follows that $u$ solves~\eqref{eq67a79b9b}.
	
	To show that $u$ is unique, suppose that $\tilde u\in C^1(\R^n\times \bar I)$ is another solution.
	Then the difference $w := u-\tilde u$ solves~\eqref{eq67a79b9b} with $f=0$ and $g=0$. 
	Notice that 
	\[
	\frac{\dd}{\dd s} w(x+ sb,t+s) = \grad w(x+sb,t+s)\cdot b + \de_t w(x+sb,t+s) = 0
	\]
	Therefore, for all $(x,t)\in \R^n\times I$, we have $w(x,t) = w(x+(t_0-t)b,t_0) = 0$.
\end{proof}

\begin{remark}
	The function $u$ defined in~\eqref{eq67a79d37} is well defined even in the case $g$ is not smooth.
	In some sense, these other functions could be regarded as \emph{weak solutions} to~\eqref{eq67a79b9b} when $g$ is not $C^1$.
\end{remark}

%%%%%%%%%%%%%%%%%%%%%%%%%%%%%%%%%%%%%%%%%%%%%%%%%%%%%%%%%%%%%%%
%%%%%%%%%%%%%%%%%%%%%%%%%%%%%%%%%%%%%%%%%%%%%%%%%%%%%%%%%%%%%%%
\newpage
\section{Laplace Equation}

%%%%%%%%%%%%%%%%%%%%%%%%%%%%%%%%%%%%%%%%%%%%%%%%%%%%%%%%%%%%%%%
\subsection{Laplace operator}\label{subs67b9d4b5}
\index{Laplace operator}
\index{operator!Laplace --}

For $U\subset\R^n$ open and $u\in C^2(U)$, define the \emph{laplacian of $u$} as
\begin{equation}
	\laplacian u \coloneqq \div(\grad u)) = \sum_{j=1}^n \frac{\de^2 u}{\de x_j^2} .
\end{equation}
The operator $\laplacian \coloneqq \sum_{j=1}^n \de_j^2$ is the \emph{Laplace operator}.

%%%%%%%%%%%%%%%%%%%%%%%%%%%%%%%%%%%%%%%%%%%%%%%%%%%%%%%%%%%%%%%
\subsection{Harmonic function}\label{subs67b9d542}
\index{harmonic!-- function}
\index{function!harmonic --}

A \emph{harmonic function} on $U\subset\R^n$ open is a function $u\in C^2(U)$ such that $\laplacian u=0$.

\begin{exercise}\label{ex68347bbd}
	Show that the only harmonic functions on $\R$ are the affine functions.
\end{exercise}

%%%%%%%%%%%%%%%%%%%%%%%%%%%%%%%%%%%%%%%%%%%%%%%%%%%%%%%%%%%%%%%
\subsection{Symmetries of Laplace operator}\label{subs67b9d638}
Let $U\subset\R^n$ open, $u\in C^2(U)$, $O\in\Orth(n)$, $b\in\R^n$, and $\lambda\in\R\setminus\{0\}$.
Define $\bar u(y) \coloneqq u(\lambda Oy+b)$.
Then $\bar u\in C^2(\lambda^{-1}O^{-1}(U-b))$ and
\begin{equation}
	\laplacian\bar u(y) = \lambda^2 (\laplacian u)(\lambda Oy+b).
\end{equation}

%%%%%%%%%%%%%%%%%%%%%%%%%%%%%%%%%%%%%%%%%%%%%%%%%%%%%%%%%%%%%%%
\subsection{How to find the fundamental solution: radial solutions}\label{subs67b9d8f3}

In this section we solve the following Exercise~\ref{ex6889d8bc}, which asks to find radial harmonic functions on $\R^n$:

\begin{exercise}\label{ex6889d8bc}
	For $n\ge2$,
	find smooth non-constant functions $u:\R^n\setminus\{0\}\to\C$ that are radial, that is, they only depend on the distance from the origin, and harmonic, that is, $\laplacian u =0 $.
\end{exercise}

In what follows, we solve Exercise~\ref{ex6889d8bc}.
The reader is invited to solve it first by themselves: it is actually easier than it looks like at first.

Such radial harmonic functions are expected to exist because of the symmetries of the laplacian that we have seen in the previous Section~\ref{subs67b9d638}:
since the laplacian has spherical symmetry, we expect to have harmonic functions with spherical symmetry.
We expect these functions to be somehow special: they are indeed, and we will see in the forthcoming sections that among them we find fundamental solutions of the laplacian.

Another type of symmetric harmonic functions are those that are homogeneous with respect to dilations.
I invite the student to try to characterize those too: they are the so-called \emph{spherical harmonics}...\\

Done with Exercise~\ref{ex6889d8bc}? Here is my take.

We consider functions of the form $u(x) = \phi(\abs{x}^2)$, with $\phi:(0,+\infty)\to\C$ smooth.
I choose to take the squared norm because in this way derivatives are easier.
The we have
\begin{align}
	\frac{\de u}{\de x_j}(x) &= \phi'(\abs{x}^2) 2 x_j ; \\
	\frac{\de^2 u}{\de x_j^2}(x) &= \phi''(\abs{x}^2) 4 x_j^2 + \phi'(\abs{x}^2) 2 ; \\
	\laplacian u (x) &= 4 \phi''(\abs{x}^2) \abs{x}^2 + 2n \phi'(\abs{x}^2) .
\end{align}
Since we want $\laplacian u (x) = 0$ for $x\neq0$, we obtain that $\phi$ must satisfy 
\begin{equation}\label{eq6889dc0f}
	\forall t\in(0,+\infty),\qquad
	4 \phi''(t) t + 2n \phi'(t) = 0 .
\end{equation}
Take $\psi = \phi'$: then $\psi$ must satisfy
\begin{equation}\label{eq6889dc59}
	\forall t\in(0,+\infty),\qquad
	\frac{\psi'(t)}{\psi(t)} = - \frac{n}{2} \frac1t .
\end{equation}

We are assuming $\psi(t) = \phi'(t)\neq0$ for all $t>0$, because if $\phi'(t)=0$ for some $t$, then $\phi$ constant, i.e., $\phi'\equiv0$, is a solution to the ODE~\eqref{eq6889dc0f}.
Since this ODE has unique solution given $\phi'$ at one point, we conclude that either $\phi$ is constant or $\phi'$ is never zero.
So, $\psi$ does satisfy~\eqref{eq6889dc59}.

The ODE~\eqref{eq6889dc59} is equivalent to
\begin{equation}
	\forall t\in(0,+\infty),\qquad
	\frac{\dd}{\dd t} \log(\psi(t))
	= - \frac{n}{2} \frac{\dd}{\dd t} \log(t) ,
\end{equation}
which is in turn equivalent to
\begin{equation}
	\exists c\in\C,\quad\forall t\in(0,+\infty),\qquad
	\log(\psi(t))  = - \frac{n}{2} \log(t) + c .
\end{equation}
Exponentiating and integrating, we obtain that non-constant solutions of~\eqref{eq6889dc0f} are
\begin{equation}
	\exists a,c\in\C\quad
	\forall t\in(0,+\infty),\qquad
	\phi(t) = a + e^c \int_1^t s^{-\frac{n}{2}} \did s .
\end{equation}
Here we need to distinguish tow cases:
\begin{align}
	n=2: && \phi(t) &= a + e^c \log(t) , \\
	n>2: && \phi(t) &= \left( a - \frac{2e^c}{2-n} \right) + \frac{2e^c}{2-n} t^{\frac{2-n}{2}} .
\end{align}

This might be pedantic, but we have shown that all functions required by Exercise~\ref{ex6889d8bc} are all functions of the form
\begin{align}
	n=2: && u(x) &= a + b \log(\abs{x}) , \\
	n>2: && u(x) &= a + \frac{b}{\abs{x}^{n-2}} ,
\end{align}
for every choice of $a,b\in\C$.

%%%%%%%%%%%%%%%%%%%%%%%%%%%%%%%%%%%%%%%%%%%%%%%%%%%%%%%%%%%%%%%
\begin{figure}\label{fig67b9fcba}
\resizebox{0.45\textwidth}{!}{
\begin{tikzpicture}[>=stealth]
    \begin{axis}[
        xmin=-0.1,xmax=2.1,
        ymin=-0.5,ymax=2.1,
        axis x line=middle,
        axis y line=middle,
        axis line style=->,
%        xlabel={$|x|$},
%        ylabel={$\Phi(x)$},
        ]
        \addplot[no marks,blue,-] 
        	expression[domain=0.01:1.9,samples=100]{-1/(2*pi)*ln(x)} 
            node[pos=0.80,anchor=north east]{$n=2$}; 
         \addplot[no marks,red,-]
         	expression[domain=0.01:1.9,samples=100]{1/3 1/x}
			node[pos=0.9,anchor=south]{$n=3$};
    \end{axis}
\end{tikzpicture}
}
\resizebox{0.45\textwidth}{!}{
\begin{tikzpicture}
    \begin{axis}[
        xmin=-0.1,xmax=2.1,
        ymin=-0.5,ymax=2.1,
        axis x line=middle,
        axis y line=middle,
        axis line style=->,
%        xlabel={$|x|$},
%        ylabel={$\Phi(x)$},
        ]
        \addplot[no marks,blue,-] 
        	expression[domain=0.01:1.9,samples=100]{x^2/4*(1-2*ln(x))}%{x^2/4(1-2*ln(x))} 
            node[pos=0.95,anchor=north east]{$n=2$}; 
         \addplot[no marks,red,-]
         	expression[domain=0.01:1.9,samples=100]{x^2/2}
			node[pos=0.90,anchor=south east]{$n=3$}; 
    \end{axis}
\end{tikzpicture}
}
\caption{Plots of $|x|\mapsto\Phi(x)$ (left) and of $R\mapsto\int_{B(0,R)}\Phi(x)\did x$ (right) for $n=2,3$.}
\end{figure}

%%%%%%%%%%%%%%%%%%%%%%%%%%%%%%%%%%%%%%%%%%%%%%%%%%%%%%%%%%%%%%%%
%\begin{figure}\label{fig67b9fcba}
%\resizebox{0.5\columnwidth}{!}{
%\begin{tikzpicture}
%    \begin{axis}[
%        xmin=-0.1,xmax=2.1,
%        ymin=-0.5,ymax=2.1,
%        axis x line=middle,
%        axis y line=middle,
%        axis line style=->,
%%        xlabel={$|x|$},
%%        ylabel={$\Phi(x)$},
%        ]
%        \addplot[no marks,blue,-] 
%        	expression[domain=0.01:1.9,samples=100]{x^2/4*(1-2*ln(x))}%{x^2/4(1-2*ln(x))} 
%            node[pos=0.95,anchor=north east]{$n=2$}; 
%         \addplot[no marks,red,-]
%         	expression[domain=0.01:1.9,samples=100]{x^2/2}
%			node[pos=0.90,anchor=south east]{$n=3$}; 
%    \end{axis}
%\end{tikzpicture}
%}
%\caption{A couple of graphs of $R\mapsto\int_{B(0,R)}\Phi(x)\did x$ for $n=2,3$.}
%\end{figure}

%%%%%%%%%%%%%%%%%%%%%%%%%%%%%%%%%%%%%%%%%%%%%%%%%%%%%%%%%%%%%%%
\subsection{Fundamental solution of the Laplace's equation}\label{subs67b9d63e}
\index{fundamental solution!-- of the Laplace's equation}
\index{Laplace!fundamental solution of the --'s equation}

Define $\Phi:\R^n\setminus\{0\}\to\R$ by
\begin{equation}\label{eq67b9d6b9}
	\Phi(x) = 
	\begin{cases}
	-\frac{1}{2\pi} \log(|x|) & \text{ if }n=2 , \\
	\frac1{n(n-2)\omega_n } \frac1{|x|^{n-2}} & \text{ if }n\ge 3 ,
	\end{cases}
\end{equation}
where $\omega_n $ is the volume of the unit ball in $\R^n$.
We will call this function $\Phi$ the \emph{Fundamental solution of the Laplace's equation}.
The most important property of this function is its role in Theorem~\ref{subs67ba081a}.
In fact, we will see fundamental solutions of linear operators from an abstract viewpoint in Section~\ref{subs6889d87f}.

%%%%%%%%%%%%%%%%%%%%%%%%%%%%%%%%%%%%%%%%%%%%%%%%%%%%%%%%%%%%%%%
%\subsection{Properties of the fundamental solution}\label{prop67b9d700}
\begin{proposition}[Properties of the fundamental solution]\label{prop67b9d700}
	Let $n\ge2$.
	The fundamental solution $\Phi:\R^n\setminus\{0\}\to\R$ of Laplace's equation,
	defined in~\eqref{eq67b9d6b9}, satisfies the following statements.
	\begin{enumerate}
	\item\label{prop67b9d700_anal}
	$\Phi$ is an analytic function $\R^n\setminus\{0\}\to\R$.
	If $n\ge3$, then $\Phi$ is strictly positive valued.
	\item\label{prop67b9d700_grad}
	For every $x\in\R^n\setminus\{0\}$,
	\begin{equation}\label{eq67b9dc62}
		\grad\Phi(x)
		= - \frac1{n\omega_n } \frac{x}{|x|^n} .
	\end{equation}
	\item\label{prop67b9d700_hess}
	For every $x\in\R^n\setminus\{0\}$,
	\begin{equation}\label{eq67b9f252}
		\Diff^2\Phi(x)
		= \frac1{\omega_n } \frac{x\otimes x}{|x|^{n+2}} - \frac{1}{n\omega_n } \frac{\Id}{|x|^n} .
	\end{equation}
	\item\label{prop67b9d700_laplacian}
	 $\laplacian\Phi=0$ on $\R^n\setminus\{0\}$.
	\item\label{prop67b9d700_estimates}
	There exists $C>0$ such that for every $x\in \R^n\setminus\{0\}$:
		\begin{enumerate}
		\item\label{prop67b9d700_2_a}
		 $0<\Phi(x)\le C |x|^{2-n}$ if $n\ge 3$;
		\item\label{prop67b9d700_2_b}
		 $\left| \frac{\de\Phi}{\de x_j}(x)\right| \le \frac{C}{|x|^{n-1}}$, for all $j\in\{1,\dots,n\}$;
		\item\label{prop67b9d700_2_c}
		 $\left| \frac{\de^2\Phi}{\de x_j\de x_k}(x)\right| \le \frac{C}{|x|^{n}}$, for all $j,k\in\{1,\dots,n\}$.
		\end{enumerate}
	\item\label{prop67b9d700_integrals}
	$\Phi\in L^1_\loc(\R^n)$ with 
	\begin{equation}\label{eq67ba10db}
		\int_{B(0,R)} \Phi(x)\did x =
		\begin{cases}
			\frac{R^2}{4} (1 - 2\log(R)) & \text{ if }n=2, \\
			\frac{R^2}{2(n-2)}  & \text{ if }n\ge3. 
		\end{cases}
	\end{equation}
		See Figure~\ref{fig67b9fcba} for a plot of these quantities.
	\end{enumerate}
\end{proposition}

\begin{proof}
	Part~\ref{prop67b9d700_anal} is clear.
	
	Statement~\ref{prop67b9d700_grad} is proven as follows.
	If $n=2$ and $x\in\R^2\setminus\{0\}$, then
	\begin{equation}\label{eq67b9dc00}
		\grad\Phi(x)
		= -\frac1{2\pi} \frac1{|x|} \frac{x}{|x|} 
		= -\frac1{2\pi} \frac{x}{|x|^2} .
	\end{equation}
	Therefore,
	$|\grad\Phi(x)| = \frac1{2\pi} \frac{1}{|x|}$.
	
	If $n\ge3$ and $x\in\R^n\setminus\{0\}$, then
	\begin{equation}\label{eq67b9dc05}
		\grad\Phi(x)
		= \frac1{n(n-2)\omega_n } (2-n) |x|^{2-n-1} \frac{x}{|x|}
		= - \frac1{n\omega_n } \frac{x}{|x|^n} .
	\end{equation}
	Therefore, $|\grad\Phi(x)| = \frac1{n\omega_n } \frac1{|x|^{n-1}}$.
	Notice that~\eqref{eq67b9dc00} and~\eqref{eq67b9dc05} are summarized in~\eqref{eq67b9dc62}, because $\alpha(2)=\pi$.
	
	For~\ref{prop67b9d700_hess}, we compute
	\begin{align*}
	\Diff^2\Phi(x)
	&= -\frac{1}{n\omega_n } \left( \frac{\Id}{|x|^n} + x\otimes \left( -n|x|^{-n-1} \frac{x}{|x|} \right) \right) \\
	&= - \frac{1}{n\omega_n } \frac{\Id}{|x|^n} + \frac1{\omega_n } \frac{x\otimes x}{|x|^{n+2}} .
	\end{align*}
	In other words,
	\begin{align*}
	\frac{\de^2\Phi}{\de x_j\de x_k}(x)
	&= \frac{\de}{\de x_j} \left( - \frac1{n\omega_n } \frac{x_k}{|x|^n} \right) \\
	&= - \frac1{n\omega_n } \left( \frac{\delta_{jk}}{|x|^n} - x_k n |x|^{-n-1} \frac12 (\sum_\ell x_\ell^2)^{1/2} 2x_k \right) \\
	&= - \frac1{n\omega_n }\frac{\delta_{jk}}{|x|^n} + \frac1{\omega_n } \frac{x_kx_j}{|x|^{n+2}} .
	\end{align*}
	So, we have~\eqref{eq67b9f252}.
	
	Statement~\ref{prop67b9d700_laplacian} is now an easy computation:
	if $x\in\R^n\setminus\{0\}$, then
	\begin{align*}
		\laplacian\Phi(x)  
		&= \trace(\Diff^2\Phi(x)) \\
		&= \frac1{\omega_n } \frac{\trace(x\otimes x)}{|x|^{n+2}} - \frac{1}{n\omega_n } \frac{\trace(\Id)}{|x|^n}  \\
		&= \frac1{\omega_n } \frac{1}{|x|^n} - \frac{1}{n\omega_n } \frac{n}{|x|^n}
		= 0 .
	\end{align*}

	The estimates~\ref{prop67b9d700_estimates} are a consequence of the explicit formulas we have computed.
	
	The last two formulas stated in part~\ref{prop67b9d700_integrals} are simply computed as follows.
	For $n=2$, we have
	\begin{align*}
		\int_{B(0,R)} \Phi(x)\did x
		&= - \frac1{2\pi} \int_0^R\int_0^{2\pi} \log(r) r \did\theta\did r \\
		&= - \int_0^R \log(r) r\did r \\
		\text{[Integration by parts]} 
		&= -\left(\log(r)\frac{r^2}{2}\right|_0^R + \int_0^R \frac1r \frac{r^2}{2} \did r \\
		&= - \log(R) \frac{R^2}{2} + \frac12 \frac{R^2}{2} 
		=  \frac{R^2}{4} (1 - 2\log(R)) .
	\end{align*}
	For $n\ge3$, we have
	\begin{align*}
		\int_{B(0,R)} \Phi(x)\did x
		&= \frac1{n(n-2)\omega_n } \int_{B(0,R)} \frac1{|x|^{n-2}} \did x \\
		&= \frac1{n(n-2)\omega_n } \int_0^R \int_{\de B(0,1)} \frac1{|x|^{n-2}} \did S^{n-1}(x) r^{n-1} \did r \\
		&= \frac{n\omega_n }{n(n-2)\omega_n } \int_0^R r \did r \\
		&= \frac1{(n-2)} \frac{R^2}{2} .
	\end{align*}
\end{proof}

%%%%%%%%%%%%%%%%%%%%%%%%%%%%%%%%%%%%%%%%%%%%%%%%%%%%%%%%%%%%%%%
\subsection{Solution to Poisson equation}\label{subs67ba081a}
\index{Poisson equation}
\index{equation!Poisson --}

\begin{theorem}[Solution to Poisson equation]\label{teo67ba14f2}
	Let $f\in C^2_c(\R^n)$ and define for $x\in\R^n$
	\begin{equation}\label{eq67ba08ac}
		u(x) = \int_{\R^n} \Phi(y) f(x-y) \did y = \int_{\R^n} \Phi(x-y) f(y) \did x  .
	\end{equation}
	Then $u\in C^2(\R^n)$ and 
	\begin{equation}\label{eq67ba14c8}
		-\laplacian u = f \text{ in }\R^n . 
	\end{equation}
\end{theorem}
\begin{proof}
	Recall that, since $\Phi\in L^1_\loc(\R^n)$ by 
	Proposition~\ref{prop67b9d700}.\ref{prop67b9d700_integrals},
	and since $f$ is continuous with bounded support, then, for every $x\in\R^n$, the integrand $y\mapsto \Phi(y)f(x-y)$ is integrable.
	It follows that $u$ is well defined and that the identity between the second and the third expressions in~\eqref{eq67ba08ac} is a simple change of variables.
	
	We need now to prove the regularity of $u$.
	For $x,y\in\R^n$, define
	\begin{equation}
		K_f(x,y) = \Phi(y) f(x-y).
	\end{equation}
	Fix $R>0$ and set $X=B(0,R)$. 
	Let $S>0$ be such that $\spt(f)\subset B(0,S)$: then, for every $x\in X$,
	\begin{equation}
		|K_f(x,y)| \le \|f\|_{L^\infty} \cdot \Phi(y) \cdot \one_{B(0,R+S)}(y) .
	\end{equation}
	Notice that the function $g_f:y\mapsto \|f\|_{L^\infty} \cdot \Phi(y) \cdot \one_{B(0,R+S)}(y)$ is integrable over $\R^n$.
	By Theorem~\ref{thm679f29de}, the function $u$ is continuous on $X$.
	Moreover, since 
	\[
	\Diff^\alpha_x K_f = K_{\Diff^\alpha f} ,
	\]
	then, again by Theorem~\ref{thm679f29de}, $u\in C^2(X)$.
	Since this holds for every $R>0$, we conclude that $u\in C^2(\R^n)$.
	
	We have also obtained that, always from Theorem~\ref{thm679f29de}, 
	for ever $\alpha\in\N^n$ with $|\alpha|\le 2$,
	\[
	\Diff^\alpha u(x) = \int_{\R^n} \Phi(y) (\Diff^\alpha f)(x-y) \did y .
	\]
	We conclude that
	\[
	\laplacian u(x) = \int_{\R^n} \Phi(y) (\laplacian f)(x-y) \did y .
	\]
	We need to compute the latter integral.
	
	For $\epsilon>0$, we set
	\[
	\laplacian u(x) = \underbrace{ \int_{B(0,\epsilon)} \Phi(y) (\laplacian f)(x-y) \did y }_{I_\epsilon}
		+ \underbrace{ \int_{\R^n\setminus B(0,\epsilon)} \Phi(y) (\laplacian f)(x-y) \did y }_{J_\epsilon} .
	\]
	From~\eqref{eq67ba10db} we obtain that, if $\epsilon<1/2$, there is $C>0$ (depending on $n$) such that 
	\[
	|I_\epsilon| 
	\le \|\laplacian f\|_{L^\infty} \int_{B(0,\epsilon)} |\Phi(y)| \did y 
	\le \begin{cases}
	C \|\laplacian f\|_{L^\infty} \epsilon^2 |\log(\epsilon)| & \text{ if }n=2 , \\
	C \|\laplacian f\|_{L^\infty} \epsilon^2 & \text{ if }n\ge3 .
	\end{cases}
	\]
	In both cases, we have 
	\begin{equation}\label{eq67ba10b2}
		\lim_{\epsilon\to0}I_\epsilon = 0 .
	\end{equation}
	
	For $J_\epsilon$, we compute
	\begin{align*}
		J_\epsilon 
		&= \int_{\R^n\setminus B(0,\epsilon)} \Phi(y) (\laplacian f)(x-y) \did y \\
		&= \int_{\R^n\setminus B(0,\epsilon)} \Phi(y) \laplacian_y (f(x-y)) \did y \\
		&\overset{\eqref{eq67a7a587}}= \int_{\R^n\setminus B(0,\epsilon)} \div_y (\Phi(y)  \grad_yf(x-y)) \did y 
			- \int_{\R^n\setminus B(0,\epsilon)} \grad_y\Phi(y) \cdot \grad_y f(x-y) \did y \\
		&\overset{\eqref{eq67a7a587}}= \int_{\R^n\setminus B(0,\epsilon)} \div_y (\Phi(y)  \grad_yf(x-y)) \did y
			+ \int_{\R^n\setminus B(0,\epsilon)} \div_y(\grad_y\Phi(y)) \cdot f(x-y) \did y  \\
		&\qquad
			- \int_{\R^n\setminus B(0,\epsilon)} \div_y(f(x-y) \grad_y\Phi(y) ) \did y \\
		&\overset{\ref{prop67b9d700}.\ref{prop67b9d700_laplacian}}= \int_{\R^n\setminus B(0,\epsilon)} \div_y (\Phi(y)  \grad_yf(x-y)) \did y
			- \int_{\R^n\setminus B(0,\epsilon)} \div_y(f(x-y) \grad_y\Phi(y) ) \did y \\
		&\overset{\eqref{eq67a7a3f9}}= - \underbrace{ \int_{\de B(0,\epsilon)}  \Phi(y)  \grad_yf(x-y) \cdot \frac{y}{|y|} \did S^{n-1}(y) }_{L_\epsilon}
			+ \underbrace{ \int_{\de B(0,\epsilon)} f(x-y) \grad_y\Phi(y)  \cdot \frac{y}{|y|} \did S^{n-1}(y) }_{K_\epsilon} .
	\end{align*}
	
	We then have that, for $\epsilon<1/2$, there is $C>0$ (depending on $n$) such that 
	\begin{align*}
	|L_\epsilon| 
	&\le \|\grad f\|_{L^\infty} \int_{\de B(0,\epsilon)}  \Phi(y)  \did S^{n-1}(y) \\
	&\le \begin{cases}
	C \|\grad f\|_{L^\infty} |\log(\epsilon)|\epsilon & \text{ if }n=2 , \\
	C \|\grad f\|_{L^\infty} \epsilon & \text{ if }n\ge3 .
	\end{cases}
	\end{align*}
	In both cases, we have 
	\begin{equation}\label{eq67ba13a8}
		\lim_{\epsilon\to0} L_\epsilon = 0 .
	\end{equation}
	
	The remaining quantity is
	\begin{align*}
	K_\epsilon 
	&= \int_{\de B(0,\epsilon)} f(x-y) \grad_y\Phi(y)  \cdot \frac{y}{|y|} \did S^{n-1}(y) \\
	&= - \frac1{n\omega_n } \int_{\de B(0,\epsilon)} f(x-y) \frac{y}{|y|^n}  \cdot \frac{y}{|y|} \did S^{n-1}(y) \\
	&= - \frac1{n\omega_n } \int_{\de B(0,\epsilon)} f(x-y) |y|^{2-n-1} \did S^{n-1}(y) \\
	&= - \frac1{n\omega_n \epsilon^{n-1}} \int_{\de B(0,\epsilon)} f(x-y) \did S^{n-1}(y) \\
	&= - \fint_{\de B(0,\epsilon)} f(x-y) \did S(y) .
	\end{align*}
	We conclude that 
	\begin{equation}\label{eq67c5daf9}
		\lim_{\epsilon\to0} K_\epsilon = - f(x) .
	\end{equation}
	
	We conclude that~\eqref{eq67ba14c8} holds.
\end{proof}

From the proof of the above Theorem~\ref{teo67ba14f2}, we obtain the following corollary
\begin{corollary}[Fundamental property of the fundamental solution]\label{cor67c6cc59}
	Let $U\subset\R^n$ open and $f\in C^0(U)$.
	For every $x\in U$,
	\begin{equation}\label{eq67c6cca3}
%		\lim_{\epsilon\to0} \int_{\de B(0,\epsilon)} f(x-y) \grad_y\Phi(y)  \cdot \frac{y}{|y|} \did S(y) 
		\lim_{\epsilon\to0} \int_{\de B(x,\epsilon)} f(y) \grad_y\Phi(x-y)  \cdot \frac{y-x}{|y-x|} \did S(y)
		= -f(x) ;
	\end{equation}
	or, equivalentely,
	\begin{equation}
		\lim_{\epsilon\to0} \int_{\de B(0,\epsilon)} f(x-y) \grad_y\Phi(y)  \cdot \frac{y}{|y|} \did S(y) 
%		= \lim_{\epsilon\to0} \int_{\de B(x,\epsilon)} f(y) \grad_y\Phi(x-y)  \cdot \frac{y-x}{|y-x|} \did S(y)
		= -f(x) .
	\end{equation}
\end{corollary}
\begin{proof}
	See~\eqref{eq67c5daf9}.
\end{proof}

\begin{remark}
	\begin{enumerate}
	\item We are inverting the Laplacian, in some sense.
	\item If $\phi\in C^2(\R^n)$ is harmonic (e.g., a harmonic polynomial), then $\laplacian (u+\phi) = \laplacian u = f$. So, the above integral formula~\eqref{eq67ba08ac} selects one solution among many.
	\item Our integral formula is only solving $\laplacian u = f$ on the whole space $\R^n$ and only for $f$ with compact support.
	There are various direction in which to extend this result: 
	to $f$ continuous (in fact, $f$ of class $C^2$ looks weird: if $u\in C^2$, then we expect $f=\laplacian u\in C^0$) [see~\ref{subs6889edd6}];
	to $f$ not with compact support (in fact, the selected $u$ does not have compact support anyway);
	to the case of $f$ defined only on an open subset $U$ [see~\ref{subs67bb7d34}].
	I hope we can see all these extensions in this course: we will check!
	\end{enumerate}
\end{remark}

\begin{exercise}
	Using the ideas in the proof of Theorem~\ref{teo67ba14f2}, prove the following statement:
	{\it 
	If $u\in C^2(\R^n)$ is such that $\laplacian u\in C^0_c(\R^n)$, then, for every $x\in\R^n$,
	\begin{equation}
		u(x) = - \int_{\R^n} \Phi(x-y) \laplacian u(y) \did y .
	\end{equation}
	}
	Can you weaken the condition ``$\laplacian u\in C^0_c(\R^n)$''?
	
	{\it Hint.} Look also at Theorem~\ref{thm67bb7d61}.
\end{exercise}

%%%%%%%%%%%%%%%%%%%%%%%%%%%%%%%%%%%%%%%%%%%%%%%%%%%%%%%%%%%%%%%
\subsection{Mean value property}\label{subs67bac7f9}
\index{mean value property for harmonic functions}
\index{harmonic!mean value property for -- functions}

See Section~\ref{sec67bac9ff} first.

\begin{theorem}[Mean value property for harmonic functions]\label{teo67bac805}
	Let $U\subset\R^n$ be open and $u\in C^2(U)$.
	The following statements are equivalent:
	\begin{enumerate}[label=(\roman*)]
	\item\label{teo67bac805_i}
		$u$ is harmonic, i.e., $\laplacian u =0$;
	\item\label{teo67bac805_ii}
		$\forall x\in U$, $\forall r>0$, such that $\bar B(x,r)\subset U$,
		\begin{equation}\label{eq67bac950}
			u(x) = \fint_{\de B(x,r)} u(y) \did S(y) .
		\end{equation}
	\item\label{teo67bac805_iii}
		$\forall x\in U$, $\forall r>0$, such that $\bar B(x,r)\subset U$,
		\begin{equation}\label{eq67bac961}
			u(x) = \fint_{B(x,r)} u(y) \did y .
		\end{equation}
	\end{enumerate}
\end{theorem}
\begin{proof}
	As in~\eqref{eq67bacbb5}, we set 
	\[
	\phi_u(x;r) := \fint_{B(x,r)} u(z) \did z ,
	\quad\text{and}\quad
	\psi_u(x;r) := \fint_{\de B(x,r)} u(z) \did S(z) .
	\]
	Then, from~\eqref{eq67bacb9a} we have,
	\begin{align}
		\label{eq67bace9e}
		\de_r\phi_u(x;r) &= \frac{n}{r} (\psi_u(x;r)- \phi_u(x;r) ) , \\
		\label{eq67bacea3}
		\de_r\psi_u(x;r) &= \frac{r}{n} \phi_{\laplacian_xu}(x;r) .
	\end{align}
	
	$\ref{teo67bac805_i}\THEN\ref{teo67bac805_ii}$:
	By~\eqref{eq67bacea3}, if $\laplacian u=0$, then $\de_r\psi_u=0$, i.e., $r\mapsto\psi_u(x;r)$ is constant and, since $B(x,r)\subset U$, we have
	\[
	u(x) 
	\overset{\eqref{eq67bad1c9}}= \lim_{\epsilon\to0} \psi_u(x;\epsilon) 
	= \psi_u(x;r) .
	\]
	
	$\ref{teo67bac805_ii}\THEN\ref{teo67bac805_iii}$:
	Let $x\in U$ and $r>0$ be such that $\bar B(x,r)\subset U$.
	Then
	\begin{align*}
		\fint_{B(x,r)} u(y) \did y
		&= \frac1{r^n\omega_n } \int_0^r \int_{\de B(x,s)} u(y) \did S(y)\did r \\
		&= \frac{n}{r^n} \int_0^r \fint_{\de B(x,s)} u(y) \did S(y) s^{n-1}\did r \\
		&\overset{\ref{teo67bac805_ii}}= u(x) \frac{n}{r^n} \int_0^r s^{n-1}\did r 
		= u(x) .
	\end{align*}
	
	$\ref{teo67bac805_iii}\THEN\ref{teo67bac805_ii}$:
	Let $x\in U$ and $r>0$ be such that $\bar B(x,r)\subset U$.
	From~\eqref{eq67bace9e}, we obtain,
	\[
	0 = \de_r\phi_u(x;r) 
	= \frac{n}{r} (\psi_u(x;r)- u(x) ) ,
	\]
	i.e., $\psi_u(x;r) = u(x)$.
	
	$\ref{teo67bac805_ii}\THEN\ref{teo67bac805_i}$:
	Let $x\in U$. 
	From~\eqref{eq67bacea3}, we have
	\[
	\laplacian u(x) 
	\overset{\eqref{eq67bad1c9}}= \lim_{r\to0} \phi_{\laplacian u}(x;r) 
	\overset{\eqref{eq67bacea3}}= \lim_{r\to0} \frac{n}{r}\de_r\psi_u(x;r)
	\overset{\ref{teo67bac805_ii}}= 0 .
	\]
\end{proof}

\begin{remark}
	If $\de B(x,r)\subset U$ but $B(x,r)\not\subset U$, then the mean value formula does not need to hold.
	Find an example.
%	To see this, take the fundamental solution $\Phi$ and any ball centered at $0$ (exercise).
\end{remark} 

\begin{exercise}
	For all $n\ge2$, compute
	\[
	\psi_\Phi(0;r) := \fint_{\de B(0,r)} \Phi(y) \did S(y) ,
	\]
	where $\Phi$ is the fundamental solution of the Laplace equation.
	Is $r\mapsto\psi_\Phi(0;r)$ constant? 
\end{exercise}

\begin{exercise}
	In the proof of Theorem~\ref{teo67bac805}, we have shown $\ref{teo67bac805_ii}\THEN\ref{teo67bac805_i}$.
	Give a direct proof of $\ref{teo67bac805_iii}\THEN\ref{teo67bac805_i}$.
\end{exercise}

\begin{exercise}\label{ex67badfb7}
	Show the equivalence $\ref{teo67bac805_iii}\IFF\ref{teo67bac805_ii}$
	assuming only $u\in C^0(U)$.
\end{exercise}


%%%%%%%%%%%%%%%%%%%%%%%%%%%%%%%%%%%%%%%%%%%%%%%%%%%%%%%%%%%%%%%
\subsection{Strong maximum principle}\label{subs67bad3fc}
\index{strong maximum principle}
\index{maximum!strong -- principle}
\index{principle!strong maximum --}

\index{strong minimum principle}
\index{minimum!strong -- principle}
\index{principle!strong minimum --}

\begin{theorem}[Strong maximum principle - First version]\label{teo67bad430}
	Suppose $U\subset\R^n$ is open and connected,
	and $u\in C^2(U;\R)$ is harmonic.
	If there exists $x\in U$ is such that $u(x) = \sup_Uu$, then $u$ is constant.
\end{theorem}
\begin{proof}
	Define
	\[
	M = \sup_Uu
	\quad\text{and}\quad
	W = \{x\in U:u(x)=M\} = u^{-1}(\{M\}) .
	\]
	Since $u$ is continuous, then $W$ is closed in $U$.
	
	We claim that $W$ is open.
	If $x\in W$, then there is $r>0$ such that $B(x,r)\subset U$, because $U$ is open.
	Therefore, since $u$ is harmonic, by Theorem~\ref{teo67bac805} we have
	\[
	M = u(x) 
	\overset{\ref{teo67bac805}.\ref{teo67bac805_iii}}= \fint_{B(x,r)} u(y) \did y
	\overset{(*)}\le M
	\]
	where the inequality $(*)$ is strict\footnotemark{} unless $u(y)=M$ for all $y\in B(x,r)$.%
	\footnotetext{
	Indeed, if $|B(x,r)\setminus W|>0$, then
	\begin{align*}
	\fint_{B(x,r)} u(y) \did y 
	&= \frac1{|B(x,r)|} \left( \int_{B(x,r)\cap W} u(y)\did y
		+  \int_{B(x,r)\setminus W} u(y)\did y \right) \\
	&< \frac1{|B(x,r)|} \left( M |B(x,r)\cap W| 
		+  M |B(x,r)\setminus W|  \right)
	= M .
	\end{align*}
	}
	It follows that $B(x,r)\subset W$.
	
	In conclusion, since $U$ is connected and $W\subset U$ is both open and closed, 
	then $W\in\{\emptyset,U\}$.
	If $W$ is not empty, then $W=U$, i.e., $u$ is constant.
\end{proof}

\begin{exercise}
	Show that, if $u\in C^0(B(0,r);\R)$ is such that $u\le M$ on $B(0,r)$, then
	$\fint_{B(x,r)} u(y) \did y\le M$, with equality if and only if $u=M$ on $B(0,r)$.
\end{exercise}

\begin{theorem}[Strong maximum principle - Second version]\label{teo67bad761}
	Suppose $U\subset\R^n$ is open and bounded,
	and $u\in C^2(U;\R)\cap C^0(\bar U;\R)$ is harmonic.
	Then
	\begin{equation}\label{eq67bad7a9}
		\max_{\bar U} u = \max_{\de U} u .
	\end{equation}
\end{theorem}
\begin{proof}
	We clearly have $\max_{\bar U} u \ge \max_{\de U} u$.
	Suppose that there exists $x\in U$ with $u(x) = \max_{\bar U}u$.
	We need to show that there is also a point on $\de U$ with the same value.
	Let $C\subset U$ be the connected component of $U$ containing $x$.
	Since $C$ is open and connected, by Theorem~\ref{teo67bad430}, $u$ is constant on $C$.
	Therefore, if $y\in\de C$, then $u(y) = u(x) = \max_{\bar U}u$.
	Since $\R^n$ is locally connected, then $\de C\subset \de U$ (see Exercise~\ref{ex67baddda}).
	Therefore, 
	\[
	\max_{\bar U}u = u(x) = \max_{\de C}u \le \max_{\de U} u .
	\]
\end{proof}

\begin{exercise}\label{ex67baddda}
	Show the following statement:
	{\it Let $X$ be a locally connected topological space and $U\subset X$ open.
	Let $C\subset U$ be a connected component of $U$.
	Then $\de C\subset\de U$.}
	
	After this, show that the hypothesis of $X$ being locally connected is necessary.
	In other words, find an example of a topological space $X$ that is not locally connected, and $U\subset X$ open that has a connected component $C\subset U$ with $\de C\cap\de U=\emptyset$.
\end{exercise}
\begin{proof}[Solution to Exercise~\ref{ex67baddda}]
	Let $x\in\de C$.
	Then $x\in\bar U$.
	If $x\in U$, then there exists a connected neighborhood $V$ of $x$ with $V\subset U$.
	Since $x\in\de C$, then $V\cap C\neq\emptyset$.
	We then have $C\cup V\subset U$ is connected.
	Since $C$ is a maximal connected subset of $U$, we conclude $V\subset C$, in contradiction with $x\in\de C$.
\end{proof}

\begin{exercise}
	State and prove the \emph{Strong Minimum Principle} for harmonic functions, both in the first and second versions.
\end{exercise}

\begin{exercise}\angryface
	Show the strong maximum and minimum principles for harmonic functions without using the mean value property, i.e., explicitly using that $\laplacian u=0$.
	
	{\it Hint}:
	First consider $u\in C^2(U;\R)$ with $\laplacian u>0$ in $U$.
	Suppose $x\in U$ is such that $u(x) = \max_Uu$.
	Then $t=0$ is a point of maximum for $t\mapsto u(x+te_j)$, 
	for each $j\in\{1,\dots,n\}$.
	Therefore $\frac{\de^2u}{\de x_j^2}(x) = \left.\frac{\diff^2}{\diff t^2}\right|_{t=0}u(x+te_j) \le 0$.
\end{exercise}

\begin{remark}
	In Theorem~\ref{teo67bad761}, we actually prove more than what we claim.
	In fact, do not need ``$u\in C^2(U;\R)\cap C^0(\bar U;\R)$ harmonic'', 
	but only $u\in C^0(\bar U;\R)$ satisfying the mean value property, as stated in
	Theorem~\ref{teo67bac805}.\ref{teo67bac805_ii}.	
	See also Exercise~\ref{ex67badfb7}.
	Maybe, in your personal notes you can rewrite these statements in the more general form.
\end{remark}

%%%%%%%%%%%%%%%%%%%%%%%%%%%%%%%%%%%%%%%%%%%%%%%%%%%%%%%%%%%%%%%
\subsection{Subharmonic functions}\label{subs67bae229}
\index{subharmonic function}
\index{function!subharmonic --}
\index{harmonic!sub-- function}

\index{superharmonic function}
\index{function!superharmonic --}
\index{harmonic!super-- function}

\begin{exercise}
	Let $U\subset\R^n$ open.
	A function $u\in C^2(U)$ is \emph{subharmonic} if and only if $\laplacian u\ge 0$ on $U$.
	A function $u\in C^2(U)$ is \emph{superharmonic} if and only if $\laplacian u\le 0$ on $U$.
	(Notice the inversion between ``sub'' and ``$\ge$'').

	State and prove modified versions of Theorem~\ref{teo67bac805} 
	and Theorem~\ref{teo67bad430}
	for subharmonic functions.
	
	After this, do the same for superharmonic functions.
\end{exercise}

%%%%%%%%%%%%%%%%%%%%%%%%%%%%%%%%%%%%%%%%%%%%%%%%%%%%%%%%%%%%%%%
\subsection{Monotonicity of Laplace's boundary value problem}\label{subs67baea83}

\begin{corollary}\label{cor67bae9fc}
	Let $U\subset\R^n$ open, bounded and connected.
	Let $u\in C^2(U;\R)\cap C^0(\bar U;\R)$ and $g\in C^0(\de U;\R)$ be such that
	\[
	\begin{cases}
		\laplacian u = 0 &\text{ in }U ,\\
		u = g &\text{ on }U .
	\end{cases}
	\]
	If $g$ is not constant and $g\ge 0$ on $\de U$, then $u>0$ on $U$.
\end{corollary}

\begin{exercise}
	Write a proof of Corollary~\ref{cor67bae9fc}.
\end{exercise}

\begin{exercise}
	Write and prove a version of Corollar~\ref{cor67bae9fc} for sub- and superharmonic functions.
\end{exercise}

%%%%%%%%%%%%%%%%%%%%%%%%%%%%%%%%%%%%%%%%%%%%%%%%%%%%%%%%%%%%%%%
\subsection{Uniqueness for the Poisson equation}\label{subs67baea7d}
\index{Poisson equation}
\index{equation!Poisson --}

\begin{theorem}[Uniqueness for the Poisson equation]\label{teo67baea93}
	Let $U\subset\R^n$ be open and bounded, $f\in C^0(U)$ and $g\in C^0(\de U)$.
	
	There exists at most one solution in $C^2(U)\cap C^0(\bar U)$ to the boundary value problem
	\begin{equation}\label{eq67baeb15}
	\begin{cases}
		- \laplacian u = f &\text{ in }U ,\\
		u = g &\text{ on }U .
	\end{cases}
	\end{equation}
\end{theorem}
\begin{proof}
	Suppose $u_1,u_2\in C^2(U)\cap C^0(\bar U)$ are solutions to~\eqref{eq67baeb15}.
	Then if $u$ is the real or imaginary part of $u_1-u_2$, then $u\in C^2(U;\R)\cap C^0(\bar U;\R)$ and $u$ solves
	\begin{equation}\label{eq67baeb46}
	\begin{cases}
		\laplacian u = 0 &\text{ in }U ,\\
		u = 0 &\text{ on }U .
	\end{cases}
	\end{equation}
	Apply the maximum principle Theorem~\ref{teo67bad761} and the minimum principle (or the maximum principle to $-u$), to conclude that $u=0$ on $U$.
	Therefore, $u_1=u_2$.
\end{proof}

%%%%%%%%%%%%%%%%%%%%%%%%%%%%%%%%%%%%%%%%%%%%%%%%%%%%%%%%%%%%%%%
\subsection{Smoothness of harmonic functions}\label{subs67baec70}

\begin{theorem}\label{thm67baed2c}
	Let $U\subset \R^n$ be open and $u\in C^0(U)$.
	Suppose that $u$ satisfies the spherical mean-value property, i.e.,
	\begin{equation}\label{eq67baed21}
		\forall x\in U\ \forall r>0\ \text{ s.t. } \bar B(x,r)\subset U :
		\qquad u(x) = \fint_{\de B(x,r)} u(y) \did S(y) .
	\end{equation}
	Then $u\in C^\infty(U)$.
\end{theorem}
\begin{proof}
	Let $\{\eta_\epsilon\}_{\epsilon>0}$ be the family of standard mollifiers.
	Recall that $\eta_\epsilon(x) = \tilde\eta(|x|/\epsilon)/\epsilon^n$ for a compactly supported function $\tilde\eta\in C^\infty(\R)$ with $\tilde\eta\ge0$, $\spt(\tilde\eta)\subset[-1,1]$, and $\int_\R\tilde\eta(t)\did t = 1$.
	In particular, it follows that $\spt(\eta_\epsilon)\subset \bar B(0,\epsilon)$ 
	and $\int_{\R^n}\eta_\epsilon(x)\did x = 1$.

	We will show that the spherical mean-value property~\eqref{eq67baed21} implies that, for every $\epsilon>0$, $u*\eta_\epsilon = u$ on 
	\[
	U_\epsilon = \{x\in U : \dist(x,\de U) > \epsilon\} .
	\]
	Since $u*\eta_\epsilon \in C^\infty(U_\epsilon)$ by Proposition~\ref{prop688c72e1}, we will conclude $u\in C^\infty(U)$.
	
	So, if $x\in U_\epsilon$, then
	\begin{align*}
		u*\eta_\epsilon(x)
		&= \int_{B(x,\epsilon)} u(y) \eta_\epsilon(x-y) \did y \\
		&= \int_0^\epsilon \int_{\de B(x,r)} u(y)\eta_\epsilon(x-y) \did S(y) \did r\\
		&= \int_0^\epsilon \tilde\eta_\epsilon(r) \left(\int_{\de B(x,r)}  \did S(y)\right) \fint_{\de B(x,r)} u(y) \did S(y) \did r\\
		&\overset{\eqref{eq67baed21}}= u(x) \int_0^\epsilon \tilde\eta_\epsilon(r) \int_{\de B(x,r)} \did S(y) \did r\\
		&= u(x) \int_0^\infty  \int_{\de B(x,r)} \eta_\epsilon(x-y) \did S(y) \did r\\
		&= u(x) \int_{\R^n} \eta_\epsilon(x-y) \did y 
		= u(x) .
	\end{align*}
\end{proof}

\begin{corollary}
	Harmonic functions are $C^\infty$ smooth.
\end{corollary}

%%%%%%%%%%%%%%%%%%%%%%%%%%%%%%%%%%%%%%%%%%%%%%%%%%%%%%%%%%%%%%%
\subsection{Analyticity of harmonic functions}\label{subs67baeff5}

\begin{proposition}[Estimates on derivatives of harmonic functions]\label{prop67bb09b8}
	Let $U\subset\R^n$ open and $u\in C^\infty(U)$ harmonic.
	For every $ x\in U$ and $r>0$ with $B( x,r)\subset U$,
	and for every $\alpha\in\N^n$ with $|\alpha|=k$ we have
	\begin{equation}\label{eq67baf145}
		 |\Diff^\alpha u( x)| \le \frac{C_k}{r^{n+k}} \|u\|_{L^1(B(x,r))} ,
		 \qquad\text{where }
		 \begin{array}{l}
		 C_0 = \frac1{\omega_n } , \text{ and }\\
		 C_k = \frac{(2^{n+1}n k)^k}{\omega_n } \text{for }k\ge1 .
		 \end{array}
	\end{equation}
%	where
%	\begin{equation}
%		C_0 = \frac1{\omega_n }
%		\quad\text{and}\quad
%		C_k = \frac{(2^{n+1}n)^k}{\omega_n } 
%		\quad\text{for }k\ge1.
%	\end{equation}
\end{proposition}

The proof of Proposition~\ref{prop67bb09b8} will come after a few lemmas.

%%%%%%%%%%%%%%%%%%%%%%%%%%%%%%%%%%%%%%%%%%%%%%%%%%%%%%%%%%%%%%%
\begin{lemma}[Derivative of harmonic functions]\label{lem67bb0a0c}
	Every partial derivative of a harmonic function is harmonic.
\end{lemma}
\begin{proof}
	Let $U\subset\R^n$ open and $u\in C^\infty(U)$ harmonic.
	If $j\in\{1,\dots,n\}$, then 
	\[
	\laplacian(\de_ju) = \sum_{k=1}^n \de_k^2\de_ju = \de_j\sum_{k=1}^n \de_k^2u 
	= \de_j\laplacian u = 0 ,
	\]
	because $\de_j\de_k=\de_k\de_j$ for all $j,k$.
	So, all first-order derivatives of $u$ are harmonic.
	Iterating, all derivatives of harmonic functions are harmonic harmonic functions.
\end{proof}

\begin{lemma}[Case $k=0$]\label{lem67bb11e4}
	Let $U\subset\R^n$ open and $u\in C^\infty(U)$ harmonic.
	For every $ x\in U$ and $r>0$ with $B( x,r)\subset U$,
	\begin{equation}\label{eq67bb11da}
		 |u(x)| \le \frac1{\omega_n r^n}  \|u\|_{L^1(B(x,r))} .
	\end{equation}
\end{lemma}
\begin{proof}
	The ball mean-value property for harmonic functions gives
	\[
	|u(x)| 
	= \left| \fint_{B(x,r)} u(y) \did y \right|
	\le \frac1{\omega_n r^n}  \|u\|_{L^1(B(x,r))} .
	\]
%	This shows~\eqref{eq67baf145} for $k=0$.
\end{proof}

%%%%%%%%%%%%%%%%%%%%%%%%%%%%%%%%%%%%%%%%%%%%%%%%%%%%%%%%%%%%%%%
\begin{lemma}[Case $k=1$: Estimate of first derivatives of harmonic functions]\label{lem67bb0fed}
	Let $U\subset\R^n$ open and $u\in C^\infty(U)$ harmonic.
	For every $ x\in U$ and $r>0$ with $B( x,r)\subset U$,
	and for every $j\in\{1,\dots,n\}$ and $\theta\in(0,1)$, we have
	\begin{equation}\label{eq67bb1919}
		|\de_ju(x)| \le \frac{1}{\omega_n  (\theta r)^n} \int_{\de B(x,\theta r)} |u(y)| \did S(y) .
	\end{equation}
	As a consequence, we also have
	\begin{equation}\label{eq67bb100e}
		 |\de_j u(x)| \le \frac{2^{n+1}n}{\omega_n } \frac{1}{r^{n+1}} \|u\|_{L^1(B(x,r))} .
	\end{equation} 
\end{lemma}
\begin{proof}
	Fix $j\in\{1,\dots,n\}$.
	By Lemma~\ref{lem67bb0a0c}, $\de_ju$ is also harmonic, and thus it has the mean-value property.
	Therefore, for every $\theta\in(0,1)$,
	\begin{align*}
		|\de_ju(x)| 
		&= \left| \fint_{B(x,\theta r)} \de_ju(y) \did y \right| \\
		&= \frac{1}{\omega_n  (\theta r)^n} \left| \int_{B(x,\theta r)} \div(ue_j)(y) \did y \right| \\
		&= \frac{1}{\omega_n  (\theta r)^n} \left| \int_{\de B(x,\theta r)} u(y) e_j\cdot \frac{y-x}{|y-x|} \did S(y) \right| \\
		&\le \frac{1}{\omega_n  (\theta r)^n} \int_{\de B(x,\theta r)} |u(y)| \did S(y) . 
%	\end{align*}
	\intertext{ This shows~\eqref{eq67bb1919}.
%	Next, take $\theta =1/2$.
	Next, we use the fact that, if $y\in\de B(x,\theta r)$, then $B(y,(1-\theta)r)\subset B(x,r)$.
	Thus, applying Lemma~\ref{lem67bb11e4} to the running estimates, we get}
%	\begin{align*}
		|\de_ju(x)| 
		&\le \frac{1}{\omega_n  (\theta r)^n} \int_{\de B(x,\theta r)} |u(y)| \did S(y) \\
		&\overset{\eqref{eq67bb11da}}\le \frac{1}{\omega_n  (\theta r)^n} \int_{\de B(x,\theta r)} \frac1{\omega_n ((1-\theta)r)^n}  \|u\|_{L^1(B(y,(1-\theta)r))} \did S(y) \\
		&\le \frac{1}{\omega_n  (\theta r)^n} \frac{1}{\omega_n  ((1-\theta) r)^n} n\omega_n  (\theta r)^{n-1} \|u\|_{L^1(B(x,r))} \\
		&= \frac{n}{\omega_n \theta (1-\theta)^n r^{n+1}}\|u\|_{L^1(B(x,r))}
%
%		|\de_ju(x)| 
%		&\le \frac{1}{\omega_n  (r/2)^n} \int_{\de B(x,r/2)} |u(y)| \did S(y) \\
%		&\overset{\eqref{eq67bb11da}}\le \frac{1}{\omega_n  (r/2)^n} \int_{\de B(x,r/2)} \frac1{\omega_n (r/2)^n}  \|u\|_{L^1(B(y,r/2))} \did S(y) \\
%		&\le \frac{2^n}{\omega_n  r^n} \frac{1}{\omega_n  (r/2)^n} n\omega_n  (r/2)^{n-1} \|u\|_{L^1(B(x,r))} \\
%		&= \frac{n2^{n+1}}{\omega_n  r^{n+1}}\|u\|_{L^1(B(x,r))}
	\end{align*}
	All in all, if we take $\theta =1/2$, we have~\eqref{eq67bb100e}.
\end{proof}

\begin{proof}[Proof of Proposition~\ref{prop67bb09b8}]
	We shall argue by induction over $k$.
	
	For $k=0$ and $k=1$, we have already Lemma~\ref{lem67bb11e4} and Lemma~\ref{lem67bb0fed}, respectively.
		
	Let $m\ge1$ and assume that \eqref{eq67baf145} holds for all $k\le m$.
	We now prove \eqref{eq67baf145} for $k=m+1$.
	Let $\alpha\in\N^n$ with $|\alpha|=m+1$.
	Then there is $\beta\in\N^n$ with $|\beta|=m$ and $j\in\{1,\dots,n\}$ such that $\Diff^\alpha = \Diff^\beta\frac{\de}{\de x_j}$.
	
	The function $\Diff^\beta u$ is harmonic by Lemma~\ref{lem67bb0a0c}.
	We apply Lemma~\ref{lem67bb0fed} to the function $\Diff^\beta u$ and get that,
	for $\theta\in(0,1)$,
	\begin{align*}
		|\Diff^\alpha u(x)| 
		&= |\de_j \Diff^\beta  u(x)| \\
		&\overset{\eqref{eq67bb1919}}\le \frac{1}{\omega_n  (\theta r)^n} \int_{\de B(x,\theta r)} |\Diff^\beta u(y)| \did S(y) \\
%		&\le \frac{1}{\omega_n  (\theta r)^n} \frac{(2^{n+1}n m)^m}{\omega_n } \frac{1}{((1-\theta)r)^{n+m}} \|u\|_{L^1(B(x,r))} \int_{\de B(x,\theta r)} \did S(y) \\
		\text{[inductive hypothesis]}
		&\le \frac{1}{\omega_n  (\theta r)^n}  \frac{C_m}{((1-\theta)r)^{n+m}} \|u\|_{L^1(B(x,r))} \int_{\de B(x,\theta r)} \did S(y) \\
%		&= \frac{1}{\omega_n  (r/m)^n} \frac{C_m}{((m-1)r/m)^{n+m}} \|u\|_{L^1(B(x,r))} n\omega_n  (r/m)^{n-1} \\
		&= C_m n \frac{1}{(1-\theta)^{n+m}\theta} \frac1{r^{n+m+1}} \|u\|_{L^1(B(x,r))} .
	\end{align*}
	
	Taking $\theta = \frac{1}{m+1}$ and $C_m = \frac{(2^nnm)^m}{\omega_n}$, we get
	\begin{align*}
	|\Diff^\alpha u(x)|
	&\le \frac{(2^nnm)^m}{\omega_n} \frac{ n(m+1)^{m+n+1} }{ m^{m+n} } \frac1{r^{n+m+1}} \|u\|_{L^1(B(x,r))} \\
	\text{[computations]}&\le \frac{(2^nn(m+1))^{m+1}}{\omega_n} \frac1{r^{n+m+1}} \|u\|_{L^1(B(x,r))} .
	\end{align*}
\end{proof}

\begin{exercise}
	Let $a,b<0$ and define $\phi:(0,1)\to\R$, $\phi(\theta) = \theta^a\theta^b$.
	Find the minimum and the point of minimum of $\phi$ on $(0,1)$.
	
	{\it Solution}:
	$\theta_m = \frac{a}{a+b}$ and $\phi(\theta_m)=\left(\frac{a}{a+b}\right)^{a+b}$.
\end{exercise}


%%%%%%%%%%%%%%%%%%%%%%%%%%%%%%%%%%%%%%%%%%%%%%%%%%%%%%%%%%%%%%%
\index{analyticity of harmonic functions}
\index{harmonic!analyticity of -- functions}
\index{theorem!analyticity of harmonic functions}

\begin{theorem}[Analyticity]\label{thm67bb666b}
	Harmonic functions are analytic.
\end{theorem}
\begin{proof}
	NOTA BENE: This proof contains a mistake. 
	Find it and correct it. The correct proof is in Section~\ref{subs688c7a1a}, page~\pageref{subs688c7a1a}.

	Let $U\subset\R^n$ be open and $u\in C^\infty(U)$ a harmonic function.
	Fix $\hat x\in U$ and set $\hat r = \frac14 \dist(\hat x,\de U)$.
	We claim that there exists $\epsilon\in(0,1)$ such that, if
	\begin{equation}\label{eq67bb678e}
		r < %\frac{\hat r}{2^{n+2} n^3 e} = 
		\epsilon \hat r,
	\end{equation}
	then the Taylor series of $u$ centered at $\hat x$ converges on $B(\hat x,r)$ to $u$,
	that is, for every $x\in B(\hat x,r)$,
	\begin{equation}\label{eq67bb6792}
		u(x) = \lim_{N\to\infty} \sum_{k=0}^N \sum_{|\alpha|=k} \frac{\Diff^\alpha u(\hat x)}{\alpha!} (x-\hat x)^\alpha .
	\end{equation}
	
	To this aim, define the reminder function
	\[
	R_N(x) = u(x) - \sum_{k=0}^{N-1} \sum_{|\alpha|=k} \frac{\Diff^\alpha u(\hat x)}{\alpha!} (x-\hat x)^\alpha .
	\]
	For every $x\in B(\hat x,r)$ there exists $t_x\in[0,1]$ such that
	\[
	R_N(x) = \sum_{|\alpha|=N} \frac{\Diff^\alpha u(\hat x+t(x-\hat x))}{\alpha!} (x-\hat x)^\alpha .
	\]
	Using Proposition~\ref{prop67bb09b8}, we make the following estimate:
	since $\hat x+t(x-\hat x)\in B(\hat x,r) \subset B(\hat x,\hat r)$,
	\begin{align*}
		|R_N(x)|
		&\le \sum_{|\alpha|=N} \frac{|\Diff^\alpha u(\hat x+t(x-\hat x))|}{\alpha!} |x-\hat x|^N \\
		&\overset{\eqref{eq67baf145}}\le  \frac{(2^{n+1}nN)^N \|u\|_{L^1(B(\hat x,\hat r))}}{\omega_n \hat r^{n+N} }  r^N \sum_{|\alpha|=N}\frac1{\alpha!}\\
		&\overset{\eqref{eq67bb6ddf}}= \frac{\|u\|_{L^1(B(\hat x, \hat r))}}{\omega_n \hat r^n} \epsilon^N (2^{n+1}nN)^N \frac{n^N}{N!} \\
		&= \frac{\|u\|_{L^1(B(\hat x, \hat r))} }{ \omega_n \hat r^n } \frac{\sqrt{2\pi N} (N/e)^N }{N!} 
			\frac{(\epsilon 2^{n+1}n^2N)^{N}}{\sqrt{2\pi N} (N/e)^N } \\
		&= \frac{\|u\|_{L^1(B(\hat x, \hat r))} }{ \omega_n \hat r^n } \frac{\sqrt{2\pi N} (N/e)^N }{N!} 
			\frac{(\epsilon 2^{n+1}n^2e)^{N}}{\sqrt{2\pi N} } ,\\
	\end{align*}
	where we have used the Multinomial Theorem
	\begin{equation}\label{eq67bb6ddf}
		n^N = (\sum_{j=1}^n 1)^N = \sum_{|\alpha|=N} \frac{N!}{\alpha!} .
	\end{equation}
	Using Stirling's formula
	\begin{equation}\label{eq67bb6e37}
		\lim_{N\to\infty} \frac{\sqrt{2\pi N} (N/e)^N }{N!}  = 1 ,
	\end{equation}
	we conclude that, if $\epsilon\in(0,1)$ is so that $\epsilon 2^{n+1}n^2e<1$, then $\lim_{N\to\infty}|R_N(x)|=0$.
\end{proof}




%%%%%%%%%%%%%%%%%%%%%%%%%%%%%%%%%%%%%%%%%%%%%%%%%%%%%%%%%%%%%%%
\subsection{Liouville's Theorem}\label{subs67bb6f4c}
\index{Liouville's Theorem}
\index{theorem!Liouville's --}

\begin{theorem}[Liouville's Theorem]\label{thm67bb782e}
	If $u:\R^n\to\C$ is harmonic and bounded, then $u$ is constant.
\end{theorem}
\begin{proof}
	Fix $ x\in\R^n$.
	Then, for all $r>0$, we have from Proposition~\ref{prop67bb09b8}, 
	\begin{align*}
		|Du( x)| 
		&\le \sum_{j=1}^n \left| \frac{\de u}{\de x_j}( x) \right| 
		\le \frac{nC_1}{r^{n+1}} \int_{B( x,r)} |u(y)| \did y \\
		&\le \frac{nC_1}{r^{n+1}} \|u\|_{L^\infty(\R^n)} \omega_n r^n 
		= \frac{nC_1\omega_n \|u\|_{L^\infty(\R^n)} }{r} .
	\end{align*}
	Since this inequality holds for every $r>0$, then $\Diff u( x)=0$.
	Since $ x$ is arbitrary, then $\Diff u\equiv0$ on $\R^n$,
	i.e., $u$ is constant.
\end{proof}

\begin{exercise}
	Prove Liouville's Theorem using only the ball mean-value property.
	
	{\it Solution}:
	Fix $\hat x\in\R^n$ and $x\in B(\hat x,1)$.
	Then 
	\begin{align*}
		|u(x)-u(\hat x)|
		&= \left| \fint_{B(x,r)} u(y) \did y - \fint_{B(\hat x,r)} u(y) \did y \right| \\
		&= \frac1{\omega_nr^n} \left| \int_{B(x,r)\setminus B(\hat x,r)} u(y) \did y - \int_{B(\hat x,r)\setminus B(x,r)} u(y) \did y \right| \\
		&\le \frac1{\omega_nr^n} \int_{B(\hat x,r+1)\setminus B(\hat x,r-1)} |u(y)| \did y \\
		&\le \|u\|_{L^\infty(\R^n)} \frac{ (r+1)^n - (r-1)^n }{ r^n } \\
		&= \|u\|_{L^\infty(\R^n)} \left( (1+1/r)^n - (1-1/r)^n \right) 
		\overset{r\to\infty}\longrightarrow 0 .
	\end{align*}
\end{exercise}

%%%%%%%%%%%%%%%%%%%%%%%%%%%%%%%%%%%%%%%%%%%%%%%%%%%%%%%%%%%%%%%
\subsection{Representation formula}\label{subs67bb761c}

Define
\[
C^2_b(\R^n) = \{u\in C^2(\R^n): \|u\|_{L^\infty}<\infty \} .
\]

\begin{theorem}
	Suppose $n\ge3$ and fix $f\in C^2_c(\R^n)$.
	The solutions $u$ of $-\laplacian u=f$ in $C^2_b(\R^n)$ are exactly all functions in the family
	\begin{equation}
		\left\{ x\mapsto \int_{\R^n} \Phi(x-y) f(y) \did y + c : c\in \C \right\}
%		\qquad\forall x\in\R^n .
	\end{equation}
	
%	If $u\in C^2(\R^n)$ and $f\in C^2_c(\R^n)$ are such that $\laplacian u=f$ on $\R^n$,
%	and if $u$ is bounded, then there exists $C\in\C$ such that
%	\begin{equation}
%		u(x) = \int_{\R^n} \Phi(x-y) f(y) \did y + C
%		\qquad\forall x\in\R^n .
%	\end{equation}
\end{theorem}
\begin{proof}
	We already know that the function $\tilde u(x) = \int_{\R^n} \Phi(x-y) f(y) \did y$ 
	is of class $C^2$ and solves $-\laplacian\tilde u=f$,
	see Theorem~\ref{teo67ba14f2}.
	
	We claim that $\tilde u\in C^2_b(\R^n)$.
	Indeed, if $\spt(f)\subset \R^n\setminus B(x,R)$, then
	\begin{align*}
		|\tilde u(x)| 
		&= \left| \int_{\spt(f)} \Phi(x-y) f(y) \did y \right| \\
		&\le \|f\|_{L^\infty} \Phi(R) |\spt(f)| .
	\end{align*}
	So, if $\spt(f)\subset B(0,L)$, then for every $x$ with $|x|\ge L$, we have
	$\tilde u(x) \le  \|f\|_{L^\infty}|\spt(f)| \Phi(|x|-L)$.
	Therefore, $\lim_{x\to\infty}|\tilde u(x)| = 0$.
	In particular, $\tilde u$ is bounded on $\R^n$.
	
	If $-\laplacian u=f$ and if $u$ is bounded, then $u-\tilde u$ is bounded and $\laplacian(u-\tilde u)=0$.
	By Liouville's Theorem~\ref{thm67bb782e}, $u-\tilde u$ is constant.
\end{proof}


%%%%%%%%%%%%%%%%%%%%%%%%%%%%%%%%%%%%%%%%%%%%%%%%%%%%%%%%%%%%%%%
\subsection{Harnack's Inequality}\label{subs67bb79dd}
\index{Harnack's Inequality}
\index{theorem!Harnack's Inequality}

\begin{theorem}[Harnack's Inequality]\label{thm67bb79ea}
	Let $U\subset\R^n$ be open and $V\Subset U$ an open and connected subset.
	There exists $C>0$ (depending on $U$ and $V$) such that,
	if $u:U\to[0,+\infty)$ is harmonic on $U$ and non-negative, then
	\begin{equation}
		\sup_V u \le C \inf_V u ,
	\end{equation}
	or, equivalentely,
	\begin{equation}
		\forall x,y\in V:
		\qquad
		u(x) \le C u(y) .
	\end{equation}
\end{theorem}%]
\begin{proof}
	Let $r=\frac14\dist(V,\de U)$.
	If $x,y\in V$ are such that $|x-y|<r$, then, by the mean value property Theorem~\ref{teo67bac805},
	\begin{align*}
		u(x) 
		&= \fint_{B(x,2r)} u(z) \did z \\
		&= \frac1{\omega_n (2r)^n} \left( \int_{B(y,r)} u(z)\did z + \int_{B(x,2r)\setminus B(y,r)} u(z) \did z \right) \\
		\text{[since $u\ge0$]}
		&\ge \frac1{2^n} \fint_{B(y,r)} u(z) \did z
		= \frac1{2^n} u(y) .
	\end{align*}
	Since $\bar V$ is compact and connected, there is a finite family of points $x_1,\dots,x_N\in V$ such that $|x_j-x_{j+1}|<r$ and $V\subset\bigcup_{j=1}^N B(x_j,r)$;
	see Exercise~\ref{ex688c8182}.
	Then, whenever $x,y\in V$,
	\[
	u(x) \ge \left( \frac1{2^n} \right)^{N+1} u(y) .
	\]
\end{proof}
\begin{remark}
	Here are a few comments to Theorem~\ref{thm67bb79ea}:
	\begin{enumerate}
		\item	If $\inf_V u=0$, then $u=0$.
		\item	If $\inf_V u=1$, then $\sup_V u \le C$, where the constant $C$ is determined by $U$ and $V$, but independent on $u$.
		\item	Example on $\R$: take $V=(1,2)$ and $U=(0,4)$ and make a picture of affine maps that must be positive on $U$.
	\end{enumerate}
\end{remark}

\begin{exercise}\label{ex688c8182}
	Let $K\subset\R^n$ be compact and connected and $r>0$.
	Show that there is a finite family of points $x_1,\dots,x_N\in K$ such that $|x_j-x_{j+1}|<r$ and $K\subset \bigcup_{j=1}^N B(x_j,r)$.
\end{exercise}

%%%%%%%%%%%%%%%%%%%%%%%%%%%%%%%%%%%%%%%%%%%%%%%%%%%%%%%%%%%%%%%
\subsection{Green's function}\label{subs67bb7d34}
\index{Green's function}
\index{corrector function}
(See section 2.2.4 in Evans' book)

\begin{theorem}[Representation formula using Green's function]\label{thm67bb7d61}
	Let $U\subset\R^n$ be open, bounded and with $C^1$ boundary $\de U$.
	Suppose that there exists a \emph{corrector function} $\phi:\bar U\times\bar U\to\R$, $\phi(x,y) = \phi^x(y)$, such that, for every $x\in U$,
	\begin{equation}\label{eq67bb7eeb}
	\begin{cases}
		\laplacian \phi^x = 0 & \text{ in }U , \\
		\phi^x(y) = \Phi(y-x) & \forall y\in\de U ,
	\end{cases}
	\end{equation}
	where $\Phi$ is the fundamental solution in $\R^n$.
	Define the \emph{Green's function of $U$} as 
	\begin{equation}
		G : \{(x,y)\in U\times U: x\neq y\} \to \R ,
		\qquad
		G(x,y) = \Phi(y-x) - \phi^x(y) .
	\end{equation}
	
	If $u\in C^2(\bar U)$, then, for every $x\in U$,
	\begin{equation}\label{eq67bb7ef3}
		u(x) = - \int_{\de U} u(y) \grad G(x,y) \cdot \nu_{U}(y) \did S(y)
			- \int_U \laplacian u(y) G(x,y) \did y .
	\end{equation}
	In particular, if $f\in C^0(U)$, $g\in C^0(\de U)$ and $u\in C^2(\bar U)$ are such that
	\begin{equation}\label{eq67bb7ef9}
	\begin{cases}
		- \laplacian u = f & \text{ in }U , \\
		u = g & \text{ on }\de U ,
	\end{cases}
	\end{equation}
	then, for all $x\in U$,
	\begin{equation}\label{eq67bb7f00}
		u(x) = - \int_{\de U} g(y) \grad G(x,y) \cdot \nu_{U}(y) \did S(y)
			- \int_U f(y) G(x,y) \did y .
	\end{equation}
\end{theorem}
\begin{proof}
	If $\epsilon>0$ is such that $B(x,\epsilon)\Subset U$, then
	\[
	\int_U \laplacian u(y) G(x,y) \did y
	= \int_{B(x,\epsilon)} \laplacian u(y) G(x,y) \did y
		+ \int_{U\setminus B(x,\epsilon)} \laplacian u(y) G(x,y) \did y .
	\]
	where both $\laplacian u$ and $G(x,\cdot)$ belong to $L^1(U)$ 
	and thus 
	\[
	\lim_{\epsilon\to0} \int_{B(x,\epsilon)} \laplacian u(y) G(x,y) \did y = 0 .
	\]
	
	Set $V_\epsilon = U\setminus B(x,\epsilon)$.
	Then
	\begin{align*}
		&\int_{V_\epsilon} \laplacian u(y) G(x,y) \did y
		= \int_{V_\epsilon} \div_y(\grad u(y) G(x,y)) \did y
			- \int_{V_\epsilon} \grad u(y) \cdot \grad_y G(x,y) \did y \\
		&= \int_{V_\epsilon} \div_y(\grad u(y) G(x,y)) \did y
			- \int_{V_\epsilon} \div_y( u(y)  \grad_y G(x,y) ) \did y \\
		&\qquad 
			+ \int_{V_\epsilon} u(y)  \underbrace{\laplacian_y G(x,y)}_{=0} \did y \\
		&= \int_{\de V_\epsilon} \grad u(y) G(x,y)\cdot\nu_{V_\epsilon}(y) \did S(y)
			- \int_{\de V_\epsilon} u(y)  \grad_y G(x,y)\cdot\nu_{V_\epsilon}(y) \did S(y) \\
		&= - \int_{\de B(x,\epsilon)} \grad u(y) G(x,y) \cdot\frac{y-x}{|y-x|} \did S(y)
			+ \int_{\de U} \grad u(y) \underbrace{G(x,y)}_{=0} \cdot\nu_U(y) \did S(y) \\
		&\qquad
			- \int_{\de U} u(y) \grad_y G(x,y) \did S(y) 
			+ \int_{\de B(x,\epsilon)} u(y) \grad_y G(x,y)\cdot\frac{y-x}{|y-x|} \did S(y) \\
		&= - \int_{\de B(x,\epsilon)} \grad u(y) G(x,y) \cdot\frac{y-x}{|y-x|} \did S(y)
			- \int_{\de U} u(y) \grad_y G(x,y) \did S(y) \\ 
		&\qquad
			+ \int_{\de B(x,\epsilon)} u(y) \grad_y \Phi(y-x)\cdot\frac{y-x}{|y-x|} \did S(y) 
			+ \int_{\de B(x,\epsilon)} u(y) \grad_y \phi^x(y)\cdot\frac{y-x}{|y-x|} \did S(y) .
	\end{align*}
	Since we have
	\begin{align}
	\label{eq67c2ccf1}
		&\lim_{\epsilon\to0}
			\int_{\de B(x,\epsilon)}  G(x,y) \grad u(y)\cdot\frac{y-x}{|y-x|} \did y 
			= 0 ,\\
	\label{eq67c2ccfa}
		&\lim_{\epsilon\to0}
			\int_{\de B(x,\epsilon)} u(y) \grad_y \Phi(y-x)\cdot\frac{y-x}{|y-x|} \did y 
			= -u(x) ,
			\qquad\text{(see~Cor.~\ref{cor67c6cc59})}
			\\
	\label{eq67c2cd00}
		&\lim_{\epsilon\to0}
			\int_{\de B(x,\epsilon)} u(y) \grad_y \phi^x(y)\cdot\frac{y-x}{|y-x|} \did y 
			=0 ,
	\end{align}
	then we conclude that
	\begin{align*}
		\int_U \laplacian u(y) G(x,y) \did y
		&= \lim_{\epsilon\to0} \int_{B(x,\epsilon)} \laplacian u(y) G(x,y) \did y
			+ \int_{U\setminus B(x,\epsilon)} \laplacian u(y) G(x,y) \did y \\
		&= - u(x) - \int_{\de U} u(y) \grad_y G(x,y) \did y .
	\end{align*}
	So, we have~\eqref{eq67bb7ef3}.
\end{proof}

\begin{remark}
	If $U$ is connected, then for each $x\in U$ there exists at most one corrector function $\phi^x$,
	by Theorem~\ref{teo67baea93}.
	At the moment, however, we do not know whether corrector function exists.
\end{remark}

\begin{exercise}
	Show~\eqref{eq67c2ccf1}, \eqref{eq67c2ccfa}, and \eqref{eq67c2cd00}, 
	from the proof of the {\it Representation formula using Green's function}, Theorem~\ref{thm67bb7d61}:
	\begin{align*}
%	\label{eq67c2ccf1}
		&\lim_{\epsilon\to0}
			\int_{\de B(x,\epsilon)}  G(x,y) \grad u(y)\cdot\frac{y-x}{|y-x|} \did y 
			= 0 ,\\
%	\label{eq67c2ccfa}
		&\lim_{\epsilon\to0}
			\int_{\de B(x,\epsilon)} u(y) \grad_y \Phi(y-x)\cdot\frac{y-x}{|y-x|} \did y 
			= - u(x) ,\\
%	\label{eq67c2cd00}
		&\lim_{\epsilon\to0}
			\int_{\de B(x,\epsilon)} u(y) \grad_y \phi^x(y)\cdot\frac{y-x}{|y-x|} \did y 
			=0 .
	\end{align*}
\end{exercise}

%%%%%%%%%%%%%%%%%%%%%%%%%%%%%%%%%%%%%%%%%%%%%%%%%%%%%%%%%%%%%%%
\subsection{Symmetry of Green's functions}\label{subs67c2ce55}

\begin{theorem}[Symmetry of Green's functions]\label{thm67c2ce50}
	Let $U\subset\R^n$ be open and bounded, with $C^1$ boundary $\de U$.
	Let $G$ be the Green's function on $U$.
	Then
	\begin{equation}\label{eq67c2ceac}
		\forall x,y\in U\text{ with } x\neq y
		\qquad
		G(x,y) = G(y,x) .
	\end{equation}
\end{theorem}
\begin{proof}
	Fox $x,y\in U$ with $x\neq y$.
	Define
	\begin{equation}
		v(z) = G(x,z) = \Phi(z-x) - \phi^x(z) 
		\quad\text{and}\quad
		w(z) = G(y,z) = \Phi(z-y) - \phi^y(z) ,
	\end{equation}
	where $\Phi$ is the fundamental solution on $\R^n$ and $\phi^y$ is a corrector function.
	Then $\laplacian v = \laplacian w = 0 $ in $U\setminus\{x,y\}$ and $u=v=0$ on $\de U$.
	
	Let $\epsilon>0$ be such that $\bar B(x,\epsilon)\cup \bar B(y,\epsilon)\subset U$ and $\bar B(x,\epsilon)\cap \bar B(y,\epsilon) = \emptyset$.
	Set $V_\epsilon = U \setminus ( \bar B(x,\epsilon)\cup \bar B(y,\epsilon) )$.
	Then\footnote{
\begin{equation}\label{eq67c2d142}
	\int_U (v\laplacian w - w \laplacian v) \did x
	= \int_{\de U} (v\grad w - w\grad v)\cdot \nu_U \did S(x) .
\end{equation}
}
	\begin{align*}
		0
		&= \int_{V_\epsilon} (v\laplacian w - w \laplacian v) \did z \\
		&\overset{\eqref{eq67c2d142}}= \int_{\de V_\epsilon} (v\grad w - w\grad v)\cdot \nu_{V_\epsilon} \did S(z) \\
		&= -\int_{\de B(x,\epsilon)} (v\grad w - w\grad v)\cdot \frac{z-x}{|z-x|} \did S(z) 
			+ \int_{\de B(y,\epsilon)} (v\grad w - w\grad v)\cdot \frac{z-y}{|z-y|} \did S(z) . \\
%		&= \int_{\de V_\epsilon} (-v \grad w\cdot \nu_{V_\epsilon} + w \grad v\cdot \nu_{V_\epsilon} ) \did S(z) \\
%		&= \int_{\de B(x,\epsilon)} (-v\grad w + w \grad v) \cdot \frac{z-x}{\epsilon} \did S(z) 
%			+ \int_{\de B(y,\epsilon)} (-v\grad w+w\grad v) \cdot \frac{z-y}{\epsilon} \did S(z) .
	\end{align*}
	Since, using also Corollary~\ref{cor67c6cc59},
	\begin{equation}\label{eq67c2d634}
	\begin{aligned}
		&\lim_{\epsilon\to 0}
		\int_{\de B(x,\epsilon)} v\grad w \cdot \frac{z-x}{|z-x|} \did S(z)
		= 0 ,\\
		&\lim_{\epsilon\to 0}
		\int_{\de B(x,\epsilon)}  w\grad v\cdot \frac{z-x}{|z-x|} \did S(z)
		= - w(x) ,
	\end{aligned}
	\end{equation}
	then we obtain $w(x) = v(y)$.
\end{proof}

\begin{exercise}
	Show~\eqref{eq67c2d634} in the proof of the {\it Symmetry of Green's functions}, Theorem~\ref{thm67c2ce50}:
	\begin{equation*}%\label{eq67c2d634}
	\begin{aligned}
		&\lim_{\epsilon\to 0}
		\int_{\de B(x,\epsilon)} v\grad w \cdot \frac{z-x}{|z-x|} \did S(z)
		= 0 ,\\
		&\lim_{\epsilon\to 0}
		\int_{\de B(x,\epsilon)}  w\grad v\cdot \frac{z-x}{|z-x|} \did S(z)
		= - w(x) .
	\end{aligned}
	\end{equation*}
\end{exercise}

%%%%%%%%%%%%%%%%%%%%%%%%%%%%%%%%%%%%%%%%%%%%%%%%%%%%%%%%%%%%%%%
\subsection{Uniqueness by Energy methods}\label{subs67c2d8ed}

\begin{theorem}[Uniqueness by Energy methods]\label{thm67c2d8f8}
	Let $U\subset\R^n$ be open, bounded and with $C^1$ boundary $\de U$.
	Fix $f\in C^0(U)$ and $g\in C^0(\de U)$.
	Then there exists at most one solution $u$ in $C^2(\bar U)$ to the boundary value problem
	\begin{equation}\label{eq67c2d992}
	\begin{cases}
		- \laplacian u = f & \text{ in }U , \\
		u = g & \text{ on }\de U .
	\end{cases}
	\end{equation}
\end{theorem}
\begin{proof}
	Suppose $u_1,u_2\in C^2(\bar U)$ are solutions to~\eqref{eq67c2d992} .
	Then $u = u_1-u_2\in C^2(\bar U)$ is so that $\laplacian u =0$ in $U$ and $u=0$ on $\de U$.
	Therefore,
	\begin{align*}
		0
		&= \int_U u\laplacian u \did x \\
		&= \int_U \div(u \grad u) \did x - \int_U \grad u\cdot\grad u \did x \\
		&= \int_{\de U} u \grad u\cdot\nu_U \did S(x) - \int_U |\grad u|^2 \did x \\
		&= - \int_U |\grad u|^2 \did x .
	\end{align*}
	Or:
	\begin{align*}
		\int_U|\grad u|^2 \did x
		&= \int_U \grad u\cdot\grad u \did x \\
		&= \int_U (\div(u\grad u) - u \laplacian u) \did x \\
		&= \int_{\de U} u\grad u \did S(x) 
		= 0 .
	\end{align*}
	Hence, $\grad u = 0 $ in $U$, that is, since $U$ is connected, $u$ is constant.
	Since $u=0$ on $\de U$, then $u=0$ in $U$.
	

\end{proof}

\begin{remark}
	What is the difference between Theorem~\ref{thm67c2d8f8} and Theorem~\ref{teo67baea93}?
\end{remark}

%%%%%%%%%%%%%%%%%%%%%%%%%%%%%%%%%%%%%%%%%%%%%%%%%%%%%%%%%%%%%%%
\subsection{Dirichlet's principle}\label{subs67c2da65}
\index{Dirichlet's principle}
\index{principle!Dirichlet's --}
\index{theorem!Dirichlet's principle}
\newcommand{\dirichletE}{\mathtt{E}}

\begin{theorem}[Dirichlet's principle]\label{thm67c2da70}
	Let $U\subset\R^n$ be open, bounded and with $C^1$ boundary $\de U$.
	Fix $f\in C^0(U)$ and $g\in C^0(\de U)$.
	
	Define the \emph{Dirichlet's energy functional}
	\begin{equation}
		\dirichletE_f : C^2(\bar U)\to\R,
		\qquad
		\dirichletE_f(u) = \int_U (\frac12 |\grad u|^2 - uf) \did x .
	\end{equation}
	Define the \emph{admissible set}
	\begin{equation}
		\scr A_g = \{ u\in C^2(\bar U) : u = g \text{ on }\de U \} .
	\end{equation}
	For all $u\in C^2(\bar U)$, the following statements are equivalent:
	\begin{enumerate}[label=(\roman*)]
	\item\label{thm67c2da70_i}
	\eqref{eq67c2d992} holds, that is,
		\begin{equation*}
		\begin{cases}
			-\laplacian u = f & \text{ in }U , \\
			u = g & \text{ on }\de U .
		\end{cases}
		\end{equation*}
	\item\label{thm67c2da70_ii}
		$u\in\scr A_g$ and $\dirichletE_f(u) \le \dirichletE_f(w)$ for all $w\in \scr A_g$, i.e., 
		$u$ is the (unique) arg-min of $\dirichletE_f$ on $\scr A_g$.
	\end{enumerate}
\end{theorem}
\begin{proof}
	$\ref{thm67c2da70_i}\THEN\ref{thm67c2da70_ii}$:
	The fact that $u\in\scr A_g$ is clear. 
	We need to show that $\dirichletE_f(u) \le \dirichletE_f(w)$ for all $w\in \scr A_g$.
	If $w\in \scr A_g$, then
	\begin{align*}
	 	\dirichletE_f(u) - \dirichletE_f(w)
		&= \int_U \left( (|\grad u|^2/2 - uf ) - ( |\grad w|^2/2 - wf ) \right) \did x \\
		&= \frac12 \int_U (|\grad u|^2 - |\grad w|^2) \did x
			- \int_U (u-w) f \did x ,
	\end{align*}
	where
	\begin{align*}
	 	-\int_U (u-w) f \did x
		&= \int_U (u-w) \laplacian u \did x \\
		&= \int_U (\div((u-w)\grad u) - \grad(u-w) \cdot \grad u) \did x \\
		&= \underbrace{\int_{\de U} (u-w) \grad u \cdot \nu_U \did S }_{ =0 \text{ because }uw,\in\scr A_g}
			- \int_U |\grad u|^2 \did x 
			+ \int_U \grad w \cdot\grad u \did x .
	\end{align*}
	Therefore,
	\begin{align*}
	 	\dirichletE_f(u) - \dirichletE_f(w)
		&= - \int_U ( |\grad u|^2/2 + |\grad w|^2/2 - \grad w\cdot \grad u ) \did x \\
		&= - \int_U |\grad u-\grad w|^2 \did x 
		\le 0 .
	\end{align*}
	
	\begin{remark}
		We see from this calculation that, if $\dirichletE_f(u) = \dirichletE_f(w)$,
		then $u=w$.
		This is indeed the proof of the uniqueness in disguised terms.
	\end{remark}
	
	$\ref{thm67c2da70_ii}\THEN\ref{thm67c2da70_i}$:
	Fix $\phi\in C^\infty_c(U)$ and define $\iota_\phi:\R\to\R$ by
	\[
	\iota_\phi(\tau) := \dirichletE_f(u + \tau\phi) .
	\]
	Notice that $u+\tau\phi\in\scr A_g$ for all $\tau\in\R$,
	because $\phi=0$ on $\de U$.
	If $\iota_\phi$ is differentiable at $0$, then $\frac{\dd}{\dd \tau}\iota_\phi(0) = 0$.
	In fact, we have
	\begin{align*}
	 	\iota_\phi(\tau)
		&= \int_U (|\grad(u+\tau\phi)|^2/2 - (u+\tau\phi)f) \did x \\
		&= \tau^2 \int_U |\grad\phi|^2/2 \did x
			+ \tau \int_U (\grad u\cdot\grad\phi - \phi f) \did x 
			+ \int_U (|\grad u|^2/2 - uf) \did x .
	\end{align*}
	So, $\iota_\phi(\tau)$ is polynomial in $\tau$ and it has a minimum at $0$ if and only if
	$\int_U (\grad u\cdot\grad\phi - \phi f) \did x = 0$,
	where
	\begin{align*}
		\int_U (\grad u\cdot\grad\phi - \phi f) \did x
		&= \underbrace{\int_{\de U} \phi\grad u\cdot \nu_U \did S }{=0\text{ because }\phi=0\text{ on }\de U}
			- \int_U \phi(f+\laplacian u) \did x \\
		&= - \int_U \phi\cdot(f+\laplacian u) \did x .
	\end{align*}
	We conclude that, if $u$ is as in \ref{thm67c2da70_ii}, then 
	\begin{equation}\label{eq67c2eb31}
		\forall\phi\in C^\infty_c(U) ,
		\qquad
		\int_U \phi\cdot(f+\laplacian u) \did x = 0 .
	\end{equation}
	By the Fundamental Theorem of Calculus of Variations~\ref{thm67c2eb77}, 
	we~\eqref{eq67c2eb31} is equivalent to $-\laplacian u=f$ on $U$.
\end{proof}


\begin{remark}
	We have not proven yet that the boundary value problem~\eqref{eq67c2d992} has any solution at all.
	Do you have an idea of how to prove it? Think about it.
\end{remark}

%%%%%%%%%%%%%%%%%%%%%%%%%%%%%%%%%%%%%%%%%%%%%%%%%%%%%%%%%%%%%%%
\newpage
\section{Heat Equation}
\index{heat equation}
\index{equation!heat --}

For $U\subset\R^n$ open and $I\subset\R$ open interval, e.g., $I=(0,+\infty)$,
$u\in C^2(U\times I)$ and $f\in C^0(U\times I)$,
the \emph{heat equation} is
\begin{equation}
	\de_t u - \laplacian u = 0 \text{ in } U\times I ,
\end{equation}
and the \emph{nonhomogeneous heat equation} is
\begin{equation}
	\de_t u - \laplacian u = f \text{ in } U\times I .
\end{equation}
In these expressions, $\de_tu$ is the derivative of $u$ in the direction of $I$ in the product $U\times I$, 
while the Laplacian $\laplacian u = \laplacian_x u$ is with respect to the space variable, that is, in the directions $U$ in $U\times I$.
In other words, if $(x,t)$ are the coordinates of $U\times I$, with $x\in U$ and $t\in I$,
then
\[
(\de_t - \laplacian) u
=\de_t u - \laplacian u 
= \frac{\de u}{\de t}(x,t) - \sum_{j=1}^n \frac{\de^2u}{\de x_j^2}(x,t).
\]
We call $\de_t - \laplacian$ the \emph{heat operator}.
The heat operator $\de_t - \laplacian$ is a parabolic linear differential operator of order 2.

%%%%%%%%%%%%%%%%%%%%%%%%%%%%%%%%%%%%%%%%%%%%%%%%%%%%%%%%%%%%%%%
\subsection{Example of solutions}\label{subs67c45976}

\begin{exercise}
	Define 
	\begin{equation}
		\psi_+ : \R^n \times (0,+\infty) \to \R ,
		\qquad
		\psi_+(x,t) = \frac{1}{t^{n/2}} \exp\left(-\frac{|x|^2}{4t} \right) .
	\end{equation}
	Show that $(\de_t-\laplacian)\psi_+=0$ in $\R^n \times (0,+\infty)$.
	Draw a graph of $x\mapsto \psi_+(x,t)$ for positive $t$ when $n=1$
	(We will see that this function represents a forward propagation: this is why we have a plus.)
\end{exercise}

\begin{exercise}
	Define 
	\begin{equation}
		\psi_- : \R^n \times (0,+\infty) \to \R ,
		\qquad
		\psi_-(x,t) = \frac{1}{t^{n/2}} \exp\left(\frac{|x|^2}{4t} \right) .
	\end{equation}
	Show that $(\de_t-\laplacian)\psi_+=0$ in $\R^n \times (0,+\infty)$.
	Draw a graph of $x\mapsto \psi_+(x,t)$ for positive $t$ when $n=1$.
	(We will see that this function represents a backward propagation: this is why we have a minus.)
\end{exercise}

%%%%%%%%%%%%%%%%%%%%%%%%%%%%%%%%%%%%%%%%%%%%%%%%%%%%%%%%%%%%%%%
\subsection{Symmetries of the heat operator}\label{subs67d16430}
Let $U\subset\R^n$ open, $I\subset\R$ an open interval, $u\in C^2(U\times I)$, $O\in\Orth(n)$ and $b\in\R^n$, $\tau\in\R$, $\lambda\in\R\setminus\{0\}$.
Define $\tilde  u(y,s) = u(\lambda Oy+b,\lambda^2s+\tau)$.
Then $\tilde u\in C^2(O^{-1}(U-b)\times (I-\tau))$ and
\begin{equation}\label{eq67c3002d}
	(\de_s - \laplacian_y) \tilde u(y,s) = \lambda^2 (\de_tu - \laplacian_x u)(\lambda Oy+b, \lambda^2 s+\tau).
\end{equation}

\begin{exercise}
	Show the identity~\eqref{eq67c3002d}.
\end{exercise}

\begin{exercise}
	Let $U\subset\R^n$ open, $I\subset\R$ an open interval, $u\in C^2(U\times I)$, $O\in\Orth(n)$ and $b\in\R^n$, $\tau\in\R$ and $\lambda,\sigma\in\R\setminus\{0\}$
	Define 
	\[
	\tilde  u(y,s) = u(\lambda Oy+b,\sigma s+\tau) .
	\]
	Compute $(\de_t - \laplacian) \tilde u$ in terms of $(\de_t - \laplacian) u$.
	Determine for which choices of transformations we have that, if $u$ is a solution to the heat equation, i.e., $(\de_t - \laplacian) u=0$, then $\tilde u$ is also a solution to the heat equation, i.e., $(\de_t - \laplacian) \tilde u=0$.
\end{exercise}

\begin{exercise}
	Show that, if $\lambda>0$ and if $(\de_t - \laplacian) u=0$, 
	then $(\de_t - \laplacian) \tilde u=0$,
	where $\tilde u(x,t) = u(\lambda x,\lambda^2 t)$.
\end{exercise}

%%%%%%%%%%%%%%%%%%%%%%%%%%%%%%%%%%%%%%%%%%%%%%%%%%%%%%%%%%%%%%%
\subsection{Fundamental solution for the heat equation}\label{subs67c6d9a8}
\index{fundamental solution for the heat equation}
\index{heat!fundamental solution for the -- equation}

The \emph{fundamental solution for the heat equation} is the function 
$\Phi:\R^n\times\R\setminus\{(0,0)\} \to [0,+\infty)$ defined by
\begin{equation}\label{eq67c3022e} %[
	\Phi(x,t) = \begin{cases}
	\frac{1}{(4\pi t)^{n/2}} \exp\left( -\frac{|x|^2}{4t} \right)
		& \text{ for $x\in\R^n$ and $t>0$,} \\
	0 & \text{otherwise, i.e., }(x,t)\in(\R^n\times(-\infty,0])\setminus\{(0,0)\} .
	\end{cases}
\end{equation}%)

\begin{exercise}
	Find the formula for the fundamental solution for the heat equation by yourself in the following way:
	look for $\Phi$ of the form 
	$\Phi(x,t) = \frac1{t^\alpha} v\left( \frac{|x|}{t^\beta} \right)$,
	or
	$\Phi(x,t) = \frac1{t^\alpha} v\left( \frac{|x|^2}{t^\beta} \right)$,
	with $(\de_t-\laplacian_x)\Phi=0$ in $\R^n\times(0,+\infty)$.
\end{exercise}

\begin{lemma}
	\begin{equation}\label{eq67c3065f}
		\int_\R e^{-x^2}\did x = \sqrt{\pi} .
	\end{equation}
\end{lemma}
\begin{proof}
	See \href{https://math.stackexchange.com/questions/154968/is-there-really-no-way-to-integrate-e-x2}{stackexchange}.
\end{proof}

\begin{exercise}
	Find a proof (by yourself or in the literature) for~\eqref{eq67c3065f}, that is,
	\begin{equation}
		\int_\R e^{-x^2}\did x = \sqrt{\pi} .
	\end{equation}
\end{exercise}

\begin{proposition}[Properties of the fundamental solution]\label{prop67c302e8}
	The function $\Phi$ defined in~\eqref{eq67c3022e} has the following properties:
	\begin{enumerate}
	\item\label{prop67c302e8_smooth}
	$\Phi\in C^\infty(\R^n\times\R\setminus\{(0,0)\})$.
	\item\label{prop67c302e8_diff}
	for every $t>0$ and $x\in\R^n$, 
	\begin{align*}
		\grad_x\Phi(x,t) 
		&= - \frac1{(4\pi t)^{n/2}} \exp\left( -\frac{|x|^2}{4t} \right) \frac{x}{2t}
		= - \Phi(x,t) \frac{x}{2t} , \\
		\de_t \Phi(x,t) 
		&= \frac1{(4\pi t)^{n/2}} \exp\left( -\frac{|x|^2}{4t} \right) 
			\left( \frac{|x|^2}{4t^2} - \frac{n}{2} \frac1{4\pi t} \right) 
		= \Phi(x,t) \left( \frac{|x|^2}{4t^2} - \frac{n}{2t} \right) , \\
		\Diff^2_x\Phi(x,t)
		&= \Phi(x,t) \left( \frac{x\otimes x}{4t^2} - \frac{\Id}{2t} \right) .
	\end{align*} 
	\item\label{prop67c302e8_heat}
	$(\de_t-\laplacian)\Phi = 0$ in $\R^n\times\R\setminus\{(0,0)\}$.
	\item\label{prop67c302e8_intr}
	for every $t>0$, $\int_{\R^n} \Phi(x,t) \did x = 1$
	\end{enumerate}
\end{proposition}
\begin{proof}
	\fbox{Proof of~\ref{prop67c302e8_smooth}:}
	To show that $\Phi$ is smooth, we proceed as follows.
	Define
	\begin{equation*}
		\scr F = \left\{ 
			\phi:\R^n\times\R\setminus\{(0,0)\} \to \R :
			\begin{array}{l}
			\exists P(x,t) \text{ polynomial with }\\
			\phi(x,t) = P(x,t^{-1/2})\exp\left( -\frac{|x|^2}{4t} \right)\text{ for $t>0$,}\\
			\text{while $\phi(x,t)=0$ for $t\le 0$}
			\end{array}
		\right\}
	\end{equation*}
	Then we have the following two facts (whose proof is left as an exercise):
	First, $\scr F\subset C^0(\R^n\times\R)$.
	Second, if $\phi\in\scr F$, then $\frac{\de\phi}{\de x_j},\frac{\de\phi}{\de t}\in\scr F$, for all $j\in\{1,\dots,n\}$.
	We conclude that $\scr F\subset C^\infty(\R^n\times\R\setminus\{(0,0)\})$.
	Since $\Phi\in\scr F$, we the proof is complete.
	
	\fbox{Proofs of~\ref{prop67c302e8_diff} and~\ref{prop67c302e8_heat}}
	are left as exercise.
	
	\fbox{Proof of~\ref{prop67c302e8_intr}:}
	Fix $t>0$. Then we have
	\begin{align*}
		\int_{\R^n} \Phi(x,t) \did x
		&= \int_{\R^n} \exp\left( -\frac{|x|^2}{4t} \right) \frac{\did x}{(4\pi t)^{n/2}} \\
		[y=\frac{x}{2\sqrt{t}}]
		&= \frac1{\pi^{n/2}} \int_{\R^n} \exp(-|y|^2) \did y \\
		&= \frac1{\pi^{n/2}} \prod_{j=1}^n \int_\R \exp(-y_j^2) \did y_j \\
		&\overset{\eqref{eq67c3065f}}= 1 .
	\end{align*}
	
\end{proof}

\begin{exercise}
	Complete the proof of part~\ref{prop67c302e8_smooth} in Proposition~\ref{prop67c302e8}.
	Specifically, define
	\begin{equation*}
		\scr F = \left\{ 
			\phi:\R^n\times\R\setminus\{(0,0)\} \to \R :
			\begin{array}{l}
			\exists P(x,t) \text{ polynomial with }\\
			\phi(x,t) = P(x,t^{-1/2})\exp\left( -\frac{|x|^2}{4t} \right)\text{ for $t>0$,}\\
			\text{while $\phi(x,t)=0$ for $t\le 0$}
			\end{array}
		\right\} .
	\end{equation*}
	Then show
	\begin{enumerate}
	\item	$\scr F\subset C^0(\R^n\times\R)$.
	\item if $\phi\in\scr F$, then $\frac{\de\phi}{\de x_j},\frac{\de\phi}{\de t}\in\scr F$, for all $j\in\{1,\dots,n\}$.
	\end{enumerate}
	Conclude that $\scr F\subset C^\infty(\R^n\times\R\setminus\{(0,0)\})$.
\end{exercise}

\begin{exercise}
	Prove parts~\ref{prop67c302e8_diff} and~\ref{prop67c302e8_heat} in Proposition~\ref{prop67c302e8}.
	Specifically, show that 
	for every $t>0$ and $x\in\R^n$, 
	\begin{align*}
		\grad_x\Phi(x,t) 
		&= - \frac1{(4\pi t)^{n/2}} \exp\left( -\frac{|x|^2}{4t} \right) \frac{x}{2t}
		= - \Phi(x,t) \frac{x}{2t} , \\
		\de_t \Phi(x,t) 
		&= \frac1{(4\pi t)^{n/2}} \exp\left( -\frac{|x|^2}{4t} \right) 
			\left( \frac{|x|^2}{4t^2} - \frac{n}{2} \frac1{4\pi t} \right) 
		= \Phi(x,t) \left( \frac{|x|^2}{4t^2} - \frac{n}{2t} \right) , \\
		\Diff^2_x\Phi(x,t)
		&= \Phi(x,t) \left( \frac{x\otimes x}{4t^2} - \frac{\Id}{2t} \right) ,
	\end{align*} 
	where $\Phi$ is the fundamental solution of the heat operator.
	Conclude that $(\de_t-\laplacian_x)\Phi=0$ in $\R^n\times\R\setminus\{(0,0)\}$.
\end{exercise}


%%%%%%%%%%%%%%%%%%%%%%%%%%%%%%%%%%%%%%%%%%%%%%%%%%%%%%%%%%%%%%%
\subsection{Solution to the Cauchy problem}\label{subs67c30b65}

\begin{theorem}\label{thm67c30be5}
	Let $g\in C^0(\R^n)\cap L^\infty(\R^n)$ and define $u:\R^n\times (0,+\infty) \to \C$ as
	\begin{equation}
		u(x,t) = \int_{\R^n} \Phi(x-y,t) g(y) \did y
		= \frac1{(4\pi t)^{n/2}} \int_{\R^n} \exp\left( -\frac{|x-y|^2}{4t} \right) g(y) \did y ,
	\end{equation}
	for $x\in\R^n$ and $t>0$,
	where $\Phi$ is the fundamental solution for the heat operator~\eqref{eq67c3022e}.
	\begin{enumerate}
	\item\label{thm67c30be5_smooth}
		$u\in C^\infty(\R^n\times (0,+\infty))$ and, for every $\alpha\in\N^{n+1}$,
		\begin{equation}\label{eq67c33d51}
			\Diff^\alpha u(x,t) = \int_{\R^n} \Diff_{x,t}^\alpha\Phi(x-y,t) g(y) \did y .
		\end{equation}
	\item\label{thm67c30be5_heat}
		$(\de_t-\laplacian)u=0$ in $\R^n\times (0,+\infty)$;
	\item\label{thm67c30be5_init}
		for each $\hat x\in\R^n$, 
		\begin{equation}\label{eq67c342bb}
			\lim_{\substack{(x,t)\to(\hat x,0) \\ x\in\R^n,t>0}} u(x,t) = g(x) .
		\end{equation}
	\end{enumerate}
	%(
	In particular, $u$ has a continuous extension $u\in C^0(\R^n\times [0,+\infty))\cap C^\infty(\R^n\times (0,+\infty))$ and %]
	\begin{equation}\label{eq67c30ca6}
	\begin{cases}
		\de_t u - \laplacian u = 0 & \text{ in }\R^n\times (0,+\infty) , \\
		u = g & \text{ on }\R^n\times\{0\} .
	\end{cases}
	\end{equation}
\end{theorem}
\begin{proof}
	\fbox{Proof of~\ref{thm67c30be5_smooth}:}
	\[
	K(x,t;y) = \frac1{(4\pi t)^{n/2}} \exp\left( -\frac{|x-y|^2}{4t} \right) g(y) .
	\]
	For $R>0$ and $\epsilon>0$, we have, for every $x\in\R^n$ with $|x|< R$ and all $t>\epsilon$,
	\[
	|K(x,t;y)|\le h_{R,\epsilon}(y)
	:= \begin{cases}
	\frac{\|g\|_{L^\infty}}{(4\pi \epsilon)^{n/2}} & \text{ if }|y|\le R \\
	\frac{\|g\|_{L^\infty}}{(4\pi \epsilon)^{n/2}} \exp\left( -\frac{(|y|-R)^2}{4 \epsilon} \right) & \text{ if }|y| > R \\
	\end{cases}
	\]
	Since $h_R\in L^1(\R^n)$, satisfy the conditions of~Theorem~\ref{thm679f29de_2}, so $u\in C^0(B(0,R)\times (\epsilon,+\infty))$.
	Since this holds for every $R>0$ and $\epsilon>0$, we conclude $u\in C^0(\R^n\times (0,+\infty))$.
	To prove the smoothness of $u$, one uses the same strategy as in the proof of Proposition~\ref{prop67c302e8}.\ref{prop67c302e8_smooth};
	see Exercise~\ref{exr67c33c9d}.
	Also~\eqref{eq67c33d51} follows from Proposition~\ref{prop67c302e8}.
	
	\fbox{Proof of~\ref{thm67c30be5_heat}:}
	From~\eqref{eq67c33d51}, we have
	\begin{equation}
		(\de_t-\laplacian)u(x,t) 
		= \int_{\R^n} (\de_t-\laplacian_x)\Phi(x-y,t) g(y) \did y 
		= 0 ,
	\end{equation}
	because $(\de_t-\laplacian_x)\Phi(x-y,t) = 0$ for every $y\in\R^n$ and $t>0$,
	thanks to Proposition~\ref{prop67c302e8}.\ref{prop67c302e8_heat}.
	
	\fbox{Proof of~\ref{thm67c30be5_init}:}
	Fix $\hat x\in\R^n$ and $\epsilon>0$.
	Since $g$ is continuous, there exists $\delta>0$ such that $|g(y)-g(\hat x)|<\epsilon$ for all $y\in B(\hat x,\delta)$.
	\begin{align*}
	 	|u(x,t)-g(\hat x)|
		&= \left| \int_{\R^n} \Phi(x-y,t) g(y) \did y - g(\hat x) \int_{\R^n} \Phi(x-y,t)\did y \right| \\
		&\le \int_{\R^n} \Phi(x-y,t) |g(y)-g(\hat x)| \did y \\
		&= \int_{B(\hat x,\delta)} \Phi(x-y,t) \underbrace{|g(y)-g(\hat x)|}_{<\epsilon} \did y \\
		&\quad + \int_{\R^n\setminus B(\hat x,\delta)} \Phi(x-y,t) |g(y)-g(\hat x)| \did y \\
		&\le \epsilon \underbrace{\int_{\R^n} \Phi(x-y,t) \did y}_{=1}
			+ \underbrace{ \int_{\R^n\setminus B(\hat x,\delta)} \Phi(x-y,t) |g(y)-g(\hat x)| \did y }_{=: J_\delta(x)} .
	\end{align*}
	If $x\in B(\hat x,\delta/2)$, then, for all $y\in \R^n\setminus B(\hat x,\delta)$,
	\begin{equation}\label{eq67c340a0}
		|y-x| 
		\ge |y-\hat x|-|\hat x-x|
		\ge |y-\hat x|-\delta/2
		\ge |y-\hat x| - |y-\hat x|/2
		= |y-\hat x|/2 .
	\end{equation}
	Hence, if $x\in B(\hat x,\delta)$, 
	\begin{align*}
	 	J_\delta(x) 
		&\le 2 \|g\|_{L^\infty} \int_{\R^n\setminus B(\hat x,\delta)} \frac{ \exp(-\frac{|x-y|^2}{4t} ) }{ (4\pi t)^{n/2} } \did y \\
		&\overset{\eqref{eq67c340a0}}\le \frac{2 \|g\|_{L^\infty}}{(4\pi)^{n/2}}
			\int_{\R^n\setminus B(\hat x,\delta)} \frac{ \exp(-\frac{|y-\hat x|^2}{16t} ) }{  t^{n/2} } \did y \\
		[z = \frac{y-\hat x}{\sqrt{t}}]
		&\le \frac{2 \|g\|_{L^\infty}}{(4\pi)^{n/2}} 
			\underbrace{\int_{\R^n\setminus B(\hat x,\delta/\sqrt{t})} \exp(-\frac{|z|^2}{16} ) \did y }_{=:E(\delta,t)} .
	\end{align*}
	Since $\delta>0$ and since the integrand in $E(\delta,t)$ is integrable over $\R^n$, then $\lim_{t\to0^+}E(\delta,t) = 0$.
	Therefore, there exists $\tau>0$ such that $E(\delta,t)<\epsilon$ for all $t\in(0,\tau)$.
	All in all, we conclude that
	\begin{equation}\label{eq67c342cf}
		\forall\epsilon>0\exists \delta,\tau>0 \forall x\in B(\hat x,\delta/2)\ \forall t\in(0,\tau)
		\qquad
		|u(x,t) - g(\hat x)| \le \epsilon + \frac{2 \|g\|_{L^\infty}}{(4\pi)^{n/2}} \epsilon .
	\end{equation}
	\eqref{eq67c342cf} is~\eqref{eq67c342bb}.
\end{proof}

\begin{exercise}\label{exr67c33c9d}
	Show part~\ref{thm67c30be5_smooth} in Theorem~\ref{thm67c30be5}.
\end{exercise}

%%%%%%%%%%%%%%%%%%%%%%%%%%%%%%%%%%%%%%%%%%%%%%%%%%%%%%%%%%%%%%%
\subsection{Approximation of the identity}\label{subs67c41af0}
From the proof of Theorem~\ref{thm67c30be5}, we can extract the following lemma:
\begin{lemma}\label{lem67c41d25}
	Let $T>0$ and $g\in C^0(\R^n\times[0,T])\cap L^\infty(\R^n\times[0,T])$.
	Then, for every $\hat x\in\R^n$,
	\begin{equation}\label{eq67c41b8a}
		\lim_{\substack{(x,t)\to(\hat x,0)\\t>0}}
		\int_{\R^n} \Phi(z,t) g( x - z,t) \did z = g(\hat x,0) .
	\end{equation}
\end{lemma}
\begin{exercise}
	Prove Lemma~\ref{lem67c41d25}.
\end{exercise}
\begin{proof}
	Define
	\begin{equation}
		u(x,t) = \int_{\R^n} \Phi(z,t) g(x - z,t) \did z 
		= \int_{\R^n} \Phi(x-z,t) g(z,t) \did z  .
	\end{equation}
	Fix $\hat x\in\R^n$ and $\epsilon>0$.
	Since $g$ is continuous, there exists $\delta>0$ such that $|g(y,t)-g(\hat x,0)|<\epsilon$ for all $y\in B(\hat x,\delta)$ and all $t\in[0,\delta]$.
	Then
	\begin{align*}
	 	|u(x,t)-g(\hat x,0)|
		&= \left| \int_{\R^n} \Phi(x-y,t) g(y,t) \did y - g(\hat x,0) \int_{\R^n} \Phi(x-y,t)\did y \right| \\
		&\le \int_{\R^n} \Phi(x-y,t) |g(y,t)-g(\hat x,0)| \did y \\
		&= \int_{B(\hat x,\delta)} \Phi(x-y,t) \underbrace{|g(y,t)-g(\hat x,0)|}_{<\epsilon} \did y \\
		&\quad + \int_{\R^n\setminus B(\hat x,\delta)} \Phi(x-y,t) |g(y,t)-g(\hat x,0)| \did y \\
		&\le \epsilon \underbrace{\int_{\R^n} \Phi(x-y,t) \did y}_{=1}
			+ \underbrace{ \int_{\R^n\setminus B(\hat x,\delta)} \Phi(x-y,t) |g(y,t)-g(\hat x,0)| \did y }_{=: J_\delta(x)} .
	\end{align*}
	If $x\in B(\hat x,\delta/2)$, then, for all $y\in \R^n\setminus B(\hat x,\delta)$,
	\begin{equation}\label{eq67c41ccb}
		|y-x| 
		\ge |y-\hat x|-|\hat x-x|
		\ge |y-\hat x|-\delta/2
		\ge |y-\hat x| - |y-\hat x|/2
		= |y-\hat x|/2 .
	\end{equation}
	Hence, if $x\in B(\hat x,\delta/2)$, 
	\begin{align*}
	 	J_\delta(x) 
		&\le 2 \|g\|_{L^\infty} \int_{\R^n\setminus B(\hat x,\delta)} \frac{ \exp(-\frac{|x-y|^2}{4t} ) }{ (4\pi t)^{n/2} } \did y \\
		&\overset{\eqref{eq67c41ccb}}\le \frac{2 \|g\|_{L^\infty}}{(4\pi)^{n/2}}
			\int_{\R^n\setminus B(\hat x,\delta)} \frac{ \exp(-\frac{|y-\hat x|^2}{16t} ) }{  t^{n/2} } \did y \\
		[z = \frac{y-\hat x}{\sqrt{t}}]
		&= \frac{2 \|g\|_{L^\infty}}{(4\pi)^{n/2}} 
			\underbrace{\int_{\R^n\setminus B(\hat x,\delta/\sqrt{t})} \exp(-\frac{|z|^2}{16} ) \did y }_{=:E(\delta,t)} .
	\end{align*}
	Since $\delta>0$ and since the integrand in $E(\delta,t)$ is integrable over $\R^n$, then $\lim_{t\to0^+}E(\delta,t) = 0$.
	Therefore, there exists $\tau\in(0,\delta)$ such that $E(\delta,t)<\epsilon$ for all $t\in(0,\tau)$.
	All in all, we conclude that
	\begin{equation}\label{eq67c342cf_bis}
		\forall\epsilon>0\exists \delta,\tau>0 \forall x\in B(\hat x,\delta/2)\ \forall t\in(0,\tau)
		\qquad
		|u(x,t) - g(\hat x,0)| \le \epsilon + \frac{2 \|g\|_{L^\infty}}{(4\pi)^{n/2}} \epsilon .
	\end{equation}
	\eqref{eq67c342cf_bis} is~\eqref{eq67c41b8a}.
\end{proof}

%%%%%%%%%%%%%%%%%%%%%%%%%%%%%%%%%%%%%%%%%%%%%%%%%%%%%%%%%%%%%%%
\subsection{Solution to the nonhomogeneous Cauchy problem}\label{subs67c34370}

A solution for the nonhomogeneous Cauchy problem~\eqref{eq67c344a9} is constructed as follows.
For every $s>0$ let $u_s$ be the solution constructed in the previous Theorem for the Cauchy problem
\begin{equation}
\begin{cases}
	\de_t u_s - \laplacian u_s = 0 & \text{ in }\R^n\times (s,+\infty) , \\
	u_s = f(\cdot,s) & \text{ on }\R^n\times\{s\} .
\end{cases}
\end{equation}
So, $u_s(x,t) = \int_{\R^n} \Phi(x-y,t-s) f(y,s) \did y$.
Then, we set $u(x,t) = \int_0^t u_s(x,t) \did s$, that is,~\eqref{eq67c34492}.
This strategy is called \emph{Duhamel's principle}.

\begin{theorem}[Solution to the nonhomogeneous Cauchy problem]\label{thm67c3437d}
	Let $f\in C^{2;1}_c(\R^n\times[0,+\infty))$ and define $u:\R^n\times(0,+\infty)\to\R$ by
	\begin{equation}\label{eq67c34492}
	\begin{aligned}
		u(x,t) &= \int_0^t \int_{\R^n} \Phi(x-y,t-s) f(y,s) \did y\did s \\
		&= \int_0^t \frac{1}{(4\pi(t-s))^{n/2}} \int_{\R^n} \exp\left( - \frac{ |x-y|^2 }{ 4(t-s) } \right) f(y,s) \did y \did s .
	\end{aligned}
	\end{equation}
	Then the following holds:
	\begin{enumerate}
	\item\label{thm67c3437d_smooth}
		$u\in C^{2;1}(\R^n\times(0,+\infty))$;
	\item\label{thm67c3437d_heat}
		$(\de_t-\laplacian)u = f$ in $\R^n\times(0,+\infty)$;
	\item\label{thm67c3437d_init}
		for every $\hat x\in\R^n$, 
		\begin{equation}\label{eq67c3448d}
			\lim_{\substack{(x,t)\to(\hat x,0) \\ x\in\R^n,t>0}} u(x,t) = 0 .
		\end{equation}
	\end{enumerate}
	In particular, $u$ has a continuous extension $u\in C^0(\R^n\times [0,+\infty))\cap C^\infty(\R^n\times (0,+\infty))$ and
	\begin{equation}\label{eq67c344a9}
	\begin{cases}
		\de_t u - \laplacian u = f & \text{ in }\R^n\times (0,+\infty) , \\
		u = 0 & \text{ on }\R^n\times\{0\} .
	\end{cases}
	\end{equation}
\end{theorem}
\begin{proof}
	\fbox{Proof of~\ref{thm67c3437d_smooth}.}
	First, notice that the integrand in the definition~\eqref{eq67c34492} of $u$ is integrable and thus $u$ is well defined.
	Moreover,
	\begin{equation}
		u(x,t) = \int_0^t \int_{\R^n} \Phi(z,r) f(x-z,t-r) \did z\did r .
	\end{equation}
	Define
	\begin{equation}
		K(x,t;z,r) := \Phi(z,r) f(x-z,t-r) \one_{[0,t]}(r).
	\end{equation}
	Then $|K(x,t;z,r)| \le \|f\|_{L^\infty} \Phi(z,r)$ for all $(x,t),(z,r)\in\R^n\times\R\setminus\{(0,0)\}$.
	So, we can apply Theorem~\ref{thm679f29de}.
	
	\fbox{Proof of~\ref{thm67c3437d_heat}.}
	With the support of Theorem~\ref{thm679f29de}, we can compute, for $x\in\R^n$, $t>0$, and $\alpha\in\N^n$ with $|\alpha|\le 2$,
	\begin{align*}
	 	\de_t u(x,t) 
		&= \int_{\R^n} \Phi(z,t) f(x-z,0) \did z
			+ \int_0^t \int_{\R^n} \Phi(z,r) \de_tf(x-z,t-r) \did z\did r , \\
		\Diff^\alpha u(x,y)
		&= \int_0^t \int_{\R^n} \Phi(z,r) \Diff^\alpha f(x-z,t-r) \did z\did r .
	\end{align*}
	Therefore, for $x\in\R^n$, $t>0$ and $\epsilon\in(0,t)$,
	\begin{align*}
		(\de_t - \laplacian)u(x,t)
		&= \int_{\R^n} \Phi(z,t) f(x-z,0) \did z \\
		&\qquad
			+ \int_0^t \int_{\R^n} \Phi(z,r) (\de_t - \laplacian)f(x-z,t-r) \did z\did r \\
		&= \underbrace{ \int_{\R^n} \Phi(z,t) f(x-z,0) \did z }_{:= K}\\
		&\qquad
			+ \underbrace{ \int_0^\epsilon \int_{\R^n} \Phi(z,r) (\de_t - \laplacian)f(x-z,t-r) \did z\did r }_{:= J_\epsilon} \\
		&\qquad
			+ \underbrace{ \int_\epsilon^t \int_{\R^n} \Phi(z,r) (\de_t - \laplacian)f(x-z,t-r) \did z\did r }_{:= I_\epsilon} .
	\end{align*}
	Now, 
	\begin{align*}
		|J_\epsilon| 
		&\le (\|\de_t f\|_{L^\infty} + \|\Diff^2f\|_{L^\infty} )
			\int_0^\epsilon \int_{\R^n} \Phi(z,r) \did z\did r \\
		&= \epsilon (\|\de_t f\|_{L^\infty} + \|\Diff^2f\|_{L^\infty} ) .
	\end{align*}
	Next,
	\begin{align*}
	 	I_\epsilon 
		&= \int_\epsilon^t \int_{\R^n} \Phi(z,r) (\de_t - \laplacian_x)f(x-z,t-r) \did z\did r \\
		&= \int_\epsilon^t \int_{\R^n} \Phi(z,r) (-\de_r - \laplacian_z)f(x-z,t-r) \did z\did r \\
		&= \int_\epsilon^t \int_{\R^n} (-\de_r(\Phi(z,r)f(x-z,t-r)) + \de_r\Phi(z,r)f(x-z,t-r) ) \did z\did r \\
		&\qquad
			- \int_\epsilon^t \int_{\R^n} \div_z\Big(\Phi(z,r)\grad_zf(x-z,t-r) - \grad_z\Phi(z,r)f(x-z,t-r)\Big) \did z\did r \\
		&\qquad
			- \int_\epsilon^t \int_{\R^n} \laplacian\Phi(z,r)f(x-z,t-r)) \did z\did r \\
		&\overset{(*)}= -K + \int_{\R^n} \Phi(z,\epsilon)f(x-z,t-\epsilon) \did z\did r ,
	\end{align*}
	where we used in $(*)$ the fact that $\Phi$ solves the homogeneous heat equation and that $f$ has compact support.
	So, we conclude that
	\begin{align*}
		(\de_t - \laplacian)u(x,t)
		&= \lim_{\epsilon\to0} \int_{\R^n} \Phi(z,\epsilon)f(x-z,t-\epsilon) \did z\did r \\
		[\text{Lemma~\ref{lem67c41d25}}]
		&= f(x,t) .
	\end{align*}
	
	\fbox{Proof of~\ref{thm67c3437d_init}.}
	We easily conclude with the following estimate:
	\begin{align*}
		|u(x,t)|
%		&\le \int_0^t \int_{\R^n} \Phi(x-y,t-s) |f(y,s)| \did y\did s \\
		&\le \|f\|_{L^\infty} \int_0^t \int_{\R^n} \Phi(x-y,t-s) \did y\did s \\
		\text{[by~\ref{prop67c302e8}.\ref{prop67c302e8_intr}]}&= \|f\|_{L^\infty} t .
	\end{align*}
\end{proof}

%%%%%%%%%%%%%%%%%%%%%%%%%%%%%%%%%%%%%%%%%%%%%%%%%%%%%%%%%%%%%%%
\subsection{Mean-value formula}\label{subs67c44996}
For $(x,t)\in R^n\times\R$ and $r>0$, define
\begin{equation}\label{eq67ceb710}
\begin{aligned}
	E(x,t;r)
	&= \{(y,s) \in \R^n\times\R : s\le t,\ \Phi(x-y,t-s) \ge \frac1{r^n} \} \\
	&= (x,t) + E(0,0;r)
	= (x,t) + \delta_r(E(0,0;1)) ,
\end{aligned}
\end{equation}
where $\delta_r(y,s) = (ry,r^2 s)$.
Notice that $E(0,0;1) = \{(y,s): \Phi(-y,-s) \ge 1\}$.
See Figure~\ref{fig67cf14e8} for a drawing of the shape of $E(x,t;r)$.
Since
\begin{align}
	\Phi(-y,-s) \ge 1
	&\IFF \frac{(4\pi|s|)^{n/2}} \exp\left(-\frac{|y|^2}{4|s|}\right) \ge 1 \\
	&\IFF \frac{n}{2} \log(4\pi|s|) \le -\frac{|y|^2}{4|s|} \\
	&\IFF -\frac{1}{4\pi} < s < 0 \qquad |y|^2 \le - 2n |s| \log(4\pi|s|) .
\end{align}

\begin{theorem}\label{thm67cf158e}
	Let $\Omega\subset\R^n\times\R$ be open and $u\in C^{2;1}(\Omega)$ such that $(\de_t-\laplacian)u=0$ in $\Omega$.
	Then, for all $(x,t)\in\Omega$ and $r>0$ such that $E(x,t;r)\subset \Omega$, we have
	\begin{equation}
		u(x,t) = \frac{1}{4r^n} \iint_{E(x,t;r)} u(y,s) \frac{|x-y|^2}{(t-s)^2} \did y \did s .
	\end{equation}
\end{theorem}
\begin{proof}
	See \cite[pag.53]{MR2597943}.
\end{proof}

\begin{figure}\label{fig67cf14e8}
\usepgfplotslibrary{fillbetween}
\begin{tikzpicture}
    \begin{axis}[
        xmin=-1,xmax=1,
        ymin=-1.25,ymax=0.25,
        axis x line=middle,
        axis y line=middle,
        axis line style=-,
        xtick=\empty,ytick=\empty,
        ]
		\addplot+ [name path=A,black,smooth,no markers,domain=-1:0,variable=\t,samples=50]  ({ sqrt(t*ln(abs(t)))} ,{ t });
		\addplot+ [name path=B,black,smooth,no markers,domain=-1:0,variable=\t,samples=50]  ({ -sqrt(t*ln(abs(t)))} ,{ t });
		\addplot[blue!50] fill between[of=A and B];
		\addplot [only marks] coordinates {(0,0)} node[above right] {$(x,t)$};
		\addplot [no marks] coordinates {(0.3,-0.5)} node {$E(x,t;r)$};
    \end{axis}
\end{tikzpicture}
\caption{The shape of the set $E(x,t;r)$ defined in~\ref{subs67c44996}.}
\end{figure}

%%%%%%%%%%%%%%%%%%%%%%%%%%%%%%%%%%%%%%%%%%%%%%%%%%%%%%%%%%%%%%%
\subsection{Strong maximum and minimum principles}\label{subs67c44a2b}
\index{closed parabolic cylinder}
\index{parabolic interior}
\index{parabolic boundary}

For $U\subset\R^n$ and $T>0$, define:
\begin{enumerate}
\item
the \emph{closed parabolic cylinder} $\bar U_T := \bar U\times[0,T]$;
\item
the \emph{parabolic interior} $U_T := U\times(0,T]$ (notice that $T$ is included);
\item
the \emph{parabolic boundary} $\Gamma_T := \bar U_T \setminus U_T = (U\times\{0\}) \cup (\de U\times [0,T) )$.
\end{enumerate}

\begin{theorem}[Strong maximum principle]\label{thm67c44b0d}
	Let $U\subset\R^n$ be open and bounded.
	Let $u\in C^{2;1}(U_T;\R)\cap C^0(\bar U_T;\R)$ be a real-valued function such that $(\de_t-\laplacian)u=0$ in $U_T$.
	Then
	\begin{equation}\label{eq67c45475}
		\max_{\bar U_T} u = \max_{\Gamma_T} u .
	\end{equation}
\end{theorem}
\begin{proof}
	(From Folland \cite[Thm 4.16]{MR1357411}).	
	In~\eqref{eq67c45475}, the inequality $\max_{\bar U_T} u \ge \max_{\Gamma_T} u$ is clear.
	We need to show $\max_{\bar U_T} u \le \max_{\Gamma_T} u$.
	
	Let $\epsilon>0$ and define $v_\epsilon(x,t) = u(x,t) + \epsilon |x|^2$.
	Let $T'\in (0,T)$.
	We claim that
	\begin{equation}\label{eq67c4531d}
		\max_{\bar U_{T'}} v_\epsilon = \max_{\Gamma_{T'}} v_\epsilon .
	\end{equation}
	First of all, notice that, in $U_T$,
	\begin{equation}\label{eq67c452da}
	\de_tv_\epsilon - \laplacian v_\epsilon 
	= \de_t u - \laplacian u - \epsilon \laplacian(|x|^2) 
	= -2\epsilon n >0 .
	\end{equation}
	Suppose that $(x,t)\in U_{T'}$ is a point of maximum for $v_\epsilon$ on $\bar U_{T'}$.
	Then $\laplacian v_\epsilon(x,t) \le 0$ and $\de_t v_\epsilon \ge 0$.
	However, this is in contradiction with~\eqref{eq67c452da}.
	We conclude that~\eqref{eq67c4531d} must hold.
	
	To prove the statement about $u$, we see that
	\begin{align*}
	 	\max_{\bar U_{T'}} u 
		&\le \max_{\bar U_{T'}} v_\epsilon \\
		&\overset{\eqref{eq67c452da}}= \max_{\Gamma_{T'}} v_\epsilon \\
		&\le \max_{\Gamma_{T'}} u + \epsilon \max_{\bar U_{T'}} |x|^2 .
	\end{align*}
	Since $U$ is bounded, then $ \max_{\bar U_{T'}} |x|^2<\infty$.
	So, letting $\epsilon\to0$, we obtain the desired inequality.
\end{proof}

\begin{exercise}
	Show the following statement:
	{\it 
	Let $U\subset\R^n$ be open, bounded and connected, $T>0$,
	and $u\in C^{2;1}(U_T;\R)\cap C^0(\bar U_T;\R)$ such that $(\de_t-\laplacian)u=0$ in $U_T$.
	If there exists $(x,t)\in U_T$ such that $u(x,t) = \max_{\bar U_T} u$, then $u$ is constant on $\bar U_t$.}
	
	[This exercise turned out to be more difficult than I expected: It is proven by Evans using the mean value formula from Theorem~\ref{thm67cf158e}.]
\end{exercise}

\begin{exercise}
	Prove the \emph{strong minimum principle} for the heat operator.
\end{exercise}

\begin{exercise}
	Prove the strong maximum principle for the heat operator using the mean-value property; see \cite{MR2597943}.
\end{exercise}

%%%%%%%%%%%%%%%%%%%%%%%%%%%%%%%%%%%%%%%%%%%%%%%%%%%%%%%%%%%%%%%
\subsection{Uniqueness on bounded domains}\label{subs67c4558e}
\begin{theorem}[Uniqueness on bounded domains]\label{thm67c455a4}
	Suppose $U\subset\R^n$ is open and bounded, and $T>0$.
	Let $g\in C^0(\Gamma_T)$ and $f\in C(U_T)$.
	Then there exist not two distinct solutions in $C^{2;1}(U_T)\cap C(\bar U_T)$ to the boundary value problem
	\begin{equation}\label{eq67c456eb}
	\begin{cases}
		(\de_t-\laplacian)u = f & \text{ in }U_T , \\
		u = g & \text{ on }\Gamma_T .
	\end{cases}
	\end{equation}
\end{theorem}
\begin{proof}
	Let $u_1,u_2\in C^{2;1}(U_T)\cap C(\bar U_T)$ be two solutions to~\eqref{eq67c456eb}.
	Then $u=u_1-u_2 \in C^{2;1}(U_T)\cap C(\bar U_T)$ satisfies $(\de_t-\laplacian)u =0$ in $U_T$ and $u = 0$ on $\Gamma_T$.
	By the strong maximum and minimum principles, see~\ref{subs67c44a2b}, both the real and the imaginary parts of $u$ are zero on $U_T$, i.e., $u_1=u_2$.
\end{proof}

%%%%%%%%%%%%%%%%%%%%%%%%%%%%%%%%%%%%%%%%%%%%%%%%%%%%%%%%%%%%%%%
\subsection{Maximum and minimum principles for the unbounded Cauchy problem}\label{subs67c45d73}

\begin{theorem}[Maximum principle for the unbounded $U$]\label{thm67c4584f}
	Let $T>0$.
	Let $u\in C^{2;1}(\R^n\times(0,T];\R)\cap C^0(\R^n\times[0,T];\R)$ be a real-valued function such that $(\de_t-\laplacian)u=0$ in $\R^n\times(0,T]$.
	Suppose also that there are $A,a>0$ such that
	\begin{equation}\label{eq67c45e15}
		u(x,t) \le A\exp(a|x|^2) .
	\end{equation}
	Then
	\begin{equation}\label{eq67c4584a}
		\max_{\R^n\times[0,T]} u = \max_{\R^n} u(\cdot,0) .
	\end{equation}
\end{theorem}
\begin{proof}
	Let $0<b\le T$ be such that
	\begin{equation}
		\frac1{4b} - a > 0 ,
		\text{ i.e., }
		0<b<\frac1{4a} .
	\end{equation}
	We will show that
	\begin{equation}\label{eq67cff4fc}
		\forall s,t\in[0,T]\text{ with $0<t-s<b$, }
		\sup_{x\in\R^n} u(x,t) \le \sup_{x\in\R^n} u(x,s) .
	\end{equation}
	
	Since both~\eqref{eq67c45e15} and the heat equation are invariant under time translations (see~\ref{subs67d16430}), we can prove~\eqref{eq67cff4fc} for $s=0$ without loss of generality.
	
	Set $g\in C^0(\R^n)$ by $g(x) = u(x,0)$.
	Fix $\hat y\in\R^n$ and $\hat t\in(0,b)$: we need to show that
	\begin{equation}\label{eq67cff5ae}
		u(\hat y,\hat t) \le \|g\|_{L^\infty(\R^n)} .
	\end{equation}
	For $\mu>0$, define, for $x\in\R^n$ and $t\in(0,b)$,
	\begin{equation}\label{eq67c4590e}
		v(x,t) = u(x,t) - \frac{\mu}{(b-t)^{n/2}} \exp\left( \frac{|x-\hat y|^2}{4(b-t)} \right) .
	\end{equation}
	So, we have 
	$v\in C^{2;1}(\R^n\times(0,b);\R)\cap C^0(\R^n\times[0,b);\R)$
	and $(\de_t-\laplacian)v = 0$ in $\R^n\times(0,b)$;
	see~\ref{subs67c45976}.
	
	We claim that there exists $r>0$ such that
	\begin{equation}\label{eq67cff772}
		\sup_{B(\hat y,r)\times[0,b)}v  \le \|g\|_{L^\infty} .
	\end{equation}
	
	By the strong maximum principle, Theorem~\ref{thm67c44b0d},
	for every $r>0$, we have
	\begin{equation}
%		\sup\left\{ v(x,t) : 
%		x\in B(\hat y,r),\ t\in[0,b)
%		\right\}
		\sup_{B(\hat y,r)\times[0,b)}v
		=
		\sup\left\{ v(x,t) : 
		\begin{array}{l}
			x\in B(\hat y,r)\text{ and }t=0,\text{ or } \\
			x\in\de B(\hat y,r)\text{ and }t\in[0,b)
		\end{array}
		\right\} .
	\end{equation}
	If $|x-\hat y|\le r$, then 
	$v(x,0) = u(x,0) - \frac{\mu}{b^{n/2}} \exp\left(  \frac{|x-\hat y|^2}{4b} \right)
	\le g(x) \le \|g\|_{L^\infty}$.
	If $|x-\hat y|=r$ and $t\in[0,b)$, then 
	\begin{align*}
	 	v(x,t) 
		&= u(x,t) - \frac{\mu}{(b-t)^{n/2}} \exp\left( \frac{r^2}{4(b-t)} \right) \\
		&\le A\exp(a|x|^2) - \frac{\mu}{b^{n/2}} \exp\left(  \frac{r^2}{4b} \right) \\
		&\le A\exp(a(|\hat y|+r)^2) - \frac{\mu}{b^{n/2}} \exp\left(  \frac{r^2}{4b} \right) \\
		&= \exp(a(|\hat y|+r)^2) \left( A - \frac{\mu}{b^{n/2}} \exp\left(  \frac{r^2}{4b} - a (|\hat y|+r)^2 \right) \right) =: J_r .
	\end{align*}
	Since $\lim_{r\to\infty} J_r = -\infty$, 
	we can choose $r$ large enough that the $J_r < \|g\|_{L^\infty}$.
	We conclude that, for $r$ large enough,~\eqref{eq67cff772} holds.
	
	In particular, from~\eqref{eq67cff772} we obtain that, for every $\mu>0$,
%	We conclude that, for every $t\in[0,b)$,
	\begin{equation}\label{eq67cff890}
		u( \hat y,\hat t) - \frac{\mu}{(b-\hat t)^{n/2}}
		= v( \hat y,\hat t)
%		\le \sup\{ u(x,t) : x\in B(\hat y,r), t\in[0,b)\} 
%		\le J_r 
		\le \|g\|_{L^\infty} .
	\end{equation}
	Taking the limit $\mu\to0$ in~\eqref{eq67cff890}, we conclude~\eqref{eq67cff5ae}.
\end{proof}

\begin{exercise}\label{ex67c45e75}
	State and prove the strong minimum principle for the unbounded Cauchy problem.
\end{exercise}

%%%%%%%%%%%%%%%%%%%%%%%%%%%%%%%%%%%%%%%%%%%%%%%%%%%%%%%%%%%%%%%
\subsection{Uniqueness for the unbounded Cauchy problem}\label{subs67c45ca2}

\begin{theorem}[Uniqueness for the unbounded Cauchy problem]\label{thm67c45cb9}
	Let $T>0$, $g\in C^0(\R^n)$ and $f\in C^0(\R^n\times[0,T])$.
	Set 
	\begin{equation}\label{eq67c45e49}
		\scr A = \left\{
		u\in C^{2;1}(\R^n\times(0,T])\cap C^0(\R^n\times[0,T]) : 
		\begin{array}{c}
		\exists A,a>0 \text{ s.t.}\\
		|u(x,t)| \le A\exp(a|x|^2)
		\end{array}
		\right\} .
	\end{equation}
	There does not exist two solutions in $\scr A$ of
	\begin{equation}\label{eq67c45d58}
	\begin{cases}
		(\de_t-\laplacian)u = f & \text{ in }\R^n\times(0,T] , \\
		u = g & \text{ on }\R^n\times\{0\} .
	\end{cases}
	\end{equation}
\end{theorem}
\begin{proof}
	If $u_1,u_2\in\scr A$ solve~\eqref{eq67c45d58}, then $u=u_1-u_2\in\scr A$
	is such that $(\de_t-\laplacian)u = 0$ in $\R^n\times(0,T]$ and
	$u=0$ for $t=0$.
	By the maximum and minimum principles~\ref{subs67c45d73},
	$u\equiv 0$.
\end{proof}
\begin{remark}
	The growth condition~\eqref{eq67c45e15} is necessary: there are counter-examples.
	See \cite[chapter 7]{zbMATH03739988}.
\end{remark}
\begin{remark}
	Notice that in the definition of $\scr A$ in~\eqref{eq67c45e49}, 
	we require the growth condition~\eqref{eq67c45e15} both as an upper bound and as a lower bound.
	You should have noticed this already in Exercise~\ref{ex67c45e75}.
\end{remark}

%%%%%%%%%%%%%%%%%%%%%%%%%%%%%%%%%%%%%%%%%%%%%%%%%%%%%%%%%%%%%%%
\subsection{Smoothness}\label{subs67cc46da}
{\color{blue}

The following blue part contains mistakes, so it is dropped from the content of the course.
I keep it here for future memory: with a bit of work, we should be able to fix it.
Notice that we have other methods to prove that solutions to the heat equation are smooth...

\TODO: Fix proof of Lemma~\ref{lem67cc828c} .


\begin{lemma}\label{lem67cc828c}
	For $(x,t)\in\R^n\times\R$ and $r>0$, define
	\begin{equation}
		C(x,t;r) = \bar B(x,r)\times [t-r^2,t] \subset\R^n\times\R .
	\end{equation}
	Let $\Omega\subset\R^n\times\R$ open.
	Fix $(\hat x,\hat t)\in \Omega$ and $\hat r>0$ such that $C:= C(\hat x,\hat t;\hat r)\subset \Omega$.
	Define $C':= C(\hat x,\hat t;\dfrac{3}{4} \hat r)$
	and $C'' := C(\hat x,\hat t;\dfrac{1}{2} \hat r)$.

	Let $\zeta\in C^\infty_c(\R^n\times(-\infty,\hat t])$ be such that
	$C' \subset \{\zeta = 1\}$ and 
	$\spt(\zeta)\subset C$.
%	 i.e., $\R^n\setminus C\subset \{\zeta=0\}$.
	
	Then, for every $u\in C^{2;1}(\Omega)$ and all $(x,t)\in C''$,
	\begin{equation}
	\begin{aligned}
		u(x,t) 
		&= \int_0^t\int_{\R^n} \bigg( \Phi(x-y,t-s) (\de_s\zeta(y,s)+\laplacian_y\zeta(y,s) )
			+ 2\Diff_y \Phi(x-y,t-s) \Diff_y\zeta(y,s) \bigg) \cdot u(y,s) \did y\did s \\
		&\qquad+ \int_0^t\int_{\R^n} \Phi(x-y,t-s) \zeta(y,s) \bigg(\de_su(y,s) - \laplacian_y u(y,s)\bigg) \did y\did s .
	\end{aligned}
	\end{equation}
\end{lemma}
%%%%%%%%%%%%%%%%%%%%%%%%%%%%%%%%%%%%%%%%%%%%%%%%%%%%%%%%%%%%%%%
\begin{center}
\begin{tikzpicture}
	% Define points
	\coordinate (hatX) at (0,0);
	\coordinate (A1) at (-2,0);
	\coordinate (B1) at (-3,0);
	\coordinate (C1) at (-4,0);

	\coordinate (A2) at (-2,-1);
	\coordinate (B2) at (-3,-9/4);
	\coordinate (C2) at (-4,-4);	

	\coordinate (A3) at (2,-1);
	\coordinate (B3) at (3,-9/4);
	\coordinate (C3) at (4,-4);

	\coordinate (A4) at (2,0);
	\coordinate (B4) at (3,0);
	\coordinate (C4) at (4,0);
	
	\coordinate (F0) at (-5,-5);
	\coordinate (F1) at (5,-5);
	\coordinate (Fi) at (-5,1);

	% Draw triangle
	\draw (A1) -- (A2) -- (A3) -- (A4) ;
	\draw (B1) -- (B2) -- (B3) -- (B4) ;
	\draw (C1) -- (C2) -- (C3) -- (C4) ;
	\draw (C1) -- (C4) ;
	
	\node[above right] at (A2) {$C''$};
	\node[above right] at (B2) {$C'$};
	\node[above right] at (C2) {$C$};
	\filldraw (hatX) circle (2pt);
	\node[above] at (hatX) {$(\hat x,\hat t)$};
	\draw (Fi) -- (F0) node[midway, left] {$\R$};
	\draw (F0) -- (F1) node[midway, below] {$\R^n$};
\end{tikzpicture}
\end{center}
%%%%%%%%%%%%%%%%%%%%%%%%%%%%%%%%%%%%%%%%%%%%%%%%%%%%%%%%%%%%%%%
\begin{proof}
	Assume $\hat t>\hat r>0$.
	Define $v:=\zeta\cdot u\in C^{2;1}_c(\R^n\times(-\infty,\hat t])$ and compute
%	\begin{align*}
%	\de_tv &= \de_t\zeta\cdot y + \zeta\de_tu \\
%	\laplacian v &= \laplacian\zeta\cdot y + 2\Diff\zeta\cdot\Diff u + \zeta\cdot\laplacian u \\
%	(\de_t-\laplacian)v &= (\de_t-\laplacian)\zeta\cdot u + \zeta (\de_t-\laplacian)u - 2 \Diff\zeta\cdot\Diff u =: \tilde f .
%	\end{align*}
	\begin{equation}
		\tilde f := (\de_t-\laplacian)v
		= (\de_t-\laplacian)\zeta\cdot u + \zeta (\de_t-\laplacian)u - 2 \Diff\zeta\cdot\Diff u .
	\end{equation}
	Set then 
	\begin{equation}
		\tilde v(x,t) := \int_0^t \int_{\R^n} \Phi(x-y,t-s) \tilde f(y,s) \did y\did s .
	\end{equation}
	By the definition of $\tilde f$ and by~\eqref{eq67c34492}, 
	both $v$ and $\tilde v$ are solutions to the Cauchy problem 
	$(\de_t-\laplacian)v=\tilde f$ in $\R^n\times(0,\hat t)$,
	with $v=0$ on $\R^n\times\{0\}$.
	Moreover, we have both $\|v\|_{L^\infty}<\infty$ and $\|\tilde f\|_{L^\infty}<\infty$.
	Therefore, $|\tilde v(x,t)|\le t \|\tilde f\|_{L^\infty}$.
	From Theorem~\ref{thm67c45cb9}, it follows that $v=\tilde v$.
	
	Since $(x,t)\in C''$, then $\zeta(x,t)=1$ and
	\begin{align*}
	u(x,t) 
	&= \zeta(x,t) u(x,t)
	= v(x,t)
	= \tilde v(x,t) \\
	&= \int_0^t \int_{\R^n} \Phi(x-y,t-s) \bigg(
	(\de_s-\laplacian_y)\zeta(y,s)\cdot u(y,s) + \zeta(y,s)\cdot (\de_s-\laplacian_y)u(y,s) 
	\\&\qquad\qquad
	- 2 \Diff\zeta(y,s)\cdot\Diff u(y,s)
	\bigg) \did y\did s .
	\end{align*}
%	{\color{Red} WRONG:
%	Notice that the integral is in fact on $C$, because $\spt(\zeta)\subset C$. }
	Moreover,
	\begin{align*}
	\int_0^t\int_{\R^n} \Phi(x-y,t-s) \Diff\zeta\cdot\Diff u \did y\did s
%	&= \int_{\hat t-\hat r^2}^{\hat t} \int_{\bar B(\hat x,\hat r)}
	&= \int_0^t \int_{\bar B(\hat x,\hat r)}
		\bigg( \div_y( \Phi(x-y,t-s) \Diff\zeta(y,s) u(y,s) ) \\
	&\qquad		- \Diff_y\Phi(x-y,t-s) \Diff\zeta(y,s) u(y,s) \\
	&\qquad\qquad		- \Phi(x-y,t-s) u(y,s) \laplacian\zeta(y,s) 
		\bigg) \did y\did s ,\\
	\end{align*}
	where, for all $s<t$,
	\begin{align*}
	 	\int_{\bar B(\hat x,\hat r)} \div_y( \Phi(x-y,t-s) \Diff\zeta(y,s) u(y,s) \did y 
%		&= \int_{\bar B(\hat x,\hat r) \setminus B(\hat x,\epsilon)} \div_y( \Phi(x-y,t-s) \Diff\zeta(y,s) u(y,s) \did y \\
		&= 0 .
	\end{align*}
\end{proof}

%%%%%%%%%%%%%%%%%%%%%%%%%%%%%%%%%%%%%%%%%%%%%%%%%%%%%%%%%%%%%%%
\begin{theorem}[Smoothness of solutions to the homogeneous heat equation]\label{thm67cc8194}
	Let $\Omega\subset\R^n\times\R$ be open and $u\in C^{2;1}(\Omega)$ such that
	$(\de_t-\laplacian)u=0$ in $\Omega$.
	Then $u\in C^\infty(\Omega)$.
\end{theorem}
\begin{proof}
	By the Lemma~\ref{lem67cc828c}, we have for $(x,t)\in C''$
	\begin{equation}
		u(x,t) = \iint_C K(x,t;y,s) u(y,s) \did y\did s ,
	\end{equation}
	where, for $(x,t)\in C''$, $K(x,t;\cdot)\in C^\infty(\bar C)$.
	{\color{Red} TRUE? Therefore, $u$ is smooth on $C''$ by Theorem~\ref{thm679f29de}, or Proposition~\ref{prop6798f690}. -> $C''$ is not open!}
\end{proof}

%%%%%%%%%%%%%%%%%%%%%%%%%%%%%%%%%%%%%%%%%%%%%%%%%%%%%%%%%%%%%%%
\begin{theorem}[Quantitative smoothness]\label{thm67cc82d1}
	Let $\Omega\subset\R^n\times\R$ be an open set.
	For every $\alpha\in\N^n$, $\ell\in\N$, there exists $C_{\alpha,\ell}\in\R$
	such that, for every $u\in C^\infty(\Omega)$ with $(\de_t-\laplacian)u=0$ in $\Omega$,
	if $C(x,t;r)\subset \Omega$, then
	\begin{equation}\label{eq67cc83af}
		\sup_{C(x,t;r/2)} \left| \frac{\de^{|\alpha|}}{\de x^\alpha} \frac{\de^\ell}{\de t^\ell} u \right|
		\le 
		\frac{C_{\alpha,\ell}}{r^{|\alpha|+2\ell+n+2}} \|u\|_{L^1(C(x,t;r))} .
	\end{equation}
\end{theorem}
\begin{proof}
	From Lemma~\ref{lem67cc828c} we get that, if $u\in C^\infty(C(0,0;1))$, then, 
	for every $(x,t)\in C(0,0;1/2)$, 
	\begin{equation}
		u(x,t) = \iint_{C(0,0;1)} K(x,t;y,s) u(y,s) \did y\did s ,
	\end{equation}
	where $K\in C^\infty(\overline{C(0,0;1/2)\times C(0,0;1)})$.
	By Proposition~\ref{prop6798f690}, or Theorem~\ref{thm679f29de},
	for every $(x,t)\in C(0,0;1/2)$ we have
	\begin{equation}
		\Diff^{\alpha,\ell}_{x,t}u(x,t) = \iint_{C(0,0;1)} \Diff^{\alpha,\ell}_{x,t}K(x,t;y,s) u(y,s) \did y\did s .
	\end{equation}
	Hence, if we take as constants
	\begin{equation}
		C_{\alpha,\ell} 
		= \sup\{ |\Diff^{\alpha,\ell}_{x,t}K(x,t;y,s)| : (x,t)\in C(0,0;1/2),\ (y,s)\in C(0,0;1) \},
	\end{equation}
	we obtain~\eqref{eq67cc83af} in this specific case.
	
	In the general case, if $(\hat x,\hat t)\in \Omega$ and $\hat r>0$ are such that $C(\hat x,\hat t;\hat r)\subset\Omega$, and if $u\in C^\infty(\Omega)$ is such that $(\de_t-\laplacian)u=0$ in $\Omega$,
	then the function
	\begin{equation}
		\hat u(x,t) := u(\hat x+\hat r x, \hat t + \hat r^2 t) 
	\end{equation}
	is such that $(\de_t-\laplacian)u=0$ in $C(0,0;1)$.
	Therefore, for every $(x,t) \in C(0,0;1/2)$,
	\begin{align*}
	 	|\Diff^{\alpha,\ell}_{x,t}u(\hat x+\hat r x,\hat t + \hat r^2 t )| \hat r^{|\alpha| + 2\ell}
		&= |\Diff^{\alpha,\ell}_{x,t} \hat u(x,t)| \\
		&\le C_{\alpha,\ell} \iint_{C(0,0;1)} \hat u(y,s) \did y\did s \\
		&= C_{\alpha,\ell} \iint_{C(0,0;1)}  u(\hat x+\hat r y,\hat t + \hat r^2 s) \did y\did s \\
		&= C_{\alpha,\ell} \iint_{C(\hat x,\hat t;\hat r)} u(\bar y,\bar s) \frac{\did \bar y}{\hat r^n} \frac{\did \bar s}{\hat r^2} .
	\end{align*}
	Hence, we conclude with~\eqref{eq67cc83af}.
\end{proof}

\begin{remark}
	Notice that this smoothness result tells us that time cannot be inverted!
\end{remark}
}

%%%%%%%%%%%%%%%%%%%%%%%%%%%%%%%%%%%%%%%%%%%%%%%%%%%%%%%%%%%%%%%
\subsection{Energy Methods}
\begin{lemma}\label{lem67cc8a6c}
	Let $U\subset\R^n$ open and bounded with $C^1$ boundary $\de U$, and let $T>0$.
	For $u:U_T\to\C$ define $e_u:[0,T]\to\R$,
	\begin{equation}
		e_u(t) = \int_U |u(x,t)|^2 \did x ,
	\end{equation}
	whenever it is well defined.
	If $u\in C^{2;1}(\bar U_T)$ is such that $(\de_t-\laplacian)u=0$ in $U_T$, then $e_u\in C^1([0,T])$ and
	\begin{equation}
		\dot e_u(t) = 2 \int_{\de U} u \Diff u\cdot\nu_U \did S 
			- 2 \int_U |\Diff u|^2 \did x .
	\end{equation}
	If $u=0$ on $\de U\times [0,T]$, then $\dot e_u(t) \le 0$ and so, $e_t$ is decreasing.
\end{lemma}
\begin{proof}
	Since $U$ is bounded and $u$ is continuous on $\bar U$, then $\|u\|_{L^\infty(\bar U_T)} < \infty$.
	Since $U$ is bounded, then constants belong to $L^2(U)$.
	Therefore, by Theorem~\ref{thm679f29de}, $e_u\in C^1((0,T))$ and
	\begin{align*}
	 	\dot e_u(t) 
		&= 2\int_U \de_t u(x,t) u(x,t) \did x \\
		&= 2 \int_U \laplacian u(x,t) u(x,t) \did x \\
		&= 2 \int_{\de U} \Diff u(x,t) \cdot \nu_U(x) \did S(x)
			- 2 \int_U |\Diff u(x,t)|^2 \did x .
	\end{align*} 
\end{proof}

%%%%%%%%%%%%%%%%%%%%%%%%%%%%%%%%%%%%%%%%%%%%%%%%%%%%%%%%%%%%%%%
\begin{theorem}[Uniqueness by Energy methods]\label{thm67cc8975}
	Let $T>0$ and $U\subset\R^n$ be an open, bounded subset with $C^1$ boundary $\de U$.
	Let $f\in C^0(U_T)$ and $g\in C^0(\Gamma_T)$.
	Then there do not exist two distinct solutions in $C^{2;1}(\bar U_T)$ to
	\begin{equation}\label{eq67cc89f9}
	\begin{cases}
		(\de_t-\laplacian)u = f & \text{ in }U_T , \\
		u = g & \text{ on }\Gamma_T .
	\end{cases}
	\end{equation}
\end{theorem}
\begin{proof}
	Suppose there are two solutions $u_1,u_2\in C^{2;1}(\bar U_T)$ to~\eqref{eq67cc89f9}.
	Then $w=u_1-u_2 \in C^{2;1}(\bar U_T)$ solves~\eqref{eq67cc89f9} with $f=0$ and $g=0$.
%	Therefore, $w\in C^\infty(U_T)$ by Theorem~\ref{thm67cc8194} 
%	and, 
	By Lemma~\ref{lem67cc8a6c}, $0 = e_w(0) \ge e_w(t) = \int_U |w(x,t)|^2 \did x \ge 0$ for all $t>0$.
	It follows that $w=0$ and thus $u_1=u_2$.
\end{proof}


%%%%%%%%%%%%%%%%%%%%%%%%%%%%%%%%%%%%%%%%%%%%%%%%%%%%%%%%%%%%%%%
\subsection{Backward uniqueness}\label{subs67cc8a9b}

\begin{theorem}[Backward uniqueness]
	Let $T>0$ and $U\subset\R^n$ be an open, bounded subset with $C^1$ boundary $\de U$.
	Let $g\in C^0(\de U\times[0,T])$.
	Suppose that $u_1,u_2\in C^2(\bar U_T)$ solve the Cauchy problem
	\begin{equation}\label{eq67cc8aec}
	\begin{cases}
		(\de_t-\laplacian)u = 0 & \text{ in }U_T , \\
		u = g & \text{ on }\de U\times[0,T] .
	\end{cases}
	\end{equation}
	If $u_1(x,T) = u_2(x,T)$ for all $x\in U$, then $u_1=u_2$ in $U_T$.
\end{theorem}
\begin{proof}
	Take $w=u_1-u_2 \in C^2(\bar U_T)$, which solves~\eqref{eq67cc8aec} with $g=0$.
	By the Lemma~\ref{lem67cc8a6c}, 
	\begin{align*}
		e_w(t) &= \int_U |w(x,t)|^2 \did x ; \\
		\dot e_w(t) &= 2 \int_U |\Diff w(x,t)|^2 \did x ; \\
		\ddot e_w(t) &= [...] = 4 \int_U | \laplacian w(x,t)|^2 \did x .
	\end{align*}
	Moreover, using the Hölder inequality,
	\begin{align*}
		\int_U | \Diff w(x,t) |^2 \did x
		&= \int_U \Diff w(x,t)\cdot \Diff w(x,t) \did x \\
		&= - \int_U \laplacian w(x,t) \cdot w(x,t) \did x \\
		&\le \left(\int_U |\laplacian w(x,t)|^2 \did x\right)^{1/2}
			\cdot \left(\int_U |w(x,t)| \did x\right)^{1/2} .
	\end{align*}
	Hence, 
	\begin{equation}
		\dot e_w(t)^2 \le \ddot e_w(t) \cdot e_w(t) .
	\end{equation}
	Suppose that $e_w$ is not zero on $[0,T]$.
	Let $(a,b)\subset[0,T]$ be a maximal interval where $e_w>0$.
	Then $e_w(a) = 0$.
	By Lemma~\ref{lem67cc8da8}, for all $[t_1,t_2]\subset(a,b)$,
	\begin{equation}
		e_w(1/2 t_1 + 1/2 t_2 ) 
		\le e_w(t_1)^{1/2} e_w(t_2)^{1/2} .
	\end{equation}
	Taking the limit to $t_1\to a$ and $t_2\to b$ we obtain $e_w(\dfrac{a+b}{2}) = 0$, in contradiction with $e_w>0$ on $(a,b)$.
	We conclude that $e_w=0$ on $[0,T]$.
\end{proof}

\begin{lemma}\label{lem67cc8da8}
	Let $I\subset\R$ be an interval and $e:I\to (0,+\infty)$ be a $C^2$ function with 
	$\dot e \le \dd e \dot e$.
	Then, for every $t_1<t_2$ belonging to $I$, and every $\tau\in(0,1)$,
	\begin{equation}\label{eq67cc8d9e}
		e((1-\tau)t_1 + \tau t_2) \le e(t_1)^{1-\tau} e(t_2)^\tau .
	\end{equation}
\end{lemma}
\begin{proof}
	Set $f(t) = \log(e(t))$, which is well defined because $e(t)>0$.
	Then $f'(t) = \frac{\dot e(t)}{e(t)}$ and 
	\begin{equation}
		f''(t) = - \frac1{e(t)^2} \dot e(t)^2 + \frac{\ddot e(t)}{e(t)} 
	\ge \frac{ - \ddot e e + \ddot e e }{e^2} = 0 .
	\end{equation}
	Therefore $f$ is convex.
	From the convexity of $f$, we get~\eqref{eq67cc8d9e}.
\end{proof}













%%%%%%%%%%%%%%%%%%%%%%%%%%%%%%%%%%%%%%%%%%%%%%%%%%%%%%%%%%%%%%%
\newpage
\section{Wave Equation}

\subsection{The wave operator}\label{subs67cd41c5}
\newcommand{\waveop}{\square}
\index{wave equation}
\index{equation!wave --}
\index{wave opeartor}

Let $U\subset\R^n$ open and $I\subset\R$ an open interval.
For $u:U\times I\to \C$, we call
\begin{itemize}
\item the (homogeneous) wave equation $\frac{\de^2u}{\de t^2} - \laplacian u = 0$ in $U\times I$;
\item the \emph{nonhomogeneous wave equation} $\frac{\de^2u}{\de t^2} - \laplacian u = f$ in $U\times I$, for some $f:U\times I\to\C$.
\end{itemize}
A common abbreviation is to denote the \emph{wave operator} by
\begin{equation}
	\waveop w = \de_t^2 - \laplacian .
\end{equation}

%%%%%%%%%%%%%%%%%%%%%%%%%%%%%%%%%%%%%%%%%%%%%%%%%%%%%%%%%%%%%%%
\subsection{Symmetries of $\waveop$}\label{subs67cd41e1}
The wave operator $\waveop$ is a linear homogeneous hyperbolic differential operator of order 2.

Define $v(x,t) = u(Ax+a, Bt+b)$ with $A\in \GL(n)$, $a\in \R^n$, $B,b\in\R$
with $AA^T = \lambda^2\Id$.
Then
\begin{align*}
	\waveop v 
	&= (B^2\de_t^2u - \lambda^2 \laplacian u)(Ax+a,Bt+b) \\
	\text{[If $B^2=\lambda^2$:]}
	&= \lambda^2 \waveop u(Ax+a,Bt+b) .
\end{align*}

\begin{exercise}
	Compute $\waveop u$ for $u(x,t) = \exp((a+bi)\cdot x+\phi t)$, where $a,b\in\R^n$ and $\phi\in\R$.
\end{exercise}

\begin{exercise}
	Show that $\waveop$ is invariant under the Lorentz group of transformations.
	Recall that the Lorentz group is the group of linear automorphisms of the Minkowski space.
	More explicitely, if $M$ is the $(n+1)\times(n+1)$ matrix
	\begin{equation}
		M = \begin{pmatrix}\Id & 0 \\ 0 & -1 \end{pmatrix} ,
	\end{equation}
	then the Minkowski space is $(\R^n\times\R,M)$
	and the Lorentz group is made of matrices $A\in \GL(\R^n\times\R)$ such that $A^TMA = M$.
	
	{\it Hint.}
	If it seems too hard, try to solve the exercise at leas for $n=1$.
\end{exercise}

%%%%%%%%%%%%%%%%%%%%%%%%%%%%%%%%%%%%%%%%%%%%%%%%%%%%%%%%%%%%%%%
\subsection{Examples}
For $a,b\in\C$,
consider the following functions $u_{a,b}:\R\times\R\to\C$,
\begin{equation}
	u_{a,b}(x,t) = \exp(ax + bt) ,
\end{equation}
\begin{exercise}
	For which $a,b\in\C$ we have $\waveop u_{a,b}=0$?
	
	{\it Hint:} $a=b$ or $a=-b$.
\end{exercise}
\begin{exercise}
	For which $a,b\in\C$ we have $u_{a,b}$ 1-periodic (i.e., $u_{a,b}(0) = u_{a,b}(1)$)
	and $\waveop u_{a,b}=0$?
	
	{\it Hint:} $a,b\in 2\pi i\Z$ with $a=b$ or $a=-b$.
\end{exercise}
\begin{exercise}
	For every $k\in\N$ and $\ell>0$, find $u:[0,\ell]\times\R\to\R$ such that 
	\begin{enumerate}
	\item	$\waveop u = 0$,
	\item	$u(0,t) = u(\ell,t) = 0$ for all $t$, and
	\item	there are $0\le x_0 < x_1 < \dots < x_k \le \ell$ such that $u(x_j,t) = 0$ for all $j$ and all $t$.
	\end{enumerate}	
	These functions are the Harmonics of the string pinched at the two ends.
\end{exercise}

\begin{exercise}
	Let $c\in\C$.
	Find $u:\R^n\times\R\to\C$ such that
	\begin{equation}
		\waveop u = c u .
	\end{equation}
%
%	{\it Hint.}
%	First take $n=1$ and look at the functions $u_{a,b}$ we saw above.
\end{exercise}

%%%%%%%%%%%%%%%%%%%%%%%%%%%%%%%%%%%%%%%%%%%%%%%%%%%%%%%%%%%%%%%
\subsection{Finite propagation speed}\label{subs67ce9a17}
\begin{theorem}[Finite propagation speed]\label{thm67ce9a24}
	Let $U\subset \R^n$ be open, $\hat x\in U$ and $R>0$.
	Define
	\begin{equation}
		K(\hat x,R) = \{ (x,t)\in\R^n\times [0,R] : |x|\le R \} .
	\end{equation}
	Suppose $\bar B(\hat x,R)\subset U$ and that $u\in C^\infty(U\times[0,R])$ is such that
	$\waveop u = 0$ in $U\times(0,R)$,
	and $u=\de_tu =0 $ on $U\times\{0\}$.
	Then $u= 0$ in $K(\hat x,R)$.
\end{theorem}
\begin{proof}
	Fix $t\ge 0$, define 
	\begin{equation}
		E(t) = \frac12 \int_{B(\hat x,R-t)} (|\de_tu(x,t)|^2 + |\Diff u(x,t)|^2 ) \did x .
	\end{equation}
	Then, if $t\in (0,R)$,
	\begin{align*}
	 	\de_t E(t) 
		&= \frac12 \left.\frac{\dd}{\dd h}\right|_{h=0}
			\bigg( \int_{B(\hat x,R-t-h)} (|\de_tu(x,t)|^2 + |\Diff_xu(x,t)|^2 ) \did x
		\\&\qquad
				+ \int_{B(\hat x,R-t)} (|\de_tu(x,t+h)|^2 + |\Diff_xu(x,t+h)|^2 ) \did x 
			\bigg) \\
		&= \frac12 \bigg( - \int_{\de B(\hat x,R-t)} (|\de_tu(x,t)|^2 + |\Diff_xu(x,t)|^2 ) \did S(x)
		\\&\qquad
			+ \int_{B(\hat x, R-t)} (2\de_tu\de_t^2u + 2\Diff u \Diff\de_t u) \did x 
			\bigg) \\
		&= \int_{\de B(\hat x,R-t)} \bigg( - \frac{|\de_tu|^2}{2} - \frac{|\Diff u|^2}{2} + \de_t u \Diff u\cdot \nu \bigg) \did S(x) \\
		&\overset{(*)}\le  -\int_{\de B(x,R-t)} (|\de_tu|-|\Diff u|)^2 \did S
		\le 0 ,
	\end{align*}
	where we have used in $(*)$ that
	\begin{equation}
		(|\de_tu|-|\Diff u|)^2
		= |\de_tu|^2 + |\Diff u|^2 - 2 |\de_tu|\cdot|\Diff u|
		\le |\de_tu|^2 + |\Diff u|^2 - 2 |\de_tu|\cdot|\Diff u\cdot \nu|,
	\end{equation}
	because $|\nu|=1$.
	It follows that $0\le E(t) \le E(0) = 0$.
	So, $u$ is constant in $K(\hat x,R)$.
\end{proof}
\begin{remark}
	Notice that we actually only need $u$ constant in $B(\hat x,R)$ at time $0$ to obtain that $u$ is constant in $K(\hat x,R)$.
\end{remark}

\begin{exercise}
	Let $\hat x\in\R^n$ and $R>0$.
	Prove the following uniqueness result: if $u_1,u_2\in C^2(K(\hat x,R))$ are such that $\waveop u_1 = \waveop u_2$ in $K(\hat x,R)$ and $u_1=u_2$ on $K(\hat x,R) \cap \R^n\times\{0\}$, then $u_1 = u_2$.
\end{exercise}


%%%%%%%%%%%%%%%%%%%%%%%%%%%%%%%%%%%%%%%%%%%%%%%%%%%%%%%%%%%%%%%
\subsection{Uniqueness of solution to the wave equation}\label{subs67d92277}

\begin{theorem}[Uniqueness of solution to the wave equation]\label{thm67da8593}
	Let $U\subset\R^n$ be an open set with $C^1$ boundary $\de U$, and fix $T>0$.
	Let $f\in C^0(U_T)$, $g\in C^0(\Gamma_T)$, and $h\in C^0(U)$.
	Then there is at most one solution in $C^2(\bar U_T)$ to 
	\begin{equation}\label{eq67d923d6}
	\begin{cases}
		\waveop u = f & \text{ in } U_T , \\
		u = g & \text{ on }\Gamma_T , \\
		\de_t u = h & \text{ on } U\times\{0\} .
	\end{cases}
	\end{equation}
\end{theorem}
\begin{proof}
	As usual, if $u_1,u_2\in C^2(\bar U \times[0,T])$ solve~\eqref{eq67d923d6}, 
	then their difference $w = u_2-u_1$ solve the same problem with $f=0$, $g=0$, and $h=0$.
	Define
	\begin{equation}
		E(t) = \int_U (\de_tw(x,t)^2 + |\grad w(x,t)|^2) \did x .
	\end{equation}
	Then
	\begin{align*}
	 	E'(t)
		&= \int_U (2\de_tw \de_t^2w + 2 \grad w\grad\de_tw) \did x \\
		&= 2 \int_U\de_t w \de_t^2w + \int_{\de U} \de_t w \grad w\cdot \nu_U \did S(x) 
			- \int_U \laplacian w \de_tw \did x \\
		&= 0 ,
	\end{align*}
	because $\de_tw = 0 $on $\de U$ for all $t$ (because $w=0$ on $\de U\times [0,T]$),
	and because $\de_t^2 w = \laplacian w $ on $U\times(0,T)$.
	Therefore, $E$ is constant.
	Since $E(0)=0$, then $E\equiv 0$.
	It follows that $w$ is constant, and thus $0$ since $w=0$ on $U\times\{0\}$.
\end{proof}



%%%%%%%%%%%%%%%%%%%%%%%%%%%%%%%%%%%%%%%%%%%%%%%%%%%%%%%%%%%%%%%
\subsection{Solution  by spherical means: case $n=1$, d'Alambert's formula}\label{subs67cd4319}
We consider the case $n=1$.
We will solve the PDE
\begin{equation}\label{eq67cd4386}
\begin{cases}
	\waveop u = 0 & \text{ in }\R\times(0,+\infty) , \\
	u = g,\ \de_tu = h & \text{ on }\R\times\{0\} ,
\end{cases}
\end{equation}
with $g$ and $h$ given functions $\R\to\C$.
Before we give the full statement that we can prove at the moment, we will see how to find such a formula.
So, we forget about the regularity of $u$, or, in other words, we assume that $u$ is $C^\infty$ (or $C^2$, which is enough to justify our reasoning).

First, we notice that we can rewrite~\eqref{eq67cd4386} as
\begin{equation}
	(\de_t+\de_x)(\de_t-\de_x) u = 0 .
\end{equation}
Therefore, if $v=\de_tu - \de_x u$, then $v$ solves the transport equation~\eqref{eq67a79b9b} (with $b=1$):
\begin{equation}
\begin{cases}
	\de_t v + \de_x v = 0 & \text{ in }\R\times(0,+\infty) , \\
	v = h-g' & \text{ on }\R\times\{0\} .
\end{cases}
\end{equation}
If $u$ is $C^2$, then $v$ is $C^1$ and we know that
\begin{equation}
	v(x,t) = v(x-t,0) = h(x-t) - g'(x-t) .
\end{equation}
Next, $u$ solves $\de_tu-\de_xu = v$ in $\R\times(0,+\infty)$.
By the solution to the nonhomogeneous transport equation given by Theorem~\ref{thm67d7b17c} (with $b=-1$ and $f=v$), we have
\begin{align*}
 	u(x,t)
	&= u(x+t,0) + \int_0^t v(x-(r-t),r) \did r \\
	&= g(x+t) + \int_0^t (h(x-r+t-r) - g'(x-r+t-r)) \did r \\
	&= g(x+t + \left. \frac{g(x+t-2r)}{2} \right|^t_0 + \int_0^t h(x+t-2r) \did r \\
	\text{[$\bar r = x+t-2r$]}
	&= \frac12 (g(x+t)+g(x-t)) - \int_{x+t}^{x-t} h(\bar r) \frac{\did \bar r}{2} \\
	&= \frac12 (g(x+t)+g(x-t)) + \frac12 \int_{x-t}^{x+t} h(r) \did r .
\end{align*}
We thus obtain \emph{D'Alambert's formula}
\begin{equation}\label{eq67cd45ee}
	u(x,t) = \frac12 (g(x+t)+g(x-t)) + \frac12 \int_{x-t}^{x+t} h(r) \did r .
\end{equation}
\index{D'Alambert's formula}

\begin{theorem}[Solution to $\waveop=0$ for $n=1$]\label{thm67cd4667}
	Let $g\in C^2(\R)$ and $h\in C^1(\R)$.
	Define $u:\R\times\R\to \C$ by D'Alambert's formula~\eqref{eq67cd45ee}.
	Then
	\begin{enumerate}
	\item
	$u\in C^2(\R\times\R)$;
	\item
	$\de_t^2u-\de_x^2u = \waveop u=0$ in $\R\times\R$;
	\item
	for every $\hat x\in\R$, $u(\hat x,0) = g(\hat x)$ and $\de_tu(\hat x,0) = h(\hat x)$.
	\end{enumerate}	
	Moreover, $u$ is the only solution to~\eqref{eq67cd4386}.
\end{theorem}
\begin{proof}
	The proof is left as an exercise.
	Uniqueness is given by Theorem~\ref{thm67da8593}.
\end{proof}

\begin{exercise}
	Prove Theorem~\ref{thm67cd4667}.
\end{exercise}

\begin{exercise}
	Try to solve the nonhomogeneous version of~\eqref{eq67cd4386}, that is,
	\begin{equation}\label{eq67d7bc33}
	\begin{cases}
		\waveop u = f & \text{ in }\R\times(0,+\infty) , \\
		u = g,\ \de_tu = h & \text{ on }\R\times\{0\} ,
	\end{cases}
	\end{equation}
	for some given $f$.
\end{exercise}

\begin{remark}\label{rmk67da8652}
	Solutions to the wave equation given by D'Alambert's formula~\eqref{eq67cd45ee} are of the form
	\begin{equation}
		u(x,t) = F(x+t) + G(x-t) ,
	\end{equation}
	with $F,G\in C^2(\R)$.
	The two functions represent a forward-moving wave %($G(x-t)$) 
	and a backward-moving wave. % ($F(x+t)$)
	Can you say which is which?
\end{remark}

\begin{remark}
	The solutions we have found with D'Alambert's formula~\eqref{eq67cd45ee} extend to negative time!
\end{remark}

\begin{remark}
	The solutions we have found with D'Alambert's formula~\eqref{eq67cd45ee} are not $C^\infty$ smooth if $g$ and $h$ are not $C^\infty$.
	This is different from the other equations we have studied so far.
	Can you spot the difference in the equation? 
	In fact, the highest order part in the Laplace and in the heat equations is the laplacian, while here the whole $\waveop$ is homogeneous of order 2.
\end{remark}

\begin{remark}
	For $n>1$, we can define $\tilde u(x,t) = u(x_1,t)$, where $u$ is a solution in dimension $1$ to the wave equation.
	It follows that $\tilde u$ is a solution to the wave equation in $\R^n\times\R$.
	Therefore, we have found non-smooth solutions to $\waveop u=0$ in all dimensions!
\end{remark}

%%%%%%%%%%%%%%%%%%%%%%%%%%%%%%%%%%%%%%%%%%%%%%%%%%%%%%%%%%%%%%%
\subsection{A Reflexion method: solution on the half-line}\label{subs67cd4e88}

We want to solve the PDE
\begin{equation}\label{eq67d7b898}
	\begin{cases}
		\waveop u =0 & \text{ in }\R_+\times(0,+\infty) , \\
		u(x,0)=g(x),\ \de_tu(x,0) = h(x) & \forall x\in\R_+ , \\
		u(0,t) = 0 & \forall t>0 .
	\end{cases}
\end{equation}
This PDE represent a vibrating string that is pinched at one end and infinite in the other direction.
A solution is given in the following Theorem~\ref{thm67d7b770},
whose proof is a direct calculation.
The formula~\eqref{eq67d7b859} is obtained by a \emph{reflection method},
that is, we extend the problem~\eqref{eq67d7b898} to a PDE on the whole line $\R$ by taking
\begin{align*}
	\tilde u(x,t) &= 
	\begin{cases}
		u(x,t) & \text{ if } x\ge 0 ,\\
		-u(-x,t) & \text{ if } x\le 0 ;
	\end{cases} \\
	\tilde g(x) &= 
	\begin{cases}
		g(x) & \text{ if } x\ge 0 ,\\
		-g(-x) & \text{ if } x\le 0 ;
	\end{cases} \\
	\tilde h(x) &= 
	\begin{cases}
		h(x) & \text{ if } x\ge 0 ,\\
		-h(-x) & \text{ if } x\le 0 .
	\end{cases} 
\end{align*}
Then one can convince themselves that $\tilde u$ solves the 1D wave equation and thus, we can take $\tilde u$ as given from the D'Alambert's formula~\eqref{eq67cd45ee}.
From there, we can obtain~\eqref{eq67d7b859}.

\begin{theorem}\label{thm67d7b770}
	Let $\R_+ = (0,+\infty)$.
	Let $g\in C^2(\bar \R_+)$ and $h\in C^1(\bar \R_+)$ be such that $g(0)=h(0) = g''(0)= 0$.
	The function $u:\R_+ \times[0,+\infty)\to\C$,
	\begin{equation}\label{eq67d7b859}
		u(x,t) = 
		\begin{cases}
			\frac12 (g(x+t)+g(x-t)) + \frac12 \int_{x-t}^{x+t} h(y)\did y
				& \text{ if }0\le t\le x ,\\
			\frac12 (g(x+t)-g(x-t)) + \frac12 \int_{-x+t}^{x+t} h(y)\did y
				& \text{ if }0\le x\le t ,\\
		\end{cases}
	\end{equation}
	belongs to $C^2(\R_+\times(0,+\infty))\cap C^0(\bar\R_+\times[0,+\infty))$ and
	it is THE solution to 
	\begin{equation}\label{eq67da86cc}
	\begin{cases}
		\waveop u =0 & \text{ in }\R_+\times(0,+\infty) , \\
		u(x,0)=g(x),\ \de_tu(x,0) = h(x) & \forall x\in\R_+ , \\
		u(0,t) = 0 & \forall t>0 .
	\end{cases}
	\end{equation}
\end{theorem}
\begin{proof}
	Left as an exercise. Uniqueness is given by Theorem~\ref{thm67da8593}.
	
	{\it Hint.} The assumption $g''(0)=0$ implies that $\tilde g$ is $C^2$.
\end{proof}

\begin{exercise}
	Prove Theorem~\ref{thm67d7b770}.
\end{exercise}

\begin{exercise}
	For $F,G\in C^2(\R)$, define
	$\tilde u(x,t) = F(x+t) + G(x-t)$. 
	Then we know that $\waveop\tilde u$, see Remark~\ref{rmk67da8652}.
	For which $F$ and $G$ we have $\tilde u(0,t) = 0$ for all $t$?
	Solve~\eqref{eq67da86cc} finding the correct $F$ and $G$.
\end{exercise}

\begin{exercise}
	Solve the string problem
	\begin{equation}\label{eq67da8881}
	\begin{cases}
		\waveop u =0 & \text{ in }[0,1]\times(0,+\infty) , \\
		u(x,0)=g(x),\ \de_tu(x,0) = h(x) & \forall x\in[0,1] , \\
		u(0,t) = u(1,t) = 0 & \forall t>0 .
	\end{cases}
	\end{equation}
	You need to find some conditions on $g$ and $h$ to make $u$ of class $C^2$: it is part of the exercise.
	
	Next, for every $n\in\N$, find $u_n\in C^2([0,1]\times\R)$ such that $\waveop u_n=0$ and $u_n(k/n,t)=0$ for all $k\in\{1,\dots,n\}$.
	These functions $u_n$ are called \emph{harmonics} of the string.
\end{exercise}

%%%%%%%%%%%%%%%%%%%%%%%%%%%%%%%%%%%%%%%%%%%%%%%%%%%%%%%%%%%%%%%
%\subsection{\TODO}
%\begin{itemize}
%%\item A reflection method
%%\item Vibrating string pinched in two points, eigenvalues, harmonics...
%\item Change of coordinates
%\end{itemize}

%%%%%%%%%%%%%%%%%%%%%%%%%%%%%%%%%%%%%%%%%%%%%%%%%%%%%%%%%%%%%%%
\subsection{Spherical means: Euler--Poisson--Darboux equation}\label{subs67cd4983}
\index{Euler--Poisson--Darboux equation}

\begin{lemma}[Euler--Poisson--Darboux equation]\label{lem67cd498e}
	Let $n\ge2$, $m\ge 2$, $u\in C^m(\R^n\times[0,+\infty))$, $g,h\in C^m(\R^n)$.
	Suppose
	\begin{equation}\label{eq67cd4bec}
	\begin{cases}
		\de_t^2u-\laplacian u = 0 & \text{ in }\R^n\times(0,+\infty) , \\
		u = g,\ \de_t u = h & \text{ on }\R^n\times\{0\} . \\
	\end{cases}
	\end{equation}
	Referring to~\ref{sec67bac9ff},
	set $U(x;r,t) = \psi_u(t,x;r)$, $G(x;r):= \psi_g(x,r)$, $H(x;r) = \psi_h(x;r)$.
	Here $\psi_{(\cdot)}$ is defined in~\eqref{eq67bacbb5}.
	Then $U(x;\cdot)\in C^m([0,+\infty)\times[0,+\infty))$ and
	\begin{equation}\label{eq67cd4b9e}
	\begin{cases}
		\de_t^2 U - \de_r^2 U - \frac{n-1}{r} \de_r U = 0 & \text{ in }(0,+\infty)\times(0,+\infty) , \\
		U = G,\ \de_t U = H & \text{ on }(0,+\infty)\times\{0\} .
	\end{cases}
	\end{equation}
\end{lemma}
\begin{proof}
	The regularity of all functions $U$, $G$ and $H$ is proven in Lemma~\ref{lem67cd4b64}.
	To show~\eqref{eq67cd4b9e}, we just need to perform the following computations, using again Lemma~\ref{lem67cd4b64}.
	\begin{align}
		\de_r\psi_u(x,t;r) 
		&= \frac{r}{n} \phi_{\laplacian u} (x,t;r) \\
		&\overset{\eqref{eq67cd4bec}}= \frac{r}{n} \phi_{\de_t^2 u} (x,t;r) \\
		&= \frac{r}{n} \de_t^2 \phi_u(x,t;r) , \\
		\de_r^2\psi_u(x,t;r)
		&= \frac1n \de_t^2\phi_u + \frac{r}{n} \de_t^2\left( \frac{n}{r}(\psi_u-\phi_u)\right) \\
		&= \de_t^2 \psi_u + \frac{1-n}{n} \de_t^2\phi_u \\
		&= \de_t^2 \psi_u + \frac{1-n}{n} \left(\frac{n}{r} \de_r\psi_u \right) \\
		&= \de_t^2 \psi_u + \frac{1-n}{r} \de_r \psi_u .
	\end{align}
\end{proof}

%%%%%%%%%%%%%%%%%%%%%%%%%%%%%%%%%%%%%%%%%%%%%%%%%%%%%%%%%%%%%%%
\subsection{Solution by spherical means: case $n=3$}\label{subs67ce94a5}
Here we show how to obtain \emph{Kirchhoff's formula}~\eqref{eq67ce93f5},
where we can assume $s=0$ without loss of generality by~\ref{subs67cd41e1}.

So, let $u\in C^2(\R^3\times[0,+\infty))$ be such that~\eqref{eq67d7f53c} holds, that is
\begin{equation}
\begin{cases}
	\de_t^2u-\laplacian u = 0 & \text{ in }\R^3\times(0,+\infty) , \\
	u = g,\ \de_t u = h & \text{ on }\R^3\times\{0\} . \\
\end{cases}
\end{equation}
Define $U$, $G$ and $H$ as in Lemma~\ref{lem67cd498e} and then set
\begin{align}
	\tilde U(x;r,t) = r U(x;r,t) ,
	&& 
	\tilde G(x;r) = r G(x;r) ,
	&&
	\tilde H(x;r) = r H(x;r) .
\end{align}
Then a direct computation shows that, for each $x\in\R^3$,
\begin{equation}
\begin{cases}
	\de_t^2 \tilde U - \de_r^2 \tilde U = 0 & \text{ for }r,t>0 \\
	\tilde U(x;r,0) = \tilde G(x;r) , \ 
	\de_t \tilde U(x;r,0) = \tilde H(x;r)
	& \text{ for }r>0 .
\end{cases}
\end{equation}
So, for each $x\in\R^3$, the function $U(x;\cdot,\cdot)$ solves the pinched string problem~\ref{subs67cd4e88}.
Since the solution to the pinched string problem~\ref{subs67cd4e88} is the only one,
we obtain, for $r\in(0,t)$,
\begin{equation}
	\tilde U (x;r,t) = \frac12 (\tilde G(x;r+t)-\tilde G(x;t-r)) + \frac12 \int_{-r+t}^{r+t} \tilde H(x;s) \did s .
\end{equation}
So,
\begin{align*}
 	u(x,t)
	&= \lim_{r\to0} U(x;r,t) 
	= \lim_{r\to0} \frac{ \tilde U(x;r,t)}{r} \\
	&= \lim_{r\to0} \left( \frac{\tilde G(x;r+t) - \tilde G(x;t-r) }{2r} 
		+ \fint_{t-r}^{t+r} \tilde H(x;s) \did s \right) \\
	&= \de_r\tilde G(x;r)|_{r=t} + \tilde H(x;t) \\
	&= (G(x;t) + t\de_r G(x;r)|_{r=t}) + t H(x;t) \\
	&\overset{\eqref{eq67bacb9a}}= \fint_{\de B(x,t)} g(y) \did S(y) 
		+ \frac{t^2}{3} \fint_{B(x,t)} \laplacian g(y) \did y 
		+ t \fint_{\de B(x,t)} h(y) \did S(y) \\
	&= \fint_{\de B(x,t)} \left( g(y) + t \grad g(y)\cdot \frac{y-x}{|y-x|} + th(y) \right) \did S(y) \\
	&= \fint_{\de B(x,t)} \left( g(y) + \grad g(y)\cdot (y-x) + th(y) \right) \did S(y) .
\end{align*}

We have thus obtained the \emph{Kirchhoff's formula}.
\index{Kirchhoff's formula}

\begin{theorem}[Kirchhoff's formula]\label{thm67ce9378}
	Let $u\in C^2(\R^3\times[0,+\infty))$ be such that $\waveop u =0$ in $\R^3\times(0,+\infty)$.
	Then, for every $s\ge 0$ and every $t>0$,
	\begin{equation}\label{eq67ce93f5}
		u(x,s+t) = \fint_{\de B(x,t)} \left( u(y,s) + \grad_yu(y,s) \cdot (y-x) + t \de_s u(y,s) \right) \did S(y) .
	\end{equation}
	
	Vice-versa, let $g\in C^2(\R^3)$ and $h\in C^1(\R^3)$, and define
	\begin{equation}\label{eq67d7f537}
		u(x,t) = \fint_{\de B(x,t)} \left( g(y) + \grad_yg(y) \cdot (y-x) + t h(y) \right) \did S(y) .
	\end{equation}
	Then $u\in C^2(\R^3\times[0,+\infty))$ and $u$ is a solution to 
	\begin{equation}\label{eq67d7f53c}
	\begin{cases}
		\waveop u = ( \de_t^2-\laplacian )u = 0 & \text{ in }\R^3\times(0,+\infty) , \\
		u = g,\ \de_t u = h & \text{ on }\R^3\times\{0\} .
	\end{cases}
	\end{equation}
\end{theorem}
\begin{proof}
	The second part of the theorem, that is, that the function~\eqref{eq67d7f537} is of class $C^2$ and that it solves the problem~\eqref{eq67d7f537}, follows from Lemma~\ref{lem67cd4b64} and direct computations.
	
	The first part of the theorem, that is, that every solution to 
	$\waveop u =0$ in $C^2(\R^3\times[0,+\infty))$ satisfies Kirchhoff's formula~\eqref{eq67ce93f5}, is the result of the discussion before the theorem.
%	
%	needs some type of uniqueness result, which we do not have at the moment.
%	For this reason, we cannot really state~\eqref{eq67ce93f5}...
%	Let's prove the first part of the theorem.
%	Without loss of generality, by~\ref{subs67cd41e1}, we can assume $s=0$.
%	So, let $u\in C^2(\R^3\times[0,+\infty))$ be such that~\eqref{eq67d7f53c} holds.
%%	 that is
%%	\begin{equation}
%%	\begin{cases}
%%		\de_t^2u-\laplacian u = 0 & \text{ in }\R^3\times(0,+\infty) , \\
%%		u = g,\ \de_t u = h & \text{ on }\R^3\times\{0\} . \\
%%	\end{cases}
%%	\end{equation}
%	Define $U$, $G$ and $H$ as in Lemma~\ref{lem67cd498e} and then set
%	\begin{align}
%		\tilde U(x;r,t) = r U(x;r,t) ,
%		&& 
%		\tilde G(x;r) = r G(x;r) ,
%		&&
%		\tilde H(x;r) = r H(x;r) .
%	\end{align}
%	Then a direct computation shows that, for each $x\in\R^3$,
%	\begin{equation}
%	\begin{cases}
%		\de_t^2 \tilde U - \de_r^2 \tilde U = 0 & \text{ for }r,t>0 \\
%		\tilde U(x;r,0) = \tilde G(x;r) \ 
%		\de_t \tilde U(x;r,0) = \tilde H(x;r)
%		& \text{ for }r>0 .
%	\end{cases}
%	\end{equation}
%	So, for each $x\in\R^3$, the function $U(x;\cdot,\cdot)$ solves the pinched string problem~\ref{subs67cd4e88}, from which we obtain, for $r\in(0,t)$,
%	\begin{equation}
%		\tilde U (x;r,t) = \frac12 (\tilde G(x;r+t)-\tilde G(x;t-r)) + \frac12 \int_{-r+t}^{r+t} \tilde H(x;s) \did s .
%	\end{equation}
%	So,
%	\begin{align*}
%	 	u(x,t)
%		&= \lim_{r\to0} U(x;r,t) 
%		= \lim_{r\to0} \frac{ \tilde U(x;r,t)}{r} \\
%		&= \lim_{r\to0} \left( \frac{\tilde G(x;r+t) - \tilde G(x;t-r) }{2r} 
%			+ \fint_{t-r}^{t+r} \tilde H(x;s) \did s \right) \\
%		&= \de_r\tilde G(x;r)|_{r=t} + \tilde H(x;t) \\
%		&= (G(x;t) + t\de_r G(x;r)|_{r=t}) + t H(x;t) \\
%		&\overset{\eqref{eq67bacb9a}}= \fint_{\de B(x,t)} g(y) \did S(y) 
%			+ \frac{t^2}{3} \fint_{B(x,t)} \laplacian g(y) \did y 
%			+ t \fint_{\de B(x,t)} h(y) \did S(y) \\
%		&= \fint_{\de B(x,t)} \left( g(y) + t \grad g(y)\cdot \frac{y-x}{|y-x|} + th(y) \right) \did S(y) \\
%		&= \fint_{\de B(x,t)} \left( g(y) + \grad g(y)\cdot (y-x) + th(y) \right) \did S(y) .
%	\end{align*}
\end{proof}



\begin{remark}
	Notice that the integral in~\eqref{eq67ce93f5} is supported on the sphere!
\end{remark}
%
%\begin{theorem}
%	Let $g\in C^2(\R^3)$ and $h\in C^1(\R^2)$, and define
%	\begin{equation}\label{eq67d7f537}
%		u(x,t) = \fint_{\de B(x,t)} \left( g(y) + \grad_yg(y) \cdot (y-x) + t h(y) \right) \did S(y) .
%	\end{equation}
%	Then $u\in C^2(\R^3\times[0,+\infty))$ and $u$ is a solution to 
%	\begin{equation}
%	\begin{cases}
%		\waveop u = ( \de_t^2-\laplacian )u = 0 & \text{ in }\R^3\times(0,+\infty) , \\
%		u = g,\ \de_t u = h & \text{ on }\R^3\times\{0\} .
%	\end{cases}
%	\end{equation}
%\end{theorem}
%\begin{proof}
%	The proof that the function~\eqref{eq67d7f537} is of class $C^2$ and that it solves the problem~\eqref{eq67d7f537}, is a direct application of Lemma~\ref{lem67cd4b64} and some computations.
%\end{proof}

%%%%%%%%%%%%%%%%%%%%%%%%%%%%%%%%%%%%%%%%%%%%%%%%%%%%%%%%%%%%%%%
\subsection{Solution by spherical means: case $n=2$}\label{subs67ce94b2}

\begin{remark}
	The issue here is that, for $n=2$, the Euler--Poisson--Darboux equation~\eqref{eq67cd4b9e} cannot be transformed into a wave equation.
	(Why?)
\end{remark}

Let $u\in C^2(\R^2\times[0,+\infty))$ be a solution to the system
\begin{equation}
\begin{cases}
	\waveop u = ( \de_t^2-\laplacian )u = 0 & \text{ in }\R^2\times(0,+\infty) , \\
	u = g,\ \de_t u = h & \text{ on }\R^2\times\{0\} .
\end{cases}
\end{equation}
We define $\tilde u\in C^2(\R^3\times[0,+\infty))$ by $\tilde u(x_1,x_2,x_3,t) = u(x_1,x_2,t)$,
$\tilde g\in C^2(\R^3)$ by $\tilde g(x_1,x_2,x_3) = g(x_1,x_2)$,
and $\tilde h\in C^1(\R^3)$ by $\tilde h(x_1,x_2,x_3) = u(x_1,x_2)$.
Therefore, $\tilde u$ solves
\begin{equation}
\begin{cases}
	\waveop \tilde u = ( \de_t^2-\laplacian )\tilde u = 0 & \text{ in }\R^3\times(0,+\infty) , \\
	\tilde u = g,\ \de_t \tilde u = h & \text{ on }\R^3\times\{0\} .
\end{cases}
\end{equation}
To avoid confusion, we denote by $B_3$ balls in $\R^3$ and by $B_2$ balls in $\R^2$.
Theorem~\ref{thm67ce9378} tells us that, 
denoting by $\tilde x= (x_1,x_2,0)\in \R^3$ the lift of the point $x=(x_1,x_2)\in\R^2$, then
\begin{align*}
	u(x_1,x_2,t)
 	&= \tilde u(x_1,x_2,0,t) 
	= \fint_{\de B_3(\tilde x,t)} \left( g(\tilde y) + \Diff g(\tilde y)\cdot (\tilde y-\tilde x) + th(\tilde y) \right) \did S(\tilde y)  \\
	&\overset{\eqref{eq67a735b3}}= \frac{1}{3\omega_3 t^2} \int_{-t}^t \int_{\de B_2(x,\sqrt{t^2-s^2})}
		\left( g(y) + t h(y) + \Diff g(y)\cdot (y-x) \right) \did S(y) \frac{1}{\sqrt{t^2-s^2}} \did s \\
	&= \frac{2}{3\omega_3 t^2} \int_0^t \int_{\de B_2(x,\sqrt{t^2-s^2})}
		\left( g(y) + t h(y) + \Diff g(y)\cdot (y-x) \right) \did S(y) \frac{1}{\sqrt{t^2-s^2}} \did s \\
	\text{[$r=\sqrt{t^2-s^2}$]}
	&= \frac{2}{3\omega_3 t^2} \int_0^t \int_{\de B_2(x,r)} \left(g(y) + th(y) + \Diff g(y)\cdot (y-x) \right) \did S(y) \frac{r}{\sqrt{t^2-r^2}} \frac{\did r}{r} \\
	&= \frac{2}{3\omega_3 t^2} \int_0^t \int_{\de B_2(x,r)} \frac{ \left(g(y) + th(y) + \Diff g(y)\cdot (y-x) \right) }{ (t^2 - |y-x|^2)^{1/2} } \did S(y) \did r \\
	&= \frac{2}{3\omega_3 t^2} \int_{B_2(x,t)} \frac{ \left(g(y) + th(y) + \Diff g(y)\cdot (y-x) \right) }{ (t^2 - |y-x|^2)^{1/2} } \did y \\
	&= \frac12 \fint_{B_2(x,t)} \frac{g + th + \Diff g\cdot (y-x) }{ ( t^2 - |y-x|^2 )^{1/2} } \did (y) .
\end{align*}

\index{Poisson's formula}
\begin{theorem}[Poisson's formula]\label{thm67e1265b}
	Let $u\in C^2(\R^2\times[0,+\infty))$ be such that $\waveop u =0$ in $\R^2\times(0,+\infty)$.
	Then, for every $s\ge 0$ and every $t>0$,
	\begin{equation}\label{eq67e12659}
		u(x,s+t) = \frac12 \fint_{B_2(x,t)} \frac{u(y,s) + t\de_t(y,s) + \grad u(y,s)\cdot (y-x) }{ ( t^2 - |y-x|^2 )^{1/2} } \did (y) .
	\end{equation}
	
	Vice-versa, let $g\in C^2(\R^2)$ and $h\in C^1(\R^2)$, and define
	\begin{equation}\label{eq67d7f537}
		u(x,t) = \frac12 \fint_{B_2(x,t)} \frac{g(y) + th(y) + \grad g(y)\cdot (y-x) }{ ( t^2 - |y-x|^2 )^{1/2} } \did (y) .
	\end{equation}
	Then $u\in C^2(\R^2\times[0,+\infty))$ and $u$ is a solution to 
	\begin{equation}\label{eq67d7f53c}
	\begin{cases}
		\waveop u = ( \de_t^2-\laplacian )u = 0 & \text{ in }\R^2\times(0,+\infty) , \\
		u = g,\ \de_t u = h & \text{ on }\R^2\times\{0\} .
	\end{cases}
	\end{equation}
\end{theorem}


\begin{exercise}
	Compute $\omega_3$, the volume of the unit ball in $\R^3$.
	
	{\it Solution:}
	\begin{align*}
	 	\omega_3
		&= 2\int_{B_2(x,1)} \sqrt{1-|x|^2} \did x \\
		&= 2 \int_0^1 \int_0^{2\pi} \sqrt{1-r^2} \did \theta r \did r \\
		&= 4\pi \int_0^1 r \sqrt{1-r^2} \did r \\
		\text{[$r=\sin t$]}
		&= 4\pi \int_0^{\pi/2} \sin t \cos t \did t
		= [...] = \frac{4\pi}{3} .
	\end{align*}
\end{exercise}

%%%%%%%%%%%%%%%%%%%%%%%%%%%%%%%%%%%%%%%%%%%%%%%%%%%%%%%%%%%%%%%
\subsection{Solution of the wave equation in all dimensions}

For the next theorem, see \cite[Th. 5.15, page 170]{MR1357411} or \cite[Th. 2.4.2, page 77]{MR2597943}.
\begin{theorem}[Odd dimensions]\label{thm67e1275d}
	Let $n\ge3$ odd, say $n= 2m-1$ for $m\ge2$, or $m=\frac{n+1}{2}$.
	Let $g\in C^{m+1}(\R^n)$, $h\in C^m(\R^n)$, and define
	\begin{equation}\label{eq67e16805}
	\begin{aligned}
		u(x,t) &= \frac1{\gamma_n} 
		\bigg[ 
		\left(\frac{\de}{\de t}\right) \left( \frac1t \frac{\de}{\de t} \right)^{\frac{n-3}{2}} \left( t^{n-2} \fint_{\de B(x,t)} g(y) \did S(y) \right) \\
		&\qquad\qquad\qquad\qquad\qquad\qquad
		+ \left(\frac1t \frac{\de}{\de t}\right)^{\frac{n-3}{2}} \left( t^{n-2} \fint_{\de B(x,t)} h(y) \did S(y) \right)
		\bigg] ,
	\end{aligned}
	\end{equation}
	where $\gamma_n$ is the product of the odd numbers from 1 to $n-2$.
	
	Then $u\in C^2(\R^n\times[0,+\infty)$ and $u$ is a solution to
	\begin{equation}
	\begin{cases}
		\waveop u = ( \de_t^2-\laplacian )u = 0 & \text{ in }\R^n\times(0,+\infty) , \\
		u = g,\ \de_t u = h & \text{ on }\R^n\times\{0\} .
	\end{cases}
	\end{equation}
\end{theorem}

For the next theorem, see \cite[Th. 5.17, page 171]{MR1357411} or \cite[Th. 2.4.3, page 80]{MR2597943}.
\begin{theorem}[Even dimensions]\label{thm67e12abf}
	Let $n\ge2$ even, say $n= 2m-2$ for $m\ge2$, or $m=\frac{n+2}{2}$.
	Let $g\in C^{m+1}(\R^n)$, $h\in C^m(\R^n)$, and define
	\begin{equation}\label{eq67e16844}
	\begin{aligned}
		u(x,t) &= \frac1{\beta_n} 
		\bigg[ 
		\left(\frac{\de}{\de t}\right) \left( \frac1t \frac{\de}{\de t} \right)^{\frac{n-2}{2}} \left( t^{n} \fint_{B(x,t)} \frac{g(y)}{(t^2-|y-x|^2)^{1/2}} \did S(y) \right) \\
		&\qquad\qquad\qquad\qquad\qquad\qquad
		+ \left(\frac1t \frac{\de}{\de t}\right)^{\frac{n-2}{2}} \left( t^{n} \fint_{B(x,t)} \frac{h(y)}{(t^2-|y-x|^2)^{1/2}}\did S(y) \right)
		\bigg] ,
	\end{aligned}
	\end{equation}
	where $\beta_n$ is the product of the even numbers from 1 to $n$.
	
	Then $u\in C^2(\R^n\times[0,+\infty)$ and $u$ is a solution to
	\begin{equation}
	\begin{cases}
		\waveop u = ( \de_t^2-\laplacian )u = 0 & \text{ in }\R^n\times(0,+\infty) , \\
		u = g,\ \de_t u = h & \text{ on }\R^n\times\{0\} .
	\end{cases}
	\end{equation}
\end{theorem}

\begin{remark}
	Notice that in both theorems~\ref{thm67e1275d} and~\ref{thm67e12abf},
	\begin{equation}
		m = \left\lfloor \frac{n}{2} \right\rfloor + 1 .
	\end{equation}
\end{remark}

\begin{exercise}
	Recover Kirchhoof's formula~\eqref{eq67d7f537} from~\eqref{eq67e16805}.
\end{exercise}

\begin{exercise}
	Recover Poisson's formula~\eqref{eq67d7f537} from~\eqref{eq67e16844}.
\end{exercise}

%%%%%%%%%%%%%%%%%%%%%%%%%%%%%%%%%%%%%%%%%%%%%%%%%%%%%%%%%%%%%%%
\subsection{Solution to the nonhomogeneous wave equation: Duhamel's principle}\label{subs67e12d15}

For the next theorem, see \cite[Th. 5.25, page 175]{MR1357411} or \cite[Th. 2.4.4, page 81]{MR2597943}.
\begin{theorem}[Nonhomogeneous equation with null initial data]\label{thm67e131b9}
	Let $n\ge2$ and $f\in C^{\left\lfloor \frac{n}{2} \right\rfloor + 1} (\R^n\times[0,+\infty))$.
	For every $s>0$, let $u_s:\R^n\times[s,+\infty)\to\C$ be the solution 
	in $C^2(\R^n\times[0,+\infty))$ to 
	\begin{equation}
	\begin{cases}
		\waveop u = ( \de_t^2-\laplacian )u = 0 & \text{ in }\R^n\times(s,+\infty) , \\
		u = 0,\ \de_t u = f(\cdot,s) & \text{ on }\R^n\times\{s\} .
	\end{cases}
	\end{equation}
	Define $u:\R^n\times[0,+\infty)\to\C$ by
	\begin{equation}
		u(x,t) = \int_0^t u_s(x,t) \did s .
	\end{equation}
	Then $u\in C^2(\R^n\times[0,+\infty))$ and $u$ is a solution to
	\begin{equation}
	\begin{cases}
		\waveop u = ( \de_t^2-\laplacian )u = f & \text{ in }\R^n\times(0,+\infty) , \\
		u = 0,\ \de_t u = 0 & \text{ on }\R^n\times\{0\} .
	\end{cases}
	\end{equation}
\end{theorem}

\begin{exercise}
	Prove Theorem~\ref{thm67e131b9}.
\end{exercise}

\begin{exercise}
	Write explicitly $u$ from Theorem~\ref{thm67e131b9} for $n=2$ and $n=3$.
\end{exercise}

\begin{theorem}[Nonhomogeneous wave equation]\label{thm67e132f7}
	Let $n\ge2$ and $m=\left\lfloor \frac{n}{2} \right\rfloor + 1$.
	Let $f\in C^m (\R^n\times[0,+\infty))$,
	Let $g\in C^{m+1}(\R^n)$, and $h\in C^m(\R^n)$.
	
	Let $u_0$ be the function given by Theorem~\ref{thm67e1275d} and~\ref{thm67e12abf},
	and $u_1$ the function given by Theorem~\ref{thm67e131b9}.
	Set $u=u_0+u_1$.
	Then $u \in C^2(\R^n\times[0,+\infty))$ and $u$ is a solution to
	\begin{equation}
	\begin{cases}
		\waveop u = ( \de_t^2-\laplacian )u = f & \text{ in }\R^n\times(0,+\infty) , \\
		u = g,\ \de_t u = h & \text{ on }\R^n\times\{0\} .
	\end{cases}
	\end{equation}
\end{theorem}

\begin{exercise}
	Prove Theorem~\ref{thm67e132f7}.
\end{exercise}

























%%%%%%%%%%%%%%%%%%%%%%%%%%%%%%%%%%%%%%%%%%%%%%%%%%%%%%%%%%%%%%%
\newpage
\part{Distributions}



%%%%%%%%%%%%%%%%%%%%%%%%%%%%%%%%%%%%%%%%%%%%%%%%%%%%%%%%%%%%%%%
%\newpage
\section{Distributions}

Here we run through the fundamentals of the theory of distributions omitting a few details.
The details can be recovered by thinking through this material, or reading Rudin's book~\cite{MR1157815}, which I strongly recommend:

• \fullcite{MR1157815}\\

Another valuable reference is Hörmander's first book of his series on linear partial differential operators~\cite{MR717035}:

• \fullcite{MR717035}

Hörmander monograph is one of the most important resources on PDE.
The text is has a very high density, and this can lead to obscurity: if you read it slowly, it will become crystalline clear!\\

An even denser reference is Chapter Four of Federer's book~\cite{MR0257325}:

• \fullcite{MR0257325}

There, you can find a more general construction of distributions, where test functions are smooth functions $\Omega\to Y$, where $\Omega$ is an open set in a Banach space and $Y$ is another Banach space. 
As for Hörmander, Federer's style is at times obscure, but extremely precise, abstract and general.
Read it slowly.\\

Never forget to check out Wikipedia:

• \url{https://en.wikipedia.org/wiki/Distribution_(mathematics)}\\

Finally, the founding father of distributional calculus was Laurent Schwartz\footnote{\url{https://en.wikipedia.org/wiki/Laurent_Schwartz}}, who won the Fields medal in 1950 for the reason\footnote{\url{https://www.mathunion.org/fileadmin/IMU/Prizes/Fields/1950/index.html}}:
\begin{quotation}
	Developed the theory of distributions, a new notion of generalized function motivated by the Dirac delta-function of theoretical physics.
\end{quotation}
He then wrote a beautiful autobiography \cite{zbMATH01571670}, which I suggest everyone to read:

• \fullcite{zbMATH01571670}.

This is the English translation. I have an Italian translation at home, but the original is in French. Probably, there is a German translation too.

%%%%%%%%%%%%%%%%%%%%%%%%%%%%%%%%%%%%%%%%%%%%%%%%%%%%%%%%%%%%%%%
\subsection{Test Functions}\label{subs67def15f}
\index{test function}

For every set $E\subset\R^n$,
we define
\begin{equation}
	\scr D(E) = \{\phi\in C^\infty_c(\R^n): \spt(\phi)\subset E \} .
\end{equation}

If $\Omega\subset\R^n$ is an open set,
then $\scr D(\Omega)$ is the space $C^\infty_c(\Omega)$ of smooth functions $\Omega\to\C$ with compact support.
We call elements of $\scr D(\Omega)$ \emph{test functions}.
We write just $\scr D$ for $\scr D(\R^n)$.

We endow $\scr D(\Omega)$ with a topology that has the following property:
For every sequence $\{\phi_j\}_{j\in\N}\subset\scr D(\Omega)$ and $\phi\in\scr D(\Omega)$,
\begin{equation}\label{eq67e16d22}
	\phi_j\overset{\scr D(\Omega)}\longrightarrow \phi
	\quad\IFF\quad
	\begin{array}{c}
		\exists K\Subset\Omega\ \forall j\in\N\ \spt(\phi_j)\subset K ,\text{ and }\\
		\forall \alpha\in\N^n\ \lim_{j\to\infty}\|\Diff^\alpha\phi_j-\Diff^\alpha\phi\|_{L^\infty} = 0 .
	\end{array}
\end{equation}

\begin{exercise}
	We can see $\scr D(\Omega)$ as a subspace of $\scr D(\R^n)$, but not as a closed subspace.
	Why?
\end{exercise}

%%%%%%%%%%%%%%%%%%%%%%%%%%%%%%%%%%%%%%%%%%%%%%%%%%%%%%%%%%%%%%%
\subsection{The topology of test functions}\label{subs67e16d33}
We will use only the 
In most of the situations,
Property~\eqref{eq67e16d22} is everything we need to know about the topology of $\scr D(\Omega)$.
However, the fact that this notion of convergence descends from a topology, is a non-trivial fact which needs precise definition of that topology.

There are three ways to construct the topology of test functions.
First of all, for $K\subset\R^n$ compact, the space $\scr D(K)$ is a Frechét space when endowed with the family of pseudonorms 
\begin{equation}
	\|u\|_\alpha = \|\Diff^\alpha u\|_{L^\infty(K)} ,
	\qquad\alpha\in\N^n . 
\end{equation}

The first way to construct the topology of $\scr D(\Omega)$ is defining the collection
$\beta$ of all convex balanced sets $W\subset\scr D(\Omega)$ such that $\scr D(K)\cap W$ is open $\scr D(K)$ for all $K\subset\Omega$ compact.
A a set $W$ is \emph{balanced} if $\lambda W \subset W$ for all $\lambda\in\C$ with $|\lambda|\le 1$.
The collection $\beta$ induces a topology $\tau$ make of unions of elements of $\{x+W:x\in\scr D(K),\ w\in\beta\}$.
Then $\tau$ makes $\scr D(\Omega)$ into a locally convex topological vector space.

The other two ways are in terms of initial and final topologies:

\begin{definition}[Initial, or projective, topology]
	Given a set $Y$ and a family of topological spaces $\{Z_i\}_{i\in I}$ and functions $f_i:Y\to Z_i$.
	The \emph{initial topology} or \emph{projective topology} induced by the family of functions $f_i$ is the coarsest (i.e., smallest) topology in $Y$ that makes all functions $f_i$ continuous.
\end{definition}
 
\begin{definition}[Final, or inductive, topology]
	Given a set $Y$ and a family of topological spaces $\{X_i\}_{i\in I}$ and functions $f_i:X_i\to Y$.
	The \emph{final topology} or \emph{inductive topology} induced by the family of functions $f_i$ is the finest (i.e., largest) topology in $Y$ that makes all functions $f_i$ continuous.
\end{definition}

So, we start with the Banach spaces
\begin{equation}
	C^m_c(K) = \{\phi:\R^n\to\C\text{ of class $C^m$ with }\spt(\phi)\subset K\} ,
\end{equation}
where the norm is 
\begin{equation}\label{eq67deff0b}
	\|\phi\|_{C^m(\Omega)} = \max\{ \|\Diff^\alpha\phi\|_{L^\infty(\Omega)} : |\alpha|\le m \} .
\end{equation}
%$\|\phi\|_m = \max\{\|\Diff^\alpha\phi\|_{L^{\infty}}:|\alpha|\le m\}$.

Then we set
\begin{equation}
	C^\infty_c(K) = \bigcap_{m\in\N} C^m_c(K)
\end{equation}
endowed with the initial topology induced by the functions $C^\infty_c(K)\into C^m_c(K)$;
and then
\begin{equation}\label{eq67e2591e}
	\scr D(\Omega) = C^\infty_c(\Omega) = \bigcup_{K\Subset\Omega} C^\infty_c(K) 
	= \bigcup_{K\Subset\Omega}\bigcap_{m\in\N} C^m_c(K) ,
\end{equation}
endowed with the final topology induced by the functions $C^\infty_c(K)\into C^\infty_c(\Omega)$.

An equivalent way is to take first
\begin{equation}
	C^m_c(\Omega) = \bigcup_{K\Subset\Omega} C^m_c(K)
\end{equation}
endowed with the final topology induced by the functions $C^m_c(K)\into C^m_c(\Omega)$,
and then
\begin{equation}\label{eq67e2593b}
	\scr D(\Omega) = C^\infty_c(\Omega) = \bigcap_{m\in\N} C^m_c(\Omega)
	= \bigcap_{m\in\N}\bigcup_{K\Subset\Omega} C^m_c(K)
\end{equation}
endowed with the initial topology induced by the functions $C^\infty_c(\Omega)\into C^m_c(\Omega)$.


\begin{exercise}
	Show that the two topologies coming from~\eqref{eq67e2591e} and~\eqref{eq67e2593b} are the same.
\end{exercise}

%%%%%%%%%%%%%%%%%%%%%%%%%%%%%%%%%%%%%%%%%%%%%%%%%%%%%%%%%%%%%%%
\subsection{Continuity of linear operators}\label{subs67e26996}

\begin{proposition}\label{prop67e26a2f}
	Let $Y$ be a locally convex space and $L:\scr D(\Omega) \to Y$ linear.
	Then the following are equivalent:
	\begin{enumerate}
	\item
	$L$ is continuous;
	\item
	if $\phi_j\to 0$ in $\scr D(\Omega)$ then $L\phi_j\to 0$ in $Y$;
	\item
	the restrictions of $L$ to every $C^\infty_c(K)\subset\scr D(\Omega)$, for $K\Subset\Omega$, are continuous.
	\end{enumerate}
\end{proposition}
\begin{proof}
	The equivalence $(2)\IFF(3)$ is clear, as it is clear $(1)\THEN(2)$.
	The implication $(3)\THEN(1)$ follows from the properties of the topology in $\scr D(\Omega)$.
	See Rudin~\cite[Thm.6.6]{MR1157815}.
\end{proof}


%%%%%%%%%%%%%%%%%%%%%%%%%%%%%%%%%%%%%%%%%%%%%%%%%%%%%%%%%%%%%%%
\subsection{Distributions}\label{subs67def3c0}
A distribution is an element of the dual space $\scr D'(\Omega)$, that is, 
a continuous linear functional $\scr D(\Omega)\to\C$.
The topology of $\scr D'(\Omega)$ is the weak* topology:
For every sequence $\{A_j\}_{j\in\N}\subset\scr D'(\Omega)$ and $A\in\scr D(\Omega)$,
\begin{equation}
	A_j\overset{\scr D'(\Omega)}\longrightarrow A
	\quad\IFF\quad
	\forall \phi\in\scr D(\Omega)\ \lim_{j\to\infty} A_j[\phi] = A[\phi] .
\end{equation}
We write just $\scr D'$ for $\scr D'(\R^n)$.


%%%%%%%%%%%%%%%%%%%%%%%%%%%%%%%%%%%%%%%%%%%%%%%%%%%%%%%%%%%%%%%
\subsection{Functions as distributions}\label{sub67def4a2}
To every $f\in L^1_\loc(\Omega)$, we associate a distribution $A_f\in \scr D'(\Omega)$ defined by, for $\phi\in\scr D(\Omega)$,
\begin{equation}
	A_f[\phi] = \int_\Omega f(x) \phi(x) \did x .
\end{equation}
Let's show that $A_f$ is a distribution.
Clearly $A_f$ is linear. We need to show it is continuous.
If $\phi_j\to \phi$ in $\scr D(\Omega)$, then there is $K\Subset\Omega$ with $\spt(\phi_j)\subset K$ and $\|\phi_j-\phi\|_{L^\infty(\Omega)} \to 0$.
Therefore, 
\begin{equation}
	|A_f[\phi_j] - A_f[\phi]|
	\le \int_\Omega |f(x)| |\phi_j(x) - \phi(x)| \did x
	\le \int_K |f(x)|\did x \|\phi_j-\phi\|_{L^\infty(\Omega)}
	\to 0 .
\end{equation}
This shows that $A_f$ is continuous.

The Fundamental Theorem of Calculus implies that, if $f,g\in L^1_\loc(\Omega)$ are such that $A_f = A_g$ as distributions, then $f=g$ almost everywhere, that is, $f=g$ in $L^1_\loc(\Omega)$.
We thus have an inclusion $L^1_\loc(\Omega)\into \scr D(\Omega)$.

We can say more about this inclusion: it is continuous.
Here we consider on $L^1_\loc(\Omega)$ the topology of local convergence in $L^1$, that is, $f_j\to f$ in $L^1_\loc(\Omega)$ if and only if, for every $K\Subset\Omega$, we have $f_j|_K\to f|_K$ in $L^1(K)$, i.e., $\|f_j-f\|_{L^1(K)} \to 0$.
So, if $f_j\to f$ in $L^1_\loc(\Omega)$ and $\phi\in\scr D(\Omega)$, 
$| A_{f_j}[\phi] - A_{f}[\phi] | \le \int_{\spt(\phi)} |f_j(x)-f(x)| |\phi(x)| \did x 
\le \|f_j-f\|_{L^1(\spt(\phi))} \|\phi\|_{L^\infty} \to 0$.

\begin{exercise}
	Find a sequence $f_j\in L^1_\loc(\R)$ such that $\|f_j\|_{L^1([0,1])}=1$ but $A_{f_j}\to 0$ in $\scr D'(\R)$.
	
	{\it Hint}:
	Take $f_j(x) = \sum_{j=1}^{2^n} (-1)^j \one_{((j-1)/2^n,j/2^n)}(x)$.
	Then $\int_0^1|f_j(x)| \did x = 1$.
	If $\phi\in\scr D(\Omega)$, then there is $L$ such that $|\phi(x)-\phi(y)|<L|x-y|$
	for every $x,y\in\R$.
	\begin{align*}
	 	\left| \int_\R f_j(x)\phi(x) \did x \right|
		&= \left| \sum_{j=1}^{2^n} (-1)^j \int_{(j-1)/2^n}^{j/2^n} \phi(x)\did x \right|\\
		&= \left| \sum_{j=1}^{2^{n-1}} \int_{0}^{1/2^n} \left(- \phi\left(\frac{j-1}{2^{n-1}} + x \right) + \phi\left(\frac{j-1}{2^{n-1}} + \frac{1}{2^n} + x \right) \right)\did x \right|\\
		&\le \sum_{j=1}^{2^{n-1}} \int_{0}^{1/2^n} \frac{L}{2^n} \did x \\
		&= \frac{L}{2^n} \frac{2^{n-1}}{2^n} = \frac{L}{2^{n+1}} .
	\end{align*}
%	\qed
\end{exercise}

\begin{exercise}
	Show that $f_j\to 0$ weakly* in $L^1_\loc(\R^n)$, if and only if $A_{f_j}\to 0$ in $\scr D'(\R^n)$.
	The weak* convergence is $\int_\R f_j g \did x \to 0$ for all $g\in L^\infty(\R^n)$ with compact support.
\end{exercise}

%Finally, we can say that this inclusion is an embedding, that is, if $A_{f_j}\to A_f$ in $\scr D'(\Omega)$, then $f_n\to f$ in $L^1_\loc(\Omega)$.
%Indeed, up to taking $f_n-f$ instead of $f_n$, we can assume $A_f=0$.
%So, $A_{f_n}\to0$ in $\scr D'(\Omega)$.
%\TODO....

%A note: a sequence $A_{f_n}$ could converge to some distribution $A$ that is not a function. 
%This does not contradict our discussion above: there, we assumed that the limit $A$ is a function.

After this, we can denote still by $f$ the distribution $A_f$, that is,
\begin{equation}
	L^1_\loc(\Omega) \subset \scr D'(\Omega) .
\end{equation}

\begin{exercise}
	Let $\{u_k\}_{k\in\N}\subset C^\infty(\Omega)$ be a sequence of harmonic functions
	and suppose that $u_k\to A$ in $\scr D'(\Omega)$.
	Show that $A$ is a harmonic function.
\end{exercise}

%%%%%%%%%%%%%%%%%%%%%%%%%%%%%%%%%%%%%%%%%%%%%%%%%%%%%%%%%%%%%%%
\subsection{Measures as distributions}\label{subs67def74c}
\newcommand{\radon}{\mathtt{Rad}}
Let $\mu$ be a Radon measure on $\Omega$: $\mu\in\radon(\Omega)$.
Then we define the distribution
\begin{equation}
	A_\mu[\phi] = \int_\Omega \phi(x) \did \mu(x) .
\end{equation}
As we saw with functions, the map $\mu\mapsto A_\mu$ is a continuous embedding:
\begin{equation}
	\radon(\Omega) \subset \scr D'(\Omega) .
\end{equation}

Examples of Radon measures:
\begin{enumerate}
\item The \emph{Dirac delta} centered at $x\in\Omega$ is the measure $\delta_x$ defined by: $\delta_x(E) = 1$ if $x\in E$, $\delta_x(E)=0$ if $x\notin E$.
On test functions, the Dirac delta acts as an evaluation: $\delta_x[\phi] = \phi(x)$.
\item If $E\subset\Omega$ closed or open, the measure $\scr L^n|_E$ is Radon.
\item Integration over embedded submanifolds of $\Omega$ are Radon measures.
\end{enumerate}

%%%%%%%%%%%%%%%%%%%%%%%%%%%%%%%%%%%%%%%%%%%%%%%%%%%%%%%%%%%%%%%
\subsection{Order of a distribution}\label{subs67defb33}
A distribution $A\in\scr D'(\Omega)$ \emph{has order (up to) $N$} is there is $C<\infty$ with
\begin{equation}
	\forall \phi\in\scr D(\Omega),\qquad
	A[\phi] \le C \|\phi\|_{C^N(\Omega)} .
\end{equation}
Recall that the norm $\|\phi\|_{C^N(\Omega)}$ was defined in~\eqref{eq67deff0b}.
Notice that if $A$ has order $N$, then it has also order $N+1$.
We say that $A$ has order exactly $N$ if it has order $N$ but not order $N-1$.

\begin{proposition}\label{prop67defc54}
	A linear functional $A:\scr D'(\Omega)\to\C$ is continuous, i.e., a distribution,
	if and only if it has locally finite order, that is, for every $K\Subset\Omega$ there are $N\in\N$ and $C<\infty$ such that, for every $\phi\in\scr D(\Omega)$ with $\spt(\phi)\subset K$, $A[\phi]\le C\|\phi\|_{C^N(\Omega)}$.
\end{proposition}
\begin{proof}
	In this proof, we use two key facts.
	First, the norms defined in~\eqref{eq67deff0b} satisfy 
	$\|\phi\|_{C^N(\Omega)} \ge \|\phi\|_{C^k(\Omega)}$ whenever $k\le N$.
	Second, if we fix $k\in\N$ and $K\Subset\Omega$, then $A$ is a continuous linear functional $(\scr D(K) , \|\cdot \|_{C^k(K)}) \to\C$,
	that is, there is $C_k$ such that $A[\phi]\le C_k\|\phi\|_{C^k(\Omega)}$ for every $\phi\in\scr D(K)$.
	This follows from the very definition of distribution.
	
	So, arguing by contradiction, assume our proposition is false,
	that is, there is $K\Subset\Omega$ such that, for every $N\in\N$ there is $\phi_N\in\scr D(\Omega)$ with $\spt(\phi_N)\subset K$, and $A[\phi_N]\ge N\|\phi_N\|_{C^N(K)}$.
	We have reached a contradiction: for every $N\ge k$, we should also have 
	$N\|\phi_N\|_{C^k(\Omega)} \le N\|\phi_N\|_{C^N(\Omega)} \le A[\phi_N] \le C_k \|\phi_N\|_{C^k(\Omega)}$.
	Since $A[\phi_N]\neq0$, then $\|\phi_N\|_{C^k(\Omega)}\neq0$,
	and thus we get $N\le C_k$ for all $N\ge k$: \Lightning.
	
	We conclude that such a sequence $\{\phi_N\}_N$ cannot exists.
\end{proof}

\begin{exercise}
	Show that, if $A\in\scr D'(\Omega)$ has finite order $N$, then $A$ extends as a continuous linear operator from $\scr D(\Omega)$ to $C^N(\Omega)$.
\end{exercise}

%%%%%%%%%%%%%%%%%%%%%%%%%%%%%%%%%%%%%%%%%%%%%%%%%%%%%%%%%%%%%%%
\subsection{Distributions of order 0}\label{subs67df019e}

If $X$ is a topological space, we define $C_c(X)$ as the space of continuous functions $X\to\C$ with compact support endowed with the $L^\infty$ norm.
The closure of $C_c(X)$ in $C(X)$ is $C_0(X)$, which is the space of continuous functions vanishing ``at infinity'', that is,
\begin{equation}
	C_0(X) = \{f\in C(X): \text{for every $\epsilon>0$ the set $\{|f|\ge\epsilon\}$ is compact}\} .
\end{equation}
The space $C_0(X)$ is a Banach space when endowed with the $L^\infty$ norm.
Its (topological) dual $C_0(X)'$ is also a Banach space when endowed with the operator norm $\|\xi\|_{ C_0(X)'} = \sup\{\xi[\phi] : \phi\in C_0(X),\ \|\phi_{L^\infty}\le 1\}$.

We define $\radon(X;\C)$ as the space of all $\C$-valued Radon measures: the precise definition goes as follows.
Let $\scr B(X)$ be the $\sigma$-algebra of all Borel sets.
A \emph{positive Radon measure} on $X$ is a measure $\lambda:\scr B(X)\to[0,+\infty]$ such that (see \cite[page 212]{MR767633}):
\begin{enumerate}
\item
	$\lambda$ is outer regular on all Borel sets, that is, if $E\in\scr B(X)$, then
	\begin{equation}
		\lambda(E) = \inf\{\lambda(U) : E\subset U,\ U\text{ open}\} ;
	\end{equation}
\item
	$\lambda$ is inner regular on all open sets, that is, if $E\subset X$ is open, then
	\begin{equation}
		\lambda(E) = \sup\{\lambda(K) : K\subset E,\ K\text{ compact}\} ;
	\end{equation}
\item
	$\lambda$ is finite on compact sets, that is, $\lambda(K)<\infty$ for all $K\subset X$ compact.
\end{enumerate}
We denote by $\radon(X;[0,+\infty])$ the space of all positive Radon measures on $X$.
A complex-valued or real-valued Radon measure is a Borel measure $\mu:\scr B(X)\to\K$ with $\K\in\{\R,\C\}$ whose total variation is $|\mu|\in\radon(X;[0,+\infty])$.
We call the space of these measures $\radon(X;\K)$ for $\K\in\{\R,\C\}$.
In fact, if $\mu\in\radon(X;\K)$, then $|\mu|(X)<\infty$.
Moreover, if $\mu\in\radon(X;\R)$, then there are $\mu_+,\mu_-\in\radon(X;[0,+\infty))$ such that $\mu=\mu_+-\mu_-$.
If $\mu\in\radon(X;\C)$, then there are $\mu_r,\mu_i\in\radon(X;\R)$ such that $\mu=\mu_r + i\mu_i$.

%Define 
%\begin{equation}
%	\scr B_b(X) = \{B\in\scr B(X):\exists K\subset X\text{ compact, with }B\subset K\} .
%\end{equation}
%A \emph{real-valued Radon measure} on $X$ is a function $\mu:\scr B_b(X)\to \R$ such that there are $\mu_+,\mu_-\in\radon(X;[0,+\infty])$ such that $\mu=\mu_+-\mu_-$.
%We denote by $\radon(X;\R)$ the space of real-valued Radon measures.
%
%A \emph{complex-valued Radon measure} on $X$ is a function $\mu:\scr B_b(X)\to \C$ such that there are $\mu_r,\mu_i\in\radon(X;\R)$ such that $\mu=\mu_r + i\mu_i$.
%
%Notice that such a real-valued and complex-valued Radon measures are not measures in general.
%To understand why, simply take the Radon measure $\mu\in\radon(\R;\R)$, $\mu(E) = \int_E x\did x$.
%Then $\mu_+(E) = \int_{[0,+\infty)\cap E} x\did x$ and $\mu_-(E) = \int_{(-\infty,0]\cap E} |x| \did x$.
%So, if we want $\mu$ to be a measure, we should have $\mu([0,+\infty))= +\infty$, $\mu([-\infty,0))=-\infty$, 
%and $\mu(\R) = \mu([0,+\infty)) + \mu([-\infty,0)) = +\infty - \infty$, which is not well defined.
%However, real-valued and complex-valued Radon measures are indeed measures on compact subsets of $X$.


Recall the following result (see \cite[Thm 7.17, page 223]{MR767633})
\begin{theorem}[Riesz Representation Theorem]\label{thm67e2d9b5}
	Let $X$ be a topological space that is locally compact\footnote{locally compact: for every $x\in X$ and every $U\subset X$ open with $x\in U$ there exists $V\subset U$ compact with $x\in{\rm interior}(K)$} and Hausdorff\footnote{Hausdorff: for every $x,y\in X$ distinct there are $U,V\subset X$ open such that $x\in U$, $y\in V$ and $U\cap V=\emptyset$.}.
	For $\mu\in\radon(X;\C)$ and $f\in C_0(X)$, define $I_\mu(f) = \int_X f\did\mu$.
	Then $I_\mu$ is an element of the dual $C_0(X)'$ and the map $\mu\mapsto I_\mu$ is an isometric equivalence of $\radon(X;\C)$ with $C_0(X)'$.
\end{theorem}


\begin{proposition}\label{prop6810ed13}
	A distribution $A\in\scr D'(\Omega)$ has order zero if and only if it is a Radon measure.
\end{proposition}
\begin{proof}
	We already know that Radon measures are distributions of order zero.
	Let's prove the other implication.
	
	If $A\in\scr D'(\Omega)$ has order zero, then there is $C$ such that $A[\phi]\le C\|\phi\|_{L^\infty}$ for all $\phi\in \scr D(\Omega)$.
	Since $ \scr D(\Omega)$ is dense in $C_0(\Omega)$, 
	it follows that $A$ continuously extends to a linear operator $C_0(\Omega)\to\C$.
	By Riesz Representation Theorem~\ref{thm67e2d9b5}, there is a Radon measure $\mu\in\radon(\Omega;\C)$ such that $A=I_\mu$.
	This means for us that $A$ is a Radon measure.
\end{proof}

\begin{exercise}\label{ex67df0243}
	Show that $\scr D(\Omega)$ is dense in $( C_0(\Omega),\|\cdot\|_{L^\infty})$.
\end{exercise}

%\begin{exercise}
%	Show that, if $\Omega$ is bounded, then $\radon(\Omega) = C^0_0(\Omega)'$.
%	Show also that $\radon(\R^n) = C^0_b(\R^n)'$, where $C^0_b(\R^n)$ is the Banach space of continuous and bounded functions on $\R^n$ endowed with the $L^\infty$ norm.
%\end{exercise}

%%%%%%%%%%%%%%%%%%%%%%%%%%%%%%%%%%%%%%%%%%%%%%%%%%%%%%%%%%%%%%%
\subsection{Distributions of order 1}\label{subs67df056f}
Here are some examples of distributions of order $1$.
\begin{enumerate}
\item	
$A[\phi] = \de_x\phi(0)$;
\item
$A[\phi] = \int_{\de B(x,r)} \grad\phi(y)\cdot\frac{(y-x)}{r} \did S(y)$.
\item
In general, if $\Sigma\subset\R^n$ is a smooth submanifold, $v:\Sigma\to\R^n$ a smooth vector field, and $\did S$ is the surface measure on $\Sigma$, then 
$A[\phi] = \int_\Sigma \grad\phi(y)\cdot v(y) \did S(y)$ is a distribution of order 1.
\item
For example,
\begin{equation}
	A_{x,r}[\phi] = \fint_{\de B(x,r)} \grad\phi(y) \cdot (y-x) \did y .
\end{equation}
\item 
	the distribution $A[\phi] = \int_0^1 \phi'(x)\did x$ is a distribution of order one on $\R$: $|A[\phi]| \le \int_0^1 |\phi'(x)|\did x \le \|\phi\|_{C^1} $.
	However, $A$ is actually of order zero: $|A[\phi]| = |\phi(1)-\phi(0)| \le 2\|\phi\|_{C^0}$.
\end{enumerate}


%%%%%%%%%%%%%%%%%%%%%%%%%%%%%%%%%%%%%%%%%%%%%%%%%%%%%%%%%%%%%%%
\subsection{Principal value}\label{subs67e1b7b7}
\newcommand{\pv}{\text{p.v.}}
The function $f:\R\setminus\{0\}\to\R$, $f(x)= \frac1x$, is not integrable in a neighborhood of $0$.
For this reason, we cannot see it as a distribution on $\R$, although it is a distribution on $\R\setminus\{0\}$, because $f\in L^1_\loc(\R\setminus\{0\})$.
However, we can define a distribution on $\R$ as follows:
\begin{equation}\label{eq67e1bd0c}
	\scr D(\R) \ni \phi \mapsto \lim_{\epsilon\to0} \int_{\R\setminus[-\epsilon,\epsilon]} \frac{\phi(x)}{x} \did x 
	= \pv \int_\R \frac{\phi(x)}{x} \did x
\end{equation}
If $\spt(\phi)\subset[-a,a]$, then
\begin{equation}
 	\int_{\R\setminus[-\epsilon,\epsilon]} \frac{\phi(x)}{x} \did x
	= \int_{[-a,a]\setminus[-\epsilon,\epsilon]} \frac{\phi(x)-\phi(0)}{x} \did x .
\end{equation}
Since
$\left|\frac{\phi(x)-\phi(0)}{x}\right| \le \|\phi'\|_{L^\infty}$,
then the latter integral converges as $\epsilon\to0$:
\begin{equation}
	\pv \int_\R \frac{\phi(x)}{x} \did x
	= \lim_{\epsilon\to0} \int_{[-a,a]\setminus[-\epsilon,\epsilon]} \frac{\phi(x)-\phi(0)}{x} \did x
	= \int_{[-a,a]} \frac{\phi(x)-\phi(0)}{x} \did x .
\end{equation}
So,~\eqref{eq67e1bd0c} defines a distribution by Proposition~\ref{prop67defc54}.


%%%%%%%%%%%%%%%%%%%%%%%%%%%%%%%%%%%%%%%%%%%%%%%%%%%%%%%%%%%%%%%
\subsection{Adjoint operators}\label{subs67df0783}
\begin{proposition}\label{prop67df0940}
	Let $\Omega_1$ and $\Omega_2$ be open subsets of $\R^n$.
	Let $\Phi:\scr D(\Omega_1)\to\scr D(\Omega_2)$ be a continuous linear operator.
	Then there is a sequentially\footnotemark continuous linear operator $\Phi^*:\scr D'(\Omega_2)\to\scr D'(\Omega_1)$ such that, for every $A\in\scr D'(\Omega_2)$ and $\phi\in \scr D(\Omega_1)$,
	\begin{equation}\label{eq67df0824}
		\Phi^*A[\phi] = A[\Phi\phi] .
	\end{equation}
	\footnotetext{Recall that $T:X\to Y$ is \emph{sequentially continuous} if for every $x_j\to x_\infty$ in $X$ we have $Tx_j\to T_x\infty$ in $Y$.
	Instead, $T$ is (topologically) continuous if for every open set $V\subset Y$ the preimage $T^{-1}(V)$ is open in $X$.
	In topological spaces, (topological) continuity implies sequential continuity.
	The converse is false in this generality.
	
	(A way to recover continuity from convergence is by means of \emph{nets}, a generalization of sequences: a \emph{sequence in $X$} is a function $\N\to X$, a \emph{net in $X$} is a function $\omega\to X$ for some ordinal $\omega$. ``Netial'' continuity implies continuity.)
	
	It remains unclear to me if sequential continuity implies continuity in the context of Proposition~\ref{subs67df0783}.
	}
\end{proposition}
\begin{proof}
	For every $A\in\scr D'(\Omega_2)$, define $\Phi^*A$ by~\eqref{eq67df0824}.
	Notice that, $\Phi^*A = A\circ\Phi$.
	Therefore, since both $A$ and $\Phi$ are linear and continuous, then $\Phi^*A$ is linear and continuous.
	We need to show that $\Phi^*$ is continuous.
	Let $A_j\to A$ in $\scr D'(\Omega_2)$.
	Then, for every $\phi\in \scr D(\Omega_1)$, we have
	$\lim_{j\to\infty} \Phi^*A_j[\phi] = \lim_{j\to\infty} A_j[\Phi\phi] = A[\Phi\phi] = \Phi^*A[\phi]$.
	This shows that $\Phi^*$ is sequentially continuous.
\end{proof}

Proposition~\ref{prop67df0940}, combined with Proposition~\ref{prop67e26a2f}, is the key tool to extend operations from functions to distributions.
We will use it all the time!

%%%%%%%%%%%%%%%%%%%%%%%%%%%%%%%%%%%%%%%%%%%%%%%%%%%%%%%%%%%%%%%
\subsection{Derivatives of distributions}\label{subs67df096d}
If $A\in\scr D'(\Omega)$ and $\alpha\in\N^n$, define, for every $\phi\in\scr D(\Omega)$,
\begin{equation}\label{eq67df24c8}
	\Diff^\alpha A[\phi] = (-1)^{|\alpha|} A[\Diff^\alpha \phi] .
\end{equation}

\begin{exercise}\label{ex67df0a00}
	Show that, if $\alpha\in\N^n$, the function $\phi\mapsto \Diff^\alpha\phi$ is a continuous linear operator $\scr D(\Omega)\to\scr D(\Omega)$.
\end{exercise}

By Exercise~\ref{ex67df0a00} and Proposition~\ref{prop67df0940},
the $A\mapsto \Diff^\alpha A$ is a continuous operator $\scr D'(\Omega)\to\scr D'(\Omega)$.

\begin{exercise}\label{ex67df246e}
	Let $f\in C^N(\Omega)$ and $\phi\in\scr D(\Omega)$.
	Show that, for every $\alpha\in\N^n$ with $|\alpha|\le N$,
	\begin{equation}
		\int_\Omega \Diff^\alpha f(x) \phi(x) \did x 
		= (-1)^{|\alpha|} \int_\Omega f(x) \Diff^\alpha \phi(x) \did x .
	\end{equation}
	In other words, $\Diff^\alpha A_f = A_{\Diff^\alpha f}$.
\end{exercise}

\begin{exercise}
	Show that, if $A\in\scr D'(\Omega)$, then $\Diff^\alpha\Diff^\beta A = \Diff^{\alpha+\beta}A = \Diff^\beta\Diff^\alpha A$ for all $\alpha,\beta\in\N^n$.
\end{exercise}

Distributions are thus infinitely differentiable: this is one of the main features of distributions.
In particular, we have derivatives of every order for each function in $L^1_\loc(\Omega)$.
Exercise~\ref{ex67df246e} shows that, if a distribution $A_f$ is a function $f$ of class $C^N$, then derivatives defined by~\eqref{eq67df24c8} are coherent with derivatives of $f$.

\begin{proposition}
	Let $f\in L^1_\loc(\R)$ be such that there is $g\in L^1(\R)$ with 
	$\Diff A_f = A_g$, i.e., $\Diff f = g$ in distributional sense.
	Then, for almost every $x\in\R$,
	\begin{equation}
		f(x) = \int_{-\infty}^x g(y)\did y .
	\end{equation}
	Consequently, up to changing $f$ on a set of measure zero, $f$ is absolutely continuous and $f'=g$.
	In particular, if $f$ and $g$ are continuous, then $f\in C^1(\R)$ and $f'=g$.
\end{proposition}
\begin{proof}
	The identity $\Diff A_f = A_g$ means that, for every $\phi\in\scr D(\R)$,
	\begin{equation}\label{eq67e9028e}
		\int_{-\infty}^\infty g(x) \phi(x) \did x
		= - \int_{-\infty}^\infty f(x) \phi'(x) \did x .
	\end{equation}
	Define
	\begin{equation}
		v(x) = \int_{-\infty}^x g(y)\did y .
	\end{equation}
	We claim that $A_f = A_v$.
	If $\phi\in\scr D(\R)$, then
	\begin{align*}
	 	A_v [\phi] 
		&= \int_{-\infty}^\infty v(x) \phi(x) \did x \\
		&= \int_{-\infty}^\infty \int_{-\infty}^x g(y) \phi(x) \did y\did x \\
		[\text{Change of variables}]
		&= \int_{-\infty}^\infty \int_0^\infty g(t) \phi(t+s) \did s\did t \\
		&= \int_0^\infty\int_{-\infty}^\infty  g(t) \phi(t+s) \did t\did s \\
		&\overset{\eqref{eq67e9028e}}= -\int_0^\infty \int_{-\infty}^\infty  f(t) \phi'(t+s) \did t\did s \\
		&= - \int_{-\infty}^\infty  f(t) (-\phi(t)) \did t 
		= A_f[\phi] .
	\end{align*}
	where we performed the following change of variables:
	$x=t+s$, $y=t$, $\did x\wedge\did y = (\did t+\did s)\wedge\did t = \did s\wedge\did t$, and $\{x\in\R,\ y\in(-\infty,x]\} = \{y\in\R,\ x\in[y,+\infty)\} = \{t\in\R,\ s\in[0,+\infty)\}$.
	We have thus obtained that $A_v=A_f$ as distributions, 
	and we know that this means that $f=v$ almost everywhere (see~\ref{sub67def4a2}),
	that is,~\eqref{eq67e9028e}.
\end{proof}

\begin{exercise}
	Show the following proposition:
	
	\begin{proposition}
		Let $\Omega\subset\R^n$ be open
		and $f\in C(\Omega)$ a continuous function.
		Suppose that, for every $j\in\{1,\dots,n\}$, there is a continuous function $g_j\in C(\Omega)$ such that $\Diff^jA_f = A_{g_j}$, i.e., $\Diff^jf=g_j$ in distributional sense.
		Then $f\in C^1(\Omega)$ and $\Diff^jf = g_j$.
	\end{proposition}
\end{exercise}

\begin{exercise}
	Let $f:\R\to\R$ be a function with bounded variation.
	For instance, the Cantor staircase function.
	Show that $\Diff A_f = A_\mu$, where $\mu\in\radon(\R)$ is the measure defined by 
	\begin{equation}
		\mu([a,b)) = f(b)-f(a)
	\end{equation}
	for all $a,b\in\R$ with $a<b$.
	For instance, if $f$ is the Cantor staircase function, then we know that,
	for almost every $x\in\R$, $f$ is differentiable at $x$ and $f'(x)=0$.
	However, $\Diff A_f\neq0$.
	
	{\it Hint:} see \cite[\S6.14]{MR1157815}.
\end{exercise}

%%%%%%%%%%%%%%%%%%%%%%%%%%%%%%%%%%%%%%%%%%%%%%%%%%%%%%%%%%%%%%%
\subsection{Intermezzo: Banach--Steinhaus Theorem}\label{subs67e9629a}
We will need to cite Banach--Steinhaus Theorem. 
Here we see a version of its statement that is less general than the original, but it is what we will need later on.
Before stating the theorem, we fix the terms used.

A family $\Gamma$ of linear functions $\gamma:X\to Y$ between topological vector spaces is 
\emph{equicontinuous} if for every open neighborhood $V$ of $0$ in $Y$ there exists an open neighborhood $U$ of $0$ in $X$ such that $\gamma(U)\subset V$ for all $\gamma\in\Gamma$.
As as short notation, for $\Gamma\subset\Lin(X;Y)$ and $E\subset X$, we define
\begin{equation}
	\Gamma(E) = \{\gamma(x):\gamma\in \Gamma,\ x\in E\}\subset Y .
\end{equation}

In the finite-dimensional case, i.e., when both $X$ and $Y$ have finite dimension,
equicontinuity is equivalent to boundedness (in the operator norm) and pre-compactness.
In the infinite dimensional case, this is not the case.

A \emph{Fréchet space} is a topological vector space $(X,\tau)$ such that
\begin{enumerate}
\item	the topology $\tau$ is generated by a complete \emph{invariant metric} (i.e., a complete distance function $d:X\times X\to(0,+\infty)$ such that $d(x+z,y+z) = d(x,y)$ for all $x,y,z\in X$);
\item	$X$ is locally convex, that is, for every $U\subset X$ open with $0\in U$ there exists $V\subset U$ open with $0\in V$ and $V$ convex.
\end{enumerate}

A subset $E\subset X$ in a topological vector space $X$ is \emph{bounded} if for every $U\subset X$ open with $0\in U$ there exists $\lambda>0$ such that $E\subset \lambda U$.
This notion of boundedness might look abstract, but you can easily show the following statement:
{\it If $E=\{x_j\}_{j\in\N}$ is a convergent sequence in $X$, then $E$ is bounded.}

\begin{exercise}
	Show that, if $E=\{x_j\}_{j\in\N}$ is a convergent sequence in $X$, then $E$ is bounded.
\end{exercise}

\begin{theorem}[(Consequence of) Banach--Steinhaus Theorem]\label{thm67e96360}
	Suppose that $X$ is a Fréchet space and $Y$ a topological vector space, 
	$\Gamma$ is a collection of continuous linear maps from $X$ to $Y$.
	If 
	\begin{equation}
		\forall x\in X\qquad \Gamma(x) \text{ is bounded in $Y$} ,
%		B = \{x\in X: \Gamma(x) \text{ is bounded in $Y$} \} .
	\end{equation}
	then $\Gamma$ is equicontinuous.
%	If $B$ is of the second category in $X$, then $B=X$ and $\Gamma$ is equicontinuous.
\end{theorem}

See \cite[\S2.1--\S2.6]{MR1157815}, and also the~\href{https://en.wikipedia.org/wiki/Uniform_boundedness_principle}{\underline{Wikipedia page}}.

\begin{corollary}\label{cor67e96e59}
	Let $X$ be a Fréchet space, $Y$ and $Z$ topological vector spaces,
	and $B:X\times Y\to Z$ a bilinear map.
	Suppose that $B$ is continuous in each entry separately, i.e., 
	for every $x\in X$ the linear map $Y\to Z$, $y\mapsto B(x,y)$, is continuous,
	and for every $y\in Y$ the linear map $X\to Z$, $x\mapsto B(x,y)$, is continuous.
	
	If $\{x_j\}_{i\in\N}\subset X$ and $\{y_j\}_{j\in\N}\subset Y$ are sequences with $\lim_{j\to\infty} x_j = x_\infty$ in $X$ and $\lim_{j\to\infty} y_j = y_\infty$ in $Y$,
	then
	\begin{equation}\label{eq67e96b48}
		\lim_{j\to\infty} B(x_j,y_j) = B(x_\infty,y_\infty) .
	\end{equation}
\end{corollary}
\begin{proof}
	For $j\in\N\cup\{\infty\}$, define $b_j:X\to Z$, $b_j(x) = B(x,y_j)$.
	Set $\Gamma = \{b_j\}_{j\in\N\cup\{\infty\}}$.
	Since $B$ is continuous in each entry, the functions $b_j$ are continuous.
	If $x\in X$, since $\Gamma(x)$ is a sequence in $Z$ convergent to $b_\infty(x)$,
	then $\Gamma(x)$ is bounded.
	By the Banach--Steinhaus Theorem~\ref{thm67e96360}, $\Gamma$ is equicontinuous.
	
	We are now ready to prove~\eqref{eq67e96b48}.
	Let $U\subset Z$ be a neighborhood of $0$ in $Z$.
	We want to show that there is $N\in\N$ such that 
	\begin{equation}\label{eq67e96dba}
		\exists N\in\N\ \forall j>N\qquad
		B(x_j,y_j)\in B(x_\infty,y_\infty) + U .
	\end{equation}
	
	As a general fact in topological vector spaces, there is $\tilde U\subset U$ neighborhood of $0$ in $Z$ such that $\tilde U-\tilde U\subset U$.
	Since $\Gamma$ is equicontinuous, there is $V\subset X$ neighborhood of $0$ in $X$ such that $\Gamma(V)\subset \tilde U$.
	Since $x_n\to x_\infty$, then there is $N\in\N$ such that $x_n-x_\infty \in V$ for all $n>N$.
	Up to taking $N$ larger, we have also $B(x_\infty,y_\infty-y_j)\in \tilde U$ for all $n>N$, 
	because $y\mapsto B(x_\infty,y)$ is continuous and thus $\lim_{j\to\infty} B(x_\infty,y_\infty-y_j) = 0$.
	So, for $j>N$ we have
	\begin{align*}
		B(x_j,y_j) - B(x_\infty,y_\infty)
		&= B(x_j-x_\infty,y_j) - B(x_\infty,y_\infty-y_j)\\
		&= b_j(x_j-x_\infty) - B(x_\infty,y_\infty-y_j) 
		\in \tilde U - \tilde U \subset U .
	\end{align*}
	We have thus obtained~\eqref{eq67e96dba}.
\end{proof}

%%%%%%%%%%%%%%%%%%%%%%%%%%%%%%%%%%%%%%%%%%%%%%%%%%%%%%%%%%%%%%%
\subsection{Product of distributions}\label{sub67df2561}
A derivative is usually defined using the Leibniz rule: $\de(fg) = f\de g + g\de f$.
The product of two distributions is not well defined... We will see special situations in which we can multiply two distributions, but there is not a general product of two distributions.

%%%%%%%%%%%%%%%%%%%%%%%%%%%%%%%%%%%%%%%%%%%%%%%%%%%%%%%%%%%%%%%
\subsection{Product of a smooth function and a distribution}\label{subs67df2598}
If $A\in\scr D(\Omega)$ and $f\in C^\infty(\Omega)$, then we define $fA$ by
\begin{equation}
	(fA)[\phi] = A[f\phi] .
\end{equation}
Again, we justify this formula with Proposition~\ref{prop67df0940}:
indeed, the map $\Phi_f:\phi\mapsto f\phi$ is a continuous linear operator $\scr D(\Omega)\to\scr D(\Omega)$.
So, $\Phi_f^*A = A\circ\Phi_f$ defines a continuous linear operator $\scr D'(\Omega)\to\scr D'(\Omega)$.

\begin{proposition}\label{prop67e267de}
	If $f_k\to f_\infty$ in $C^\infty(\Omega)$ and $A_k\to A_\infty$ in $\scr D'(\Omega)$,
	then $f_kA_k\to f_\infty A_\infty$ in $\scr D'(\Omega)$.
\end{proposition}
\begin{proof}
	Consider the bilinear map $B:C^\infty(\Omega)\times \scr D'(\Omega)\to \scr D'(\Omega)$,
	$B(f,A) = fA$.
	We have seen that $B$ is continuous in each entry separately.
	Since $C^\infty(\Omega)$ is a Fréchet space, we conclude thanks to Corollary~\ref{cor67e96e59} of the Banach--Steinhaus Theorem.
\end{proof}
%\begin{proof}
%	This statement is fundamentally a consequence of \href{https://en.wikipedia.org/wiki/Uniform_boundedness_principle}{\underline{Banach--Steinhaus Theorem}}.
%	See \cite[Theorem 6.18]{MR1157815}.
%\end{proof}

\begin{corollary}\label{cor67e265dd}
	If $\phi_j\to\phi_\infty$ in $\scr D(\Omega)$ and $A_j\to A$ in $\scr D'(\Omega)$,
	then $\lim_{j\to\infty} A_j[\phi_j] = A[\phi_\infty]$.
\end{corollary}
%\begin{proof}
%	\TODO
%\end{proof}

\begin{exercise}
	Prove Corollary~\ref{cor67e265dd} using Proposition~\ref{prop67e267de}.
\end{exercise}

\begin{exercise}[Generalized Leibniz Rule]
	 Show that, if $u\in\scr D'(\Omega)$ and $f\in C^\infty(\Omega)$, then, for every $\alpha\in\N^n$,
	 \begin{equation}\label{eq67f2c2e1}
	 	\Diff^\alpha(fu) = \sum_{\beta\le \alpha} \binom{\alpha}{\beta} \Diff^\beta f \cdot \Diff^{\alpha-\beta} u .
	 \end{equation}
	 {\it Hint:} First of all, understand this formula when $u$ is a smooth function.
	 Then consider the case $|\alpha|=1$ (just one derivative).
\end{exercise}

%%%%%%%%%%%%%%%%%%%%%%%%%%%%%%%%%%%%%%%%%%%%%%%%%%%%%%%%%%%%%%%
\subsection{Locality}\label{subs67e26aff}
We say that two distributions $A_1,A_2\in\scr D'(\Omega)$ are equal on an open set $\omega\subset\Omega$, that is, $A_1=A_2$ in $\omega$, if $A_1\phi=A_2\phi$ for all $\phi\in\scr D(\omega)$.

\begin{proposition}\label{prop67eaabf1}
	Let $\scr U$ be an open cover of an open set $\Omega\subset\R^n$ and let $\{A_\omega\}_{\omega\in\scr U}$ be a collection of distributions with:
	\begin{enumerate}
	\item $A_\omega\in\scr D'(\omega)$ for all $\omega\in\scr U$, and,
	\item $A_{\omega_1} = A_{\omega_2}$ in $\omega_1\cap\omega_2$ for all $\omega_1,\omega_2\in \scr U$.
	\end{enumerate}
	Then there exists a unique $A\in\scr D'(\Omega)$ with $A = A_\omega$ in $\omega$ for all $\omega\in\scr U$.
\end{proposition}
\begin{proof}
	Let $\{\psi_j\}_{j\in\N}$ be a partition of unity subordinated to $\scr U$.
	More precisely:
	\begin{enumerate}
	\item $\psi_j\in C^\infty_c(\Omega)$ for all $j$;
	\item for every $x\in\Omega$ the set $\{j:\psi_j(x)\neq0\}$ is finite; 
	\item $\sum_j\psi_j(x) = 1$ for all $x\in \Omega$;
	\item for every $j\in\N$ there is $\omega_j\in\scr U$ such that $\spt(\psi_j)\subset\omega_j$.
	\end{enumerate}
	We fix the subcover $\{\omega_j\}_{j\in\N}\subset\scr U$.
	
	Define $A:\scr D(\Omega)\to\C$ as follows: if $\phi\in\scr D(\Omega)$, then
	$A[\phi] = \sum_{j\in\N} A_{\omega_j}[\psi_j\phi]$.
	Notice that, since $\spt(\phi)$ is compact and the local finiteness of the cover $\{\omega_j\}_j$, the sum is a finite sum.
	
	Clearly $A$ is linear.
	If $\phi_k\to0$ in $\scr D(\Omega)$, then there is $K\Subset\Omega$ with $\spt(\phi_k)\subset K$ for all $k$.
	So, there is $N\in\N$ with $K\subset\bigcup_{j=1}^N\omega_j$ and 
	$A[\phi_k] = \sum_{j=1}^N A_{\omega_j}[\psi_j\phi_k]$ for all $k$.
	We conclude that $\lim_{k\to\infty}A[\phi_k] = \sum_{j=1}^N \lim_{k\to\infty}A_{\omega_j}[\psi_j\lim_{k\to\infty}\phi_k] = 0$.
	
	By Proposition~\ref{prop67e26a2f}, we obtain that $A\in\scr D'(\Omega)$.
	
	If $\omega\in\scr U$, then, for every $\phi\in\scr D(\omega)$, we have
	$A[\phi]= \sum_{j\in\N} A_{\omega_j}[\psi_j\phi]
	= \sum_{j\in\N} A_{\omega}[\psi_j\phi]
	=  A_{\omega}[\sum_{j\in\N}\psi_j\phi]
	= A_\omega [\phi]$,
	where we used the fact that $\spt(\psi_j\phi)\subset\spt(\psi_j)\cap\spt(\phi)\subset\omega_j\cap\omega$,
	that $A_{\omega_j}=A_\omega$ on $\omega_j\cap\omega$,
	that there is a finite set $J\subset\N$ such that $\spt(\psi_j)\cap\spt(\phi)\neq\emptyset$, 
	and that, for every $x\in\spt(\phi)$, $\sum_{j\in\N}\psi_j(x) = \sum_{j\in J}\psi_j(x) = 1$.
	
	Finally, we need to show uniqueness.
	Suppose there is $\bar A\in\scr D'(\Omega)$ such that $\bar A=A_\omega$ on $\omega$ for every $\omega\in\scr U$.
	Then, for every $\phi\in \scr D(\Omega)$,
	$\bar A[\phi] = \bar A[\sum_j \psi_j\phi]
	= \sum_j \bar A[\psi_j\phi]
	= \sum_j A_{\omega_j}[\psi_j\phi]
	= A[\phi]$.
	Hence, $\bar A=A$.
\end{proof}

\begin{exercise}
	Show that a partition of unity exists for every open cover of an open subset of $\R^n$.
	In other words, if $\Omega\subset\R^n$ is open and $\scr U$ is a collection of open subsets of $\Omega$ such that $\bigcup\scr U = \Omega$, 
	then there is a countable family of functions $\{\psi_j\}_{j\in\N}\subset C^\infty_c(\Omega)$ 
	such that
	\begin{enumerate}
%	\item $\psi_j\in C^\infty_c(\Omega)$ for all $j$;
	\item $0\le \psi_j(x)\le 1$ for all $j$ and all $x\in\Omega$
	\item for every $x\in\Omega$ the set $\{j:\psi_j(x)\neq0\}$ is finite; 
	\item $\sum_j\psi_j(x) = 1$ for all $x\in \Omega$;
	\item for every $j\in\N$ there is $\omega_j\in\scr U$ such that $\spt(\psi_j)\subset\omega_j$.
	\end{enumerate}
	
	Moreover, show that, if $K\subset \Omega$ is compact, then there is a finite set $J\subset\N$ such that $\sum_{j\in J}\psi_j(x) = 1$ for all $x\in K$.
	
	{\it Hint:} Rudin has a proof.
\end{exercise}

\begin{exercise}
	Let $A\in\scr D'(\Omega)$ and $\scr U$ an open cover of $\Omega$.
	Show that, if $\bar A\in\scr D'(\Omega)$ is such that $\bar A = A$ on $\omega$, for every $\omega\in\scr U$, then $\bar A=A$.

	The meaning of this exercise is as follows:
	If we start from some $A\in\scr D'(\Omega)$, then we can localize $A$ to each open subset of some given cover of $\Omega$, and then we can recover $A$ from Proposition~\ref{prop67eaabf1}.
\end{exercise}

\begin{corollary}\label{cor67ea99a6}
	If $A\in \scr D'(\Omega)$ is such that for every $x\in\Omega$ there is $\omega\subset\Omega$ open with $x\in\omega$ and $A=0$ on $\omega$, then $A=0$ on $\Omega$.
\end{corollary}
\begin{proof}
	Define $\scr U$ as the open cover of $\Omega$ given by the open sets $\omega\subset\Omega$ with $A=0$ on $\omega$.
	Proposition~\ref{prop67eaabf1} claims that there exists a unique distribution in $\scr D'(\Omega)$ that is equal to $A$ on each $\omega\in\scr U$.
	Since $0$ is one of such distributions, uniqueness implies $A=0$.
\end{proof}

\begin{remark}
	Corollary~\ref{cor67ea99a6} might look confusionally trivial.
	To understand it better, it is useful to think of a situation (outside the world of distributions, of course) where locality fails.
	Here is one example.
	
	Consider the space $\scr C = \{\alpha\in\Omega^1(\R^2\setminus\{0\}) : \dd \alpha = 0 \}$ of closed 1-forms on the plane punctured plane.
	We consider (continuous) linear functionals $\scr C\to\C$.
	Let $\bb S^1 = \{z\in\R^2: |z|=1\}$ the the unit circle.
	Define $A[\alpha] := \int_{\bb S^1} \alpha = \int_0^{2\pi} \langle \alpha(e^{it}) | ie^{it} \rangle \did t$ 
	(we identify $\C$ with $\R^2$ for notational purposes).
	Notice that $A$ is not zero: for example, $A[x\dd y - y \dd x]\neq0$.
	
	Now, we claim that, for every $x\in\R^2$, there exists $\omega\subset\R^2$ open neighborhood of $x$ such that $A[\alpha]=0$ for all $\alpha\in\scr C$ with $\spt(\alpha)\subset\omega$.
	Indeed, if $x\notin\bb S^1$, then we can just take $\omega$ with $\omega\cap\bb S^1=\emptyset$.
	If $x\in\bb S^1$, then we can take $\omega = B(x,1/2)$, the ball of radius $1/2$ and center $x$.
	If $\alpha$ is a closed 1 form on $\R^2\setminus \{0\}$, the integral of $\alpha$ over each contractible loop is $0$.
	So, if $\spt(\alpha)\subset\omega$, and if $y,z\in\R^2$ are the two extremal points of the arc $\bb S^1\cap \omega$, then the integral of $\alpha$ over $\bb S^1\cap \omega$ is equal to the integral of $\alpha$ along another path from $y$ to $z$ that does not intersect the support of $\alpha$ (such as the boundary of $\omega$ itself).
	It follows that the integral of $\alpha$ along $\bb S^1\cap \omega$ is zero.
	This shows the claim.
\end{remark}

%%%%%%%%%%%%%%%%%%%%%%%%%%%%%%%%%%%%%%%%%%%%%%%%%%%%%%%%%%%%%%%
\subsection{Support of a distribution}\label{subs67e1bd60}
The support of $A\in\scr D'(\Omega)$ is defined as the complement of the set
\begin{equation}
	\Omega\setminus\spt(A) = \bigcup\{U\subset\Omega\text{ open, such that }A\phi=0\ \forall\phi\in\scr D(\Omega) \}.
\end{equation}

\begin{proposition}
	Let $A\in\scr D'(\Omega)$ and $\phi\in\scr D(\Omega)$.
	If $\spt\phi\subset\Omega\setminus\spt A$, then $A\phi=0$.
\end{proposition}
\begin{proof}
	Let $\scr U=\{\omega\subset\Omega:A=0\text{ in }\omega\}$.
	Then $\scr U$ is a cover of $\Omega\setminus\spt A$, by definition of support of a distribution.
	Let $\{\psi_j\}_{j\in\N}$ be a partition of unity subordinated to $\scr U$.
	Since $\spt\phi$ is a compact subset of $\Omega\setminus\spt A$,
	there is $N$ such that $\phi = \sum_{j=1}^N \psi_j \phi$.
	So,
	$A\phi = \sum_{j=1}^N A[\psi_j \phi] = 0$.
	
	We can give another proof using Proposition~\ref{prop67eaabf1}:
	Let $V = \Omega\setminus\spt A$.
	Then $V$ is open and $\scr U$ is an open cover of $V$.
	By the very definition of $\scr U$, each restriction $A_\omega$ of $A$ to $\omega\in\scr U$ is zero.
	By Proposition~\ref{prop67eaabf1}, these $A_\omega\in\scr D'(\omega)$ (which are zero) are the restrictions of a unique distribution on $V$.
	Since $0$ is a distribution on $V$ so that $0=A_\omega$ on $\omega$ for each $\omega\in\scr U$, then $A=0$ on $V$ by uniqueness.	
\end{proof}

A few examples of support:
\begin{itemize}
\item
If $A_f\in\scr D'(\Omega)$ is the distribution associated to $f\in C(\Omega)$, then $\spt(A_f) = \spt(f)$.
\item
$\spt(\delta_p) = \{p\}$.
\end{itemize}

\begin{exercise}
	Show the following statement:
	if $A\in\scr D'(\Omega)$ and $f\in C^\infty(\Omega)$ are such that $\spt A\subset\{f=1\}$,
	then $fA = A$.
\end{exercise}

\begin{exercise}
	Show the following statement:
	if $A\in\scr D'(\Omega)$ and $f\in C^\infty(\Omega)$, then $\spt(fA) \subset \spt(f) \cap\spt(A)$.
	Is equality true?
	
	{\it Hint for question:} Try with $A=\delta_0$. % with $0\in\de\spt(f)$.
\end{exercise}

%%%%%%%%%%%%%%%%%%%%%%%%%%%%%%%%%%%%%%%%%%%%%%%%%%%%%%%%%%%%%%%
\subsection{Derivatives of the Dirac delta}\label{subs67e27313}

I don't think we will ever use the following statement, but it is important to know it.

\begin{proposition}\label{prop67fe6632}
	Suppose $A\in\scr D'(\Omega)$, $p\in\Omega$ and $\spt(A) = \{p\}$.
	Then there is $N\in\N$ and constants $c_\alpha\in\C$ such that
	\begin{equation}\label{eq67e273a7}
		A = \sum_{|\alpha|\le N} c_\alpha \Diff^\alpha \delta_p ,
	\end{equation}
	where $\delta_p\phi=\phi(p)$ is the Dirac delta, see~\ref{subs67def74c}.
	
	Conversely, if $A$ is as in~\eqref{eq67e273a7}, then $\spt(A)=\{p\}$, unless $c_\alpha=0$ for all $\alpha$.
\end{proposition}
\begin{proof}
	See \cite[Thm.6.21]{MR1157815}.
\end{proof}
%\begin{proof}
%	Recall what is the Taylor expansion of a smooth function:
%	There are constants $\{T_\alpha \} _{\alpha\in\N}$ (depending only on $\alpha$ and $n$)
%	such that, if $\phi\in C^\infty(\Omega)$ and $N\in\N$, then
%	\begin{equation}\label{eq67ea37e5}
%		\phi(x) = \sum_{|\alpha|\le N} T_\alpha \Diff^\alpha\phi(p) (x-p)^\alpha + \phi_N(x) ,
%	\end{equation}
%	where $\phi_N\in C^\infty(\Omega)$ is a function such that $\Diff^\alpha \phi_N(p) = 0$ for $|\alpha|\le N$.
%	The identity~\eqref{eq67ea37e5} is called \emph{Taylor expansion of degree $N$} of $\phi$ at $p$.
%	The constants $T_\alpha$ can be computed, but we do not need their explicit value.
%	Recall also that $(\Diff^\alpha (x-p)^\beta)|_{x=p}$ is  $\alpha!$ if $\alpha=\beta$, zero otherwise.
%
%	Let $r,R>0$ be such that $B(p,r) \subset \bar B(p,R)\subset\Omega$.
%	Let $f\in C^\infty(\Omega)$ be such that $B(p,r)\subset\{f=1\}$ and $\spt(f) \subset B(p,R)$.
%	Since $\bar B(p,R)$ is compact and contained in $\Omega$, there are $N\in\N$ and $C>0$ such that
%	\begin{equation}
%		A[\phi] \le C \|\phi\|_{C^N}
%	\end{equation}
%	for all $\phi\in\scr D(B(p,R))$.
%	For $\alpha\in\N^n$ with $|\alpha|\le N$, define $c_\alpha = \frac{A[x^\alpha\cdot f]}{\alpha!}$.
%	We claim that the equality~\eqref{eq67e273a7} holds with these constants.
%	
%	For this, define $A_N = \sum_{|\alpha|\le N} c_\alpha \Diff^\alpha \delta_p$.
%	
%	Let $\phi\in\scr D(\Omega)$.
%	Let $\phi_N(x)$ be the polynomial of degree $N$ coming from the Taylor expansion of $\phi$ at $p$.
%	In other words,
%	\begin{equation}
%		\phi_N(x) = \sum_{|\alpha|\le N} T_\alpha \Diff^\alpha\phi(p) x^\alpha .
%	\end{equation}
%	Notice that $A_N[\phi] = A_N[f\cdot\phi_N]$.
%	\begin{align*}
%		A_N[\phi]  
%		&= \sum_{|\alpha|\le N} c_\alpha \Diff^\alpha \delta_p[\phi] \\
%		&= \sum_{|\alpha|\le N} c_\alpha (-1)^{|\alpha|}  \delta_p[\Diff^\alpha\phi] \\
%		&\overset{\eqref{eq67ea37e5}}= \sum_{|\alpha|\le N} c_\alpha (-1)^{|\alpha|}  T_\alpha \Diff^\alpha\phi(p) \alpha! \\
%		&= \sum_{|\alpha|\le N}  (-1)^{|\alpha|}  T_\alpha \Diff^\alpha\phi(p) A[x^\alpha\cdot f] \\
%		&= A\left[ \sum_{|\alpha|\le N}  (-1)^{|\alpha|}  T_\alpha \Diff^\alpha\phi(p) x^\alpha\cdot f \right] \\
%		&= A[f\cdot\phi_N] .
%	\end{align*}
%	Therefore, $A[\phi] - A_N[\phi] = A[\phi] - A[f_r\cdot\phi_N] = A[\phi - f_r\cdot\phi_N]$.
%	\begin{align}
%		A[\phi] - A_N[\phi]
%	\end{align}
%\end{proof}

%%%%%%%%%%%%%%%%%%%%%%%%%%%%%%%%%%%%%%%%%%%%%%%%%%%%%%%%%%%%%%%
\subsection{Distributions as derivatives of functions}\label{subs67e2746f}

I don't think we will ever use the following statement, but it is important to know it.

\begin{theorem}\label{thm67eaad41}
	For ever $A\in\scr D'(\Omega)$ there are continuous functions $g_\alpha:\Omega\to\C$ for $\alpha\in\N^n$ such that the infinite sum $\sum_\alpha g_\alpha$ is locally finite and
	\begin{equation}
		A = \sum_{\alpha\in\N^n} \Diff^\alpha g_\alpha .
	\end{equation}
\end{theorem}
\begin{proof}
	See \cite[Thm.6.28]{MR1157815}.
\end{proof}

\begin{exercise}\label{ex67eaae04}
	Use Theorem~\ref{thm67eaad41} to show the following:
	if $A\in\scr D'(\R)$, then there exists $g\in C(\R)$ and $N\in\N$ such that $A = g^{(N)}$ ($N$-th distributional derivative of $g$).
\end{exercise}

\begin{remark}
	Here is an example of the situation described by Theorem~\ref{thm67eaad41}.
	On $\R$, consider the Dirac delta at zero $\delta_0$.
	The theorem claims that there is a continuous function $g:\R\to\C$ and $N$ such that $g^{(N)}=\delta_0$, see also Exercise~\ref{ex67eaae04}.
	In fact, if $g(x) = \int_{-\infty}^x \one_{[0,+\infty)}(y) \did y$,
	then (Exercise!) $g^{(2)}=\delta_0$.
\end{remark}

%%%%%%%%%%%%%%%%%%%%%%%%%%%%%%%%%%%%%%%%%%%%%%%%%%%%%%%%%%%%%%%
\subsection{Convolution of functions}\label{subs67ea3f7c}
If $f,g:\R^n\to\C$ are measurable functions, we define
\begin{equation}\label{eq67e2849f}
	f*g (x) = \int_{\R^n} f(y) g(x-y) \did y ,
\end{equation}
whenever the integral is well defined, i.e., whenever $y\mapsto f(y) g(x-y)$ is integrable.
Notice that, as soon as $y\mapsto f(y) g(x-y)$ is integrable, then not only $f*g(x)$ is well defined, but also
\begin{equation}
	f*g(x) = \int_{\R^n} f(y) g(x-y) \did y = \int_{\R^n} f(x-y) g(y) \did y
	= g*f(x) .
\end{equation}

\begin{exercise}
	Show that, if $f\in L^1(\R^n)$ and $g\in L^\infty(\R^n)$, then $f*g(x)$ is well defined for all $x\in\R^n$ and in fact $f*g\in C^0_b(\R^n)$.
\end{exercise}
\begin{exercise}[Young's inequality]\label{ex67fbf157}
	Show that, if $f\in L^1(\R^n)$ and $g\in L^{p}(\R^n)$ with $p\in[1,\infty]$,
	then $f*g\in L^p(\R^n)$ and 
	\begin{equation}\label{eq6817203b}
		\|f*g\|_{L^p} \le \|f\|_{L^1} \|g\|_{L^p}
	\end{equation}
	
	{\it Hint.}
	By Hölder inequality, with $\frac1p+\frac1{p'}=1$,
	$\int |f(y)g(y-x)| \did y = \int |f(y)|^{1/p'}\cdot |f(y)|^{1/p} |g(y-x)| \did y
	\le (\int |f(y)| \did y)^{1/p'} \cdot (\int |f(y)| |g(y-x)|^{p} \did y)^{1/p}$.
	Therefore, $\int (f*g(x))^p \did x
	\le (\int |f(y)| \did y)^{p/p'} \cdot \int \int |f(y)| |g(y-x)|^{p} \did y \did x
	\le (\int |f(y)| \did y)^{p/p'} \cdot \int |g(y)|^p \did y \cdot \int |f(y)| \did y$.
\end{exercise}

\begin{exercise}\label{ex67ea4f16}
	Show that, if $f,g\in C^0(\R^n)$, then
	\begin{equation}
		\spt(f*g) \subset \spt(f) + \spt(g) .
	\end{equation}
	Can you find a case where equality holds? And where equality does not hold?
\end{exercise}

%%%%%%%%%%%%%%%%%%%%%%%%%%%%%%%%%%%%%%%%%%%%%%%%%%%%%%%%%%%%%%%
\subsection{Translations and inversion of a distribution}\label{sub67df26c0}
We will extend the definition of convolution to distributions.
To do so, we need two linear operators on distributions: here they are.

For $x\in\R^n$, define $\tau_x:\R^n\to\R^n$ by $\tau_x(y) = y-x$.
For $x\in\R^n$, define the continuous linear operator $\tau_x:\scr D\to\scr D$ by
\begin{equation}
	\tau_x\phi = \phi\circ\tau_x ,
\end{equation}
for all $\phi\in\scr D$.
In other words, $\tau_x\phi = \tau_x^*\phi$ is the pull-back via $\tau_x$.
By Proposition~\ref{prop67df0940}, we have a continuous linear operator $\tau_x:\scr D'\to \scr D'$ 
\begin{equation}
	\tau_xA[\phi] = A[\phi\circ\tau_x] .
\end{equation}
In other words, $\tau_x = (\tau_x^*)_*$ is the push forward induced by the pull back induced by $\tau_x$.
\[
\xymatrix{
\R^n \ar[rr]^{\tau_x}_{y\mapsto y-x} && \R^n \\
\scr  D && \scr D \ar[ll]_{\tau_x^*}^{\phi\mapsto \phi\circ\tau_x} \\
\scr D' \ar[rr]^{(\tau_x^*)_*}_{A\mapsto (\phi\mapsto A[\phi\circ\tau_x])} && \scr D'
}
\]
I think at this point the amount of confusion surpass the amount of information.

Next, define the continuous linear operator $\scr D\to\scr D$, $\phi\mapsto \check\phi$, by
\begin{equation}
	\check \phi : x\mapsto \phi(-x) .
\end{equation}
The notation $\check\phi$ will clash with the inverse of the Fourier transform: I do not know how to avoid this clash at the moment. 
For now, we follow Rudin's notation.

\begin{exercise}
	Show the relations
	\begin{align}
	\tau_y \tau_z &= \tau_{y+z} ;\\
	\label{eq67f29d18}
	(\tau_x\phi)^{\vee} &= \tau_{-x}\check\phi ;\\
	\label{eq67e2886e}
	\tau_x (\Diff^\alpha \phi)^\vee &= (-1)^{|\alpha|} \Diff^\alpha(\tau_x\check\phi) .
	\end{align}
\end{exercise}


%%%%%%%%%%%%%%%%%%%%%%%%%%%%%%%%%%%%%%%%%%%%%%%%%%%%%%%%%%%%%%%
\subsection{Convolution of a distribution}\label{subs67df2648}
Let $\phi\in C^\infty_c(\R^n) = \scr D$ and $u\in\scr D'$.
Define $\phi*u : \R^n\to\C$ by
\begin{equation}\label{eq67e28487}
	(u*\phi)(x) = u[\tau_x\check\phi] .
\end{equation}
Note that $\tau_x\check\phi(y) = \check\phi(y-x) = \phi(x-y)$.
Notice that, if $u$ is a function, then
\begin{equation}
	(u*\phi)(x)
	\overset{\eqref{eq67e28487}}= u[\tau_x\check\phi]
	= \int u(y) (\tau_x\check\phi)(y) \did y
	= \int u(y) \phi(-(y-x)) \did y
	\overset{\eqref{eq67e2849f}}= u*\phi(x) .
\end{equation}

\begin{exercise}
	Show that, if $u\in\scr D'$ and $\phi\in\scr D$, then
	\begin{equation}\label{eq67e28b61}
		u[\phi] = (u*\check\phi)(0) .
	\end{equation}
\end{exercise}

\begin{exercise}
	Show that, if $u\in\scr D'$ and $\phi\in\scr D$, then 
	\begin{equation}\label{eq67f2b0f5}
		\spt(u*\phi) \subset \spt(u) + \spt(\phi)
		= \{x+y: x\in\spt(u),\ y\in\spt(\phi) \} .
	\end{equation}
	{\it Solution.}
	Notice that, if $x\in\R^n$ and $\phi:\R^n\to\C$, then $\tau_x\check\phi(y) = \phi(x-y) \neq 0$ if and only if $x-y\in\spt(\phi)$, if and only if $y\in x-\spt(\phi)$.
	Therefore, $\spt(\tau_x\check\phi) = x-\spt(\phi)$.
	
	Let $x\in\R^n\setminus ( \spt(u) + \spt(\phi))$.
	Notice that, if $y \in \spt(\tau_x\check\phi)\cap \spt(u)$,
	then there is $z\in\spt(\phi)$ with $y=x-z\in\spt(u)$,
	thus $x = y+z \in \spt(u) + \spt(\phi)$.
	Since $x\notin \spt(u) + \spt(\phi)$, then $\spt(\tau_x\check\phi)\cap \spt(u)=\emptyset$.
	We conclude that $u*\phi(x) = u[\tau_x\check\phi] = 0$.
\end{exercise}

\begin{exercise}
	Show that, if $u\in\scr D'$, $\phi\in\scr D$ and $v\in\R^n$, then
	\begin{equation}\label{eq67ebace9}
		u*(\tau_v\phi) = \tau_v(u*\phi) .
	\end{equation}
\end{exercise}

\begin{exercise}
	Show that $\phi\mapsto u*\phi$ is linear.
\end{exercise}

\begin{exercise}
	Show that $\delta_0*\phi = \phi$ for every $\phi\in\scr D$.
	What is $\delta_v*\phi$?
\end{exercise}

\begin{proposition}\label{prop67ea5631}
	If $u\in\scr D'$ and $\phi\in\scr D$, then
	$u*\phi\in C^\infty(\R^n)$ and, for every $\alpha\in\N^n$,
	\begin{equation}\label{eq67e28779}
		\Diff^\alpha(u*\phi) = u*(\Diff^\alpha \phi) = (\Diff^\alpha u)*\phi .
	\end{equation}
\end{proposition}
\begin{proof}
	If $v\in\R^n$ and $v\in\R^n$, then
	\begin{align*}
	\Diff_v(u*\phi)(x)
	&= \lim_{h\to0} \frac{u*\phi(x+hv) - u*\phi(x)}{h} \\
	&= \lim_{h\to0} \frac{u[\tau_{x+hv}\check\phi] - u[\tau_x\check\phi]}{h} \\
	&= \lim_{h\to0} u_y\left[ \frac{ \phi(x+hv-y) - \phi(x-y) }{ h } \right] .
	\end{align*}
	Since $\phi(x+hv-y) - \phi(x-y)\to \Diff_v\phi(x-y)$ in $\scr D_y$ as $h\to0$, then 
	\begin{align*}
	\lim_{h\to0} u_y\left[ \frac{ \phi(x+hv-y) - \phi(x-y) }{ h } \right]
	= u_y[\Diff_v\phi(x-y)] 
	= u*\Diff_v\phi(x) .
	\end{align*}
	This shows that, for every $x\in\R^n$,
	\begin{equation}
		\Diff^j(u*\phi)(x) = (u*\Diff^j\phi)(x) .
	\end{equation}
	Iterating, we get the first part of~\eqref{eq67e28779}.
	
	Next, 
	\begin{align*}
	 	u*(\Diff^\alpha \phi)(x)
		&= u[\tau_x (\Diff^\alpha \phi)^\vee ]\\
		&\overset{\eqref{eq67e2886e}}= (-1)^{|\alpha|} u[\Diff^\alpha(\tau_x\check\phi)]\\
		&= \Diff^\alpha u[\tau_x\check\phi] \\
		&= ((\Diff^\alpha u)*\phi)(x) .
	\end{align*}
	We thus have completed the proof of~\eqref{eq67e28779}.
\end{proof}



\begin{proposition}\label{prop67f2992e}
	If $u\in\scr D'$ and $\phi,\psi\in\scr D$, then
	\begin{equation}\label{eq67e288dc}
		u*(\phi*\psi) = (u*\phi)*\psi .
	\end{equation}
\end{proposition}
\begin{proof}
	We use Lemma~\ref{lem67ea55c8} and Proposition~\ref{prop67ea5631}:
	\begin{align*}
	 	u*(\phi*\psi)(x)
		&= u[\tau_x(\phi*\psi)^\vee] \\
		&= u_y[(\phi*\psi)^\vee(y-x)] \\
		&= u_y[(\phi*\psi)(x-y)] \\
		&= u_y\left[ \int_{\R^n} \phi(x-y-z) \psi(z) \did z \right] \\
		&= u_y\left[ \lim_{h\to0} \sum_{z\in\Z^n} h^n \phi(x-y-hz) \psi(hz)  \right] \\
		[\text{by Lemma~\ref{lem67ea55c8}}]
		&=\lim_{h\to0} u_y\left[ \sum_{z\in\Z^n} h^n \phi(x-y-hz) \psi(hz)  \right] \\
		[\text{since the sum is finite}]
		&= \lim_{h\to0} \sum_{z\in\Z^n} h^n u_y[\phi(x-y-hz)] \psi(hz) \\
		[\phi(x-y-hz) = \tau_{x-hz}\check\phi(y)]
		&= \lim_{h\to0} \sum_{z\in\Z^n} h^n (u*\phi)(x-hz)] \psi(hz) \\
		[\text{by Lemma~\ref{lem67ea55c8} and Proposition~\ref{prop67ea5631}}]
		&= \int_{\R^n} (u*\phi)(x-z) \psi(z) \did z \\
		&= (u*\phi)*\psi(x) .
	\end{align*}
\end{proof}

\begin{lemma}\label{lem67ea55c8}
	Let $\phi\in \scr D$ and $\psi\in\scr E$.
	For $h>0$, define
	\begin{equation}
		\rho_h(x) = \sum_{z\in\Z^n} h^n \phi(x-hz) \psi(hz) .
	\end{equation}
	Then $\rho_h\in\scr E$ for all $h>0$ and
	$\lim_{h\to0} \rho_h = \phi*\psi$ in $\scr E$.
	
	Moreover, if $\psi\in\scr D$, then $\lim_{h\to0} \rho_h = \phi*\psi$ in $\scr D$.
\end{lemma}
\begin{proof}
	Clearly $\rho_h$ is smooth.
	The support of $\rho_h$ is contained in $\spt(\phi) + \spt(\psi)$ (cfr.~Exercise~\ref{ex67ea4f16}).
	Since the function $(x,y)\mapsto \phi(x-y)\psi(y)$ is continuous,
	then $\rho_h\to \phi*\psi$ uniformly on compact sets (see Exercise~\ref{ex67f3a281}).
	The derivatives of $\rho_h$ have the same form, i.e.
	\begin{equation}
		\Diff^\alpha\rho_h(x) = \sum_{z\in\Z^n} h^n \Diff^\alpha_x\phi(x-hz) \psi(hz) .
	\end{equation}
	So, $\Diff^\alpha\rho_h\to (\Diff^\alpha\phi)*\psi = \Diff^\alpha(\phi*\psi)$ uniformly.
%	By~\eqref{eq67e16d22}, 
	We obtain $\rho_h\to \phi*\psi$ in $\scr E$.
	
	If $\psi\in\scr D$, then $(x,y)\mapsto \phi(x-y)\psi(y)$ is uniformly continuous with compact support.
	Reasoning as above, we obtain $\rho_h\to \phi*\psi$ in $\scr D$.
\end{proof}

\begin{exercise}\label{ex67ea516c}
	In this exercise, you show that Riemann sums converge to the integral.
	
	Let $f:\R^n\to\C$ be a continuous and integrable function.
	(Integrable: $\int_{\R^n}|f(z)|\did z < \infty$).
	For $h>0$, define
	\begin{equation}
		F_h = \sum_{z\in\Z^n} h^n f(hz) .
	\end{equation}
	Show that $\lim_{h\to0}F_h = \int_{\R^n}f(z)\did z$.
\end{exercise}

\begin{exercise}\label{ex67ea5356}
	[To do while listening to Paganini's Caprice No. 24].
	Variation over Exercise~\ref{ex67ea516c}: 
	Let $f:\R^n\times\R^n\to\C$ be a uniformly continuous and integrable function.
	Define $F:\R^n\to\C$ by 
	\begin{equation}
		F(x) = \int_{\R^n} f(x,z) \did z .
	\end{equation}
	For $h>0$ and $x\in\R^n$, define
	\begin{equation}
		F_h(x) = \sum_{z\in\Z^n} h^n f(x,hz) .
	\end{equation}
	Show that $F_h\to F$ uniformly in $x$ as $h\to0$.
\end{exercise}

\begin{exercise}\label{ex67f3a281}
%	[To do while listening to Paganini's Caprice No. 24].
	Variation over Exercise~\ref{ex67ea5356}: 
	Let $f:\R^n\times\R^n\to\C$ be a \sout{uniformly} continuous and integrable function.
	Define $F:\R^n\to\C$ by 
	\begin{equation}
		F(x) = \int_{\R^n} f(x,z) \did z .
	\end{equation}
	For $h>0$ and $x\in\R^n$, define
	\begin{equation}
		F_h(x) = \sum_{z\in\Z^n} h^n f(x,hz) .
	\end{equation}
	Show that $F_h\to F$ uniformly \underline{on compact sets} in $x$ as $h\to0$.
\end{exercise}



%%%%%%%%%%%%%%%%%%%%%%%%%%%%%%%%%%%%%%%%%%%%%%%%%%%%%%%%%%%%%%%
\subsection{Smooth approximation of a distribution}\label{subs67df2a23}

Let $\rho\in C^\infty_c(\R^n)$ be such that $\spt(\rho)=B(0,1)$, $0\le \rho\le 1$, $\check\rho=\rho$, and $\int\rho\did x = 1$.
Define $\rho_\epsilon(x) = \rho(x/\epsilon)/\epsilon^n$.
We call the family $\{\rho_\epsilon\}_{\epsilon>0}$ an \emph{approximation of the identity on $\R^n$}, or \emph{a family of mollifiers}.
%The \emph{standard family of mollifiers} the one given by
For example, one can take
\begin{equation}
	\rho(x) = \begin{cases}
	k \exp\left(\frac{1}{|x|^2-1}\right) & \text{ if }|x|<1 \\
	0 & \text{ otherwise} ,
	\end{cases}
\end{equation}
where $k$ normalizes the integral.

\begin{exercise}
	Show that the function 
	\begin{equation}
		\rho(x) = \begin{cases}
		\exp\left(\frac{1}{|x|^2-1}\right) & \text{ if }|x|<1 \\
		0 & \text{ otherwise} ,
		\end{cases}
	\end{equation}
	is $C^\infty$-smooth on $\R^n$ and compute $\int_\R^n \rho(x)\did x$.
	Show also that $\rho$ is not analytic.
	Does it exist a family of analytic mollifiers?
	
	Finally, why have I put the square in the definition of $\rho$? Do you think we could do without?
	
	{\it Hint:} I have put the square to help you. The square is itself a hint.
\end{exercise}

\begin{proposition}\label{prop67e28cab}
	Let $\{\rho_\epsilon\}_{\epsilon>0}$ be an approximation of the identity on $\R^n$,
	$\phi\in\scr D$ and $u\in\scr D'$.
	Then
	\begin{align}
	\label{eq67e28a4f}
		\lim_{\epsilon\to0} \phi*\rho_\epsilon &= \phi \text{ in }\scr D , \\
	\label{eq67e28a56}
		\lim_{\epsilon\to0} u*\rho_\epsilon &= u \text{ in }\scr D' .\\
	\end{align}
\end{proposition}
\begin{proof}
	The first identity was an exercise, but here is my solution.
	Notice that $\spt(\rho_\epsilon) = B(0,\epsilon)$.
	Since $\phi\in C^1_c$, then there is $L$ such that $|\phi(x)-\phi(y)|\le L|x-y|$ for all $x,y$ (see Exercise~\ref{ex67ea5cf7}).
	Then, for every $x\in\R^n$,
	\begin{align*}
	 	|\phi*\rho_\epsilon(x) - \phi(x) | 
		&= \left|\int \phi(y) \rho_\epsilon(x-y) \did y - \phi(x) \int \rho_\epsilon(x-y) \did y \right| \\
		&\le \int |\phi(y) - \phi(x)| \rho_\epsilon(x-y) \did y \\
		&\le L\epsilon \int \rho_\epsilon(x-y) \did y 
		= L\epsilon.
	\end{align*}
	This shows that $\phi*\rho_\epsilon\to \phi$ uniformly.
	Since this holds for every $\phi\in\scr D$,
	we also have that $\Diff^\alpha(\phi*\rho_\epsilon) = \phi*\Diff^\alpha\rho_\epsilon \to \Diff^\alpha\phi$ uniformly, for every $\alpha\in\N^n$.
	Therefore, $\phi*\rho_\epsilon\to \phi$ in $\scr D$.
	
	For the second identity, 
	\begin{align}
	\lim_{\epsilon\to0} u*\rho_\epsilon[\phi]
	&\overset{\eqref{eq67e28b61}}= \lim_{\epsilon\to0} ((u*\rho_\epsilon)*\check\phi)(0) \\
	&\overset{\eqref{eq67e288dc}}= \lim_{\epsilon\to0} (u*(\rho_\epsilon*\check\phi))(0) \\
	[\text{because $\check\rho=\rho$}]
	&= \lim_{\epsilon\to0} (u*(\check\rho_\epsilon*\check\phi))(0) \\
	[\text{because $\check f*\check g = (f*g)^\vee$}]
	&= \lim_{\epsilon\to0} (u*(\rho_\epsilon*\phi)^\vee)(0) \\
	&\overset{\eqref{eq67e28b61}}= \lim_{\epsilon\to0} u[\rho_\epsilon*\phi] \\
	&\overset{\eqref{eq67e28a4f}}= u[\phi] .
	\end{align}
\end{proof}

\begin{exercise}\label{ex67ea5cf7}
	Let $\Omega\subset\R^n$ convex and $\phi\in C^1(\Omega)$ such that $L = \|\grad\phi\|_{L^\infty}<\infty$
	Show that, for every $x,y\in \Omega$, $|\phi(x)-\phi(y)|\le L|x-y|$.
	
	Question: what happens if we drop the hypothesis of $\Omega$ being convex?
\end{exercise}

\begin{corollary}
	The space $C^\infty(\R^n)$ is dense in $\scr D'$ (with respect to the topology of $\scr D'$).
\end{corollary}

%\begin{corollary}
%	If $\Omega\subset\R^n$ is open, then 
%	the space $C^\infty(\Omega)$ is dense in $\scr D'(\Omega)$ (with respect to the topology of $\scr D'(\Omega)$).
%\end{corollary}
%\begin{proof}
%	\TODO
%\end{proof}

\begin{exercise}
	Show that, if $\Omega\subset\R^n$ is open, then 
	the space $C^\infty(\Omega)$ is dense in $\scr D'(\Omega)$ (with respect to the topology of $\scr D'(\Omega)$).
\end{exercise}

\begin{exercise}
	Is $\scr D(\Omega)$ dense in $\scr D'(\Omega)$? (Try at least for $\Omega=\R^n$).
	
	{\it Hint:}
	Take $A[\phi] = \int \phi\did x$ and try to approximate $A$ with functions in $C^\infty_c(\Omega)$.
\end{exercise}

%%%%%%%%%%%%%%%%%%%%%%%%%%%%%%%%%%%%%%%%%%%%%%%%%%%%%%%%%%%%%%%
\subsection{Constancy theorem}\label{subs67df2829}

\begin{theorem}[Constancy theorem]\label{thm68346c5c}
	If $u\in\scr D'(\Omega)$ is such that $\de_ju = 0$ for all $j\in\{1,\dots,n\}$,
	then $u$ is a constant function, that is, there is $c\in\C$ such that $u[\phi] = c\int_\Omega\phi(x)\did x$ for all $\phi\in\scr D(\Omega)$.
\end{theorem}
\begin{proof}
	Let $\{\rho_\epsilon\}_{\epsilon>0}$ be an approximation of the identity on $\R^n$.
	Then $\de_j(u*\rho_\epsilon) = (\de_ju)*\rho_\epsilon = 0$, therefore $u*\rho_\epsilon = c_\epsilon \in\C$.
	Moreover, by Proposition~\ref{prop67e28cab},
	$u[\phi] = \lim_{\epsilon\to0} u*\rho_\epsilon[\phi]
	=  \lim_{\epsilon\to0} c_\epsilon \int \phi(y)\did y$.
	Therefore, the is $c\in\C$ with $c=\lim_{\epsilon\to0} c_\epsilon$ and $u[\phi] = c\int \phi(y)\did y$, that is, $u$ is a constant. 
\end{proof}




%%%%%%%%%%%%%%%%%%%%%%%%%%%%%%%%%%%%%%%%%%%%%%%%%%%%%%%%%%%%%%%
%%%%%%%%%%%%%%%%%%%%%%%%%%%%%%%%%%%%%%%%%%%%%%%%%%%%%%%%%%%%%%%
%\newpage
%%\section{Distributions with compact support}\label{sec67e28d62}

%%%%%%%%%%%%%%%%%%%%%%%%%%%%%%%%%%%%%%%%%%%%%%%%%%%%%%%%%%%%%%%
\subsection{The space of all smooth functions as a Fréchet space}\label{subs67eb97bc}
We define $\scr E(\Omega)$ as the vector space $C^\infty(\Omega)$ of smooth functions $\Omega\to\C$, endowed with the family of quasinorms
$p_{N,K}:\scr E(\Omega)\to [0,+\infty)$, for $K\Subset\Omega$ and $N\in\N$, where
\begin{equation}
	p_{N,K} (f) = \sup\{|\Diff^\alpha f(x)|:x\in K,\ |\alpha|\le N\} = \|f\|_{C^N(K)} .
\end{equation}
The quasinorms determine a topology on $\scr E(\Omega)$, where $f_j\to f_\infty$ in 
\begin{equation}
	f_j\to f_\infty \text{ in }\scr E(\Omega)
	\quad\IFF\quad
	\begin{array}{c}
	\forall K\Subset \Omega,\ \forall N\in\N\\
	\lim_{j\to\infty} p_{N,K} (f_j-f_\infty) = 0 .
	\end{array}
\end{equation}
The topological vector space $\scr E(\Omega)$ is a so-called \emph{Fréchet space}.

\begin{exercise}
	Check whether the topology on $\scr E(\Omega)$ is the the initial topology induced by the maps $C^\infty(\Omega)\into C^m(\Omega)$.
	
	{\it Hint:} I actually don't know if this is true. So, please send me an email with the result, if you don't mind.
\end{exercise}

Notice that $\scr D(\Omega)\subset\scr E(\Omega)$ as sets, but also $\scr D(\Omega)\into\scr E(\Omega)$ continuously.
Indeed, if $\phi_j\to0$ in $\scr D(\Omega)$, then $\phi_j\to0$ in $\scr E(\Omega)$, 
and thus the immersion is continuous by Proposition~\ref{prop67e26a2f}.

However, $\scr D(\Omega)$ is dense in $\scr E(\Omega)$ with the topology of $\scr E(\Omega)$.

\begin{exercise}
	Show that $\scr D(\Omega)$ is dense in $\scr E(\Omega)$.
\end{exercise}

\begin{exercise}
	Find a sequence $\phi_j\in\scr D(\Omega)$ that converges to some $f$ in $\scr E(\Omega)$ but it does not converge in $\scr D(\Omega)$.
	
	{\it Just as a note:} there is a notion of ``Cauchy sequence'' for topological vector spaces (see Rudin's book).
	With such a notions available, one could check that $\scr D(\Omega)$ is complete in its own topology, but that its completion in the topology of $\scr E(\Omega)$ is $\scr E(\Omega)$. 
	This should clarify the situation.
\end{exercise}

%%%%%%%%%%%%%%%%%%%%%%%%%%%%%%%%%%%%%%%%%%%%%%%%%%%%%%%%%%%%%%%
\subsection{The dual space of $\scr E(\Omega)$...}\label{subs67e2723b}
Let $\scr E'(\Omega)$ be the topological dual of $\scr E(\Omega)$, that is, the space of continuous linear functionals $\scr E(\Omega)\to\C$.
A linear functional $A:\scr E(\Omega)\to\C$ is continuous if and only if, whenever $f_j\to f_\infty$ in $\scr E(\Omega)$, then $A[f_j] \to A[f_\infty]$ in $\C$.
%\TODO: Why?

The topology on $\scr E'(\Omega)$ is the usual weak* topology,
that is, point-wise convergence:
\begin{equation}
	E_j\to E_\infty \text{ in }\scr E'(\Omega)
	\quad\IFF\quad
	\forall f\in\scr E(\Omega)\ \lim_{j\to\infty} E_j[f] = E_\infty[f] .
\end{equation}

%%%%%%%%%%%%%%%%%%%%%%%%%%%%%%%%%%%%%%%%%%%%%%%%%%%%%%%%%%%%%%%
\subsection{...is made of distributions...}\label{subs67eb9b4b}

Since $\scr D(\Omega)\subset\scr E(\Omega)$, then each $E\in\scr E(\Omega)$ is also a linear functional $E:\scr D(\Omega)\to\C$.
Since $\scr D(\Omega)\into \scr E(\Omega)$ is continuous, or, otherwise said, since $\phi_j\to\phi_\infty$ in $\scr D(\Omega)$ implies $\phi_j\to\phi_\infty$ in $\scr E(\Omega)$, then also the restriction of $E$ to $\scr D(\Omega)$ is continuous.
In other words, $E\in\scr D(\Omega)$.

Moreover, if $E_j\to E_\infty$ in $\scr E'(\Omega)$, then clearly $E_j\to E_\infty$ in $\scr D'(\Omega)$.
So, we can say $\scr E'(\Omega)\subset\scr D'(\Omega)$.

%%%%%%%%%%%%%%%%%%%%%%%%%%%%%%%%%%%%%%%%%%%%%%%%%%%%%%%%%%%%%%%
\subsection{...with compact support.}\label{subs67eb9b5f}
It remains to characterize the distributions $A\in\scr D'(\Omega)$ that belong to $\scr E'(\Omega)$:

\begin{proposition}\label{prop67eb9ee1}
	\begin{equation}
		\scr E'(\Omega) = \{A\in\scr D'(\Omega)\text{ with compact support}\} .
	\end{equation}
\end{proposition}
\begin{proof}
	\ci
	From~\ref{subs67eb9b4b}, we know that $\scr E'(\Omega)\subset\scr D'(\Omega)$.
	We need to show that, if $A\in\scr E'(\Omega)$, then $\spt(A)$ is compact.
	Arguing by contradiction, suppose this is not true, that is, that there is $A\in\scr E'(\Omega)$ with $\spt(A)$ not compact.
	Let $\{K_j\}_{j\in\N}$ be a countable increasing sequence of compact sets $K_j\subset\Omega$ 
	such that $\Omega = \bigcup_{j\in\N} K_j$.
	Since $\spt(A)$ is not compact, then, for every $j\in\N$ there is $\phi_j\in\scr D(\Omega)$ with $\spt(\phi_j)\cap K_j=\emptyset$ and $A[\phi_j]=1$.

	Notice that $\phi_j\to0$ in $\scr E(\Omega)$:
	indeed, if $K\subset\Omega$ is compact, then there is $k\in\N$ such that $K\subset K_k$ and thus $\phi_j|_K=0$ for all $j>k$;
	hence, for every $K\subset\Omega$ compact and every $N\in\N$, 
	$\lim_{j\to\infty} p_{N,K}(\phi_j)= 0$.
	
	But we have assumed that $A\in\scr E'(\Omega)$, and therefore we should have $\lim_{j\to\infty} A[\phi_j] = 0$, in contradiction with $A[\phi_j]=1$ for all $j$.
	We conclude that $A$ must have compact support.
	
	\ic
	This inclusion does not make sense when taken literally:
	if $A\in\scr D'(\Omega)$, then $A$ is a linear map $\scr D(\Omega)\to\C$, while $\scr E(\Omega)$ is a larger space than $\scr D(\Omega)$.
	So, $A$ is NOT a function $\scr E(\Omega)\to\C$, taken as it is.
	However, we need to consider that $\scr D(\Omega)$ is dense in $\scr E(\Omega)$:
	so, if there is a continuous linear functional $E:\scr E(\Omega)\to\C$ whose restriction to $\scr D(\Omega)$ is $A$, then $E$ is uniquely determined by $A$.
	
	This is to say that the correct interpretation of the inclusion $\ic$ is:
	{\it every $A\in\scr D'(\Omega)$ with compact support extends uniquely to a continuous linear functional $\bar A:\scr E(\Omega)\to\C$.}
	
	Let $A\in\scr D'(\Omega)$ be with compact support.
	Fix some $\psi\in C^\infty_c(\Omega)$ such that $\spt A \subset{\rm interior}\{\psi=1\}$, which exists because $\spt(A)$ is compact.
	Define $A_\psi:\scr E(\Omega)\to\C$ by $A_\psi f = A[\psi f]$.
	Notice that, firstly, $A_\psi\in\scr E'(\Omega)$:
	indeed, if $f_j\to f_\infty$ in $\scr E(\Omega)$, then $\psi f_j\to \psi f_\infty$ in $\scr D(\Omega)$, and thus $A_\psi[f_j] \to A_\psi[f_\infty]$ in $\C$.
	
	Moreover, if $\phi\in\scr D(\Omega)$, then $\spt((1-\psi)\phi) \cap\spt(A)=\emptyset$ and thus
	$A[\phi] = A[\psi\phi + (1-\psi)\phi] = A[\psi\phi] = A_\psi[\phi]$.
	We conclude that the restriction of $A_\psi$ to $\scr E(\Omega)$ is $A$.
%	
%	 and, secondly, $A_\psi$ does not depend on the choice of $\psi$.
%	Indeed, if $\psi_1,\psi_2$ are such that $\spt A \subset\{\psi=1\}$,
%	then $\spt(\psi_1-\psi_2)\cap \spt A=\emptyset$ and thus
%	$A_{\psi_1} f - A_{\psi_2}f = A[(\psi_1-\psi_2)f] = 0$.
%	This shows that $A$ extends to $A\in\scr E'(\Omega)$. 
%	??
%	
\end{proof}

\begin{exercise}
	In the proof of Proposition~\ref{prop67eb9ee1}, we have defined the extension of $A\in\scr D'(\Omega)$ to $\scr E(\Omega)$ as $A_\psi[f] = A[\psi f]$,
	where $\psi\in C^\infty_c(\Omega)$ such that $\spt A \subset{\rm interior}\{\psi=1\}$.
	It looks like this extension depends on the choice of $\psi$.
	Does it?
%	Is $A_\psi$ the extension of $A$ by continuation?
%	Moreover, if $\{\phi_j\}_{j\in\N}\subset\scr D(\Omega)$ is a sequence converging in $\scr E(\Omega)$ to some $f\in\scr E(\Omega)$, is it true that $A[\phi_j]
%	Show that the $A_\psi$ does not depend on the choice of $\psi$, directly.
\end{exercise}

\begin{exercise}\label{ex67f29a5b}
	Show that all distributions in $\scr E'(\Omega)$ have finite order.
	More precisely, if $u\in\scr E'(\Omega)$, then there are $N\in\N$ and $C\in\R$ such that, for every $f\in\scr E(\Omega)$,
	\begin{equation}
		|u[f]| \le C \|f\|_{C^N(\spt(u))} .
	\end{equation}
\end{exercise}

%%%%%%%%%%%%%%%%%%%%%%%%%%%%%%%%%%%%%%%%%%%%%%%%%%%%%%%%%%%%%%%
\subsection{Convolution of distributions with compact support}
We have started with the convolution $u*v$ of functions $u$ and $v$.
We have extended this operation to convolution $u*v$ where $u\in\scr D'$ is a distribution and $v\in\scr D$ is a smooth function with compact support.
As it happens with the product, convolution is not defined for arbitrary pairs of distributions.
However, we can still push our definition of convolution $u*v$ to the case when $u,v\in\scr D'$ and {\it at least one of the distributions has compact support}.
This is done in three steps.

\boxed{\textsc{Step 1:}}
If $u\in\scr E'$ and $f\in\scr E$, then we define
\begin{equation}\label{eq67f21bf6}
	u * f (x) = u[\tau_x\check f] ,
\end{equation}
exactly as we did for distributions in $\scr D'$.
So, if $u\in\scr D'$ has compact support (i.e., $u\in\scr E'$), then $u*\phi$ is defined not only for $\phi\in\scr D = C^\infty_c(\R^n)$, but also for $\phi\in\scr E = C^\infty(\R^n)$.

\begin{proposition}\label{prop67f21c33}
	If $u\in\scr E'$ and $f\in\scr E$, then $u*f\in\scr E$.
	Moreover
	\begin{align}
		\label{eq67f298b5}
		\spt(u*f) &\subset \spt(u) + \spt(f), \text{ and } \\
		\label{eq67f29911}
		\Diff^\alpha(u*f) &= u*(\Diff^\alpha f) = (\Diff^\alpha u)*f,
		\qquad \forall\alpha\in\N^n .
	\end{align} 
	In particular, if $f\in\scr D$, then $u*f\in\scr D$.
\end{proposition}
\begin{proof}
	The proof is the same as for Proposition~\ref{prop67ea5631}.
	In fact, I wonder if there is a way to merge the two proofs into one.
\end{proof}

%%%%%%%%%%%%%%%%%%%%%%%%%%%%%%%%%%%%%%%%%%%%%%%%%%%%%%%%%%%%%%%
\boxed{\textsc{Step 2:}}
We characterize continuous linear operators $\scr D\to\scr E$ that commute with translations.
First of all, notice that, if $u\in\scr D'$, then $L:\phi\mapsto u*\phi$ is a linear function $\scr D\to\scr E$.
Moreover, by~\eqref{eq67ebace9}, $L$ commutes with translation, i.e., $L\circ\tau_v = \tau_v\circ L$ for all $v\in\R^n$.
The next Lemma~\ref{lem67ebad94} shows that $L$ is in fact continuous.
Theorem~\ref{thm67ebada4} will finally prove that all continuous linear operators $\scr D\to\scr E$ that commute with translations are of this form.

\begin{lemma}\label{lem67ebad94}
	If $\phi_j\to\phi_\infty$ in $\scr D$ and $u\in\scr D'$, 
	then $u*\phi_j\to u*\phi_\infty$ in $\scr E$.
%	
%	Similarly, if $\phi_j\to\phi_\infty$ in $\scr E$ and $u\in\scr E'$, 
%	then $u*\phi_j\to u*\phi_\infty$ in $\scr E$.
\end{lemma}
\begin{proof}
%	The proof of the second part is very similar to the proof of the first part,
%	so we will leave it as an exercise.
%	
	Let $\phi_j\to\phi_\infty$ in $\scr D$.
%	We can assume $\phi_\infty=0$: 
	We need to show that, if $\phi_j\to\phi_\infty$ in $\scr D$, then $u*\phi_j\to u*\phi_\infty$ in $\scr E$,
	that is, $\Diff^\alpha(u*\phi_j)\to \Diff^\alpha(u*\phi_\infty)$ uniformly on compact sets, for every $\alpha\in\N^n$.
	In fact, since $\Diff^\alpha(u*\phi_j) = u*(\Diff^\alpha\phi_j)$ by~\eqref{eq67e28779}, 
	we only need to show that, if $\phi_j\to\phi_\infty$ in $\scr D$, then $u*\phi_j\to u*\phi_\infty$ uniformly on compact sets.
%	Since $\Diff^\alpha(u*\phi_j) = u*(\Diff^\alpha\phi_j)$ and since $\{\phi_j\}_{j\in\N}$ is an arbitrary sequence, we then conclude that,
%	Since $\Diff^\alpha(u*\phi_j) = u*(\Diff^\alpha\phi_j)$ and since $\{\phi_j\}_{j\in\N}$ is an arbitrary sequence, we only need to show that convergence for $\alpha=0$.
	
	Let $\phi_j\to\phi_\infty$ in $\scr D$.
	Fix a compact set $K\subset\R^n$.
	Let $\tilde K\subset\R^n$ be a compact such that $\spt(\phi_j)\subset \tilde K$ for all $j$.
	Notice that, if $x\in K$, then $x-y\in\tilde K$ if and only if $y\in K-\tilde K = \{a-b: a\in K,\ b\in\tilde K\}$.
	So, for every $x\in K$ and $\phi\in\scr D$ with $\spt(\phi)\subset\tilde K$, 
	then $\tau_x\check\phi$ is supported in $K-\tilde K$, which is compact.
	
	By Proposition~\ref{prop67defc54}, there are $C\in\R$ and $N\in\N$ such that $|u[\phi]| \le C\|\phi\|_{C^N}$ for all $\phi\in\scr D$ with $\spt(\phi)\subset K-\tilde K$.
%	Notice that, if $\phi$ is supported in $\tilde K$ and $x\in K$, then 
%	$\tau_x\check\phi(y) =  \phi(x-y) = 0$ if $x-y\notin \tilde K$, that is, $y\notin \tilde K - x$.
%	
	Then, for every $x\in K$,
	\begin{align*}
	|u*\phi_j(x) - u*\phi_\infty(x)|
	&= |u[\tau_x(\check \phi_j)] - u[\tau_x(\check \phi_\infty)]| \\
	&\le C \|\tau_x(\check \phi_j) - \tau_x(\check \phi_\infty)\|_{C^N} 
	= C\| \phi_j - \phi_\infty\|_{C^n} .
	\end{align*}
	Therefore, $u*\phi_j\to u*\phi_\infty$ uniformly on $K$.
%	$|u*\phi_j\to u(x)| = |u[\tau_x(\check \phi_j)]| \le C \|\tau_x(\check \phi_j)\|_{C^N} = C\|\phi_j\|_{C^n}$.
%	Therefore, 
\end{proof}
\begin{lemma}\label{lem67f21e2a}
	If $\phi_j\to\phi_\infty$ in $\scr E$ and $u\in\scr E'$, 
	then $u*\phi_j\to u*\phi_\infty$ in $\scr E$.
\end{lemma}
\begin{proof}
	The proof is very similar to the proof of Lemma~\ref{lem67ebad94},
	so we will leave it as an exercise, see Exercise~\ref{ex67f21e52}.
\end{proof}

\begin{exercise}\label{ex67f21e52}
	Show that, if $\phi_j\to\phi_\infty$ in $\scr E$ and $u\in\scr E'$, 
	then $u*\phi_j\to u*\phi_\infty$ in $\scr E$.
\end{exercise}

\begin{lemma}\label{lem67f2202d}
	If $\phi_j\to\phi_\infty$ in $\scr D$ and $u\in\scr E'$,
	then $u*\phi_j\to u*\phi_\infty$ in $\scr D$.
\end{lemma}
\begin{proof}
	If $\phi_j\to\phi_\infty$ in $\scr D$, then there is a compact $K\subset\R^n$ such that $\spt(\phi_j)\subset K$ for all $j$ and $\Diff^\alpha\phi_j\to\Diff^\alpha\phi_\infty$ uniformly on $K$ for all $\alpha\in\N^n$.
	Then $\spt(u*\phi_j)\subset \spt(u)+K = \tilde K$ by Proposition~\ref{prop67f21c33}, for all $j$.
	By Lemma~\ref{lem67f21e2a}, $u*\phi_j\to u*\phi_\infty$ in $\scr E$.
	In particular, $\Diff^\alpha(u*\phi_j)\to \Diff^\alpha(u*\phi_\infty)$ uniformly on $\tilde K$ for all $\alpha\in\N^n$.
	This shows that $u*\phi_j\to u*\phi_\infty$ in $\scr D$.
\end{proof}

\begin{theorem}[{\cite[Thm.4.2.1]{MR717035}}]\label{thm67ebada4}
	Let $L$ be a linear map from $\scr D = C^\infty_c(\R^n)$ to $C^0(\R^n)$.
	The following statements are equivalent:
	\begin{enumerate}
	\item $L$ is (sequentially) continuous, i.e., if $\phi_j\to\phi_\infty$ in $\scr D$ then $L\phi_j\to L\phi_\infty$ in $C^0(\R^n)$,
		and commutes with translations, i.e., $L[\tau_x \phi] = \tau_x L\phi$ for all $x\in\R^n$;
	\item
	there is $u\in \scr D'$ such that $L\phi = u*\phi$ for all $\phi\in\scr D$.
	\end{enumerate}
	Moreover, if the above conditions are met, then 
	the distribution $u$ is unique,
	and $L$ is in fact a continuous linear operator $\scr D\to\scr E$.
\end{theorem}
\begin{proof}
	$(1)\THEN(2)$:
	Define $u:\scr D\to\R$ by $u[\phi] = L[\check\phi](0)$.
	Since $u$ is the composition of the continuous linear functions $\phi\mapsto\check\phi$, $\psi\mapsto L[\psi]$, $f\mapsto f(0)$, 
	then $u$ is linear and continuous, that is, $u\in\scr D'$.
	Moreover,
	\begin{align*}
		u*\phi(x) 
		&= u[\tau_x\check\phi] \\
		&\overset{\eqref{eq67f29d18}}= u[(\tau_{-x}\phi)^\vee] \\
		&= L[((\tau_{-x}\phi)^\vee)^\vee](0) \\
		&= L[\tau_{-x}\phi](0) \\
		&= \tau_{-x}L[\phi](0) \\
		&= L[\phi](x) .
	\end{align*}
	For the uniqueness, notice that, if $u_1,u_2\in\scr D'$ are such that $L\phi=u_1*\phi=u_2*\phi$ for all $\phi\in\scr D$,
	then, for all $\phi\in\scr D$, $u_1[\phi] = u_1*\phi(0) = u_2*\phi(0) = u_2[\phi]$,
	that is, $u_1=u_2$.
	
	$(2)\THEN(1)$:
	By Proposition~\ref{prop67ea5631}, we know that $\phi\mapsto u*\phi$ is a linear map $\scr D\to\scr E$.
	Next, by Lemma~\ref{lem67ebad94}, we know that this map is continuous.
	Finally, the fact that this map commutes with translations has been proven in~\eqref{eq67ebace9} (which was an exercise).
\end{proof}

%%%%%%%%%%%%%%%%%%%%%%%%%%%%%%%%%%%%%%%%%%%%%%%%%%%%%%%%%%%%%%%
\boxed{\textsc{Step 3:}}
If $u_1,u_2\in\scr D'$ are distributions, one of which has compact support, then,
 for every $\phi\in\scr D$,
\begin{equation}\label{eq67ebaf4c}
	L[\phi] = 
	u_1 * (u_2*\phi) 
\end{equation}
is a well defined element of $\scr E$ and the map $\phi\mapsto L\phi$ is continuous.
Indeed, we have two cases:
\begin{enumerate}
\item 
	If $u_1\in\scr E'$, then $u_2\in\scr D'$; so, $u_2*\phi$ is defined by~\eqref{eq67e28487}, $u_2*\phi\in\scr E$ by Proposition~\ref{prop67ea5631}, and thus~\eqref{eq67ebaf4c} is defined by~\eqref{eq67f21bf6}.
	Moreover, if $\phi_j\to\phi_\infty$ in $\scr D(\Omega)$, then $u_2*\phi_j\to u_2*\phi_\infty$ in $\scr E$ by Lemma~\ref{lem67ebad94},
	and thus $u_1 * (u_2*\phi_j) \to u_1 * (u_2*\phi_\infty)$ in $\scr E$ by Lemma~\ref{lem67f21e2a}.
\item 
	If $u_2\in\scr E'$, then $u_1\in\scr D'$; so, $u_2*\phi$ is defined by~\eqref{eq67f21bf6} and $u_2*\phi\in\scr D$ by Proposition~\ref{prop67f21c33}, and thus~\eqref{eq67ebaf4c} is defined by~\eqref{eq67e28487}.
	Moreover, if $\phi_j\to\phi_\infty$ in $\scr D(\Omega)$, then $u_2*\phi_j\to u_2*\phi_\infty$ in $\scr D$ by Lemma~\ref{lem67f2202d},
	and thus $u_1 * (u_2*\phi_j) \to u_1 * (u_2*\phi_\infty)$ in $\scr E$ by Lemma~\ref{lem67ebad94}.
\end{enumerate}
%Moreover, thanks to Lemma~\ref{lem67ebad94}, $\phi\mapsto L\phi$ is continuous.
We conclude that, by Theorem~\ref{thm67ebada4}, there is a unique $u\in\scr D'$ such that $L\phi = u*\phi$ for all $\phi\in\scr D$.
So, we have the definition:
If $u_1,u_2\in\scr D'$ one of which has compact support, then
\begin{equation}
	u = u_1 * u_2 \in\scr D'
	\quad\IFF\quad
	\forall\phi\in\scr D,\ 
	u*\phi = u_1 * (u_2*\phi) .
\end{equation}

%%%%%%%%%%%%%%%%%%%%%%%%%%%%%%%%%%%%%%%%%%%%%%%%%%%%%%%%%%%%%%%
\subsection{Properties of convolutions}\label{subs67ebb2f1}
We summarize the properties of convolutions.
\begin{proposition}\label{prop67f2a0bd}
	\begin{enumerate}
	\item
	If $u_1,u_2\in\scr D'$, one of which has compact support, 
	then $u_1 * u_2 \in\scr D'$ is well defined.
	\item
	If $u_1,u_2\in\scr D'$, one of which has compact support,
	then $u_1*u_2 = u_2*u_1$.
	\item
	If $u_1,u_2\in\scr D'$, one of which has compact support,
	then $\spt(u_1*u_2) \subset \spt(u_1) + \spt(u_2)$.
	\item
	If  If $u_1,u_2,u_3\in\scr D'$, two of which have compact support,
	then $(u_1*u_2)*u_3 = u_1*(u_2*u_3)$.
	\item
	If $u_1,u_2\in\scr D'$, one of which has compact support,
	and $\alpha\in\N^n$,
	then $\Diff^\alpha(u*v) = (\Diff^\alpha u)*v = u*(\Diff^\alpha v)$.
	\item
	If $u_1,u_2\in\scr D'$, one of which has compact support, and one of which is smooth,
	then $u_1 * u_2 \in\scr E$.
	\end{enumerate}
\end{proposition}
\begin{proof}
	The proof is left as an exercise.
	I only write the proof of $\spt(u_1*u_2) \subset \spt(u_1) + \spt(u_2)$.
	
	Since convolution is commutative, we can assume $u_2\in\scr E'$.
	Take $\phi\in\scr D$ such that $\spt(\phi) \cap (\spt(u_1) + \spt(u_2)) = \emptyset$.
	Then $(u_1*u_2)[\phi] 
	\overset{\eqref{eq67e28b61}}= (u_1*u_2)*\check\phi(0)
	= u_1*(u_2*\check\phi)(0)$
	where $\spt(u_2*\check\phi) \overset{\eqref{eq67f2b0f5}}\subset \spt(u_2) + \spt(\check\phi)  = \spt(u_2) -\spt(\phi)$.
%	Therefore, $\spt(u_2*\check\phi)\cap \spt(u_1) \subset \spt(u_2) -\spt(\phi)\cap \spt(u_1) = \emptyset$.
	It follows that $\spt(u_1*(u_2*\check\phi)) \overset{\eqref{eq67f2b0f5}}\subset \spt(u_1) + \spt(u_2*\check\phi) \subset \spt(u_1) + \spt(u_2) - \spt(\phi)$.
	Since $\spt(\phi) \cap (\spt(u_1) + \spt(u_2)) = \emptyset$, then $0\notin \spt(u_1*(u_2*\check\phi))$, that is $u_1*(u_2*\check\phi)(0)=0$.
\end{proof}

\begin{exercise}
	Prove Proposition~\ref{prop67f2a0bd}.
	Some parts have already been proven, others instead have been proven only partially.
	Put all the pieces together.
\end{exercise}

\begin{exercise}
	Show that $\delta_0*u = u$ for every $u\in\scr D'$.
	What is $\delta_v*u$?
\end{exercise}

%%%%%%%%%%%%%%%%%%%%%%%%%%%%%%%%%%%%%%%%%%%%%%%%%%%%%%%%%%%%%%%
\subsection{Singular support}
The singular support of a distribution $u\in\scr D'(\Omega)$, denote by $\singspt(u)$,
is the set defined by:
\begin{equation}
	\Omega\setminus \singspt(u) 
	= \bigcup \{U\subset\Omega\text{ open, such that }u|_V\in C^\infty(V)\} .
\end{equation}

\begin{lemma}
	Let $u\in\scr D'(\Omega)$.
	The restriction of $u$ to $\Omega\setminus\singspt(u)$ is a smooth function.
\end{lemma}

\begin{exercise}
	Show that $\singspt(u) \subset \spt(u)$.
\end{exercise}

Saying that a distribution $u$ is smooth is equivalent to say that $\singspt(u)=\emptyset$.
The smoothness statements for convolutions in Proposition~\ref{prop67ea5631} and Proposition~\ref{prop67f21c33}, generalize in the following statement:

\begin{proposition}\label{prop67f2a337}
	If $u,v\in\scr D'$, one of which has compact support, then
	\begin{equation}\label{eq67f2b99a}
		\singspt(u*v) \subset \singspt(u) + \singspt(v) .
	\end{equation}
\end{proposition}
\begin{proof}
	Let $A,B\subset\R^n$ open such that $\singspt(u)\subset A$ and $\singspt(v)\subset B$.
	Then there are smooth functions $a,b\in C^\infty(\R^n)$ such that 
	\begin{align*}
	 	&\singspt(u) \subset\interior\{a=1\} \subset\spt(a) \subset A ,\\
		&\singspt(v) \subset\interior\{b=1\} \subset\spt(b) \subset B .
	\end{align*}
	Then 
	\begin{equation}
		u*v
		= ((au)*(bv)) + ((1-au)*(bv)) + ((au)*(1-bv)) + ((1-au)*(1-bv)) .
	\end{equation}
	Notice that all these convolutions are well defined.
	Moreover, in all but the first one, one of the factors of the convolutions is a smooth function.
	Hence,
	\begin{equation}
		\singspt(u*v) 
		\subset \singspt((au)*(bv))
		\subset \spt((au)*(bv))
		\subset \spt(au) + \spt(bv) 
		\subset A + B .
	\end{equation}
	Since $A$ and $B$ are arbitrary, we obtain~\eqref{eq67f2b99a}.
\end{proof}

\begin{exercise}
	Show that, if $E\in\scr D'$ is such that $\singspt(E)\subset\{0\}$, then $\singspt(E*u) \subset \singspt(u)$ for all $u$ that can be convoluted with $E$.
\end{exercise}

%%%%%%%%%%%%%%%%%%%%%%%%%%%%%%%%%%%%%%%%%%%%%%%%%%%%%%%%%%%%%%%
\subsection{Linear differential operators with constant coefficients}
A \emph{linear differential operator with constant coefficients of order $m\in\N$}
is a differential operator of the form
\begin{equation}
	P = \sum_{|\alpha|\le m} c_\alpha \Diff^\alpha ,
\end{equation} 
with some $\{c_\alpha\}_{\alpha}\subset\C$.
Our four PDE are of this type:
\begin{equation}
	\de_t - b \cdot \grad,\qquad
	-\laplacian, \qquad
	\de_t-\laplacian, \qquad
	\de_t^2-\laplacian.
\end{equation}

Such a differential operator defines a linear operator $P:\scr D'\to\scr D'$.
Suddenly, we can try to solve in $u\in\scr D'$ the equation
\begin{equation}
	Pu = f
\end{equation}
for some $f\in\scr D'$.

%%%%%%%%%%%%%%%%%%%%%%%%%%%%%%%%%%%%%%%%%%%%%%%%%%%%%%%%%%%%%%%
\subsection{Fundamental solution}\label{subs6889d87f}
A \emph{fundamental solution} of a linear differential operator with constant coefficients $P$ is a distribution $E\in\scr D'$ such that
\begin{equation}
	PE = \delta_0 ,
\end{equation}
where $\delta_0$ is the Dirac delta centered at $0$.

\begin{exercise}
	Show that the function $\Phi:\R^n\setminus\{0\}\to\R$,
	\begin{equation}%\label{eq67b9d6b9}
		\Phi(x) = 
		\begin{cases}
		-\frac{1}{2\pi} \log(|x|) & \text{ if }n=2 , \\
		\frac1{n(n-2)\omega_n } \frac1{|x|^{n-2}} & \text{ if }n\ge 3 ,
		\end{cases}
	\end{equation}
	where $\omega_n $ is the volume of the unit ball in $\R^n$,
	is a fundamental solution of $P=-\laplacian$.
\end{exercise}

\begin{exercise} %(
	Show that the function 
	$\Phi:\R^n\times\R\setminus\{(0,0)\} \to [0,+\infty)$ defined by
	\begin{equation}%\label{eq67c3022e}
		\Phi(x,t) = \begin{cases}
		\frac{1}{(4\pi t)^{n/2}} \exp\left( -\frac{|x|^2}{4t} \right)
			& \text{ for $x\in\R^n$ and $t>0$,} \\
		0 & \text{otherwise, i.e., }(x,t)\in(\R^n\times(-\infty,0])\setminus\{(0,0)\} ,
		\end{cases}
	\end{equation}
	is a fundamental solution of $P=\de_t-\laplacian$.
\end{exercise}%)

\begin{exercise}
	Find a fundamental solution for the wave operator $P = \waveop = \de_t^2-\laplacian$
	in dimension 1, 2 and 3.
\end{exercise}

%%%%%%%%%%%%%%%%%%%%%%%%%%%%%%%%%%%%%%%%%%%%%%%%%%%%%%%%%%%%%%%
\subsection{Use of fundamental solutions}\label{subs6889edd6}
A fundamental solution is useful for the following very simple reason.
If $P$ is a differential operator and $E\in\scr D'$ is such that $PE=\delta_0$, then, 
for every $f\in\scr E'$,
\begin{equation}
	f = f*\delta_0 = f*PE = P(f*E) .
\end{equation}
So, $f*E$ is a solution to $Pu=f$.

%%%%%%%%%%%%%%%%%%%%%%%%%%%%%%%%%%%%%%%%%%%%%%%%%%%%%%%%%%%%%%%
\subsection{Hypoellipticity}
A linear operator $P$ is \emph{hypoelliptic} if, for every $\Omega\subset \R^n$ open and  $u\in\scr D'(\Omega)$, 
$Pu\in C^\infty(\Omega)$ implies $u\in C^\infty(\Omega)$.

\begin{theorem}
	Let $P$ be a linear differential operator with constant coefficients.
	Then the following are equivalent:
	\begin{enumerate}[label=(\roman*)]
	\item Some fundamental solution $E$ of $P$ has $\singspt(E)\subset\{0\}$.
	\item Every fundamental solution $E$ of $P$ has $\singspt(E)\subset\{0\}$.
	\item $P$ is hypoelliptic.
	\end{enumerate}
\end{theorem}
\begin{proof}
	The implications $(iii)\THEN(ii)\THEN(i)$ are clear.
	We need to show $(i)\THEN(iii)$.
	Let $\Omega\subset\R^n$ open and $u\in\scr D'(\Omega)$ with $Pu\in C^\infty(\Omega)$.
	Fix $x\in\Omega$ and let $\psi\in C^\infty_c(\Omega)$ such that $x\in\interior\{\psi=1\}$.
	Then $\psi u\in\scr E'$.
	Moreover, by the Generalized Leibniz Rule~\eqref{eq67f2c2e1}, 
	$P(\psi u) = \psi Pu + R$ where $\spt(R) \subset\spt(\Diff \psi)$.
	So
	\begin{equation}
		\psi u = \delta_0 *(\psi u) = (PE) * (\psi u) = E*(P(\psi u))
		= E*(\psi Pu) + E*R .
	\end{equation}
	and
	\begin{align*}
		\singspt(\psi u) 
		&\subset \singspt(E*(\psi Pu)) \cup \singspt(E*R) \\
		&\subset \emptyset \cup (\singspt(E) + \singspt(R)) \\
		&\subset \spt(R) 
		\subset \spt(\Diff \psi) .
	\end{align*}
	Since $x\notin\spt(\Diff \psi)$, we obtain that $x\notin \singspt(\psi u)$.
	In other words, $\psi u$ is smooth in a neighborhood of $x$.
	Since $\psi$ is 1 in a neighborhood of $x$, we get that $u$ is smooth in a neighborhood of $x$.
	We conclude that $u$ is smooth in $x$.
\end{proof}


%%%%%%%%%%%%%%%%%%%%%%%%%%%%%%%%%%%%%%%%%%%%%%%%%%%%%%%%%%%%%%%
\subsection{Extra topics that are not covered}

\begin{itemize}
\item	kernel theorem, 
\item	convolution operators, 
\item	existence of fundamental solution
\end{itemize}





%%%%%%%%%%%%%%%%%%%%%%%%%%%%%%%%%%%%%%%%%%%%%%%%%%%%%%%%%%%%%%%
%%%%%%%%%%%%%%%%%%%%%%%%%%%%%%%%%%%%%%%%%%%%%%%%%%%%%%%%%%%%%%%
\newpage
\section{Schwartz distributions and Fourier transform}\label{sec67e28d5a}

\subsection{Schwartz test functions}
The \emph{Schwartz space} $\scr S$ of \emph{Schwartz functions} is the space of functions $f:\R^n\to\C$ that are smooth and such that, for every $\alpha\in\N^n$ and every $N\in\N$ there exists $C_{\alpha,N}$ with
\begin{equation}\label{eq67f2c68b}
	|\Diff^\alpha f(x)| \le C_{\alpha,N} (1+|x|)^{-N} ,
	\qquad\forall x\in\R^n .
\end{equation}
In other words, a smooth function $f$ belongs to $\scr S$ if $f$ and all its derivatives decrease at infinity faster than any (inverse of) polynomial.
An example is $f(x) = \exp(-|x|^2)$.

The condition~\eqref{eq67f2c68b} can be expressed in different ways.
For instance, one can require the upper bound 
$|\Diff^\alpha f(x)| \le C_{\alpha,N} (1+|x|^2)^{-N/2}$,
or $|\Diff^\alpha f(x)| \le C_{\alpha,N} (1+|x|^2)^{-N}$.

Another way to express the condition~\eqref{eq67f2c68b} is as follows:
a function $f\in C^\infty(\R^n)$ belongs to $\scr S$ if and only if, for every $\alpha,\beta\in\N^n$, the seminorms
\begin{equation}
	p_{\alpha,\beta}(f) = \sup_{x\in\R^n} |x^\alpha\Diff^\beta f(x)| 
\end{equation}
are finite.
An equivalent family of seminorms is
\begin{equation}
	p_{\alpha,N}(f) = \sup_{x\in\R^n} (1+|x|)^N |\Diff^\alpha f(x)| .
\end{equation}

We endow $\scr S$ with the family of seminorms $\{p_{\alpha,\beta}\}_{\alpha,\beta\in\N^n}$, so that $\scr S$ becomes a Fréchet space.
This means in particular that convergence in $\scr S$ is defined as follows:
\begin{equation}\label{eq67fbfa93}
	f_j\to f_\infty \text{ in }\scr S
	\quad\IFF\quad
	\forall\alpha,\beta\in\N^n\ 
		\lim_{j\to\infty}p_{\alpha,\beta}(f_j-f_\infty) =0 .
\end{equation}

Notice that, set-wise,
\begin{equation}
	\scr D \subset \scr S \subset \scr E .
\end{equation}
As usual, one can check that these embeddings are continuous.

\begin{exercise}
	Show that $\scr S\subset L^1$.
	
	{\it Hint:}
	Use~\eqref{eq67f2c68b} for $N$ large enough.
\end{exercise}

\begin{exercise}\label{ex68035dc9}
	Show that $f_j\to f_\infty$ in $\scr S$ if and only if $x^\alpha \Diff^\beta f_j \to x^\alpha \Diff^\beta f_\infty$ in $L^\infty$, for every $\alpha,\beta\in\N^n$.
\end{exercise}

\begin{exercise}
	Show that, if $f_j\to f_\infty$ in $\scr S$, then, for every $\alpha\in\N^n$, $\Diff^\alpha f_j\to\Diff^\alpha f_\infty$ in $L^1(\R^n)$.
	
	{\it Hint:}
	Take first $\alpha=0$. 
	Up to substituting $f_j$ with $f_j-f_\infty$, we can also assume $f_\infty=0$.
	So we need to show that, if $f_j\to0$ in $\scr S$, then $f_j\to 0$ in $L^1(\R^n)$.
	The convergence $f_j\to0$ in $\scr S$ implies that, for every $N>0$ and every $\epsilon>0$, there is $J\in\N$ such that $|f_j(x)| (1+|x|)^N < \epsilon$ for all $x\in\R^n$ and all $j>J$.
\end{exercise}

\begin{exercise}\label{ex67fc0379}%( 
	Define $g_\epsilon:\R^n\to[0,+\infty)$ by
	\begin{equation}
		g_\epsilon(z) = \frac1{\epsilon^n} e^{-\pi|z|^2/\epsilon^2} = \frac1{\epsilon^n}  \exp(-\pi|z|^2/\epsilon^2) .
	\end{equation}
	Show that, if $f\in\scr S$, then $f*g_\epsilon\to f$ in $\scr S$ as $\epsilon\to0$.
	
	{\it Hint:} Go back to the proof of the first statement in Proposition~\ref{prop67e28cab}.
	The now $g_\epsilon$ does not have compact support, but $\int_{\R^n\setminus B(0,1)} g_\epsilon(x)\did x$ is arbitrarily small as $\epsilon\to0$.
\end{exercise}%]]


%%%%%%%%%%%%%%%%%%%%%%%%%%%%%%%%%%%%%%%%%%%%%%%%%%%%%%%%%%%%%%%
\subsection{Fourier transform}
For $u\in L^1(\R^n)$ and $\xi\in\R^n$, define the \emph{Fourier transform}
\begin{equation}\label{eq67fbe754}
	\hat u (\xi) = \scr F(u)(\xi) 
	= %\frac1{(2\pi i)^n} 
	\int_{\R^n} e^{-2\pi i x\cdot \xi} u(x) \did x .
\end{equation}
Since $|e^{-2\pi i x\cdot \xi} u(x)|\le |u(x)|$, we readily have that $\hat u(x)$ is well defined for every $x\in\R^n$ and that
\begin{equation}\label{eq67fbf439}
	\|\hat u\|_{L^{\infty}} \le \| u \|_{L^1} .
\end{equation}
In particular, we have that, if $u_j\to u_\infty$ in $L^1$, then $\hat u_j\to \hat u_\infty$ uniformly on $\R^n$.

%%%%%%%%%%%%%%%%%%%%%%%%%%%%%%%%%%%%%%%%%%%%%%%%%%%%%%%%%%%%%%%
\subsection{Literature on the Fourier transform}
There are a lot of resources on the Fourier transforms.
We follow in particular these three:
\begin{enumerate}
\item the short account in Section 0.D, page 14, in \fullcite{MR1357411}; % Folland
\item Chapter VII in \fullcite{MR717035}; % Hörmander
\item Chapter 7 in \fullcite{MR1157815}.% Rudin 
\end{enumerate}

Be aware that the Fourier transform has slightly different definitions in the literature.
Beside our formula~\eqref{eq67fbe754}, there are also
\begin{gather*}
% 	\int_{\R^n} e^{2\pi i x\cdot \xi} u(x) \did x , \\
	\int_{\R^n} e^{-i x\cdot \xi} u(x) \did x , \\
	\frac1{(2\pi)^{n/2}}\int_{\R^n} e^{-2\pi i x\cdot \xi} u(x) \did x .
\end{gather*}
These differences will lead to differences in the formulas, usually in the multiplicative constants.
It is good to do the proofs with one formula, 
but also to try the others.

%%%%%%%%%%%%%%%%%%%%%%%%%%%%%%%%%%%%%%%%%%%%%%%%%%%%%%%%%%%%%%%
\subsection{First properties of the Fourier transform}

\begin{exercise}
	Prove the following properties:
	\begin{enumerate}
	\item
	If $u\in L^1(\R^n)$, $a\in\R^n$, and $u_a(x) = u(x+a)$
	then $\scr F(u_a)(\xi) = e^{2\pi i a\cdot \xi} \scr F(u)(\xi)$.
	\item
	If $u\in L^1(\R^n)$, $T:\R^n\to\R^n$ is linear and invertible, 
	then 
	\begin{equation}\label{eq67fccb0f}
		\scr F(u\circ T)(\xi) = |\det T|^{-1} \scr F(u)((T^{-1})^*\xi) .
	\end{equation}
	\item
	If $T$ is a rotation of $\R^n$, then $\scr F(u\circ T) = \scr F(u)\circ T$.
	\end{enumerate}
	{\it Hint:} See \cite[Proposition (0.21)]{MR1357411}
\end{exercise}

\begin{exercise}
	Compute $\scr F(\bar u)$ in terms of $\scr F(u)$.
	(Here $\bar\cdot$ denotes the complex conjugate, that is, for $x,y\in\R$, $ \overline{x+iy} = x-iy$.)
\end{exercise}

\begin{exercise}
	Compute $\scr F(u(-x))$ in terms of $\scr F(u)$.
\end{exercise}

\begin{exercise}\label{ex67fc118a}
	Fix $f\in\scr S$ and define $g(x) = \overline{f(-x)}$.
	Show that $\hat g (\xi) = \overline{\hat f(\xi)}$.
\end{exercise}

\begin{exercise}
	Compute the Fourier transform of the function 
	\begin{equation}
		u(x) = A e^{-a|x|^2} ,
	\end{equation}
	for every $A\in\C$ and $a>0$.
	
	{\it Solution:}
	Let's start with $n=1$, $A=1$ and $a=\pi$, that is, $u(x) = e^{-\pi x^2}$.
	Then, for $\xi\in\R$,
	\begin{align}
		\hat u(\xi)
		&= \int_\R e^{-2\pi i\xi x} e^{-\pi x^2} \did x
		= \int_\R e^{-2\pi i\xi x - \pi x^2} \did x .
	\end{align}
	Since $(x+i\xi)^2 = x^2-\xi^2 + 2ix\xi$, then 
	$-2\pi i\xi x - \pi x^2 
	= -\pi (x^2 + 2ix\xi) 
	= -\pi((x+i\xi)^2 + \xi^2)$ and thus
	\begin{align}
		\hat u(\xi)
		= \int_\R e^{-\pi((x+i\xi)^2 + \xi^2)} \did x
		= e^{-\pi\xi^2} \int_\R e^{-\pi(x+i\xi)^2} \did x .
	\end{align}
	Now we see that we are integrating the holomorphic function $f(z) = e^{-\pi z^2}$ along the curve $\gamma(t) = t+i\xi$.
	Applying Cauchy's theorem, we can integrate along the other curve $\eta(t) = t$, and thus 
	\begin{align}
		\hat u(\xi) 
		= e^{-\pi\xi^2} \int_\R e^{-\pi x^2} \did x
		= e^{-\pi\xi^2} ,
	\end{align}
	because $\int_\R e^{-\pi x^2} \did x = 1$.
	So, $\hat u = u$, or $\scr F(u)=u$.
	
	We extend this result to $\R^n$ for $n\ge1$: 
	\begin{align}
		\scr F(e^{-\pi|x|^2})(\xi) 
		&= \int_{\R^n} e^{-2\pi i\xi\cdot x} e^{-\pi|x|^2} \did x \\
		&= \int_{\R^n} e^{-2\pi i \sum_j\xi_jx_j - \pi \sum_j x_j^2} \did x \\
		&= \prod_{j=1}^n \int_\R e^{-2\pi i \xi_jx_j - \pi  x_j^2} \did x_j \\
		&= \prod_{j=1}^n e^{-\pi\xi_j^2}
		= e^{-\pi|\xi|^2} .
	\end{align}
	Again, we have,
	\begin{equation}
		\scr F(e^{-\pi|x|^2} ) = e^{-\pi|\xi|^2} ,
	\end{equation}
	that is, the function $u(x) = e^{-\pi|x|^2}$ is a fixed point of $\scr F:\scr S(\R^n)\to\scr S(\R^n)$.
	
	We now go back to the general $u:\R^n\to\R$, $u(x) = Ae^{-a|x|^2}$.
	By linearity of $\scr F$, we have $\scr F(Ae^{-a|x|^2}) = A \scr F(e^{-a|x|^2})$.
	We apply~\eqref{eq67fccb0f} with $Tx = \sqrt{a/\pi}x$, so that 
	\begin{align}
		\scr F(e^{-a|x|^2})(\xi) 
		&= \scr F( e^{-\pi|x|^2}\circ T )(\xi) \\
		&= (a/\pi)^{-n/2} \scr F(e^{-\pi|x|^2})(\xi/\sqrt{a/\pi}) \\
		&= \left(\frac{\pi}{a}\right)^{n/2} e^{-\frac{\pi^2}{a} |x|^2 }.
	\end{align}
	We now wrap up the solution to get
	\begin{equation}\label{eq67fe64f8}
		\scr F(Ae^{-a|x|^2}) = A \left(\frac{\pi}{a}\right)^{n/2} e^{-\frac{\pi^2}{a} |x|^2 } .
	\end{equation}
\end{exercise}

Recall from Exercise~\ref{ex67fbf157}, that the convolution of two functions $f,g\in L^1(\R^n)$ is a well defined function $f*g\in L^1(\R^n)$.

\begin{exercise}
	Show that, for every $f,g\in L^1(\R^n)$,
	\begin{equation}\label{eq67fc11d5}
		\scr F(f*g) = \scr F(f) \cdot \scr F(g) .
	\end{equation}
	Or, otherwise stated, $(f*g)^{\wedge} = \hat f \hat g$.
\end{exercise}

\begin{exercise}
	Show that, for every $f,g\in L^1(\R^n)$,
	\begin{equation}
		\int_{\R^n} f(x) \hat g(x) \did x 
		= \int_{\R^n} \hat f(\xi) g(\xi) \did \xi .
	\end{equation}
	
	{\it Hint:}
	Notice that, by~\eqref{eq67fbf439}, $f\hat g\in L^1(\R^n)$ and $\hat f,g\in L^1(\R^n)$.
	So, unpack the definition of $\hat g$ and use Fubini.
\end{exercise}

%%%%%%%%%%%%%%%%%%%%%%%%%%%%%%%%%%%%%%%%%%%%%%%%%%%%%%%%%%%%%%%
\subsection{The Fourier transform preserves the Schwartz class}
\begin{proposition}
	If $f \in\scr S$, then $\hat f\in\scr S$ and, for every $\alpha\in\N^n$,
	\begin{gather}
	\label{eq67fbf9a6}
		\scr F(\Diff^\alpha f) = (2\pi i\xi)^\alpha \scr F(f) ,\\
	\label{eq67fbf9ac}
		\scr F( (-2\pi ix)^\alpha f) = \Diff^\alpha\scr F(f) .
	\end{gather}
\end{proposition}
\begin{proof}
	For $f \in\scr S$ and $j\in\{1,\dots,n\}$, we have $\de_jf\in\scr S$ and thus
	\begin{align*}
	 	\scr F(\de_jf)(\xi)
		&= \int_{\R^n} \exp(-2\pi i\xi\cdot x) \frac{\de}{\de x_j} f(x) \did x \\
		&= \int_{\R^n} \frac{\de}{\de x_j} [\exp(-2\pi i\xi\cdot x) f(x)] \did x 
			- \int_{\R^n} (-2\pi i\xi_j) \exp(-2\pi i\xi\cdot x) f(x)) \did x \\
		&= (2\pi i\xi_j) \hat f(\xi) .
	\end{align*}
	Iterating, we obtain~\eqref{eq67fbf9a6}.
	
	Since $\scr F(\Diff^\alpha f)\in L^\infty(\R^n)$ for every $\alpha\in\N^n$,
	we conclude that $\hat f$, which a priori only belongs to $L^\infty(\R^n)$, in fact satisfies
	\begin{equation}\label{eq68035c23}
		p_{0,N}(\hat f) = \sup_{\xi\in\R^n} (1+|\xi|)^N |\hat f(\xi)| < \infty ,
	\end{equation}
	for every $N\in\N$.
	
	We don't know yet that $\hat f$ is differentiable.
	For this, we use \ref{thm679f29de_3} in Theorem~\ref{thm679f29de} to compute, for $f\in\scr S$ and $j\in\{1,\dots,n\}$,
	\begin{align}
		\de_j\scr F(f)(\xi)
		&= \frac{\de}{\de \xi_j} \int_{\R^n} e^{-2\pi i \xi\cdot x} f(x) \did x \\
		&= \int_{\R^n} \frac{\de}{\de \xi_j} [ e^{-2\pi i \xi\cdot x} f(x) ] \did x \\
		&= \int_{\R^n} (-2\pi ix_j) e^{-2\pi i \xi\cdot x} f(x) \did x \\
		&= \scr F((-2\pi ix_j) f)(\xi) .
	\end{align}
	Iterating, we obtain~\eqref{eq67fbf9ac}.
	We conclude that $\hat f$ is smooth and that, for every $\alpha\in\N^n$, since $\Diff^\alpha\hat f$ is the Fourier transform of a function in $\scr S$, then, by~\eqref{eq68035c23},
	\begin{equation}
		p_{\alpha,N}(\hat f) 
		= p_{0,N}(\Diff^\alpha\hat f) < \infty .
	\end{equation}
	We conclude that $\hat f\in\scr S$.
\end{proof}

\begin{corollary}[Riemann--Lebesgue Lemma]
	If $f\in L^1(\R^n)$, then $\hat f\in C^0_0(\R^n)$, that is, $\hat f$ is continuous and tends to zero at infinity.
\end{corollary}
\begin{proof}
	We know that $\scr S$ is dense in $L^1(\R^n)$, see Exercise~\ref{ex67fbf2c0}.
	Let $f\in L^1(\R^n)$ and $f_j\in\scr S$ a sequence with $f_j\to f$ in $L^1(\R^n)$.
	By~\eqref{eq67fbf439}, $\hat f_j\to\hat f$ in $L^\infty(\R^n)$, that is $\hat f_j\to \hat f$ uniformly.
	It follows that $\hat f$ is continuous.
	
	We need to show that $\hat f$ tends to zero at infinity.
	For every $\epsilon>0$ there is $j\in\N$ with $\|f_j-f\|_{L^\infty}<\epsilon$.  
	Then, there is $R>0$ such that $|f_j(\xi)|<\epsilon$ for every $\xi$ with $|\xi|>R$.
	Therefore, if $|\xi|>R$, then $|f(\xi)| \le |f(\xi)-f_j(\xi)| + |f_j(\xi)| < 2\epsilon$.
	This shows that $\lim_{\xi\to\infty} \hat f(\xi) = 0$.
\end{proof}

\begin{exercise}\label{ex67fbf2c0}
	Show that $\scr S$ is dense in $L^1(\R^n)$.
	
	{\it Hint:} 
	Given $f\in L^1(\R^n)$, consider $f_j(x) = \psi(x/j) f*\rho_{1/j}(x)$,
	where $\{\rho_\epsilon\}_{\epsilon>0}$ is a family of mollifiers, and $\psi\in C^\infty_c(\R^n)$ is a function valued in $[0,1]$ with $B(0,1) \subset\{\psi =1\}$.
	You then need to show that $f_j\in\scr S$ and that $f_j\to f$ in $L^1(\R^n)$.
	Use the fact that, for every $\epsilon>0$ there exists $R>0$ such that $\int_{\R^n\setminus B(0,R)} |f(x)| \did x < \epsilon$ (this is a direct consequence of integrability).
\end{exercise}

%%%%%%%%%%%%%%%%%%%%%%%%%%%%%%%%%%%%%%%%%%%%%%%%%%%%%%%%%%%%%%%
\subsection{Fourier inversion theorem}

\begin{theorem}\label{thm67fcdb1f}
	The Fourier transform $\scr F:\scr S\to\scr S$ is a linear automorphism of $\scr S$.
	In particular, if $f\in\scr S$, then 
	\begin{equation}\label{eq67fbff5a}
		\scr F^{-1}(f)(x) = \int_{\R^n} e^{2\pi i x\cdot\xi} f(\xi) \did \xi .
	\end{equation}
\end{theorem}
\begin{proof}
	We know that $\scr F$ is linear.
	We need to show that it is onto and into, and that both $\scr F$ and $\scr F^{-1}$ are continuous.
	
	First we show that $\scr F$ is continuous.
	If $\alpha,\beta\in\N^n$, then 
	\begin{align}
		\xi^\alpha \Diff^\beta \hat f(\xi)
		&\overset{\eqref{eq67fbf9ac}}= \xi^\alpha \scr F((-2\pi i x)^\beta f) \\
		&\overset{\eqref{eq67fbf9a6}}= \frac{1}{(2\pi i)^\alpha} \scr F(\Diff^\alpha (-2\pi i x)^\beta f ) \\
		&= \frac{(-2\pi i)^\beta}{(2\pi i)^\alpha} \scr F(\Diff^\alpha  (x^\beta f) ) .
	\end{align}
	If $f_j\to0$ in $\scr S$, then $\Diff^\alpha  (x^\beta f_j)\to 0$ in $\scr S$, for every $\alpha,\beta\in\N^n$.
	It follows that $\Diff^\alpha  (x^\beta f_j)\to 0$ in $L^1(\R^n)$, and thus 
	$\xi^\alpha \Diff^\beta \hat f_j\to 0$ in $L^\infty(\R^n)$, for every $\alpha,\beta\in\N^n$.
	This means that $\hat f_j\to 0$ in $\scr S$ (see Exercise~\ref{ex68035dc9}).
	So, $\scr F:\scr S\to\scr S$ is continuous.
	
	Denote by $\scr G(f)(x)$ the quantity on the right-hand side of~\eqref{eq67fbff5a}.
	Notice that $\scr G(f)(x)= \scr F(f)(-x)$: so, we know that $\scr G$ is a continuous linear operator $\scr S\to\scr S$.
	We will show that $\scr G(\scr F(f)) = f$, for all $f\in\scr S$.
	
	For $x\in\R^n$ and $\epsilon>0$, define
	\begin{equation}
		\phi_{x,\epsilon}(\xi) = e^{2\pi i x\cdot\xi - \pi \epsilon^2|\xi|^2}
		= \exp( 2\pi i x\cdot\xi - \pi \epsilon^2|\xi|^2 ).
	\end{equation}
	If we set $g_\epsilon(z) = e^{-\pi|z|^2/\epsilon^2} = \exp(-\pi|z|^2/\epsilon^2)$, 
	then one can compute
	\begin{equation}
		\hat\phi_{x,\epsilon}(y) = g_\epsilon(x-y) .
	\end{equation}
	Then we have
	\begin{align}
		\scr G(e^{ - \pi \epsilon^2|\xi|^2} \hat f )(x) 
		&= \int_{\R^n} e^{ - \pi \epsilon^2|\xi|^2} e^{2\pi i x\cdot\xi}  \hat f(\xi)\did \xi \\
		&= \int_{\R^n} \phi_{x,\epsilon}(\xi) \hat f(\xi) \did \xi \\
		&= \int_{\R^n} \hat \phi_{x,\epsilon}(y) f(y) \did y \\
		&= \int_{\R^n} g_\epsilon(x-y)  f(y) \did y \\
		&= g_\epsilon * f(x) .
	\end{align}
	On the one hand, $g_\epsilon * f\to f$ in $\scr S$ as $\epsilon\to0$ by Exercise~\ref{ex67fc0379}, hence point-wise.
	On the other hand, $e^{ - \pi \epsilon^2|\xi|^2} \hat f\to \hat f$ in $L^1(\R^n)$ as $\epsilon\to0$.
	Therefore, using also the continuity of $\scr G$, we conclude that $\scr G(\hat f) = f$.	
	
	We complete the proof with a bit of algebra.
	Define $\cal I:\scr S\to\scr S$, $\scr I(f)(x) = f(-x)$.
	We have just proven
	\begin{equation}
		\scr F\circ\cal I \circ \scr F = \Id_{\scr S} .
	\end{equation}
	This shows that $\scr F$ is also surjective, with inverse $\scr G=\cal I\circ\scr F$.
\end{proof}

\begin{remark}
	Notice that, if you write down the formula for $\scr G(\scr F(f))$, you cannot conclude using Fubini!
	In fact, we have somehow proved that
	\begin{equation}
		\int_{\R^n} \int_{\R^n} e^{2\pi i (x-y)\cdot\xi} f(y) \did y\did \xi = f(x) .
	\end{equation}
	This identity does not have a meaning as integrals, because the integrand function is not in $L^1(\R^n\times\R^n)$.
	However, it does look a lot the convolution of $f$ with $\delta_0$, in which case we would have proven
	\begin{equation}\label{eq67fc0575}
		\int_{\R^n} e^{2\pi i z\cdot\xi} \did \xi = \delta_0(z) .
	\end{equation}
	In the latter identity, the integral is not an integral:
	we will make sense of~\eqref{eq67fc0575} distributionally.
\end{remark}

\begin{remark}
	The inverse of the Fourier transform applied to $f$ is also denoted as $\check f=\scr F^{-1}(f)$.
	This is a problem in my notes, because I have already used this notation for the function $\check f(x) = f(-x)$.
	This clash of notation does not have a solution yet.
	I am just renouncing to use the second meaning in this section, so that, in the context of Fourier transform, $\check f$ is only the inverse transform of $f$.
\end{remark}

\begin{exercise}
	Define $\cal I:\scr S\to\scr S$, $\cal I(f)(x) = f(-x)$.
	Show that
	\begin{equation}
		\scr F^2 = \cal I
		\quad\text{ and }\quad
		\scr F^4 = \Id_{\scr S} .
	\end{equation}
	
	{\it Hint:}
	Forget about the Fourier transform, but only use $\scr F\cal I\scr F=1$.
\end{exercise}

%%%%%%%%%%%%%%%%%%%%%%%%%%%%%%%%%%%%%%%%%%%%%%%%%%%%%%%%%%%%%%%
\subsection{Plancherel Theorem}

The space $L^2(\R^n)$ is the space of complex valued measurable functions $f:\R^n\to\C$ 
such that $\int_{\R^n} |f(x)|^2 \did x <\infty$.
The vector space $L^2(\R^n)$ is a Hilbert space when endowed with the sesquilinear form
\begin{equation}
	\langle f,g \rangle = \int_{\R^n} f(x) \bar{g(x)} \did x ,
	\quad\forall f,g\in L^2(\R^n) .
\end{equation}

\begin{theorem}
	The Fourier transform extends uniquely to a unitary automorphism of $L^2(\R^n)$.
	More precisely, the Fourier transform has a continuous extension $\scr F:L^2(\R^n)\to L^2(\R^n)$ with 
	\begin{gather}
		\|\scr F(f)\|_{L^2} = \|f\|_{L^2} \quad\forall f\in L^2(\R^n) , \\
		\langle \scr F(f),\scr F(g) \rangle = \langle f,g \rangle \quad\forall f,g\in L^2(\R^n) .
	\end{gather}
\end{theorem}
\begin{proof}
	By Exercise~\ref{ex67fc0c47} and Theorem~\ref{thm67fcdb1f}, we know that $\scr F:\scr S\to\scr S$ is continuous with respect to the topology of $L^2(\R^n)$.
	Since $\scr S$ is dense in $L^2(\R^n)$, then $\scr F$ admits a unique continuous extension $L^2(\R^n)\to L^2(\R^n)$.
	The inverse of $\scr F$ has the same properties, and thus the extension of $\scr F$ is a continuous invertible linear operator.
	
	We only need to show $\|\scr F(f)\|_{L^2} = \|f\|_{L^2}$ for every $f\in\scr S$.
	Fix $f\in\scr S$ and define $g(x) = \overline{f(-x)}$.
	Notice that $\hat g (\xi) = \overline{\hat f(\xi)}$; 
	see also Exercise~\ref{ex67fc118a}.
	Then:
	\begin{align}
		\|\scr F(f)\|_{L^2} 
		&= \int_{\R^n} \hat f(\xi) \overline{\hat f(\xi)} \did \xi \\
		[\text{by Exercise~\ref{ex67fc118a}}]
		&= \int_{\R^n} \hat f(\xi) \hat g (\xi) \did \xi \\
		&\overset{\eqref{eq67fc11d5}}= \int_{\R^n} (f*g)^{\wedge}(\xi) \did \xi \\
		&= \int_{\R^n} e^{2\pi i 0\cdot \xi}(f*g)^{\wedge}(\xi) \did \xi \\
		&\overset{\eqref{eq67fbff5a}}= f*g(0) \\
		&= \int_{\R^n} f(x) \overline{f(-x)} \did x \\
		&= \|f\|_{L^2} .
	\end{align}
\end{proof}

\begin{exercise}\label{ex67fc0c47}%(((
	(Maybe this already appeared before).
	Show that $\scr D$ is dense in $L^p(\R^n)$ for all $p\in[1,\infty)$.
	Deduce that $\scr S$ is dense in $L^p(\R^n)$ for all $p\in[1,\infty)$.
	
	But then, show also that, if $f_j\to f_\infty$ in $\scr D$ or $\scr S$, then $f_j\to f_\infty$ in $L^p(\R^n)$ for all $p\in[1,\infty)$.
\end{exercise}%]]]

%%%%%%%%%%%%%%%%%%%%%%%%%%%%%%%%%%%%%%%%%%%%%%%%%%%%%%%%%%%%%%%
\subsection{Schwartz distributions}
The topological dual $\scr S'$ is the space of \emph{Schwartz distributions}, or \emph{tempered distributions}.
The convergence of sequences in $\scr S'$ is pointwise, that is, 
$A_n\to A_\infty$ in $\scr S'$ if and only if $A_n[\phi] \to A_\infty[\phi]$ for every $\phi\in\scr S$.

\begin{exercise}
	Show that
	\begin{equation}
		\scr D' \supset \scr S' \supset \scr E' .
	\end{equation}
\end{exercise}

\begin{exercise}
	Show the following properties of tempered distributions:
	\begin{enumerate}
	\item
	If $u\in\scr S'$ and $\alpha\in\N^n$, then $\Diff^\alpha u\in\scr S'$.
	\item
	If $u\in\scr S'$ and $f\in C^\infty(\R^n)$ is such that, 
	for all $\alpha\in\N^n$, $\Diff^\alpha f$ grows at most polynomially at infinity,
	then $fu\in\scr S'$.
	\item
	If $u\in\scr S'$ and $f\in\scr S$, then $u*f\in\scr S$.
%	\item
%	(Hard, but you have all the elements to do it:)
%	If $u,v\in\scr S'$, then $u*v$ is a well defined tempered distribution,
%	and this convolution has the usual properties listed in Proposition~\ref{prop67f2a0bd}.
	\end{enumerate}
\end{exercise}

\begin{exercise}
	Show that, if $h:\R^n\to\C$ is a measurable function that grows at most polynomially, then $f\mapsto \int_{\R^n} h(x)f(x)\did x$ defines a tempered distribution.
\end{exercise}

%%%%%%%%%%%%%%%%%%%%%%%%%%%%%%%%%%%%%%%%%%%%%%%%%%%%%%%%%%%%%%%
\subsection{The Fourier transform of tempered distributions}
We define the Fourier transform of a tempered distribution $u\in\scr S'$
as $\hat u=\scr F(u)$, where
\begin{equation}
	\hat u[f] = u[\hat f] ,
	\qquad\forall f\in\scr S.
\end{equation}

\begin{exercise}
	Show the following properties of the Fourier transform of tempered distributions:
	\begin{enumerate}
	\item
	If $u\in\scr S'$, then $\hat u\in\scr S'$.
	\item
	If $u\in\scr S'$ is actually in $\scr S$, i.e., $u[f] = \int_{\R^n} u(x) f(x) \did x$ for all $f\in\scr S$, then $\hat u$ as distribution is equal to $\hat u$ as function.
	\item
	$\scr F:\scr S'\to\scr S'$ is a continuous, invertible linear operator, with inverse
	$\scr F^{-1}u[f] = u[\scr F^{-1}(f)]$.
	\item
	If $u\in\scr S'$ and $f\in\scr S$, then 
	\begin{equation}\label{eq680350a7}
		\scr F(u*f) = \hat u \hat f .
	\end{equation}
	\item
	If $f \in\scr S'$, then $\hat f\in\scr S'$ and, for every $\alpha\in\N^n$,
	\begin{gather}
	\label{eq67fe65d6}
		\scr F(\Diff^\alpha f) = (2\pi i\xi)^\alpha \scr F(f) ,\\
	\label{eq67fe65da}
		\scr F( (-2\pi ix)^\alpha f) = \Diff^\alpha\scr F(f) .
	\end{gather}
	\end{enumerate}
\end{exercise}

%\begin{proposition}
%	If $f \in\scr S'$, then $\hat f\in\scr S'$ and, for every $\alpha\in\N^n$,
%	\begin{gather}
%	\label{eq67fe65d6}
%		\scr F(\Diff^\alpha f) = (2\pi i\xi)^\alpha \scr F(f) ,\\
%	\label{eq67fe65da}
%		\scr F( (-2\pi ix)^\alpha f) = \Diff^\alpha\scr F(f) .
%	\end{gather}
%\end{proposition}
%\begin{proof}
%	\TODO
%\end{proof}

\begin{exercise}\label{ex67fcc1d2}
	Compute $\scr F(\delta_0)$.
\end{exercise}

\begin{exercise}\label{ex6834885c}
	For $a\in\R^n$, compute $\scr F(\delta_a)$ and $\scr F(e^{ia\cdot x})$.
	Then compute $\scr F^{-1}(\delta_a)$ and $\scr F^{-1}(e^{ia\cdot x})$.
	
	{\it Solution.}
	Using only definitions, we have,
	for all $a\in\R^n$ and $\phi\in\scr S$,
	$\scr F(\delta_a)[\phi]
	= \delta_a[\scr F\phi]
	= \scr F\phi(a)
	= \int_{\R^n} \exp(-2\pi ia\cdot x) \phi(x) \did x$.
	This identity exactly means that, as distributions over $\R^n$,
	$\scr F(\delta_a) = \exp(-2\pi ia\cdot x)$.
	
	There is not much else to compute.
	Notice that $\cal I\delta_a[\phi] = \delta_a[\cal I\phi] = \phi(-a) =\delta_{-a}[\phi]$.
	So,
	\begin{align}
		\scr F(\delta_a) &= \exp(-2\pi ia\cdot x) , \\
		\delta_{-a} &= \cal I \delta_a = \scr F^2(\delta_a) = \scr F_x(\exp(-2\pi ia\cdot x)) ,\\ 
		\scr F_x(e^{ia\cdot x}) &= \delta_{\frac{a}{2\pi}} , \\
		\scr F^{-1}\delta_a &= \exp(2\pi ia\cdot x) , \label{eq683488c2}\\
		\scr F^{-1}(e^{ia\cdot x}) &= \cal I\scr F(e^{ia\cdot x}) = \cal I\delta_{\frac{a}{2\pi}} = \delta_{-\frac{a}{2\pi}} .
	\end{align}
\end{exercise}

\begin{exercise}
	Compute $\scr F(1)$.
\end{exercise}

%{\color{Red} \TODO
%If you have done Exercise~\ref{ex67fcc1d2}, you know what is $\scr F(1)$.
%However, let's try to compute it directly.
%There are two ways.
%
%First method: we can apply directly the definition:
%\begin{align}
%	\scr F(1)[f] = 1[\hat f] = \int_{\R^n}\hat f(\xi)\did \xi = \scr F^{-1}(\hat f)(0) = f(0) .
%\end{align}
%So, $\scr F(1) = \delta_0$.
%
%Second method: set $\phi_\epsilon(x) = \exp(-\epsilon|x|^2)$ and notice that $\phi_\epsilon \cdot 1 \to 1$ in $\scr S'$.
%Since $\phi_\epsilon\in\scr S$, then
%\begin{align}
%	\scr F(1) 
%	= \lim_{\epsilon\to0} \scr F(\phi_\epsilon)
%	= \lim_{\epsilon\to0} 
%\end{align}
%}

\begin{exercise}
	Compute $\scr F(p(x))$, where $p(x) = \sum_{|\alpha|\le N} c_\alpha x^\alpha$ is a polynomial of degree $N$.
\end{exercise}

\begin{remark}
	It is a fact that, if $\hat u$ has compact support, then $u$ is analytic, see \cite[Thm 7.1.14]{MR717035}
	It is a recurrent theme that regularity of $u$ is proportional to integrability of $\hat u$ (and viceversa, of course).
	There are two sorts of ``equilibrium points'' of this behavior: $\scr S$ and $L^2(\R^n)$.
\end{remark}


%%%%%%%%%%%%%%%%%%%%%%%%%%%%%%%%%%%%%%%%%%%%%%%%%%%%%%%%%%%%%%%
\subsection{Applications to PDE: Harmonic polynomials}
If $u\in\scr S'$ is such that 
\begin{equation}
	-\laplacian u = 0 ,
\end{equation}
then 
\begin{equation}
	0 = \scr F(-\laplacian u)
	\overset{\eqref{eq67fe65d6}}= - \sum_{j=1}^n(2\pi i\xi_j)^2 \hat u
	= 4\pi^2 |\xi|^2 \hat u .
\end{equation}
Therefore, $\spt(\hat u) \subset\{0\}$, that is, by Proposition~\ref{prop67fe6632},
$\hat u = \sum_{|\alpha|\le N} c_\alpha \Diff^\alpha \delta_0$ for some constants $c_\alpha\in\C$.
It follows that $u$ is a polynomial:

\begin{proposition}
	Harmonic tempered distributions are harmonic polynomials.
\end{proposition}

\begin{exercise}
	The condition $|\xi|^2 v=0$ implies $\spt(v) \subset\{0\}$.
	Can you characterize the distributions $v\in\scr S'$ that satisfy $|\xi|^2 v=0$?
\end{exercise}

%%%%%%%%%%%%%%%%%%%%%%%%%%%%%%%%%%%%%%%%%%%%%%%%%%%%%%%%%%%%%%%
\subsection{Applications to PDE: Heat equation}%(
See also \cite[\S4.A,p.143]{MR1357411} and \cite[p.192]{MR2597943}.

Let's consider the heat equation
\begin{equation}\label{eq680346f3}
\begin{cases}
	\de_t u - \laplacian u = 0 & \text{ in }\R^n\times(0,+\infty) ,\\
	u = g & \text{ on } \R^n\times\{0\} .
\end{cases}
\end{equation}
where we can see $u$ as a function $[0,+\infty)\to \scr S'$, and $g\in\scr S'$.

We apply the Fourier transform in the spatial variable, so that~\eqref{eq680346f3} becomes
\begin{equation}\label{eq6803470c}
\begin{cases}
	\de_t \hat u - (2\pi i)^2 |\xi|^2 \hat u = 0 & \text{ in }\R^n\times(0,+\infty) ,\\
	\hat u = \hat g & \text{ on } \R^n\times\{0\} .
\end{cases}
\end{equation}
A solution to the ODE $\de_t\hat u = -4\pi^2 |\xi|^2 \hat u$ with $\hat u(0)=\hat g$ is
\begin{equation}
	\hat u(t) =  \exp(-4\pi^2|\xi|^2 t) \hat g .
\end{equation}
Notice that, for each $t>0$, $\xi\mapsto \exp(-4\pi^2|\xi|^2 t)$ is an element of $\scr S$,
so the $\hat u(t)\in\scr S'$.
It follows that, for $t>0$, 
\begin{align}
	u(t) 
	&= \scr F^{-1}_\xi (\exp(-4\pi^2|\xi|^2 t) \hat g) \\
	&\overset{\eqref{eq680350a7}}= \scr F^{-1}_\xi (\exp(-4\pi^2|\xi|^2 t) ) * \scr F^{-1}_\xi(\hat g) \\
	&\overset{\eqref{eq67fe64f8}}= \left( \frac1{(4\pi t)^{n/2}} e^{-\frac{|x|^2}{4t}} \right) * g .
\end{align}

We have obtained again a representation formula for the solution of the heat equation, as we did in Theorem~\ref{thm67c30be5}.

%%%%%%%%%%%%%%%%%%%%%%%%%%%%%%%%%%%%%%%%%%%%%%%%%%%%%%%%%%%%%%%
\subsection{Applications to PDE: Wave equation}%(
See also \cite[\S5.D,p.177]{MR1357411} and \cite[p.194]{MR2597943}.

Let's consider the wave equation
\begin{equation}\label{eq6803521d}
\begin{cases}
	\de_t^2 u - \laplacian u = 0 & \text{ in }\R^n\times(0,+\infty) ,\\
	u = g,\ \de_tu = h & \text{ on } \R^n\times\{0\} .
\end{cases}
\end{equation}
where we can see $u$ as a function $[0,+\infty)\to \scr S'$, and $g,h\in\scr S'$.

We apply the Fourier transform in the spatial variable, so that~\eqref{eq6803521d} becomes
\begin{equation}\label{eq68035243}
\begin{cases}
	\de_t^2 \hat u - (2\pi i)^2 |\xi|^2 \hat u = 0 & \text{ in }\R^n\times(0,+\infty) ,\\
	\hat u(0) = \hat g\ \de_t\hat u(0) = \hat h & \text{ on } \R^n\times\{0\} .
\end{cases}
\end{equation}
A solution to the ODE $\de_t^2\hat u = -4\pi^2 |\xi|^2 \hat u$  is
\begin{equation}
	\hat u(t) =  \exp(2\pi i|\xi| t) A + \exp(-2\pi i|\xi| t) B .
\end{equation}
The initial data $\hat u(0)= \hat g$ and $\de_t\hat u(0) = \hat h$, give
\begin{align}
	\hat u(t) 
	&= \frac{ \exp(2\pi i|\xi| t) }{ 2 } \left(\hat g + \frac{\hat h}{2\pi i|\xi|} \right) 
		+ \frac{ \exp(-2\pi i|\xi| t) }{ 2 } \left(\hat g - \frac{\hat h}{2\pi i|\xi|} \right) \\
	&= \cos(2\pi i|\xi| t) \hat g + \frac{\sin(2\pi i|\xi| t)}{2\pi i|\xi|} \hat h .
\end{align}
These formulas give another representation for the solutions of the wave equation.
Since we know what they must be for $n=1,2,3$, then we know that $u(t)$ must be as in 
\ref{subs67cd4319}, \ref{subs67ce94b2} and \ref{subs67ce94a5}, respectively.

\begin{remark}
	Notice that the function $\xi\mapsto\frac1{|\xi|}$ belongs to $L^1_\loc(\R^n)$ for $n\ge2$.
	Since it also decays at infinity, this function is a tempered distribution.
	So, it has a well defined Fourier transform.
	You can try to compute it: I don't know what it should be, or even if there is a closed formula for it.
\end{remark}

\begin{exercise}
	Assuming $g,h\in\scr S$, write $u(t)$ from the formula 
	\begin{equation}
		\hat u(t)
		= \cos(2\pi i|\xi| t) \hat g + \frac{\sin(2\pi i|\xi| t)}{2\pi i|\xi|} \hat h .
	\end{equation}
\end{exercise}

%%%%%%%%%%%%%%%%%%%%%%%%%%%%%%%%%%%%%%%%%%%%%%%%%%%%%%%%%%%%%%%
\subsection{Applications to PDE: Bessel potentials}
(See \cite[\S4.3,p.191]{MR2597943})

%We investigate the PDE in $\R^n$:
%\begin{equation}
%	-\laplacian u + u = f .
%\end{equation}
%
%\TODO

\begin{exercise}[Bessel Potentials]
	Using the Fourier transform, give a representation to the solutions of 
	\begin{equation}
		-\laplacian u + u = f .
	\end{equation}
	{\it Hint:}
	See \cite[\S4.3,p.191]{MR2597943}.
\end{exercise}

%%%%%%%%%%%%%%%%%%%%%%%%%%%%%%%%%%%%%%%%%%%%%%%%%%%%%%%%%%%%%%%
\subsection{Applications to PDE: Eigenfunctions of the Laplacian}
%We investigate the PDE
%\begin{equation}
%	-\laplacian u = \lambda u ,
%\end{equation}
%for $\lambda\in\C$.
%
%\TODO

\begin{exercise}[Eigenfunctions of the Laplacian]
	Using the Fourier transform, give a representation to the solutions of
	\begin{equation}
		-\laplacian u = \lambda u ,
	\end{equation}
	for $\lambda\in\C$.
	For which $\lambda$ there exists a solution? 
	(see also \ref{subs6890901c})
\end{exercise}


%%%%%%%%%%%%%%%%%%%%%%%%%%%%%%%%%%%%%%%%%%%%%%%%%%%%%%%%%%%%%%%
%%%%%%%%%%%%%%%%%%%%%%%%%%%%%%%%%%%%%%%%%%%%%%%%%%%%%%%%%%%%%%%
\newpage
\part{Sobolev Spaces}




%%%%%%%%%%%%%%%%%%%%%%%%%%%%%%%%%%%%%%%%%%%%%%%%%%%%%%%%%%%%%%%
\section{Sobolev spaces}

%%%%%%%%%%%%%%%%%%%%%%%%%%%%%%%%%%%%%%%%%%%%%%%%%%%%%%%%%%%%%%%
\subsection{Definition of sobolev space}
Let $\Omega\subset\R^n$ open.
We have seen in~\ref{sub67def4a2} that functions $u\in L^1_\loc(\Omega)$ are distributions, and as such they have all derivatives $\Diff^\alpha u$.
We are interested in those functions whose distributional derivatives are in fact integrable functions.

For $m\in\N$ and $p\in[1,+\infty]$, define
\begin{equation}\label{eq680b26ee}
	W^{m,p}(\Omega) = \left\{u\in L^1_\loc(\Omega): \Diff^\alpha u\in L^p(\Omega)\ \forall\alpha\in\N^n\text{ with }|\alpha|\le m \right\} .
\end{equation}
Of these spaces, we also have a ``loc'' version:
\begin{equation}\label{eq680b26f4}
	W^{m,p}_\loc(\Omega) = \left\{u\in L^1_\loc(\Omega): \Diff^\alpha u\in L^p_\loc(\Omega)\ \forall\alpha\in\N^n\text{ with }|\alpha|\le m \right\} .
\end{equation}

%%%%%%%%%%%%%%%%%%%%%%%%%%%%%%%%%%%%%%%%%%%%%%%%%%%%%%%%%%%%%%%
\subsection{Sobolev spaces are Banach spaces}
We endow $W^{m,p}(\Omega)$ with the norm
\begin{equation}\label{eq680b27c5}
	\|u\|_{W^{m,p}(\Omega)}
	= \|u\|_{W^{m,p}}
	= \sum_{|\alpha|\le m} \|\Diff^\alpha u\|_{L^p(\Omega)} .
\end{equation} 
In fact, we may well take in place of $\|u\|_{W^{m,p}(\Omega)}$ other norms, such as
\begin{equation}\label{eq680b27ca}
	(\sum_{|\alpha|\le m} \|\Diff^\alpha u\|_{L^p(\Omega)}^q)^{1/q},
	\text{ for $q\in[1,+\infty)$, or }
	\max_{|\alpha|\le m} \|\Diff^\alpha u\|_{L^p(\Omega)} .
\end{equation}
See Exercise~\ref{eq680b28e2}.

\begin{exercise}\label{eq680b28e2}
	For $x\in\R^N$ and $q\in[1,+\infty)$, define
	\begin{equation}
		\|x\|_q = (\sum_{j=1}^N |x_j|^q )^{1/q} ,
		\text{ and }\|x\|_\infty = \max_{j\le N} |x_j| .
	\end{equation}
	Show that for every $q_1,q_2\in[1,+\infty]$ there is $L$ such that
	\begin{equation}
		\frac1L \|x\|_{q_2} \le \|x\|_{q_1} \le L \|x\|_{q_2} .
	\end{equation}
	Deduce that all norms on $W^{m,p}(\Omega)$
	\begin{equation}\label{eq680b27ca}
		(\sum_{|\alpha|\le m} \|\Diff^\alpha u\|_{L^p(\Omega)}^q)^{1/q},
		\text{ for $q\in[1,+\infty)$, or }
		\max_{|\alpha|\le m} \|\Diff^\alpha u\|_{L^p(\Omega)} .
	\end{equation}
	are biLipschitz equivalent to $\|u\|_{W^{m,p}(\Omega)}$.
\end{exercise}

\begin{lemma}\label{lem680b2dfd}
	Let $\Omega\subset\R^n$ be an open set and $p\in[1,\infty]$.
	Let $\{u_j\}_{j\in\N}\subset L^p(\Omega)$ be a sequence of functions, and $v,w\in L^p(\Omega)$.
	Fix $\ell\in\{1,\dots,n\}$.
	Suppose that $u_j\to v$ in $L^p(\Omega)$ and that $\de_\ell u_j\to w$ in $L^p(\Omega)$.
	Then $\de_\ell v = w$. 
\end{lemma}
\begin{proof}
	First of all, we claim that $u_j\to v$ in $\scr D'(\Omega)$.
	Indeed, if $\phi\in\scr D(\Omega)$, then
	\begin{align*}
	 	\left| \int_\Omega u_j(x) \phi(x) \did x - \int_\Omega v(x)\phi(x) \did x \right|
		&\le \int_\Omega |u_j(x)-v(x)| |\phi(x)| \did x \\
		&\overset{\eqref{Hölder}}\le \|u_j-v\|_{L^p(\Omega)} \|\phi\|_{L^{p'}(\Omega)}
		\rightarrow 0 .
	\end{align*}
	Therefore, as distributions, $u_j\to v$.
	It follows that $\de_\ell u_j\to \de_\ell v$ in $\scr D'(\Omega)$, see Exercise~\ref{ex67df0a00}.
	
	Since $\de_\ell u_j\to w$ in $L^p(\Omega)$, we also have $\de_\ell u_j\to w$ in $\scr D'(\Omega)$.
	Therefore, $\de_\ell v = w$, for the uniqueness of the limit in $\scr D'(\Omega)$.
\end{proof}

\begin{proposition}
	The normed vector space $( W^{m,p}(\Omega) , \|\cdot\|_{W^{m,p}(\Omega)} )$ is a Banach space.
\end{proposition}
\begin{proof}
	The proof that $\|\cdot\|_{W^{m,p}(\Omega)}$ is in fact a norm is left as an exercise, see Exercise~\ref{ex680b2d5c}.
	To show that $W^{m,p}(\Omega)$ is complete, let $\{u_j\}_{j\in\N}\subset W^{m,p}(\Omega)$ be a Cauchy sequence.
	It follows that, for ever $\alpha\in\N^n$ with $|\alpha|\le m$, the sequence $\{\Diff^\alpha u_j\}_{j\in\N}\subset L^p(\Omega)$ is a Cauchy sequence.
	Since $L^p(\Omega)$ are Banach spaces, it follows that there are $u_\alpha\in L^p(\Omega)$ such that $\Diff^\alpha u_j\to u_\alpha$ in $L^p(\Omega)$.
	Iterating Lemma~\ref{lem680b2dfd}, we obtaind that $\Diff^\alpha u_0 = u_\alpha$,
	and thus $u_0$, the limit of $u_j$ in $L^p(\Omega)$, is the limit of $u_j$ in $W^{m,p}(\Omega)$.
\end{proof}

\begin{exercise}\label{ex680b2d5c}
	Show that $\|\cdot\|_{W^{m,p}(\Omega)}$ is a norm on $W^{m,p}(\Omega)$.
\end{exercise}

\begin{exercise}
	Show that $W^{m,p}(\Omega)$ is isometric to a closed subspace of $L^p(\Omega)^N$, where $N=\#\{\alpha\in\N^n:|\alpha|\le m\}$.
\end{exercise}

\begin{exercise}\label{ex680b37f5}
	Inspired by Lemma~\ref{lem680b2dfd}, we can say with no doubt that, if $u_j\to u$ in $W^{m,p}(\Omega)$, then $u_j\to u$ in $\scr D'(\Omega)$.
	We may wonder if the converse is also true, that is:
	is it true that, if $u_j\to u$ in $\scr D'(\Omega)$ and $u\in W^{m,p}(\Omega)$,
	then $u_j\to u$ in $W^{m,p}(\Omega)$?
	If this was a video, you would pause it to think about the question yourself.
	However, this is not a video: can you stop reading?
	
	The answer is not hard. The point is to find a sequence of functions $v_j$ on $\R$ that converge distributionally to $0$, but have $L^p$ norm equal to 1.
	For instance, take $v_j(x) = \sum_{m=0}^{2^j} \one_{[m2^{-j},(m+1)2^{-j})}(x)$.
	Then $|v_j| = \one_{[0,1)}$, and $\int_\R v_j\phi\did x\to0$ for every $\phi\in\scr D(\R)$ (check it!).
	We have thus a sequence $v_j\to 0$ in $\scr D'(\R)$ with $\|v_j\|_{L^p}=1$ for all $j$.
	Take $u_j(x) = \int_{-\infty}^x v_j(t)\did t$: check that $u_j(2)=0$.
	So, $u_j$ is an absolutely continuous function $\R\to\R$ with compact support.
	Moreover, $u_j\to 0$ distributionally (check it!) but $u_j$ does not converge to $0$ in $W^{1,1}(\R)$.
	So, the answer is no.
\end{exercise}

\begin{exercise}\label{ex681772db}
	An example of weird Sobolev functions.
	For $\beta\in\R$, define $u_\beta:\R^n\to\R$,
	\begin{equation}
		u_\beta(x) = |x|^\beta .
	\end{equation}
	\begin{enumerate}[leftmargin=*]
	\item
	Show that $u_\beta\in L^p_\loc(\R^n)$ whenever $\beta>-n$.
	\item
	Show that $u_\beta\in W^{1,p}_\loc(\R^n)$ for all $\beta>1-n$.
	\item
	Take an enumeration $\{q_k\}_{k\in\N} = \Q^n\cap B(0,1)$ of points with rational coordinates inside the unit ball.
	Define
	\begin{equation}
		w_\beta(x) = \sum_{k=1}^\infty \frac{u_\beta(x-q_k)}{2^k} .
	\end{equation}
	Notice that, since $u_\beta\ge0$, then $w_\beta(x) \in[0,+\infty]$ is well defined for every $x\in\R^n$.
	Show that, if $\beta>1-n$, then $w_\beta\in W^{1,p}_\loc(\R^n)$.
	\end{enumerate}
	I want to remark that, if $n\ge2$, then $1-n<0$ and thus we can take $\beta<0$:
	in these cases, $w_\beta$ has a ``pole'' at every point in $\Q^n\cap B(0,1)$.
\end{exercise}

%%%%%%%%%%%%%%%%%%%%%%%%%%%%%%%%%%%%%%%%%%%%%%%%%%%%%%%%%%%%%%%
\subsection{Smooth approximation of Sobolev functions: characterization of Sobolev spaces}
We are going to prove the following theorem:
\begin{theorem}
	Let $\Omega\subset\R^n$, $p\in[1,+\infty)$ and $m\in\N$.
	A distribution $u\in\scr D'(\Omega)$ belongs to $W^{m,p}(\Omega)$ if and only if there is a sequence $\{u_j\}_{j\in\N}\subset C^\infty(\Omega)\cap W^{m,p}(\Omega)$ such that $\lim_{j\to\infty} \|u_j-u\|_{W^{m,p}(\Omega)} = 0$.
	
	In other words, $W^{m,p}(\Omega)$ is the closure of $C^\infty(\Omega)\cap W^{m,p}(\Omega)$ in the norm $\|\cdot\|_{W^{m,p}(\Omega)}$.
\end{theorem}

%%%%%%%%%%%%%%%%%%%%%%%%%%%%%%%%%%%%%%%%%%%%%%%%%%%%%%%%%%%%%%%
\subsection{Smooth approximation of Sobolev functions: local approximation by convolution}

Recall that we say that $w\in C^\infty(\bar\Omega)$ if there is an open set $V\subset\R^n$ with $\bar\Omega\subset V$, and there is a function $\tilde w\in C^\infty(V)$ such that $\tilde w|_{\Omega}=w$.
Notice that, if $\Omega$ is bounded, then $\bar \Omega$ is compact and thus, if  $w\in C^\infty(\bar\Omega)$ then $w\in W^{p,m}(\Omega)$ for all $p\in[1,+\infty]$ and $m\in\N$.
We will show three approximation results 

\begin{theorem}\label{thm680b49bf}
	Let $\Omega\subset\R^n$ open.
	Let $\{\rho_\epsilon\}_{\epsilon>0}$ be a standard family of mollifiers.
	For $\epsilon>0$, define
	\begin{equation}
	\begin{aligned}
		\Omega_\epsilon &= \{x\in\Omega:\dist(x,\de\Omega)>\epsilon\} \\
		&= \{x\in\Omega: \bar B(x,\epsilon)\subset\Omega\}
	\end{aligned}
	\end{equation}
	For $u\in L^1_\loc(\Omega)$, define $u_\epsilon:\Omega_\epsilon\to\C$,
	$u_\epsilon = u*\rho_\epsilon$;
	more precisely,
	\begin{equation}\label{eq680b3795}
		u_\epsilon = ((u\one_\Omega) * \rho_\epsilon)|_{\Omega_\epsilon} .
	\end{equation}
	
	If $u\in W^{m,p}_\loc(\Omega)$ and $V\Subset\Omega$ is open, then $u_\epsilon \to u$ in $W^{m,p}(V)$.
%	For every $u\in W^{m,p}(\Omega)$ there is a sequence $u_j\in C^\infty(\bar\Omega)$
%	with $u_j\to u$ in $W^{m,p}(\Omega)$.
\end{theorem}
\begin{proof}
	Recall that $V\Subset\Omega$ means that the closure of $V$ is compact and contained in $\Omega$.
	Notice also that there is $\delta>0$ such that $V\subset\Omega_\epsilon$ for all $\epsilon\in(0,\delta)$.
	So, the convergence ``$u_\epsilon \to u$ in $W^{m,p}(V)$'' makes sense eventually for $\epsilon\to0$.
	
	From Proposition~\ref{prop67e28cab}, we already know that the functions $u_\epsilon$ defined in~\eqref{eq680b3795} converge to $u\one_\Omega$ in $\scr D'(\R^n)$ as $\epsilon\to0$.
	But be careful, this does not imply convergence in $W^{m,p}(V)$, see Exercise~\ref{ex680b37f5}.
	We need some more work, and use some specific property of $\rho_\epsilon$.
	
	Since, for every $\alpha\in\N^n$, we have $\Diff^\alpha u_\epsilon = (\Diff^\alpha u)*\rho_\epsilon$, we only need to show that, if $u\in L^p_\loc(\Omega)$, then $u_\epsilon\to u$ in $L^p(V)$ for every $V\Subset \Omega$.
	This is the content of Lemma~\ref{lem680f1ffe}.
\end{proof}

\begin{lemma}\label{lem680f1ffe}
	Let $\Omega\subset\R^n$ be an open set and $u\in L^p_\loc(\Omega)$.
	If $V\Subset\Omega$, then $\lim_{\epsilon\to0}\|u_\epsilon-u\|_{L^p(V)} = 0$, 
	where $u_\epsilon$ is defined as in~\eqref{eq680b3795}.
\end{lemma}
\begin{proof}
	Let $\bar\epsilon>0$ be such that $B(V,\bar\epsilon) = \bigcup_{x\in V} B(x,\bar\epsilon) \Subset \Omega$.
	If $x\in V$ and $\epsilon\in(0,\bar\epsilon)$, then
	\begin{align*}
	 	u_\epsilon(x)
		&= \int_{B(0,\epsilon)} u(x-y) \rho_\epsilon(y) \did y \\
		&= \int_{B(0,\epsilon)} u(x-y) \rho_\epsilon(y)^p \cdot \rho_\epsilon(y)^{1-1/p} \did y \\
		&\overset{\eqref{Hölder}}\le \left(\int_{B(0,\epsilon)} |u(x-y)|^p \rho_\epsilon(y) \did y \right)^{1/p}
		\left( \int_{B(0,\epsilon)}\rho_\epsilon(y)\did y\right)^{1/p'} \\
		&= \left(\int_{B(0,\epsilon)} |u(x-y)|^p \rho_\epsilon(y) \did y \right)^{1/p} .
	\end{align*}
	Therefore, 
	\begin{align}
	 	\int_V|u_\epsilon(x)|^p \did x
		&\le \int_V \int_{B(0,\epsilon)} |u(x-y)|^p \rho_\epsilon(y) \did y\did x \\
		&= \int_{B(0,\epsilon)} \int_V |u(x-y)|^p \rho_\epsilon(y) \did x \did y \\
		&\le \int_{B(V,\epsilon)} |u(x)|^p \did x . \label{eq680f1f43}
	\end{align}
	Next, fix $\eta>0$.
	By Theorem~\ref{thm680f2023}, there exists $g\in C(\Omega)$ such that $\|g-u\|_{L^p(B(V,\bar\epsilon))} < \eta$.
	For $\epsilon\in(0,\bar\epsilon)$, we have
	\begin{align*}
		\|u-u_\epsilon\|_{L^p(V)}
		&\le \|u-g\|_{L^p(V)} + \|g-g_\epsilon\|_{L^p(V)} + \|g_\epsilon-u_\epsilon\|_{L^p(V)} \\
		&\overset{\eqref{eq680f1f43}}\le \|g-u\|_{L^p(B(V,\bar\epsilon))} + \|g-g_\epsilon\|_{L^\infty(V)} \scr L^n(V) + \|g-u\|_{L^p(B(V,\bar\epsilon))} \\
		&\le 2\eta + \|g-g_\epsilon\|_{L^\infty(V)} \scr L^n(V) ,
	\end{align*}
	Since $g$ is continuous and $V$ is compact, we know that $\|g-g_\epsilon\|_{L^\infty(V)}\to0$ as $\epsilon\to0$.
	Therefore, there is $\bar{\bar\epsilon}$ such that, if $\epsilon\in(0,\bar{\bar\epsilon})$, then $\|u-u_\epsilon\|_{L^p(V)}<3\eta$.
\end{proof}

%%%%%%%%%%%%%%%%%%%%%%%%%%%%%%%%%%%%%%%%%%%%%%%%%%%%%%%%%%%%%%%
\subsection{Smooth approximation of Sobolev functions: interior approximation}

\begin{lemma}\label{lem680f3acc}
	Let $\Omega\subset\R^n$ be an open set.
	Then there are a countable set $\scr N\subset\Omega$ and $r:\scr N\to(0,\infty)$ such that
%	, for every $x\in\Omega$,
%	\begin{equation}\label{eq680f216e}
%		0<\# \{\ell\in\N: x\in B(x_\ell,r_\ell) \} 
%		\le \# \{\ell\in\N: x\in B(x_\ell,2r_\ell) \} \le 2^{5n} < \infty .
%	\end{equation}
	\begin{enumerate}[label=(\ref{lem680f3acc}.\alph*)]
	\item\label{lem680f3acc_1}
	for every $x\in\scr N$, $B(x,2r_x)\subset \Omega$,
	\item\label{lem680f3acc_2}
	$\Omega\subset\bigcup_{x\in\scr N} B(x,r_x)$,
	\item\label{lem680f3acc_3} 
	for every $y\in\Omega$, $\# \{x\in\scr N: y\in B(x,2r_x) \} \le 2^{5n} < \infty$.
	\end{enumerate}
	
	Moreover, are also functions $\{\phi_x\}_{x\in\scr N}$ such that, for each $x\in\scr N$, $\phi_x\in C^\infty_c(B(x,r_x))$, $0\le \phi_x\le 1$,
	and such that $\sum_{x\in\scr N}\phi_x(y) = 1$ for every $y\in \Omega$.
\end{lemma}
\begin{proof}
	If $\Omega=\R^n$, the proof is left as an exercise.
	Suppose that $\Omega\neq\R^n$, that is, $\de\Omega\neq\emptyset$.
	Define $\delta:\Omega\to(0,+\infty)$, $\delta(x) = {\rm dist}(x,\de\Omega)$,
	and, for $k\in\Z$,
	\[
	\Omega_k = \{x\in\Omega: 2^k \le \delta(x) < 2^{k+1}\} .
	\]
	For each $k\in\Z$, let $\scr N_k\subset\Omega_k$ be a maximal $2^{k-3}$-separated set,
	that is, a subset of $\Omega_k$ such that, if $x,y\in\scr N_k$, then $B(x,2^{k-3})\cap B(y,2^{k-3})=\emptyset$, and for every $z\in\Omega_k$ there is $x\in\scr N_k$ with $|z-x|< 2\cdot 2^{k-3} = 2^{k-2}$.
	The sets $\scr N_k$ are countable, thus $\scr N=\bigcup_{k\in\Z}\scr N_k$ is also countable.
	We claim that the set of points $\scr N$ with the radii $r_x = 2^{k-2}$ for $x\in\scr N_k$ is the wanted set.
	
	The requirement~\ref{lem680f3acc_1} is clear.
	It is also clear that
	\begin{equation}
		\Omega 
		\subset\bigcup_{x\in\scr N} B(x,r_x)
		= \bigcup_{k\in\Z} \bigcup_{x\in\scr N_k} B(x,2^{k-2}) ,
	\end{equation}
	and thus~\ref{lem680f3acc_2}.
	
	We need to show the upper bound in~\ref{lem680f3acc_3}.
	Let $y\in\Omega$ and set
	$A(y) = \{x\in\scr N: y\in B(x,2r_x)\}$.
	Keep in mind that, if $x\in\scr N_k$, then $r_x=2^{k-2}$ and $2r_x = 2^{k-1}$.
	Let $j\in\Z$ such that $y\in\Omega_j$.
	
	Notice that
	\begin{equation}\label{eq680f38fd}
		x\in\scr N_k 
		\quad\THEN\quad
		B(x,2^{k-1}) \subset \Omega_{k-1} \cup \Omega_k \cup \Omega_{k+1} .
	\end{equation}
	Indeed, if $z\in B(x,2^{k-1})$, then
	\begin{align*}
	 	\delta(z) &\le \delta(x) + 2^{k-1} < 2^{k+1} + 2^{k-1} \le 2^{k+2} , \\
		\delta(z) &\ge \delta(x) - 2^{k-1} \ge 2^k - 2^{k-1} = 2^{k-1} .
	\end{align*}
	
	From~\eqref{eq680f38fd} we get that, if $x\in\scr N_k$ and $y\in B(x,2^{k-1})$, then $k\in\{j-1,j,j+1\}$.
	It follows that $A(y) \subset \scr N_{j-1}\cup\scr N_j \cup\scr N_{j+1}$.
	
	We thus have that $A(y)\subset B(y,2^{(j+1)-1}) = B(y,2^j)$, 
	and that balls of radius $2^{j-4}$ centered at points of $A(y)$ are pairwise disjoint.
	By Exercise~\ref{ex680f2822}, we conclude that $\#A(y) \le 2^{5n}$.
	
	The construction of the functions $\phi_x$ is left as an exercise.
\end{proof}

\begin{exercise}\label{ex680f2822}
	Let $j\in\Z$.
	Show that, if $A\subset B(0,2^j)$ is such that, for every $a_1,a_2\in A$ distinct, $B(a_1,2^{j-4})\cap B(a_2,2^{j-4}) = \emptyset$, then $\#A \le 2^{5n}$.
	
	{\it Hint:}
	The union of the pairwise disjoint balls $\{B(a,2^{j-4})\}_{a\in A}$ is a subset of $B(0,2^j)$, so its volume...
\end{exercise}

\begin{exercise}
	Construct the functions $\phi_x$ in Lemma~\ref{lem680f3acc}.
\end{exercise}

\begin{theorem}\label{thm6810811d}
	Let $\Omega\subset\R^n$ open.
	If $u\in W^{k,p}(\Omega)$, then there is a sequence $\{u_j\}_{j\in\N}\subset W^{k,p}(\Omega)\cap C^\infty(\Omega)$ such that $u_j\to u$ in $W^{k,p}(\Omega)$.
\end{theorem} 
\begin{proof}
	We use the following notation: if $B$ is the ball $B(x,r)$, then $2B$ is the ball $B(x,2r)$.
	
	From Lemma~\ref{lem680f3acc}, we get a partition of unity $\{\phi_\ell\}_{\ell\in\N}$
	and a countable family of balls $\{B_\ell\}_{\ell\in\N}$ 
	such that $\Omega = \bigcup_{\ell\in\N} B_\ell$, 
	$\spt(\phi_\ell)\subset B_\ell$,
	and $\#\{\ell:y\in 2B_\ell\}<\infty$ for all $y\in\Omega$.
	
	By Theorem~\ref{thm680b49bf},
	for every $\ell\in\N$ and $k\in\N$, there is $\epsilon_k>0$ such that the function
	\begin{equation}
		u_{k,\ell} = (\phi_\ell u)*\rho_{\epsilon_k} 
	\end{equation}
	is supported in $2B_\ell$ and $\|u_{k,\ell} - \phi_\ell u\|_{W^{m,p}(2B_\ell)} < \frac{1/k}{2^\ell}$.
	Define
	\begin{equation}
		u_k = \sum_{\ell\in\N} u_{k,\ell} .
	\end{equation}
	Notice that $u_k\in C^\infty(\Omega)$, because each $u_{k,\ell}$ is smooth and the sum is locally finite.
	Moreover,
	\begin{align*}
	 	\|u_k-u\|_{W^{m,p}(\Omega)}
		&\le \sum_{\ell\in\N} \| u_{k,\ell} - \phi_\ell u\|_{W^{m,p}(\Omega)} \\
		[\text{because }\spt(u_{k,\ell} - \phi_\ell u)\subset 2B_\ell]
		&= \sum_{\ell\in\N} \| u_{k,\ell} - \phi_\ell u\|_{W^{m,p}(2B_\ell)} \\
		&\le \sum_{\ell\in\N} \frac{1/k}{2^\ell} 
		= \frac{2}{k} .
	\end{align*}
	Therefore, $u_k\to u$ in $W^{m,p}(\Omega)$.
\end{proof}
%\begin{proof}
%	For every $x\in\Omega$, let $r_x\in(0,1]$ be such that $\bar B(x,2r_x)\subset \Omega$.
%	Then, $\{B(x,r_x)\}_{x\in\Omega}$ covers $\Omega$.
%	
%	We need to extract a cover $\scr B = \{B(x_\ell,r_\ell)\}_{\ell\in\N} \subset \{B(x,r_x)\}_{x\in\Omega}$ such that, for every $x\in\Omega$,
%	\begin{equation}
%		0< \# \{\ell\in\N: x\in B(x_\ell,2r_\ell) \} < \infty .
%	\end{equation}
%	
%	Let $\{\phi_\ell\}_{\ell\in\N}$ be a partition of unity subordinated to this cover.
%	So, for each $\ell\in\N$ there is $V_\ell = B(x_\ell,r_{x_{\ell}})$ such that $\spt(\phi_\ell)\subset V_\ell\Subset \Omega$.
%	
%	Fix $j\in \N$.
%	By Theorem~\ref{thm680b49bf},
%	for each $\ell\in\N$, there is $\epsilon\in (0,r_{x_\ell})$ such that 
%	\begin{equation}
%	\| \phi_\ell u - (\phi_\ell u)*\rho_\epsilon \|_{W^{k,p}(\Omega)} 
%	= \| \phi_\ell u - (\phi_\ell u)*\rho_\epsilon \|_{W^{k,p}(B(x,2r_x))} 
%	\le \frac{1/j}{2^\ell} .
%	\end{equation}
%	Define $u_{\ell,j}=(\phi_\ell u)*\rho_\epsilon\in C^\infty_c(\Omega)$ for such $\epsilon$.
%	Then, 
%	\TODO
%\end{proof}

%%%%%%%%%%%%%%%%%%%%%%%%%%%%%%%%%%%%%%%%%%%%%%%%%%%%%%%%%%%%%%%
\subsection{Smooth approximation of Sobolev functions: global approximation}

For a proof of the following theorem, see \cite[\S5.3.3]{MR2597943}.
\begin{theorem}
	Let $\Omega\subset\R^n$ open with $\de U$ of class $C^1$.
	If $u\in W^{k,p}(\Omega)$, then there is a sequence $\{u_j\}_{j\in\N}\subset W^{k,p}(\Omega)\cap C^\infty(\bar\Omega)$ such that $u_j\to u$ in $W^{k,p}(\Omega)$.
\end{theorem}
%\begin{proof}
%	\TODO
%\end{proof}

%%%%%%%%%%%%%%%%%%%%%%%%%%%%%%%%%%%%%%%%%%%%%%%%%%%%%%%%%%%%%%%
\subsection{The closure of smooth functions with compact support}
%We may wonder what is the closure of $C^\infty_c(\Omega)$ in $W^{m,p}(\Omega)$.
Using Theorem~\ref{thm680f2023} and some smoothing, one can show that $C^\infty_c(\Omega)$ is dense in $L^p(\Omega)$, see Exercise~\ref{ex680f4988}
However, it is NOT true in general that $C^\infty_c(\Omega)$ is dense in $W^{m,p}(\Omega)$.
 
\begin{exercise}
	Show that, for every $p\in[1,\infty)$, the constant function $u\equiv1$ belongs to $W^{1,p}((0,1))$, but it is not the limit in $W^{1,p}$ of functions $C^\infty_c((0,1))$.
\end{exercise}

So, for $\Omega\subset\R^n$ open, $m\in\N$ and $p\in[1,+\infty]$, we define
\begin{equation}
	W^{m,p}_0(\Omega) = \{u\in W^{m,p}(\Omega): \exists\{u_j\}_{j\in\N}\subset C^\infty_c(\Omega)\text{ such that }u_j\to u\text{ in }W^{m,p}(\Omega)\} .
\end{equation}

\begin{exercise}\label{ex680f4988}
	Show that $C^\infty_c(\Omega)$ is dense in $L^p(\Omega)$.
\end{exercise}

\begin{theorem}\label{thm6810ea3d}
	For every $m\in\N$ and $p\in[1,\infty)$,
	$W^{m,p}_0(\R^n) = W^{m,p}(\R^n)$.
\end{theorem}
\begin{proof}
	Given $u\in W^{m,p}(\R^n)$, we need to find a sequence $\{u_j\}_{j\in\N}\subset C^\infty_c(\R^n)$ that converges to $u$ in $W^{m,p}(\R^n)$.
	By Theorem~\ref{thm6810811d} and a diagonal argument, we can assume $u\in C^\infty(\R^n)$.
	
	Let $\zeta:\R^n\to[0,1]$ be a smooth function such that 
	\begin{equation}
		B(0,1) \subset \{\zeta = 1\} \subset \spt(\zeta) \subset \bar B(0,2) .
	\end{equation}
	For $R>0$ and $x\in\R^n$, define $\zeta_R(x) = \zeta(x/R)$.
	Notice that, for every $R>0$, $\beta\in\N^n$ and $x\in\R^n$,
	\begin{equation}
		|\Diff^\beta\zeta_R(x)| = |R^{-|\beta|} \Diff^\beta\zeta(x/R)| 
		\le \frac1{R^{|\alpha|}} \|\Diff^\beta\zeta\|_{L^\infty(\R^n)} .
	\end{equation}
	We shall approximate $u$ with $\zeta_Ru$ as $R\to\infty$.
	Notice that, for $R>0$ and $\alpha\in\N^n$, we have
	\begin{align}
		\|\Diff^\alpha u - \Diff^\alpha(\zeta_Ru)\|_{L^p(\R^n)}
		&\le \| \Diff^\alpha u - \zeta_R\Diff^\alpha u\|_{L^p(\R^n)}
			+ \sum_{\beta\le\alpha}\binom{\alpha}{\beta} \| \Diff^\beta\zeta_R \cdot \Diff^{\alpha-\beta}u\|_{L^p(\R^n)} \\
		&\hspace{-3cm}\le \left( \int_{\R^n\setminus B(0,R)} |\Diff^\alpha u(x)|^p \did x \right)^{1/p}
			+ \sum_{\beta\le\alpha}\binom{\alpha}{\beta} \frac1{R^{|\beta|}} \|\Diff^\alpha\zeta\|_{L^\infty(\R^n)} \|\Diff^{\alpha-\beta}u\|_{L^p(\R^n)} .
	\end{align}
	Since 
	\begin{equation}
		\lim_{R\to\infty} \int_{\R^n\setminus B(0,R)} |\Diff^\alpha u(x)|^p \did x = 0,
	\end{equation}
	we conclude that 
	$\lim_{R\to\infty} \|u-\zeta_Ru\|_{W^{m,p}(\R^n)} = 0$.
\end{proof}

%%%%%%%%%%%%%%%%%%%%%%%%%%%%%%%%%%%%%%%%%%%%%%%%%%%%%%%%%%%%%%%
\subsection{Sobolev inequalities: $1\le p < n$}

\begin{theorem}\label{thm681085c3}
	Let $n\ge2$.
	If $u\in C^1_c(\R^n)$, 
	\begin{equation}\label{eq681085c8}
		\left( \int_{\R^n} |u(x)|^{\frac{n}{n-1}} \did x \right)^{\frac{n-1}{n}}
		\le \int_{\R^n} |\grad u(x)| \did x .
	\end{equation}
\end{theorem}
\begin{proof}
	Let $u\in C^1_c(\R^n)$.
	For $j\in\{1,\dots,n\}$, $x\in\R^n$ and $y\in\R$ we denote by $u(x|y)$ the evaluation
	$u(x|_jy) = u(x_1,\dots,x_{j-1},y,x_{j+1},\dots,x_n)$.
	Notice that, for every $x\in\R^n$ and $j\in\{1,\dots,n\}$,
	\begin{equation}
		u(x) = \int_{-\infty}^{x_j} \de_ju(x|_jy_j) \did y_j .
	\end{equation}
	Therefore, for every $x\in\R^n$,
	\begin{equation}\label{eq6810e5a6}
		|u(x)|^{\frac{n}{n-1}} 
		\le \prod_{j=1}^n \left( \int_{-\infty}^{\infty} |\de_ju(x|_jy_j)| \did y_j \right)^{\frac1{n-1}} 
		\le \prod_{j=1}^n \left( \int_{\R} |\grad u(x|_jy_j)| \did y_j \right)^{\frac1{n-1}} .
	\end{equation}
	We claim that, for each $\ell\in\{1,\dots,n\}$, we have
	\begin{equation}\label{eq6810e541}
		\int_{\R^\ell} |u(x)|^{\frac{n}{n-1}}  \did x_{\le \ell}
		\le \left( \int_{\R^\ell} |\grad u(x)| \did x_{\le \ell} \right)^{\frac{\ell}{n-1}} 
		 \prod_{j=\ell+1}^n \left( \int_{\R^\ell} \int_\R |\grad u(x|_jy_j)| \did y_j \did x_{\le \ell}\right)^{\frac{1}{n-1}} .
	\end{equation}
	where $\did x_{\le \ell} = \did x_1\cdots\did x_\ell$.
	To prove~\eqref{eq6810e541} for $\ell=1$, we integrate~\eqref{eq6810e5a6} in $x_1$:
	\begin{align}
		&\int_\R |u(x)|^{\frac{n}{n-1}} \did x_1 \\
		&\overset{\eqref{eq6810e5a6}}\le \int_\R \prod_{j=1}^n \left( \int_{\R} |\grad u(x|_jy_j)| \did y_j \right)^{\frac1{n-1}} \did x_1 \\
		&= \left( \int_{\R} |\grad u(x|_1y_1)| \did y_1 \right)^{\frac1{n-1}}
			\int_\R \prod_{j=2}^n \left( \int_{\R} |\grad u(x|_jy_j)| \did y_j \right)^{\frac1{n-1}} \did x_1 \\
		&\overset{\eqref{General Hölder}}\le \left( \int_{\R} |\grad u(x)| \did x_1 \right)^{\frac1{n-1}}
			\prod_{j=2}^n \left( \int_\R  \int_{\R} |\grad u(x|_jy_j)| \did y_j  \did x_1 \right)^{\frac1{n-1}} .
	\end{align}
	Next, given~\eqref{eq6810e541} for $\ell = k\in\{1,\dots,n-1\}$, we prove the same for $\ell = k+1$ integrating in $x_{k+1}$:
	\begin{align*}
	 	&\int_\R \int_{\R^k} |u(x)|^{\frac{n}{n-1}}  \did x_{\le k} \did x_{k+1} \\
		&\overset{\eqref{eq6810e541}}\le \int_\R \left( \int_{\R^k} |\grad u(x)| \did x_{\le k} \right)^{\frac{k}{n-1}} 
		 \prod_{j=k+1}^n \left( \int_{\R^k} \int_\R |\grad u(x|_jy_j)| \did y_j \did x_{\le k}\right)^{\frac{1}{n-1}} \did x_{k+1} \\
		&= \left( \int_{\R^k} \int_\R |\grad u(x|_{k+1}y_{k+1})| \did y_{k+1} \did x_{\le k}\right)^{\frac{1}{n-1}} \times\\
		&\quad\times \int_\R \left( \int_{\R^k} |\grad u(x)| \did x_{\le k} \right)^{\frac{k}{n-1}} 
		 \prod_{j=k+2}^n \left( \int_{\R^k} \int_\R |\grad u(x|_jy_j)| \did y_j \did x_{\le k}\right)^{\frac{1}{n-1}} \did x_{k+1} \\
		 &\overset{\eqref{General Hölder}}\le\left( \int_{\R^k} \int_\R |\grad u(x|_{k+1}y_{k+1})| \did y_{k+1} \did x_{\le k}\right)^{\frac{1}{n-1}} \times\\
		&\quad\times \left( \int_\R\int_{\R^k} |\grad u(x)| \did x_{\le k} \did x_{k+1}\right)^{\frac{k}{n-1}} 
		 \prod_{j=k+2}^n \left( \int_\R\int_{\R^k} \int_\R |\grad u(x|_jy_j)| \did y_j \did x_{\le k} \did x_{k+1}\right)^{\frac{1}{n-1}}  ,
	\end{align*}
	which gives~\eqref{eq6810e541} for $\ell = k+1$. 
	
	Finally, notice that~\eqref{eq6810e541} for $\ell=n$ is~\eqref{eq681085c8}.
\end{proof}

\begin{corollary}\label{cor6810ee47}
	Let $n\ge 2$.
	Then~\eqref{eq681085c8} holds for all $u\in W^{1,1}(\R^n)$.
\end{corollary}
\begin{proof}
	Let $u\in W^{1,1}(\R^n)$.
	By Theorem~\ref{thm6810ea3d}, there is a sequence $\{u_j\}_{j\in\N}\subset C^1_c(\R^n)$ such that $u_j\to u$ in $W^{1,1}(\R^n)$.
	Notice that, 
	\begin{equation}
		\|u_j-u_k\|_{L^{\frac{n}{n-1}}} 
		\overset{\eqref{eq681085c8}}\le \|\grad(u_j-u_k)\|_{L^1}
		\le \|u_j-u_k\|_{W^{1,1}} .
	\end{equation}
	Therefore, $\{u_j\}_j$ is a Cauchy sequence in $L^{\frac{n}{n-1}}(\R^n)$.
	Since $u_j\to u$ in $L^1$ already, then $u_j$ must converge to $u$ in $L^{\frac{n}{n-1}}(\R^n)$ too.
	We conclude that
	\begin{equation}
		\|u\|_{L^{\frac{n}{n-1}}} 
		= \lim_{j\to\infty} \|u_j\|_{L^{\frac{n}{n-1}}}
		\overset{\eqref{eq681085c8}}\le \lim_{j\to\infty} \|\grad u_j\|_{L^1}
		= \|\grad u\|_{L^1} .
	\end{equation}
\end{proof}

\begin{exercise}
	In Theorem~\ref{thm681085c3} we required $n\ge 2$. What happens when $n=1$?
\end{exercise}

%%%%%%%%%%%%%%%%%%%%%%%%%%%%%%%%%%%%%%%%%%%%%%%%%%%%%%%%%%%%%%%
\subsection{A remark on the isoperimetric inequality}

The inequality~\eqref{eq681085c8} has a geometric meaning: it is indeed equivalent to the \emph{isoperimetric inequality}.
More precisely, the characteristic function $\one_E$ of a measurable $E\subset\R^n$ belongs to $L^1_\loc(\R^n)$ and thus it defines a distribution: $\one_E\in\scr D'(\R^n)$.
Define the \emph{perimeter} of $E$ as the \emph{total variation} of $\Diff E$, that is,
\begin{equation}
	\operatorname{Per}(E) 
	= \sup\{\sum_{j=1}^n\de_j\one_E[\phi_j] : \phi_j\in\scr D,\ \|\phi_j\|_{L^\infty} \le 1 \} .
\end{equation}
In order to put this formula into perspective, we define for $u\in\scr D'$
\begin{equation}
	|\Diff u|(\R^n)
	= \sup\{\sum_{j=1}^n\de_ju[\phi_j] : \phi_j\in\scr D,\ \|\phi_j\|_{L^\infty} \le 1 \} \in [0,+\infty] .
\end{equation}
It is clear that $|\Diff u|(\R^n)<\infty$ if and only if $|\Diff u|$ is a Radon measure, as seen in Proposition~\ref{prop6810ed13}.
In this case, we say that $u$ has \emph{bounded variation}.
The space of \emph{functions with bounded variation}, \emph{BV functions} for short,
is 
\begin{equation}
	\operatorname{BV}(\R^n)
	= \{u\in L^1(\R^n) : |\Diff u|(\R^n)<\infty \} .
\end{equation}
If $u\in C^\infty_c(\R^n)$, then $|\Diff u|(\R^n) = \int_{\R^n} |\grad u| \did x$.
One can show that~\eqref{eq681085c8} extends to BV functions: if $u\in \operatorname{BV}(\R^n)$, then
\begin{equation}\label{eq6810ee38}
	\|u\|_{L^{\frac{n}{n-1}}} \le |\Diff u|(\R^n) .
\end{equation}
The proof is by smooth approximation, similarly to the proof of Corollary~\ref{cor6810ee47}.
If we apply~\eqref{eq6810ee38} to $\one_E$, we get
\begin{equation}\label{eq6810eea6}
	|E|^{\frac{n-1}{n}} \le \operatorname{Per}(E) ,
\end{equation}
where $|E|$ is the volume of $E$.
This is the isoperimetric inequality in $\R^n$. 
Equality in~\eqref{eq6810eea6} is attained only when $E$ is an Euclidean ball.

In fact, one can show that from~\eqref{eq6810eea6}, using the coarea inequality one can get~\eqref{eq681085c8} for $u\in C^1_c(\R^n)$.



%%%%%%%%%%%%%%%%%%%%%%%%%%%%%%%%%%%%%%%%%%%%%%%%%%%%%%%%%%%%%%%
\subsection{Gagliardo–Nirenberg–Sobolev inequality}

The \emph{Sobolev conjugate} of $p\in[1,n)$ is the number $p^*$ such that
\begin{equation}
	\frac{1}{p} - \frac{1}{p^*} = \frac{1}{n} ,
	\text{ or, equivalently, }
	 \frac{1}{p^*} = \frac{1}{p} - \frac{1}{n} .
\end{equation}
In other words,
\begin{equation}
	p^* = \frac{np}{n-p} .
\end{equation}
Notice that, \eqref{eq681085c8} is \eqref{eq6810881c} for $p=1$.

\begin{theorem}[Gagliardo–Nirenberg–Sobolev inequality]\label{thm68108820}
	Let $n\ge2$ and $p\in[1,n)$ with Sobolev conjugate $p^* = \frac{np}{n-p}$.
	There exists $C\in\R$ such that, for all $u\in C^1_c(\R^n)$,
	\begin{equation}\label{eq6810881c}
		\|u\|_{L^{p^*}(\R^n)} \le C \|\grad u\|_{L^p(\R^n)} .
	\end{equation}
	In fact, we can take $C=\frac{p(n-1)}{n-p} = \frac{p^*}{1^*}$.
\end{theorem}
\begin{proof}
	If $p=1$, then~\eqref{eq6810881c} is~\eqref{eq681085c8}.
	Therefore, we assume $p>1$.
	
	Fix $u\in C^1_c(\R^n)$.
	Let $\gamma>1$ and set $v=|u|^\gamma$.
	By Exercise~\ref{ex68123356}, $v\in C^1_c(\R^n)$.
	Therefore, by Corollary~\ref{cor6810ee47}, we have
	\begin{align}
		\left( \int_{\R^n} |u|^{\gamma\frac{n}{n-1}} \did x \right)^{\frac{n-1}{n}}
		&= \left( \int_{\R^n} v^{\frac{n}{n-1}} \did x \right)^{\frac{n-1}{n}} \\
		&\overset{\eqref{eq681085c8}}\le \int_{\R^n} |\grad v(x)|\did x \\
		&= \gamma \int_{\R^n} |u|^{\gamma-1} |\grad u(x)|\did x \\
		&\overset{\eqref{Hölder}}\le \gamma \left( \int_{\R^n} |u|^{\frac{p}{p-1}(\gamma-1)}\did x \right)^{\frac{p-1}{p}}
			\left( \int_{\R^n} |\grad u(x)|^p \did x \right)^{1/p} .
	\end{align}
	We want $\gamma\frac{n}{n-1} = \frac{p}{p-1}(\gamma-1)$, that is $\gamma = \frac{p(n-1)}{n-p}$.
	Notice that, if $p\in(1,n)$, then $\gamma>1$.
	If $u\equiv0$, then~\eqref{eq6810881c} trivially holds.
	So, we can assume $u\neq0$, and thus the above estimates, with this $\gamma$, gives
	\begin{equation}
		\left( \int_{\R^n} |u|^{\gamma\frac{n}{n-1}} \did x \right)^{\frac{n-1}{n}-\frac{p-1}{p}} 
		\le \gamma \left( \int_{\R^n} |\grad u(x)|^p \did x \right)^{1/p} .
	\end{equation}
	A direct computations gives $\gamma\frac{n}{n-1} = p^*$ and $\frac{n-1}{n}-\frac{p-1}{p} = \frac1{p^*}$.
	So, we have~\eqref{eq6810881c}.
\end{proof}

\begin{exercise}\label{ex68123356}
	Show that, if $u\in C^1_c(\R^n)$ and $\gamma>1$, then $|u|^\gamma\in C^1_c(\R^n)$ and $\grad(|u|^\gamma) = \gamma |u|^{\gamma-1} \grad u$.
	
	{\it Hint:}
	on the set $\{u\neq 0\}$, there the statement is trivial.
	What remains to show is that, if $u(x)=0$, then $|u|^\gamma$ is differentiable at $x$ with derivative equal to $0$.
\end{exercise}

\begin{corollary}\label{cor681fa13d}
	Let $n\ge 2$ and $p\in[1,n)$ with Sobolev conjugate $p^* = \frac{np}{n-p}$.
	Then~\eqref{eq6810881c} holds for all $u\in W^{1,p}(\R^n)$.
\end{corollary}
\begin{proof}
	Let $u\in W^{1,p}(\R^n)$.
	By Theorem~\ref{thm6810ea3d}, there is a sequence $\{u_j\}_{j\in\N}\subset C^1_c(\R^n)$ such that $u_j\to u$ in $W^{1,p}(\R^n)$.
	Notice that, 
	\begin{equation}
		\|u_j-u_k\|_{L^{p^*}} 
		\overset{\eqref{eq6810881c}}\le \|\grad(u_j-u_k)\|_{L^p}
		\le \|u_j-u_k\|_{W^{1,p}} .
	\end{equation}
	Therefore, $\{u_j\}_j$ is a Cauchy sequence in $L^{p^*}(\R^n)$.
	Since $u_j\to u$ in $L^1$ already, then $u_j$ must converge to $u$ in $L^{p^*}(\R^n)$ too.
	We conclude that
	\begin{equation}
		\|u\|_{L^{p^*}} 
		= \lim_{j\to\infty} \|u_j\|_{L^{p^*}}
		\overset{\eqref{eq6810881c}}\le \lim_{j\to\infty} \|\grad u_j\|_{L^p}
		= \|\grad u\|_{L^p} .
	\end{equation}
\end{proof}

%%%%%%%%%%%%%%%%%%%%%%%%%%%%%%%%%%%%%%%%%%%%%%%%%%%%%%%%%%%%%%%
\subsection{Sobolev inequalities: $n < p < \infty$. Morrey's inequality}

\begin{theorem}[Morrey's inequality]\label{thm681726ad}
	Let $n\ge2$ and $p\in(n,\infty)$.
	There is $C\in\R^n$ such that, for every $o\in\R$, $R>0$ and $u\in W^{1,p}(B(o,4R))$,
	and for two every Lebesgue points $x,y\in B(o,R)$ of $u$ (thus, for a.e.~$x,y\in B(o,R)$),
	\begin{equation}\label{eq68177bc8}
%	\text{for a.e.~} x,y\in B(o,R), \qquad
		|u(x)-u(y)| \le C |x-y|^{1-n/p} \|\grad u\|_{L^p(B(o,4R))} .
	\end{equation}
	In particular, if $\Omega\subset\R^n$ is open, then every $u\in W^{1,p}_\loc(\Omega)$ has a locally $(1-n/p)$-Hölder continuous representative.
	
	The constant $C$, which depends on $p$ and $n$, can be taken to be $C=\frac{ 2^n (n\omega_n)^{\frac{p-1}{p}} }{ \omega_{n-1} } \left(\frac{p-1}{p-n}\right)^{\frac{p-1}{p}}$.
\end{theorem}

We shall prove Theorem~\ref{thm681726ad} after a few intermediate statements.

\begin{proposition}\label{prop6816457e}
	For every $n\ge 2$ there is $C\in\R$ such that, for every $u\in C^1(\R^n)$, every $x\in\R^n$ and every $r>0$, 
	\begin{equation}\label{eq681641a0}
		\fint_{B(x,r)} |u(y)-u(x)| \did y 
		\le C \int_{B(x,r)} \frac{ |\grad u(y)| }{ |y-x|^{n-1} } \did y .
	\end{equation}
	The constant $C$ can be taken equal to $C=\frac{1}{\omega_n n}$.
\end{proposition}
\begin{proof}
	\begin{align*}
	\int_{B(x,r)} |u(y)-u(x)|\did y
	&= \int_0^r \int_{\de B(x,s)} |u(y)-u(x)| \did S(y) \did s \\
	&= \int_0^r \int_{\de B(0,1)} |u(x+sz)-u(x)| s^{n-1}\did S(z) \did s \\
	&= \int_0^r \int_{\de B(0,1)} \left| \int_0^s \grad u(x+sz)\cdot z \did t \right| s^{n-1}\did S(z) \did s \\
	&\le \int_0^r \int_{\de B(0,1)}  \int_0^s |\grad u|(x+sz) \did t s^{n-1}\did S(z) \did s \\
	&= \int_0^r  \int_0^s \int_{\de B(0,1)} \frac{|\grad u|(x+sz)}{t^{n-1}} t^{n-1} \did S(z) \did t s^{n-1} \did s \\
	&= \int_0^r  \int_{B(x,s)} \frac{|\grad u|(y)}{|y-x|^{n-1}} \did y s^{n-1} \did s \\
	&\le \int_0^r  \int_{B(x,r)} \frac{|\grad u|(y)}{|y-x|^{n-1}} \did y s^{n-1} \did s \\
	&= \frac{r^n}{n} \int_{B(x,r)} \frac{|\grad u|(y)}{|y-x|^{n-1}} \did y  .
	\end{align*}
	Since the volume of $B(x,r)$ is $\omega_nr^n$, then we obtain~\eqref{eq681641a0} with $C=\frac{1}{\omega_n n}$. 
\end{proof}

\begin{exercise}
	Compute left and right hand sides of~\eqref{eq681641a0} for $u(y)=|y|$.
\end{exercise}

\begin{remark}
	Proposition~\ref{prop6816457e} is important because the integral kernel in the right-hand side of~\eqref{eq681641a0} is the famous \emph{Riesz potential}.
\end{remark}

\begin{exercise}
	Show that,
	for every $n\ge2$, for every $u\in C^1(\R^n)$, every $x\in\R^n$ and every $r>0$,
	\begin{equation}\label{eq68356a65}
		\left| \fint_{B(x,r)} u(y) \did y -u(x) \right|
		\le \fint_{B(x,r)} |u(y)-u(x)| \did y .
	\end{equation}
	
	{\it Solution.}
	Since $\fint_{B(x,r)}\did y =1$, we have
	\begin{equation}
		\left| \fint_{B(x,r)} u(y) \did y -u(x) \right|
		= \left| \fint_{B(x,r)} (u(y)  -u(x)) \did y \right|
		\le \fint_{B(x,r)} |u(y)-u(x)| \did y .
	\end{equation}
\end{exercise}

\begin{proposition}\label{prop683569eb}
	For every $n\ge2$ and $p\in(n,+\infty)$, there is $C(p,n)\in\R$ such that, for every $u\in C^1(\R^n)$, every $x\in\R^n$ and every $r>0$, 
	\begin{equation}\label{eq6816524f} 
		\fint_{B(x,r)} |u(y)-u(x)| \did y 
		\le C(p,n) r^{1-\frac{n}{p}} \|\grad u\|_{L^p(B(x,r))} .
	\end{equation}
	The constant $C(p,n)$ can be taken equal to $C(p,n)=(n\omega_n)^{-\frac1p} \left(\frac{p-1}{p-n}\right)^{\frac{p-1}{p}}$.
\end{proposition}
\begin{proof}
	We simply apply the Hölder inequality to~\eqref{eq681641a0}:
	\begin{multline}
		\fint_{B(x,r)} |u(y)-u(x)| \did y 
		\overset{\eqref{eq681641a0}}\le \frac{1}{\omega_n n} \int_{B(x,r)} \frac{ |\grad u(y)| }{ |y-x|^{n-1} } \did y \\
		\overset{\eqref{Hölder}}\le \frac{1}{\omega_n n} 
			\left( \int_{B(x,r)} |\grad u(y)|^p \did y \right)^{\frac1p}
			\left( \int_{B(x,r)} \frac{1}{ |y-x|^{(n-1)\frac{p}{p-1}} } \did y  \right)^{\frac{p-1}{p}} .
	\end{multline}
	Notice that, using polar coordinates,
	\begin{align}
		\int_{B(x,r)} \frac{1}{ |y-x|^{(n-1)\frac{p}{p-1}} } \did y
		&= \int_0^r n\omega_n s^{n-1} \frac{1}{s^{(n-1)\frac{p}{p-1}} } \did s \\
		&= n\omega_n  \int_0^r s^{-\frac{n-1}{p-1}} \did s \\
		&= n\omega_n \frac{p-1}{p-n} r^{\frac{p-n}{p-1}} ,
	\end{align}
	where $p>n$ ensures that $-\frac{n-1}{p-1}>-1$ and thus that the above integral is finite.
	Therefore,
	\begin{equation}
		\left( \int_{B(x,r)} \frac{1}{ |y-x|^{(n-1)\frac{p}{p-1}} } \did y  \right)^{\frac{p-1}{p}}
		= \left(n\omega_n\frac{p-1}{p-n}\right)^{\frac{p-1}{p}} r^{\frac{p-n}{p}} ,
	\end{equation}
	from which we conclude~\eqref{eq681641a0} with $C(p,n) = (n\omega_n)^{\frac{p-1}{p}-1} \left(\frac{p-1}{p-n}\right)^{\frac{p-1}{p}} = (n\omega_n)^{-\frac1p} \left(\frac{p-1}{p-n}\right)^{\frac{p-1}{p}}$.
	Finally, since We have obtained
\end{proof}

\begin{remark}\label{rem68165c8c}
	For $W\subset\R^n$ bounded and with positive volume, and $u\in L^1_\loc(\R^n)$, define
	\begin{equation}
		u_W = \fint_W u(y)\did y = \frac1{|W|} \int_W u(y) \did y ,
	\end{equation}
	where $|W|=\scr L^n(W)$ is the volume of $W$.
%	If $B= B(x,r)$, then~\eqref{eq6816524f} can be rewritten as
%	\begin{equation}
%		|u(x) - u_{B(x,r)}| \le  C(p,n) r^{1-\frac{n}{p}} \|\grad u\|_{L^p(B(x,r))} .
%	\end{equation}
	In the proof of Proposition~\ref{prop68165fd2}, we will use an idea that is worth to keep in mind and understand correctly. 
	Suppose we have two points $x,z\in\R^n$ with $r=|x-z|$, and set $W=B(x,r)\cap B(z,r)$.
	Clearly we have $|u(x)-u(z)|\le |u(x) - u_W| + |u_W - u(z)|\le \fint_W |u(y)-u(x)| \did y + \fint_W |u(y)-u(z)| \did y$.
	We now do the following:
	\begin{align}
		\fint_W |u(y)-u(x)| \did y
		&= \frac{|B(x,r)|}{|W|} \frac{1}{|B(x,r)|} \int_W |u(y)-u(x)| \did y \\
		&\le \frac{|B(x,r)|}{|W|} \frac{1}{|B(x,r)|} \int_{B(x,r)} |u(y)-u(x)| \did y \\
		&= \frac{|B(x,r)|}{|W|} \fint_{B(x,r)} |u(y)-u(x)| \did y .
	\end{align}
	Now we have that $\frac{|B(x,r)|}{|W|} = \frac{|B(x,r)|}{|B(x,r)\cap B(z,r)|}$ is a constant that depends only on $n$.
	Apart of computing it explicitely, see Exercise~\ref{ex681659a2}, we can convince ourselves that it is a constant because, clearly, for every $v\in\R^n$, $\lambda>0$ and $O\in\Orth(n)$,
	\begin{equation}
%		\frac{|B(x,r)|}{|B(x,r)\cap B(z,r)|}
		 \frac{|B(x+v,r)|}{|B(x+v,r)\cap B(z+v,r)|}
		= \frac{|B(x,\lambda r)|}{|B(x,\lambda r)\cap B(z,\lambda r)|}
		= \frac{|B(Ox,r)|}{|B(Ox,r)\cap B(Oz,r)|} .
	\end{equation}
	So, $\frac{|B(x,r)|}{|B(x,r)\cap B(z,r)|} = \frac{|B(0,1)|}{|B(0,1)\cap B(e_1,1)|}$.
\end{remark}
\begin{exercise}\label{ex681659a2}
	Compute $\frac{|B(x,r)|}{|B(x,r)\cap B(z,r)|}$.

	{\it Solution.}
	First check that 
	$B(0,1)\cap B(e_1,1) = \{(y_1,y')\in\R^n : y_1\in(0,1),\ y'\in\R^{n-1},\ |y'|^2\le y_1^2,\ |y'|^2 \le 1-y_1^2\}$.
	Then check $|B(0,1)\cap B(e_1,1)| = 2\int_0^{1/2} \omega_{n-1} y_1^{n-1} \did y_1 = \frac{\omega_{n-1}}{n2^{n-1}}$.
	Finally, $\frac{|B(x,r)|}{|B(x,r)\cap B(z,r)|} = \frac{|B(0,1)|}{|B(0,1)\cap B(e_1,1)|} = \frac{n\omega_n2^{n-1}}{\omega_{n-1}}$.
\end{exercise}

\begin{proposition}[Morrey's inequality for smooth functions]\label{prop68165fd2}
	For every $n\ge2$ and $p>n$,
	there is $C\in\R$ such that, 
	for every $o\in\R$, $R>0$ and $u\in C^1(B(o,3R))$,
%	for every $u\in C^1(\R^n)$,
	\begin{equation}\label{eq68166038}
	\forall x,y\in B(o,R)\qquad
		|u(x)-u(y)| \le C |x-y|^{1-n/p} \|\grad u\|_{L^p(B(o,3R))} .
	\end{equation}
	The constant $C$, which depends on $p$ and $n$, can be taken to be $C=\frac{ 2^n (n\omega_n)^{\frac{p-1}{p}} }{ \omega_{n-1} } \left(\frac{p-1}{p-n}\right)^{\frac{p-1}{p}}$.
\end{proposition}
\begin{proof}
	Let $x,y\in B(o,R)$, and set $r=|x-y|\le 2R$ and $W=B(x,r)\cap B(y,r)$.
	Using the argument described in Remark~\ref{rem68165c8c}, with constant $C=\frac{|B(0,1)|}{|B(0,1)\cap B(e_1,1)|} = \frac{n\omega_n2^{n-1}}{\omega_{n-1}}$ 
	(see Exercise~\ref{ex681659a2} for the explicit value of $C$),
	\begin{align}
		|u(x)-u(y)|
		&\le |u(x) - u_W| + |u_W - u(y)| \\
		&\le \fint_W |u(z)-u(x)| \did z + \fint_W |u(z)-u(y)| \did z \\
		&\le C \fint_{B(x,r)} |u(z)-u(x)| \did z + C \fint_{B(y,r)} |u(z)-u(y)| \did z \\
		&\overset{\eqref{eq6816524f}}\le C C(p,n) r^{1-\frac{n}{p}} \|\grad u\|_{L^p(B(x,r))} + C C(p,n) r^{1-\frac{n}{p}} \|\grad u\|_{L^p(B(y,r))} \\
		&\overset{(*)}\le 2 C C(p,n) |x-y|^{1-\frac{n}{p}} \|\grad u\|_{L^p(B(o,3R))} .
%		&= 2 C C(p,n) |x-z|^{1-\frac{n}{p}} \|\grad u\|_{L^p(B(o,3R))}
	\end{align}
	In $(*)$ we have used that $B(x,r) \subset B(x,2R)\subset B(o,3R)$.
	Since $C(p,n)=(n\omega_n)^{-\frac1p} \left(\frac{p-1}{p-n}\right)^{\frac{p-1}{p}}$,
	\begin{equation}\label{eq68165dde}
		2 C C(p,n) = \frac{ 2^n (n\omega_n)^{\frac{p-1}{p}} }{ \omega_{n-1} } \left(\frac{p-1}{p-n}\right)^{\frac{p-1}{p}} .
	\end{equation}
\end{proof}

\begin{remark}
	The constant we have found is not necessarily optimal.
	Finding optimal constants in Sobolev inequalities may be quite tricky.
	See Remark~\ref{rem6816609d} for another constant.
\end{remark}

\begin{remark}\label{rem6816609d}
	The proof presented here has been taken from Evans \cite[\S5.6.2]{MR2597943}.
	Parviainen in \cite{Parviainen} proves Proposition~\ref{prop68165fd2} with a slightly different approach. 
	First, they don't go through the Riesz potential as in Proposition~\ref{prop6816457e}.
	Second, they use cubes instead of balls.
	Third, the constant they get in~\eqref{eq68166038} is $C=\frac{2np}{p-n}$.
\end{remark}

\begin{exercise}
	For $n\ge2$ and $p>n$, which of these constants is better in Morrey's inequality?
	\begin{equation}
		A = \frac{ 2^n (n\omega_n)^{\frac{p-1}{p}} }{ \omega_{n-1} } \left(\frac{p-1}{p-n}\right)^{\frac{p-1}{p}},
		\qquad
		B = \frac{2np}{p-n} .
	\end{equation}
\end{exercise}

\begin{exercise}\label{ex68172771}
	Recall that a point $x\in\R^n$ is a \emph{Lebesgue point} of $u\in L^1_\loc(\R^n)$ if
	\begin{equation}
		\lim_{r\to0} \fint_{B(x,r)} |u(y)-u(x)| \did y = 0 .
	\end{equation}
	Recall also that, if $u\in L^1_\loc(\R^n)$, then almost every $x\in\R^n$ is a Lebesgue point of $u$; see \cite[Theorem (3.20),p.93]{MR767633}.
	
	Let $\{\rho_\epsilon\}_{\epsilon>0}$ be a family of standard mollifiers on $\R^n$.
	For $u\in L^1_\loc(\R^n)$, set $u_\epsilon = u*\rho_\epsilon \in C^\infty(\R^n)$.
	Show that %, if $x\in\R^n$ is a Lebesgue point of $u$, then
	\begin{equation}\label{eq681727df}
		\forall x\in\R^n\text{ Lebesgue point of $u$,}\qquad
		\lim_{\epsilon\to0} u_\epsilon(x) = u(x) .
	\end{equation}
	Note that here we do not consider $u$ ``up to a set of measure zero'', but really a fixed function $u:\R^n\to\C$.
	
	{\it Solution.}
	Notice that, for every $y\in\R^n$,  $|\rho_\epsilon(y)| = \left|\frac{\rho(y/\epsilon)}{\epsilon^n}\right| \le \frac{\|\rho\|_{L^\infty}}{\epsilon^n}$.
	Recall also that $\spt(\rho_\epsilon)\subset \bar B(0,\epsilon)$.
	Therefore,
	\begin{align}
		|u_\epsilon(x) - u(x)|
		&= \left| \int_{B(x,\epsilon)} u(x-y) \rho_\epsilon(y) \did y - u(x) \int_{B(x,\epsilon)}\rho_\epsilon(y) \did y \right| \\
		&\le \int_{B(x,\epsilon)} |u(x-y) - u(x) |  \rho_\epsilon(y) \did y \\
		&\le \frac{\|\rho\|_{L^\infty}}{\epsilon^n} \int_{B(x,\epsilon)} |u(x-y) - u(x) | \did y \\
		&= \|\rho\|_{L^\infty}\omega_n \fint_{B(x,\epsilon)} |u(x-y) - u(x) | \did y .
	\end{align}
	Therefore, if $x$ is a Lebesgeue point of $u$, then 
	$\lim_\epsilon u_\epsilon(x) = u(x)$.
	
	See \cite[Theorem 8.15]{MR767633} for a more general statement (for mollifications that do not have compact support).
\end{exercise}

\begin{proof}[\underline{Proof of Theorem~\ref{thm681726ad}: Morrey's inequality for Sobolev functions}]
	Let $u\in W^{1,p}(B(o,4R))$ and $x,y\in B(o,R)$ be two Lebesgue points of $u$.

	Let $\{\rho_\epsilon\}_{\epsilon>0}$ be a family of standard mollifiers on $\R^n$
	and denote by $u_\epsilon = u*\rho_\epsilon \in C^\infty(\R^n)$.
	By Exercise~\ref{ex68172771}, by the Morrey's inequality for smooth functions~\ref{prop68165fd2}, 
	and by standard properties of mollifiers, 
	we have
	\begin{align}
		|u(x)-u(y)|
		&\overset{\eqref{eq681727df}}= \lim_{\epsilon\to0} |u_\epsilon(x) - u_\epsilon(y)| \\
		&\overset{\eqref{eq68166038}}\le \limsup_{\epsilon\to0}  C |x-y|^{1-n/p} \|\grad u_\epsilon\|_{L^p(B(o,3R))} \\
		&= C |x-y|^{1-n/p} \limsup_{\epsilon\to0}\|(\grad u)_\epsilon\|_{L^p(B(o,3R))} \\
		&\overset{\eqref{eq6817203b}}\le C |x-y|^{1-n/p} \|\grad u\|_{L^p(B(o,4R))} .
	\end{align}
	In the last step we use Young's inequality~\eqref{eq6817203b} as follows:
	if $\epsilon<R$, then there is a cut-off function $\zeta\in C^\infty_c(B(o,4R)$ such that $\bar B(o,3R+\epsilon) \subset \{\zeta=1\}$ and $0\le \zeta\le 1$.
	Then $(\zeta\grad u)_\epsilon = (\grad u)_\epsilon$ in $B(o,3R)$.
	So,
	\begin{align}
		\|(\grad u)_\epsilon\|_{L^p(B(o,3R))} 
		&= \|(\zeta\grad u)_\epsilon\|_{L^p(B(o,3R))}
		\le \|(\zeta\grad u)_\epsilon\|_{L^p(\R^n)} \\
		&= \|\rho_\epsilon*(\zeta\grad u)\|_{L^p(\R^n)} 
		\le \|\rho_\epsilon*|\zeta\grad u|\|_{L^p(\R^n)} \\
		&\overset{\eqref{eq6817203b}}\le \|\rho_\epsilon\|_{L^1(\R^n)} \|\zeta\grad u\|_{L^p(\R^n)} 
		\le \|\grad u\|_{L^p(B(o,4R))} .
	\end{align}
\end{proof}

\begin{exercise}
	Let $n\ge2$ and $\Omega\subset\R^n$ open.
	Show that, if $u\in W^{m,p}_\loc(\Omega)$ for some $m\ge1$ and $p>n$, then $u\in C^{m-1}(\Omega)$.
	In particular, show that, for every $p>n$,
	\begin{equation}
		\bigcap_{m\ge1} W^{m,p}_\loc(\Omega) = C^\infty(\Omega) .
	\end{equation}
\end{exercise}

%%%%%%%%%%%%%%%%%%%%%%%%%%%%%%%%%%%%%%%%%%%%%%%%%%%%%%%%%%%%%%%
\subsection{Difference quotients}
Let $\Omega\subset\R^n$ open.
For $\epsilon>0$, define
\begin{equation}
	\Omega_{\epsilon} = \{x\in\Omega: \dist(x,\de\Omega)>\epsilon\}
	= \{x\in\Omega: \bar B(x,\epsilon) \subset\Omega\}. 
\end{equation}

For $u:\Omega\to\C$, $j\in\{1,\dots,n\}$ and $h\neq0$, define $\Delta_j^hu:\Omega_{|h|}\to\C$ as
\begin{equation}
	\Delta_j^hu(x) = \frac{u(x+he_j) - u(x)}{h} ,
\end{equation}
where $e_j$ is the $j$-th element of the standard basis of $\R^n$.

\begin{exercise}
	Let $\Omega\subset\R^n$ open and $h\neq0$.
	Show that, if $u,v:\Omega\to\C$, then, for every $j\in\{1,\dots,n\}$ and $x\in\Omega_{|h|}$,
	\begin{equation}\label{eq681734d8}
		\Delta_j^h(uv)(x) = \Delta_j^hu(x) v(x+he_j) + u(x)\Delta_j^hv(x) . 
	\end{equation}
\end{exercise}
\begin{exercise}
	Let $\Omega\subset\R^n$ open and $h\neq0$.
	Show that, if $u,v:\Omega\to\C$ and $\spt(v)\subset\Omega_{|h|}$, then,
	for every $j\in\{1,\dots,n\}$,
	\begin{equation}\label{eq681734d2}
		\int_\Omega \Delta_j^h u(x)  v(x) \did x 
		= - \int_\Omega u(x)\Delta_j^{-h}v(x) \did x .
	\end{equation}
	Notice that, since $\spt(v)\subset\Omega_{|h|}$, 
	we have: first, if $x\in\spt(v)$, then $x+he_j\in\Omega$;
	second, if $x\in\spt(\Delta_j^{-h}v)$, then $x\in\spt(v)$ or $x-he_j\in\spt(v)$, which implies $x\in\Omega$.
\end{exercise}

\begin{exercise}
	Let $\Omega\subset\R^n$ open and $j\in\{1,\dots,n\}$.
	Show that, if $\phi\in C^1_c(\Omega)$, then
	\begin{equation}\label{eq6817354b}
		\lim_{h\to0} \|\Delta_j^h\phi - \de_j\phi\|_{L^\infty(\Omega)} = 0 .
	\end{equation}
\end{exercise}

\begin{exercise}
	Let $\Omega\subset\R^n$ open.
	Show that, if $u\in L^1_\loc(\Omega)$, then, for every $j\in\{1,\dots,n\}$ and every $\phi\in C^\infty_c(\Omega)$,
	\begin{equation}\label{eq6817480f}
		\lim_{j\to 0} \int_{\Omega} \Delta_j^h u(x) \phi(x) \did x 
		= \de_ju[\phi] ,
	\end{equation}
	where we see $u\in\scr D'(\Omega)$ as a distribution with distributional derivative $\de_ju\in\scr D'(\Omega)$.
	In other words, $\Delta_j^h u\to\de_ju$ in $\scr D'(\Omega)$.
	
	{\it Solution.}
%	First of all, we see the following easy result.
	Fix $u$, $j$ and $\phi$.
	Since $\spt(\phi)$ is compact, there is $\epsilon>0$ so that $\phi\in C^\infty_c(\Omega_\epsilon)$.
	Since $\spt(\phi)\Subset\Omega_\epsilon$, then there is $\delta>0$ such that $B(\spt(\phi),\delta) = \bigcup_{x\in \spt(\phi)} B(x,\delta) \Subset\Omega_\epsilon$.
	Therefore, 
	\begin{equation}
	\begin{aligned}
		&\hspace{-1cm}
		\limsup_{h\to 0} \left| \int_{\Omega} \Delta_j^h u(x) \phi(x) \did x - \de_ju[\phi] \right| \\
		&\overset{\eqref{eq681734d2}}= \limsup_{h\to 0} \left| - \int_{\Omega}  u(x) \Delta_j^{-h}\phi(x) \did x + \int_{\Omega}  u(x) \de_j\phi(x) \did x \right| \\
		&\le \limsup_{h\to 0}  \int_{\Omega}  |u(x)|\cdot |-\Delta_j^{-h}\phi(x)+ \de_j\phi(x)| \did x  \\
		&\overset{\eqref{Hölder}}\le \limsup_{h\to 0} 
			\|u\|_{L^1(B(\spt(\phi),\delta))} \cdot
			\|-\Delta_j^{-h}\phi(x)+ \de_j\phi(x)\|_{L^{\infty}(\Omega)} \\
%		&\le \limsup_{h\to 0} \|u\|_{L^p(\Omega)}  \|-\Delta_j^{-h}\phi(x)+ \de_j\phi(x)\|_{L^\infty(\Omega)} \|1\|_{L^{p'}(B(\spt(\phi),\epsilon))} \\
		&\overset{\eqref{eq6817354b}}=0 .
	\end{aligned}
	\end{equation}
%	We have proven that $\Delta_j^h u\to\de_ju$ in $\scr D'(\Omega)$.
\end{exercise}


\begin{proposition}\label{prop68176bc1}
	Let $\Omega\subset\R^n$ be an open set and $p\in(1,+\infty]$.
	For every $u\in L^p(\Omega)$ and $j\in\{1,\dots,n\}$,
	the following are equivalent:
	\begin{enumerate}[label=(\roman*)]
	\item\label{prop68176bc1_1}
	$\de_ju\in L^p(\Omega)$; % and $\|\de_ju\|_{L^p(\Omega)} \le C$;
	\item\label{prop68176bc1_2}
	there exists $C\in\R$ such that, for every $\epsilon>0$, 
	$\limsup_{h\to0} \|\Delta_j^h u \|_{L^p(\Omega_\epsilon)} \le C$.
	\end{enumerate}
	Moreover, for every $\epsilon>0$ and $h$ with $|h|<\epsilon$,
	\begin{equation}\label{eq68176675}
		\|\Delta_j^h u \|_{L^p(\Omega_\epsilon)} \le \|\de_ju\|_{L^p(\Omega)} .
	\end{equation}
%	For $p=\infty$, the implication $\ref{prop68176bc1_2}\THEN\ref{prop68176bc1_1}$ and the estimate~\eqref{eq68176675} hold 
%	(but not the implication $\ref{prop68176bc1_1}\THEN\ref{prop68176bc1_2}$).
\end{proposition}
\begin{proof}
	\framebox{$\ref{prop68176bc1_1}\THEN\ref{prop68176bc1_2}$: case $p\in(1,\infty)$.}
	Let $u\in C^1(\Omega)$ and $\epsilon>0$.
	For $|h|<\epsilon$ and $p\in[1,+\infty]$, we have
	\begin{equation}\label{eq68176488}
	\begin{aligned}
		\int_{\Omega_\epsilon} |\Delta_j^hu(x)|^p \did x
		&= \int_{\Omega_\epsilon} \left| \frac{u(x+he_j)-u(x)}{h} \right|^p \did x \\
		&\overset{\eqref{eq68175caa}}= \int_{\Omega_\epsilon} \left| \int_0^1 \de_ju(x+th e_j) \did t \right|^p \did x \\
		&\overset{\eqref{Hölder}}\le \int_{\Omega_\epsilon}  \int_0^1 |\de_ju(x+th e_j)|^p \did t \did x \\
		&= \int_0^1 \int_{\Omega_\epsilon}   |\de_ju(x+th e_j)|^p  \did x \did t \\
		&\le \int_0^1 \int_{\Omega_{\epsilon-|h|}}   |\de_ju(x)|^p  \did x \did t 
		= \int_{\Omega_{\epsilon-|h|}}   |\de_ju(x)|^p  \did x .
	\end{aligned}
	\end{equation}
	
	We have shown~\eqref{eq68176675} for smooth functions.
	Next, let $u\in L^p(\Omega)$ with $\de_ju\in L^p(\Omega)$.
	Let $\{\rho_\eta\}_{\eta>0}$ be a standard family of mollifiers and define $u_k = u*\rho_{1/k}:\Omega_{1/k}\to\C$.
	For $1/k<\epsilon-|h|$, we have $\Omega_\epsilon\subset \Omega_{\epsilon-|h|}\subset \Omega_{1/k}$ and thus, from~\eqref{eq68176488},
	\begin{equation}\label{eq68176a0c}
		\int_{\Omega_\epsilon} |\Delta_j^hu_k(x)|^p \did x 
		\le \int_{\Omega_{\epsilon-|h|}}   |\de_ju_k(x)|^p  \did x .
	\end{equation}
	Moreover, we know that $u_k\to u$ and $\de_ju_k\to \de_ju$ in $L^p(\Omega_{\epsilon-|h|})$.
	We thus have, for $|h|<\epsilon$,
	\begin{align}
		&\|\Delta^h_ju_k - \Delta^h_ju\|_{L^p(\Omega_\epsilon)} \\
		&= \left( \int_{\Omega_\epsilon} \left| \frac{u_k(x+he_j) - u_k(x)}{h} - \frac{u(x+he_j)-u(x)}{h} \right|^p \did x \right)^{1/p} \\
		&\overset{\eqref{Minkowski}}\le \left( \int_{\Omega_\epsilon} \left| \frac{u_k(x+he_j) - u(x+he_j)}{h} \right|^p \did x \right)^{1/p} 
		+ \left( \int_{\Omega_\epsilon} \left| \frac{u(x) - u_k(x)}{h} \right|^p \did x \right)^{1/p} \\
		&\le \frac2h \|u_k-u\|_{L^p(\Omega_{\epsilon-|h|})} . 
	\end{align}
	Therefore, for each $h\neq0$ with $|h|<\epsilon$ fixed, 
	$\Delta^h_ju_k\to \Delta^h_ju$ in $L^p(\Omega_{\epsilon-|h|})$ as $k\to\infty$.
	This convergence allows us to extend the estimate~\eqref{eq68176a0c} to the limit.
	We obtain~\eqref{eq68176675} and thus~\ref{prop68176bc1_2}.
	
	%
	\framebox{$\ref{prop68176bc1_1}\THEN\ref{prop68176bc1_2}$: case $p=\infty$.}
	Let $u\in C^1(\Omega)$ and $\epsilon>0$.
	For $|h|<\epsilon$ and $x\in\Omega_\epsilon$, we have
	\begin{equation}\label{eq683573e1}
	\begin{aligned}
		\sup_{x\in\Omega_\epsilon} |\Delta^h_ju(x)|
		&= \sup_{x\in\Omega_\epsilon} \left| \frac{u(x+he_j) - u(x)}{h} \right| \\
		&= \sup_{x\in\Omega_\epsilon} \left| \int_0^1 \frac{\grad u(x+the_j) \cdot he_j }{h} \did t \right| \\
		&\le \sup_{x\in\Omega} |\de_j u(x)| 
		= \|\de_j u\|_{L^\infty(\Omega)} .
	\end{aligned}
	\end{equation}
	We have shown~\eqref{eq68176675} for smooth functions.
	Next, let $u\in L^\infty(\Omega)$ with $\de_ju\in L^\infty(\Omega)$.
	Like in the previous case, define $u_k = u*\rho_{1/k}:\Omega_{1/k}\to\C$, so that $u_k\in C^1(\Omega_{1/k})$.
	For $1/k<\epsilon-|h|$, we have $\Omega_\epsilon\subset \Omega_{\epsilon-|h|}\subset \Omega_{1/k}$.
%	 and thus, from~\eqref{eq683573e1},
	[...]
	
	Using Exercise~\ref{ex683577c9}, we have that, for fixed $h$ and $j$, there is a full measure set $E\subset\Omega_{\epsilon-|h|}$, 
	such that, for every $x\in E$, we have
	$\lim_{k\to\infty} u_k(x) = u(x)$, and $\lim_{k\to\infty} u_k(x+he_j) = u(x+he_j)$. 
	Moreover, $\|\de_ju_k\|_{L^\infty(\Omega_{1/k})} \le \|\de_j u\|_{L^\infty(\Omega)}$.
%	and $\lim_{k\to\infty} \de_ju_k(x) = \de_ju_k(x)$.
	Therefore, for every $x\in E$,
	\begin{equation}
	\begin{aligned}
		|\Delta^h_ju(x)|
		&= \left| \frac{u(x+he_j) - u(x)}{h} \right| \\
		&= \lim_{k\to\infty} \left| \frac{u_k(x+he_j) - u_k(x)}{h} \right| \\
		&\overset{\eqref{eq683573e1}}\le \limsup_{k\to\infty} \|\de_j u_k\|_{L^\infty(\Omega_{1/k})} \\
		&\le  \|\de_j u\|_{L^\infty(\Omega)} .
	\end{aligned}
	\end{equation}
	We obtain~\eqref{eq68176675} and thus~\ref{prop68176bc1_2} for $p=\infty$ too.
	
	%
	\framebox{$\ref{prop68176bc1_2}\THEN\ref{prop68176bc1_1}$.}
	Fix $\epsilon>0$.
	If $p\in(1,+\infty)$, then $L^p(\Omega_\epsilon)$ is the dual of $L^{p'}(\Omega_\epsilon)$, where $p'=\frac{p}{p-1}$.
	If $p=\infty$, then the same is true with $p'=1$.
	
	Since $\limsup_{h\to0} \|\Delta_j^h u \|_{L^p(\Omega_\epsilon)} \le C<\infty$,
	we can apply the Banach–Alaoglu Theorem~\ref{thm68173edd}.
	Therefore, there is a sequence $\{h_k\}_{k\in\N}\subset(0,\epsilon)$ with $\lim_{k\to\infty}h_k=0$
	and there is $v_\epsilon\in L^p(\Omega_\epsilon)$ such that
	\begin{equation}\label{eq681743bc}
		\|v_\epsilon\|_{L^p(\Omega_\epsilon)} \le C 
	\end{equation}
	and, for every $\phi\in L^{p'}(\Omega)$, 
	\begin{equation}\label{eq681742e1}
		\lim_{k\to\infty} \int_{\Omega_\epsilon} \Delta_j^{h_k} u(x) \phi(x) \did x 
		= \int_{\Omega_\epsilon} v_\epsilon(x) \phi(x) \did x .
	\end{equation}
	
	In particular, $C^\infty_c(\Omega_\epsilon)\subset L^{p'}(\Omega_\epsilon)$.
	Combining~\eqref{eq6817480f} with~\eqref{eq681742e1}, we get
	$v_\epsilon=\de_ju$, that is, the distributional derivative $\de_ju$ of $u$ on $\Omega_\epsilon$ is in fact a function in $L^p(\Omega_\epsilon)$.
	
	By the locality of distirbutions, we obtain that $\de_ju$ is a function on $\Omega$
	that satisfies, by~\eqref{eq681743bc}, $\|\de_ju\|_{L^p(\Omega_\epsilon)} \le C$ for all $\epsilon>0$.
	Therefore, $\|\de_ju\|_{L^p(\Omega)} \le C$.
\end{proof}

Recall that the dual of $L^1$ is $L^\infty$, although the dual of $L^\infty$ is not $L^1$

\begin{exercise}
	Suppose $u\in C^1(\Omega)$, $x\in\Omega$, $v\in\R^n$ such that $x+tv\in\Omega$ for all $t\in[0,1]$.
	Then
	\begin{equation}\label{eq68175caa}
		u(x+v)-u(x) = \int_0^1 \grad u(x+tv)\cdot v \did t .
	\end{equation}
\end{exercise}

	
\begin{exercise}\label{ex683577c9}
	Let $\Omega\subset\R^n$,
	$\{\rho_\eta\}_{\eta>0}$ a standard family of mollifiers,
	and $u\in L^1_\loc(\Omega)$. 
	Define $u_\epsilon = u*\rho_\epsilon : \Omega_\epsilon\to\C$.
	Show that, for every $x\in\Omega_\epsilon$,
	\begin{align}
		u_\epsilon(x) &\le \| u \|_{L^\infty(B(x,\epsilon))} .
	\end{align}
	Moreover, for almost every $x\in\Omega$, $\lim_{\epsilon\to0} u_\epsilon(x) = u(x)$.
	
	Find an example where $u\in L^\infty(\Omega)$ but, for every $h>0$, 
	$\liminf_{\epsilon\to0} \|u_\epsilon - u\|_{L^\infty(\Omega_h)} > 0$.
	
	{\it Hint for the example.}
	Take $\Omega = (-1,1)\subset\R$ (or $\Omega=\R$), and $u=\one_{(0,1)}$...
\end{exercise}


\newcommand{\weakstarto}{\overset{*}{\rightharpoonup}}
\begin{theorem}[Banach–Alaoglu Theorem]\label{thm68173edd}
	Let $(V,\|\cdot\|)$ be a normed space, and let $(V',\|\cdot\|_*)$ be the dual space endowed with the operator norm
	\begin{equation}
		\|\alpha\|_* = \sup\{\alpha[x]:x\in V,\ \|x\|\le 1\} ,
		\qquad\forall\alpha\in V' .
	\end{equation}
	If $\{\alpha_k\}_{k\in\N}\subset V'$ is a bounded sequence, 
	that is, $\sup_{k\in\N} \|\alpha_k\|_* < \infty$, 
	then there exists a unique $\alpha_\infty\in V'$ which is the weak* limit of $\alpha_k$, that is $\alpha_k\weakstarto\alpha_\infty$.
	More explicitly, for every $x\in V$,
	\begin{equation}
		\lim_{k\to\infty} \alpha_k[x] = \alpha_\infty[x] .
	\end{equation}
	Moreover, 
	\begin{equation}\label{eq6817438c}
		\|\alpha_\infty\|_* \le \liminf_{k\to\infty} \|\alpha_k\|_* .
	\end{equation}
\end{theorem}
\begin{proof}
	See \url{https://en.wikipedia.org/wiki/Banach–Alaoglu_theorem} .
\end{proof}

\begin{theorem}[Characterization of Sobolev spaces with differential quotients]\label{thm681736aa}
	Let $\Omega\subset\R^n$ be an open set and $u\in L^1_\loc(\Omega)$.
	For every $p\in(1,+\infty]$, 
	the following are equivalent
	\begin{enumerate}[label=(\roman*)]
	\item\label{thm681736aa_1}
	$u\in W^{1,p}(\Omega)$;
	\item\label{thm681736aa_2}
	$u\in L^p(\Omega)$ and there exists $C>0$ such that, for every $j\in\{1,\dots,n\}$ and every $\epsilon>0$, 
	$\limsup_{h\to0} \|\Delta_j^h u \|_{L^p(\Omega_\epsilon)} \le C$.
	\end{enumerate}
\end{theorem}
\begin{proof}
	This is a direct consequence of Proposition~\ref{prop68176bc1}.
\end{proof}

%%%%%%%%%%%%%%%%%%%%%%%%%%%%%%%%%%%%%%%%%%%%%%%%%%%%%%%%%%%%%%%
\subsection{Differentiability a.e. for $n<p\le\infty$}

For $\Omega\subset\R^n$ open, a function $u:\Omega\to\C$ is \emph{differentiable} at $x\in \Omega$ if there exists a linear function $\alpha:\R^n\to\C$ such that
\begin{equation}\label{eq681770fb}
	\lim_{y\to x} \frac{ |u(y)-u(x) - \alpha[y-x]| }{ |y-x| } = 0 .
\end{equation}
If $u$ is differentiable at $x$, the linear function $\alpha$ in~\eqref{eq681770fb} is unique and is of the form $\alpha[y-x] = \grad_{cl} u(x)\cdot(y-x)$ for a vector $\grad_{cl} u(x)\in \C^n$.
We call this vector $\grad_{cl} u(x)$ the \emph{classical gradient} of $u$ at $x$.
If $u$ is $C^1$, then it is clear that $\grad_{cl} u(x)=\grad u(x)$ for every $x$, where $\grad u(x)$ is the gradient we have used so far.
When $u\in L^1_\loc(\Omega)$, then we have a distributional gradient $\grad u \in \scr D'(\Omega)^n$, which may be an element of $L^p(\Omega)$ in the case of Sobolev functions,
but we don't know a priori that $\grad u(x) = \grad_{cl} u(x)$.

In fact, it can be that $u\in W^{1,p}(\Omega)$ is not differentiable anywhere.
Indeed, recall from your course in Analysis, or simply prove it from~\eqref{eq681770fb}, 
that if $u$ is differentiable at $x$, then $u$ is continuous at $x$.
It follows that the function constructed in Exercise~\ref{ex681772db}, which is nowhere continuous, is nowhere differentiable although it has a weak gradient in $L^p$.

\begin{theorem}\label{thm68357e9c}
	Let $n\ge2$ and $p\in(n,+\infty]$.
	If $\Omega\subset\R^n$ is open and $u\in W^{1,p}_\loc(\Omega)$, 
	then, for almost every $x\in\Omega$, $u$ is differentiable at $x$ and $\grad u(x) = \grad_{cl} u(x)$.
\end{theorem}
\begin{proof}
	First, assume $p\in(n,+\infty)$.
	By Theorem~\ref{thm681726ad}, we can assume $u$ continuous.
	By Exercise~\ref{ex68177e18}, for almost every $x\in\Omega$, we have
	\begin{equation}\label{eq68177e4d}
		\lim_{r\to0} \fint_{B(x,r)} |\grad u(y) - \grad u(x)|^p \did x = 0 .
	\end{equation}
	Let $x\in\Omega$ be a point with~\eqref{eq68177e4d}.
	Let $R>0$ be such that $B(x,4R)\Subset\Omega$.
	Define $v:B(x,4R)\to\C$ by
	\begin{equation}
		v(y) = u(y) - u(x) - \grad u(x) \cdot(y-x) ,
		\qquad\forall y\in B(x,4R) .
	\end{equation}
	Clearly, we have $v\in W^{1,p}(B(x,4R))\cap C^0(B(x,4R))$ and 
	$\grad v(y) = \grad u(y) - \grad u(x)$.
	Therefore, applying Theorem~\ref{thm681726ad} to $v$,
	we obtain for every $y\in B(x,R)$,
	\begin{align}
		\frac{ |u(y) - u(x) - \grad u(x) \cdot(y-x) | }{ |x-y| }
		&= \frac{ |v(y)-v(x)| }{ |x-y| } \\
		&\overset{\eqref{eq68177bc8}}\le C |x-y|^{-\frac{n}{p}} \|\grad u(y) - \grad u(x)\|_{L^p(B(o,4R))} \\
		&= C \left( \frac{ \omega_n^{1/n} 4R }{ |x-y| } \right)^{n/p}
			\left( \fint_{B(x,4R)} |\grad u - \grad u(x)|^p \did x \right)^{\frac1p} .
	\end{align}
	If we take $R=|x-y|$, this estimate combined with~\eqref{eq68177e4d}, gives
	\begin{equation}
		\lim_{y\to x} \frac{ |u(y) - u(x) - \grad u(x) \cdot(y-x) | }{ |x-y| } = 0 ,
	\end{equation}
	that is, $\grad u(x) = \grad_{cl}u(x)$.
	
	Finally, if $p=\infty$, then we clearly have $W^{1,\infty}_\loc (\Omega)\subset W^{1,p}_\loc(\Omega)$ for all $p\in(n,\infty)$.
	So, we apply the result we have just proven.
\end{proof}

\begin{exercise}[A variant of Lebesgue differentiation theorem]\label{ex68177e18}
	Let $\Omega\subset\R^n$ open.
	Show that, if $f\in L^p(\Omega)$, then for almost every $x\in\Omega$ we have
	\begin{equation}%\label{eq68177e4d}
		\lim_{r\to0} \fint_{B(x,r)} |f(y)-f(x)|^p \did y = 0 .
	\end{equation}
	
	{\it Solution.}
	Look at \cite[\S1.5.7\&\S1.1.8]{MR0290095}.
	It can be proven for $f\in L^p(\mu)$ with $(X,d,\mu)$ a doubling metric measure space.
\end{exercise}





%%%%%%%%%%%%%%%%%%%%%%%%%%%%%%%%%%%%%%%%%%%%%%%%%%%%%%%%%%%%%%%
\subsection{$p=\infty$: Lipschitz functions}
Let $\Omega\subset\R^n$ open.
A function $u:\Omega\to\C$ is $L$-Lipschitz for some $L\in\R$ if
\begin{equation}
	\forall x,y\in\Omega,\qquad
	|u(x)-u(y)| \le L |x-y| .
\end{equation}
Lipschitz functions are clearly continuous

\begin{theorem}\label{thm68178f84}
	Let $\Omega\subset\R^n$ open and convex.
	For every $u\in L^\infty(\Omega)$ and $L\in\R$,
	the following are equivalent
	\begin{enumerate}[label=(\roman*)]
	\item\label{thm68178f84_1}
	$u$ has a $L$-Lipschitz representative, in the sense that, for almost every $x,y\in\Omega$, $|u(x)-u(y)| \le L |x-y|$;
	\item\label{thm68178f84_2}
	$u\in W^{1,\infty}(\Omega)$ and $\|\grad u\|_{L^\infty(\Omega)}\le L$.
	\end{enumerate}
	As as consequence, if $u\in W^{1,\infty}(\Omega)$, then $u$ is $\|\grad u\|_{L^\infty(\Omega)}$-Lipschitz.
\end{theorem}
\begin{proof}
	\framebox{$\ref{thm68178f84_1}\THEN\ref{thm68178f84_2}$.}
	If $u:\Omega\to\C$ is $L$-Lipschitz, then, for every $x\in\Omega$, $j\in\{1,\dots,n\}$ and $h\neq0$ with $|h|<\dist(x,\de\Omega)$,
	\begin{equation}
		|\Delta_j^hu(x)|
		= \left| \frac{u(x+he_j)-u(x)}{h} \right| 
		\le L .
	\end{equation}
	Therefore, by Theorem~\ref{thm681736aa}, we obtain $u\in W^{1,\infty}(\Omega)$.
	Moreover, using also Proposition~\ref{prop68176bc1}, we have, for almost every $x\in\Omega$,
	\begin{equation}
		|\grad u(x)|
		\le \max\{|\de_ju(x)|:j\in\{1,\dots,n\}\}
		\le L .
	\end{equation}
	So, $\|\grad u\|_{L^\infty(\Omega)}\le L$.
	
	\framebox{$\ref{thm68178f84_2}\THEN\ref{thm68178f84_1}$.}
	We can use Morrey's inequality, Theorem~\ref{thm681726ad}, to show that if $u\in W^{1,\infty}(\Omega)$ then $u$ is Lipschitz.
	We just need to notice that $u\in W^{1,p}(B)$ for every $B\Subset\Omega$ and $p\in(n,\infty)$.
	Then, we take the limit $p\to\infty$ in~\eqref{eq68177bc8}, where the constant $C=C(p,n)$ remains bounded.
	However, in this way it is not evident that $u$ is $\|\grad u\|_{L^\infty(\Omega)}$-Lipschitz.
	
	So, we can apply another argument, just by usual mollification.
	One easily sees that $u_\epsilon = u*\rho_\epsilon\in C^\infty(\Omega_\epsilon)$
	is a smooth function with $|\grad u_\epsilon(x)|\le \|\grad u\|_{L^\infty(\Omega)}$
	for every $x\in \Omega_\epsilon$.
	Since $u_\epsilon$ is smooth, we can estimate for every $x,y\in\Omega$, using convexity,
	\begin{equation}
		|u_\epsilon(x)-u_\epsilon(y)|
		= \left| \int_0^1 \grad u_\epsilon(x+t(y-x))\cdot (y-x) \did t\right|
		\le \|\grad u\|_{L^\infty(\Omega)} |x-y| .
	\end{equation}
	Since $u_\epsilon\to u$ uniformly on compact sets, then $\lim_{\epsilon\to0}|u_\epsilon(x)-u_\epsilon(y)| = |u(x)-u(y)|$.
\end{proof}

\begin{exercise}
	In Theorem~\ref{thm68178f84} we used convexity of the set $\Omega$.
	Give an example of open set $\Omega$ that is connected but not convex where Theorem~\ref{thm68178f84} fails.
	What can we say in any case?
\end{exercise}

\begin{theorem}[Rademacher Theorem]
	If $\Omega\subset\R^n$ is an open set and $u:\Omega\to\C$ is Lipschitz,
	then $u$ is differentiable almost every where in $\Omega$.
\end{theorem}
\begin{proof}
	This is a consequence of Theorem~\ref{thm68178f84} and Theorem~\ref{thm68357e9c}.
\end{proof}


%%%%%%%%%%%%%%%%%%%%%%%%%%%%%%%%%%%%%%%%%%%%%%%%%%%%%%%%%%%%%%%
\subsection{Compactness theorems: Ascoli--Arzelà}
One of our main tools to prove compactness is the following standard result
\begin{theorem}[Ascoli--Arzelà]\label{thm681f9a96}
	Let $K$ be a compact metric space and $\scr F = \{f_k\}_{k\in\N}\subset C(K)$ be a sequence of continuous functions $K\to\C$.
	Suppose that
	\begin{enumerate}
	\item
	$\scr F$ is (equi)bounded, that is, there is $C$ such that $|f(x)| \le C$ for all $f\in\scr F$ and all $x\in K$.
	\item
	$\scr F$ is equicontinuos, that is, for every $\epsilon>0$ there exists $\delta>0$ such that, if $x,y\in K$ and $d(x,y)<\delta$, then $|f(x)-f(y)|<\epsilon$ for all $f\in\scr F$.
	\end{enumerate}
	Then there exists a subsequence $\{f_{k_j}\}_{j\in\N}\subset\scr F$ that converges uniformly on $K$.
\end{theorem}

%%%%%%%%%%%%%%%%%%%%%%%%%%%%%%%%%%%%%%%%%%%%%%%%%%%%%%%%%%%%%%%
\subsection{Compactness theorems: Rellich--Kondrachov for $W^{1,p}_0(\R^n)$}

Recall that a linear operator $L:A\to B$ between Banach spaces is \emph{compact} if it maps bounded subsets of $A$ to pre-compact subsets of $B$.
In other words, $L$ is a compact operator if, for every bounded sequence $\{a_k\}_{k\in\N}\subset A$, there is a subsequence $\{La_{k_j}\}_{j\in\N}$ converging in $B$.
We shall prove that for certain $p$ and $q$, the ``identity map'' $W^{1,p}_0(\Omega) \to L^q(\Omega)$ is a compact linear operator.

\begin{remark}
	The proof is taken from \cite[\S5.5.7]{MR2597943}.
	For a similar proof, see \cite[Theorem 4.26]{MR2759829}.
	For more general statements, see \cite[Theorem 6.3]{zbMATH02208228}.
\end{remark}

\begin{lemma}\label{lem681f714e}
	Let $n\ge1$ and $\{\rho_\epsilon\}_{\epsilon>0}$ be a standard family of mollifiers.
	Let $u\in W^{1,1}_\loc(\R^n)$ and define $u_\epsilon = u*\rho_\epsilon$.
	For every $\Omega\subset\R^n$, we have
	\begin{equation}\label{eq681f7158}
		\int_\Omega |u_\epsilon(x)-u(x)| \did x
		\le \epsilon \int_{B(\Omega,\epsilon)} |\grad u|(x) \did x ,
	\end{equation}
	where, we recall, $B(\Omega,\epsilon) = \bigcup_{x\in\Omega} B(x,\epsilon)$.
\end{lemma}
\begin{proof}
	First, assume $u\in C^1(\R^n)$.
	Notice that, for every $x\in\R^n$, 
	\begin{align}
		|u_\epsilon (x) - u(x)|
		&= \left| \int_{\R^n} \rho_\epsilon(y) (u(x-y)-u(x)) \did y \right| \\
		&= \left| \int_{\R^n} \rho_\epsilon(y) \int_0^1 \grad u(x-ty)\cdot (-y) \did t \did y \right| \\
		&\le \int_{\R^n}\int_0^1 \rho_\epsilon(y) |\grad u|(x-ty) |y| \did t \did y \\
	[\spt(\rho_\epsilon) \subset B(0,\epsilon)]
		&\le \epsilon \int_{\R^n}\int_0^1 \rho_\epsilon(y) |\grad u|(x-ty) \did t \did y .
	\end{align}
	Therefore,
	\begin{align}
		\int_\Omega |u_\epsilon (x) - u(x)|\did x 
		&\le \epsilon \int_\Omega \int_{\R^n}\int_0^1 \rho_\epsilon(y) |\grad u|(x-ty) \did t \did y \did x \\
		&\le  \epsilon \int_{\R^n}\int_0^1 \rho_\epsilon(y) \left(\int_\Omega|\grad u|(x-ty)\did x\right) \did t \did y \\
		&\le \epsilon \int_{\R^n}\int_0^1 \rho_\epsilon(y) \did t \did y \left(\int_{B(\Omega,\epsilon)}|\grad u|(x-ty)\did x\right)\\
		&= \epsilon \int_{B(\Omega,\epsilon)}|\grad u|(x-ty)\did x .
	\end{align}
	
	Second, consider $u\in W^{1,1}_\loc(\R^n)$ arbitrary.
	For every $\eta\in(0,1)$, we have~\eqref{eq681f7158} for the smooth function $u_\eta$.
	On the one hand, since $(u_\eta)_\epsilon = (u*\rho_\eta)*\rho_\epsilon = (u*\rho_\epsilon)*\rho_\eta$, we have
	\begin{equation}
		\lim_{\eta\to0} \| (u_\eta)_\epsilon(x) - u_\eta \|_{L^1(\Omega)}
		= \lim_{\eta\to0} \| ( u_\epsilon(x) - u)_\eta \|_{L^1(\Omega)}
		\overset{\eqref{??}}= \| u_\epsilon(x) - u \|_{L^1(\Omega)} .
	\end{equation}
	On the other hand, 
	\begin{equation}
		\lim_{\eta\to0} \| \grad u_\eta \|_{L^1(B(\Omega,\epsilon))}
		\overset{\eqref{??}}= \| \grad u \|_{L^1(B(\Omega,\epsilon))} .
	\end{equation}
	Therefore, we obtain~\eqref{eq681f7158} for $u$ too.
\end{proof}

\begin{exercise}
	Extend \eqref{eq681f7158} of Lemma~\ref{lem681f714e} to all $u\in W^{1,1}(\R^n)$.
	What can we say for $u\in W^{1,1}_0(\Omega)$, when $\Omega\subset\R^n$ is open.
\end{exercise}

\begin{lemma}[Interpolation inequality for $L^p$ norms]\label{lem681f73e4}
	Let $1\le p\le q\le r\le \infty$ and $\theta\in[0,1]$ such that
	\begin{equation}\label{eq681f7532}
		\frac1q = \frac\theta p + \frac{1-\theta}r .
	\end{equation}
	Then, whenever $\mu$ is a measure and $u$ is $\mu$-measurable function,
	\begin{equation}\label{eq681f76d1}
		\|u\|_{L^q(\mu)} \le \|u\|_{L^p(\mu)}^{\theta} \cdot \|u\|_{L^r(\mu)}^{1-\theta} .
	\end{equation}
	In particular, if $u\in L^p(\mu)\cap L^r(\mu)$, then $u\in L^q(\mu)$ for all $q\in[p,r]$.
\end{lemma}
\begin{proof}
	The identity~\eqref{eq681f7532} is equivalent to
	\begin{equation}
		1 = \frac{q\theta}{p} + \frac{q(1-\theta)}{r}
		= \frac{1}{p/(q\theta)} + \frac{1}{r/(q(1-\theta))} ,
	\end{equation}
	that is, $\frac{p}{q\theta}$ and $\frac{r}{q(1-\theta)}$ are Hölder conjugate exponents.
	Since they are both belong to $[1,+\infty]$, we can apply the Hölder inequality:
	\begin{align}
		\int |u|^q \did \mu
		&= \int |u|^{\theta q} |u|^{(1-\theta)q} \did \mu  \\
		&\overset{\eqref{Hölder}}\le \left( \int |u|^{\theta q \frac{p}{q\theta}} \did \mu \right)^{\frac{q\theta}{p}}
			\left( \int |u|^{(1-\theta)q\frac{r}{q(1-\theta)}} \did \mu \right)^{\frac{q(1-\theta)}{r}} \\
		&= \left( \|u\|_{L^p(\mu)}^{\theta} \cdot \|u\|_{L^r(\mu)}^{1-\theta} \right)^{q} .
	\end{align}
	So, we get~\eqref{eq681f76d1} by exponentiating by $\frac1q$.
\end{proof}

\begin{lemma}\label{lem681f796a}
	Let $n\ge1$ and $\{\rho_\epsilon\}_{\epsilon>0}$ be a standard family of mollifiers.
	Let $u\in W^{1,1}_\loc(\R^n)$ and define $u_\epsilon = u*\rho_\epsilon$.
	Then, for every $p\in [1,n)$ and $q\in[1,p^*)$, there exists $C\in\R$ such that
	\begin{equation}\label{eq681f7bc5}
		\| u_\epsilon - u\|_{L^q(\R^n)}
		\le C \epsilon^\theta \|\grad u\|_{L^1(\R^n)}^\theta \cdot \|\grad u\|_{L^{p}(\R^n)}^{1-\theta} ,
	\end{equation}
	where $\theta\in [0,1]$ is such that $\frac1q = \theta + \frac{1-\theta}{p^*}$.
\end{lemma}
\begin{proof}
	First, we apply the Interpolation inequality for $L^p$ norms, Lemma~\ref{lem681f73e4}, to get
	\begin{equation}
		\| u_\epsilon - u\|_{L^q(\R^n)} 
		\le \| u_\epsilon - u\|_{L^1(\R^n)}^\theta \cdot \| u_\epsilon - u\|_{L^{p^*}(\R^n)}^{1-\theta} .
	\end{equation}
	Second, we apply the bound~\eqref{eq681f7158} from Lemma~\ref{lem681f714e} to the first term $\| u_\epsilon - u\|_{L^1(\R^n)}$,
	and the Gagliardo–Nirenberg–Sobolev Inequality~\eqref{eq6810881c} from Theorem~\ref{thm68108820} or Corollary~\ref{cor681fa13d} to the second term $\| u_\epsilon - u\|_{L^{p^*}(\R^n)}$.
	We thus get
	\begin{equation}
	\| u_\epsilon - u\|_{L^1(\R^n)}^\theta \cdot \| u_\epsilon - u\|_{L^{p^*}(\R^n)}^{1-\theta}
	\le \left( \epsilon \|\grad u\|_{L^1(\R^n)} \right)^\theta
		\cdot \left( C_{GNS} \|\grad u\|_{L^p(\R^n)} \right)^{1-\theta} .
	\end{equation}
	We have thus obtained~\eqref{eq681f7bc5}.
\end{proof}

%\begin{exercise}
%	Extend~\eqref{eq681f7bc5} from Lemma~\ref{lem681f796a} to all $u\in W^{1,p}(\R^n)$.
%\end{exercise}

\begin{remark}
	In the inequalities~\eqref{eq681f7158} and~\eqref{eq681f7bc5}, it might happen that the right-hand side is $+\infty$.
	The inequalities are still true, although they don't provide any extra information.
\end{remark}

\begin{proposition}\label{prop681f7c91}
	Let $n\ge2$, $p\in[1,n)$ and $q\in [1,p^*)$.
	Suppose that $\{u_k\}_{k\in\N}\subset W^{1,p}(\R^n)$ is a sequence 
	such that there is $R>0$ with $\spt(u_k)\subset B(0,R)$ for all $k\in\N$.
	Suppose also that there exists $M\in\R$ such that $\|u_k\|_{W^{1,p}(\R^n)}\le M$ for all $k\in\N$.
	Then, there exists a subsequence $\{u_{k_j}\}_{j\in\N}$ that is converging in $L^q(\R^n)$.
\end{proposition}
\begin{proof}
%	Since $C^1_c(\R^n)$ is dense in $W^{1,p}(\R^n)$ by~\ref{},
%	we can assume that $\{u_k\}_{k\in\N}\subset C^1_c(\R^n)$.
%	Indeed, if $\{u_k\}_{k\in\N}\subset W^{1,p}(\R^n)$, 
%	we can take a sequence $\{v_k\}_{k\in\N}\subset C^1_c(\R^n)$
%	with $\|u_k-v_k\|_{W^{1,p}(\R^n)}<1/k$ and $\spt(v_k)\subset B(0,R+1)$.
%	So, if $\{v_{k_j}\}_{j\in\N}$ is a subsequence converging to some $v$ in $L^q(\R^n)$,
%	then 
	
	Let $\{\rho_\epsilon\}_{\epsilon>0}$ be a standard family of mollifiers.
	and define $u_k^\epsilon = u_k*\rho_\epsilon$.

	We claim that there are $C\in\R$ and $\theta\in[0,1]$ such that, for every $\epsilon>0$ and $k\in\N$,
	\begin{equation}\label{eq681f91c4}
		\|u_k^\epsilon - u_k\|_{L^q(\R^n)} \le C\epsilon^\theta.
	\end{equation}
	Indeed, if we apply~\eqref{eq681f7bc5} from Lemma~\ref{lem681f796a}, we get
	\begin{align}
		\|u_k^\epsilon - u_k\|_{L^q(\R^n)}
		&\overset{\eqref{eq681f7bc5}}\le C_{\eqref{eq681f7bc5}} \epsilon^\theta 
			\|\grad u\|_{L^1(\R^n)}^\theta \cdot \|\grad u\|_{L^{p}(\R^n)}^{1-\theta} \\
		&\overset{\eqref{Hölder}}\le C_{\eqref{eq681f7bc5}} \epsilon^\theta 
			\|\grad u\|_{L^p(\R^n)}^\theta\cdot \scr L^n(B(0,R))^{\theta/p'} \cdot \|\grad u\|_{L^{p}(\R^n)}^{1-\theta} \\
		&= C_{\eqref{eq681f7bc5}} \epsilon^\theta \scr (\omega_nR^n)^{\theta/p'} \cdot \|\grad u\|_{L^{p}(\R^n)} \\
		&\le C_{\eqref{eq681f7bc5}}(\omega_nR^n)^{\theta/p'} M \epsilon^\theta .
	\end{align}
	So, we have proven~\eqref{eq681f91c4}.


	Notice that, for every $\epsilon>0$, $k\in\N$, and $x\in\R^n$,
	\begin{equation}
	\begin{aligned}
		|u_k^\epsilon(x)| 
		&\le \int_{\R^n} \rho_\epsilon(y) |u(x-y)| \did y \\
		&\overset{\eqref{Hölder}}\le \|\rho_\epsilon\|_{L^{p'}(\R^n)} \|u\|_{L^p(\R^n)} \\
		&= \frac{ \|\rho_1\|_{L^p(\R^n)} }{ \epsilon^{n/p} } \|u\|_{L^p(\R^n)} \\
		&\le \frac{ \|\rho_1\|_{L^p(\R^n)} }{ \epsilon^{n/p} } M ,
	\end{aligned}
	\end{equation}
	and, similarly,
	\begin{equation}\label{eq681f848e}
	\begin{aligned}
		|\grad u_k^\epsilon(x)|
		&\le \int_{\R^n} \rho_\epsilon(y) |\grad u(x-y)| \did y \\
		&\overset{\eqref{Hölder}}\le \|\rho_\epsilon\|_{L^{p'}(\R^n)} \|\grad u\|_{L^p(\R^n)} \\
		&\overset{\eqref{eq681f83e0}}= \frac1{\epsilon^{n\frac{p'-1}{p'}}} \|\rho_1\|_{L^p(\R^n)} \|\grad u\|_{L^p(\R^n)} \\
		&= \frac{ \|\rho_1\|_{L^p(\R^n)} }{ \epsilon^{n/p} } \|\grad u\|_{L^p(\R^n)} \\
		&\le \frac{ \|\rho_1\|_{L^p(\R^n)} }{ \epsilon^{n/p} } M .
	\end{aligned}
	\end{equation}
	It follows that, for each $\epsilon\in(0,1)$ fixed, the family $\{u_k^\epsilon\}_{k\in\N} \subset C^0_c(B(0,R+1))$ is bounded and equicontinuous.
	By Ascoli--Arzelà Theorem~\ref{thm681f9a96}, 
	$\{u_k^\epsilon\}_{k\in\N}$ is pre-compact in $C^0_c(B(0,R+1))$.
	
	We apply a diagonal argument.
	First, let $\{u_{k^1_j}^{1}\}_{j\in\N}$ be a subsequence of $\{u_k^1\}_{k\in\N}$ that is converging uniformly on $\bar B(0,R+1)$ to a function $v_1\in C^0_c(\bar B(0,R+1))$.
	Next, for every $m\in\N_{\ge2}$, there is a subsequence $\{u_{k^m_j}^{1/m}\}_{j\in\N}$ of $\{u_{k^{m-1}_j}^{1/m}\}_{j\in\N}$ that is converging uniformly on $\bar B(0,R+1)$.
	to some $v_m\in C^0_c(\bar B(0,R+1))$.
	Notice that, for every $m\in\N_{\ge1}$,
	\begin{equation}
		\limsup_{j\to\infty} \|u_{k^{m-1}_j}^{1/m} - v_m\|_{L^p(\R^n)}
		\le \limsup_{j\to\infty} \|u_{k^{m-1}_j}^{1/m} - v_m\|_{L^\infty(\bar B(0,R+1))} \scr L^n(\Omega)^{1/p}
		= 0 .
	\end{equation}
	Therefore, $\{u_{k^m_j}^{1/m}\}_{j\in\N}$ is converging to $v_m$ also in $L^q(\R^n)$.
	
	We claim that $\{u_{k^m_m}\}_{m\in\N}$ is a Cauchy sequence in $L^q(\R^n)$.
	Indeed, let $\delta>0$.
	Then there is $L\in\N$ such that 
	\begin{equation}\label{eq681f97da}
		C(1/L)^\theta < \delta .
	\end{equation}
	Next, since $\{u_{k^L_j}^{1/L}\}_{j\in\N}$ is a Cauchy sequence in $L^q(\R^n)$,
	then there is $N\in\N$ such that, for every naturals $a,b>N$,
	\begin{equation}\label{eq681f9606}
		\forall a,b>N,\qquad
		\|u_{k^L_a}^{1/L} - u_{k^L_b}^{1/L}\|_{L^q(\R^n)} < \delta.
	\end{equation}
	We also assume $N>L$.
	Since we have taken always subsequences, if $m>L$ then 
	$\{k_a^m\}_{a>N} \subset \{k_a^L\}_{a>N}$.
	Therefore,~\eqref{eq681f9606} implies
	\begin{equation}\label{eq681f96fc}
		\forall a,b>N,\qquad
		\|u_{k^a_a}^{1/L} - u_{k^b_b}^{1/L}\|_{L^q(\R^n)} < \delta.
	\end{equation}
%	$k_a^m > k_a^L$ for all $a$.
	It follows that, for every $a,b>N$,
	\begin{align}
		\|u_{k^a_a} - u_{k^b_b}\|_{L^q(\R^n)}
		&\le \|u_{k^a_a} - u_{k^a_a}^{1/L} \|_{L^q(\R^n)}
			+ \| u_{k^a_a}^{1/L} - u_{k^b_b}^{1/L} \|_{L^q(\R^n)}
			+ \| u_{k^b_b}^{1/L} - u_{k^b_b} \|_{L^q(\R^n)} \\
		&\overset{\eqref{eq681f91c4}\&\eqref{eq681f96fc}}\le 2C(1/L)^\theta + \delta \\
		&\overset{\eqref{eq681f97da}}\le 3\delta .
	\end{align}
	We have proven our claim, that $\{u_{k^m_m}\}_{m\in\N}$ is a Cauchy sequence in $L^q(\R^n)$.
\end{proof}

\begin{exercise}
	Let $\{\rho_\epsilon\}_{\epsilon>0}$ be a standard family of mollifiers on $\R^n$.
	Show that, for every $p\in[1,+\infty]$,
	\begin{equation}\label{eq681f83e0}
		\|\rho_\epsilon\|_{L^p(\R^n)} = \frac1{\epsilon^{n\frac{p-1}{p}}} \|\rho_1\|_{L^p(\R^n)} .
	\end{equation}
	
	{\it Solution.}
	\begin{align}
		\int_{\R^n} \rho_\epsilon(y)^p \did y
		&= \int_{\R^n} \left( \frac{\rho_1(y/\epsilon)}{\epsilon^n} \right)^p \did y \\
		[z=y/\epsilon,\ \did z = \did y/\epsilon^n]
		&= \frac{1}{\epsilon^{n(p-1)}} \int_{\R^n} \rho_1(z)^p \did z .
	\end{align}
\end{exercise}

\begin{exercise}
	In Proposition~\ref{prop681f7c91}, can we take the subsequence $\{u_{k_j}\}_{j\in\N}$ independent of $q$?
\end{exercise}

\begin{exercise}
	Proposition~\ref{prop681f7c91} uses a general fact in metric spaces.
	Let $(X,d)$ be a complete metric space and $\{x^m_k\}_{k\in\N,m\in\N\cup\{\infty\}}\subset X$.
	Suppose that:
	\begin{enumerate}
	\item
	for every $k\in\N$, $\{x^m_k\}_{m\in\N}$ converges to $x^\infty_k$, uniformly in $k$;
	explicitly, for every $\epsilon>0$ there exists $N\in\N$ such that, for every $m>N$ and for every $k\in\N$, $d(x^m_k,x^\infty_k) < \epsilon$;
	\item
	For every $m\in\N$ (but not for $m=\infty$), the set $\{x^m_k\}_{k\in\N}$ is pre-compact in $X$.
	\end{enumerate}
	Show that $\{x^\infty_k\}_{k\in\N}$ is pre-compact in $X$.
\end{exercise}


%%%%%%%%%%%%%%%%%%%%%%%%%%%%%%%%%%%%%%%%%%%%%%%%%%%%%%%%%%%%%%%
\subsection{Extension domains}
An open set $\Omega\subset\R^n$ is called an~\emph{extension domain} if, for every $p\in[1,+\infty]$, there exists $T_p:W^{1,p}(\Omega) \to W^{1,p}(\R^n)$ such that
\begin{enumerate}
\item	$T_pu|_{\Omega} = u$, for all $u\in W^{1,p}(\Omega)$;
\item	$T_p$ is bounded, that is, there is $C_p$ such that $\|T_pu\|_{W^{1,1}(\R^n)} \le C_p \|u\|_{W^{1,p}(\Omega)}$, for all $u\in W^{1,p}(\Omega)$.
\end{enumerate}

An example of a set that is NOT and extension domain, is the so called \emph{slit disk}:
\begin{equation}
	\Omega = \{(x,y)\in \R^2: x^2+y^2<1\} \setminus \{(x,0):x\ge0\} .
\end{equation}

However, there holds the following result
\begin{theorem}\label{thm681fa38a}
	If $\Omega\subset\R^n$ is an open set with $C^1$ boundary, then it is an extension domain.
\end{theorem}
In fact, Theorem~\ref{thm681fa38a} can be pushed to Lipschitz boundary.

\begin{proposition}\label{prop681fa9b9}
	Suppose $\Omega\subset\R^n$ is an extension domain and $\Omega'\subset\R^n$ is an open set such that $B(\Omega,r)\subset \Omega'$, for some $r>0$.
	Then, for every $p\in[1,+\infty]$ there exists a continuous extension operator $T_p:W^{1,p}(\Omega)\to W^{1,p}_0(\Omega')$.
\end{proposition}
\begin{proof}
	Let $\zeta\in C^\infty_c(\Omega')$ such that $0\le\zeta\le 1$ and 
	\begin{equation}
		\bar\Omega \subset\interior\{\zeta=1\} \subset\spt(\zeta)\subset\Omega' .
	\end{equation}
	Since $B(\Omega,r)\subset \Omega'$, we can take $\zeta$ with $\|\grad\zeta\|_{L^\infty(\R^n)} \le 2/r$.
	Let $\tilde T_p::W^{1,p}(\Omega) \to W^{1,p}(\R^n)$ be a bounded extension operator given by $\Omega$ being and extension domain.
	
	We claim that $T_p:u\mapsto \zeta \tilde T_pu$ is the wanted operator.
	To prove our claim, we need to show that $u\mapsto \zeta u$ defines a continuous operator $W^{1,p}(\R^n)\to W^{1,p}_0(\Omega')$.
	Clearly we have
	\begin{equation}
		\|\zeta u\|_{L^p(\Omega')} \le \|u\|_{L^p(\R^n)} .
	\end{equation}
	Moreover,
	\begin{align}
		\|\grad(\zeta u)\|_{L^p(\Omega')}
		&\le \|\zeta\grad u\|_{L^p(\R^n)} + \|u\grad\zeta\|_{L^p(\R^n)} \\
		&\le \|\grad u\|_{L^p(\R^n)} + \frac2r \|u\|_{L^p(\R^n)} \\
		&\le (1+2/r) \|u\|_{W^{1,p}(\R^n)} .
	\end{align}
	The claim is proven.
\end{proof}

\begin{proposition}\label{prop681fa951}
	Let $\Omega\subset\Omega'\subset\R^n$ and $p\in[1,+\infty)$.
	Define $T_p:W^{1,p}_0(\Omega)\to L^p(\R^n)$ as
	$T_pu = \one_\Omega u$, i.e., $T_p$ extends functions to zero outside $\Omega$.
	Then $T_p$ is a continuous operator $W^{1,p}_0(\Omega)\to W^{1,p}(\Omega')$.
	In fact,
	\begin{equation}\label{eq681fa8e1}
		\forall u\in W^{1,p}(\Omega)\qquad
		\|\one_\Omega u\|_{W^{1,p}(\Omega')} = \|u\|_{W^{1,p}(\Omega)}
	\end{equation}
\end{proposition}
\begin{proof}
	If $u\in C^1_c(\Omega)$, then $\|T_pu\|_{W^{1,p}(\Omega')} = \|u\|_{W^{1,p}(\Omega)}$.
	Since $C^1_c(\Omega)$ is dense in $W^{1,p}_0(\Omega)$, then we have~\eqref{eq681fa8e1}.
\end{proof}

\begin{remark}
	By Propositions~\ref{prop681fa9b9} and~\ref{prop681fa951}, if $\Omega$ is a bounded extension domain, we can assume that the extension operator takes values in the space of functions with compact support in some fixed neighborhood of $\Omega$.
\end{remark}

%%%%%%%%%%%%%%%%%%%%%%%%%%%%%%%%%%%%%%%%%%%%%%%%%%%%%%%%%%%%%%%
\subsection{Extension of Sobolev inequalities to Extension domains}
Many properties of $W^{1,p}(\R^n)$ extend to extension domains.

%%%%%%%%%%%%%%%%%%%%%%%%%%%%%%%%%%%%%%%%%%%%%%%%%%%%%%%%%%%%%%%
\subsection{Poincaré inequality for $W^{1,p}_0(\Omega)$}\label{subs68208317}

\begin{theorem}\label{thm68207b9a}
	Let $n\ge2$ and $\Omega\subset\R^n$ a bounded open set.
	For every $p\in[1,+\infty)$ there is $C\in\R$ such that, 
	\begin{equation}\label{eq68207b94}
		\forall u\in W^{1,p}_0(\Omega),\qquad
		\|u\|_{L^p(\Omega)} \le C\|\grad u\|_{L^p(\Omega)} .
	\end{equation}
\end{theorem}
\begin{remark}
	The inequality in~\eqref{eq68207b94} can also be written as
	\begin{equation}
		\forall u\in W^{1,p}_0(\Omega),\qquad
		\int_\Omega |u|^p \did x \le C^p \int_\Omega |\grad u|^p\did x .
	\end{equation}
\end{remark}

\begin{lemma}\label{lem68207e48}
	Let $1\le a<b<\infty$.
	If $\Omega\subset\R^n$ is a measurable set, then
	\begin{equation}\label{eq68207e42}
		\forall u\in L^1_\loc(\Omega),\qquad
		\|u\|_{L^a(\Omega)} \le |\Omega|^{\frac{b-a}{ba}} \|u\|_{L^b(\Omega)} ,
	\end{equation}
	where $|\Omega|=\scr L^n(\Omega)$ is the volume of $\Omega$.
\end{lemma}
\begin{proof}
	Since $b/a>1$, we can apply the Hölder inequality:
	\begin{align}
		\int_\Omega |u|^a \did x
		\overset{\eqref{Hölder}}\le \left( \int_\Omega |u|^{a\frac{b}{a}} \did x \right)^{\frac{a}{b}} 
			\cdot \left( \int_\Omega 1^{\frac{\frac{b}{a}}{\frac{b}{a}-1}} \right)^{\frac{\frac{b}{a}-1}{\frac{b}{a}}}
		= \left( \int_\Omega |u|^{b} \did x \right)^{\frac{a}{b}} 
			\cdot |\Omega|^{\frac{b-a}{b}} .
	\end{align}
\end{proof}

\begin{proof}[Proof of Theorem~\ref{thm68207b9a}]
	We shall prove~\eqref{eq68207b94} assuming $u\in C^1_c(\Omega)$, to obtain the general statement by approximation in $W^{1,p}_0(\Omega)$.

	Fix $p\in[1,+\infty)$.
	Recall that, if $q\in[1,n)$, then $q^* = \frac{nq}{n-q}$.
%	
%	First, assume $p<n$ and let $u\in W^{1,p}_0(\Omega)$.
%	Then there is a sequence $\{u_j\}_{j\in\N}\subset C^\infty_c(\Omega)$ such that $u_j\to u$ in $W^{1,p}(\Omega)$.
%	Since
%	\begin{align}
%		\|u\|_{L^p(\Omega)}
%		&= \lim_{j\to\infty} \|u_j\|_{L^p(\Omega)}
%		&\overset{\eqref{eq6810881c}}\le \liminf_{j\to\infty} 
%	\end{align}	
	Since $\lim_{q\to n^-} q^* = +\infty$, there is some $q\in[1,n)$ such that $q^*>p$.
	Let $u\in C^1_c(\Omega)$
	Therefore, using Theorem~\ref{thm68108820},
%	Corollary~\ref{cor681fa13d} for $W^{1,p}_0(\Omega)$ (\TODO)
	\begin{align}
		\|u\|_{L^p(\Omega)}
		&\overset{\eqref{eq68207e42}}\le |\Omega|^{\frac{q^*-p}{q^*p}} \|u\|_{L^{q^*}(\Omega)} \\
		&\overset{\eqref{eq6810881c}}\le |\Omega|^{\frac{q^*-p}{q^*p}} \| \grad u\|_{L^q(\Omega)} .
	\end{align}
	If $p<n$, then we can take $q=p$ already (because $p^*>p$), and~\eqref{eq68207b94} is proven with $C=|\Omega|^{\frac{p^*-p}{p^*p}} = |\Omega|^{\frac1n}$.
	If $p\ge n$, then, for each appropriate $q\in(1,n)$ we have $q<p$ and thus
	\begin{equation}
		\| \grad u\|_{L^q(\Omega)} 
		\overset{\eqref{eq68207e42}}\le |\Omega|^{\frac{p-q}{pq}} \| \grad u\|_{L^q(\Omega)} .
	\end{equation}
	This implies \eqref{eq68207b94} with $C= |\Omega|^{\frac{q^*-p}{q^*p}+\frac{p-q}{pq}}$.
\end{proof}

%%%%%%%%%%%%%%%%%%%%%%%%%%%%%%%%%%%%%%%%%%%%%%%%%%%%%%%%%%%%%%%
\subsection{Extra that might be added in the future}
\begin{itemize}
\item Poincaré inequality for extension domains
\end{itemize}


%%%%%%%%%%%%%%%%%%%%%%%%%%%%%%%%%%%%%%%%%%%%%%%%%%%%%%%%%%%%%%%
%%%%%%%%%%%%%%%%%%%%%%%%%%%%%%%%%%%%%%%%%%%%%%%%%%%%%%%%%%%%%%%
\newpage
\part{Application of Sobolev spaces theory to PDE}
\section{Elliptic PDEs}

%%%%%%%%%%%%%%%%%%%%%%%%%%%%%%%%%%%%%%%%%%%%%%%%%%%%%%%%%%%%%%%
\subsection{Setting}\label{subs6817a5f7}
Let $n\ge 2$ and $\Omega\subset\R^n$ open.
We are interested in functions $u:\Omega\to\C$ such that
\begin{equation}\label{eq68179d48}
	Lu(x) = \div(A(x)\grad u(x)) + b(x)\cdot \grad u(x) + c(x) u(x) 
	\overset{!}{=} f(x) ,
\end{equation}
where $A(x)\in \C^{n\times n}$ is a $n\times n$-complex matrix, $b(x)\in \C^n$ and $c(x),f(x)\in\C$, for each $x\in\Omega$.

How do we interpret the formula~\eqref{eq68179d48}?
We may use distributional calculus: in such a case, we need to make sense of the products
$A\grad u = \sum_{jk} A_{jk}\de_j u$,
$b\cdot\grad u =  \sum_j b_j\de_j u$,
and $cu$.
If the coefficients $A$, $b$ and $c$ are smooth, then we can consider $Lu=f$ for $u,f\in\scr D'(\Omega)$;
see~\ref{subs67df2598}.
Another possibility is that we consider $u\in W^{1,1}_\loc(\Omega)$:
in such a case, all derivatives of $u$ are $L^1_\loc$ functions and thus, if the coefficients $A$, $b$ and $c$ are bounded in $L^\infty$, then the products are still of class $L^1_\loc$ and thus $\div(A\grad u) = \sum_{jk} \de_j(A_{jk} \de_ku) \in\scr D'$.


\begin{definition}\label{def6817a806}
	The \emph{standard conditions on $L$} are:
	\begin{enumerate}
	\item 
	$A\in L^\infty(\Omega;\C^{n\times n})$, such that $A(x)$ is symmetric for every $x\in\Omega$ and 
	there are $0<\lambda\le \Lambda<\infty$ with 
		\begin{equation}\label{eq68179d12}
			\forall x\in\Omega,\ \forall\xi\in\R^n
			\qquad
			\lambda |\xi|^2 \le 
			\langle \xi,A(x)\xi \rangle = \sum_{i,j=1}^n A_{ij}(x) \xi_i\xi_j
			\le \Lambda |\xi|^2 .
		\end{equation}
	\item
	$b\in L^\infty(\Omega;\C^n)$, $c\in L^\infty(\Omega;\C)$.
	\end{enumerate}
\end{definition}

Condition~\eqref{eq68179d12} is called~\emph{ellipticity}.
So, $L$ as in~\eqref{eq68179d48} is \emph{elliptic} if $A$ satisfy~\eqref{eq68179d12}.

Notice that, under these conditions, the formula for $Lu(x)$ written in~\eqref{eq68179d48} is not well founded.
Indeed, even taking distributional derivatives, if $A$ is only in $L^\infty$, then the product $A(x)\cdot\grad u(x)$ is not a well defined distribution.
However, we have written $L$ in \emph{divergence form} so that we can say the following:
If $u\in W^{1,1}_\loc(\Omega)$, then we define
\begin{equation}\label{eq6817a091}
\begin{gathered}
	Lu = f \text{ in }\Omega , \\
	\IFF\quad \\ 
	\forall\phi\in C^\infty_c(\Omega) 
	\quad\int_\Omega \bigg(\langle A\grad u,\grad\phi \rangle + b\cdot\grad u \phi + cu \phi \bigg) \did x = \int_\Omega f \phi \did x , \\
	\IFF\quad \\ 
	\forall\phi\in C^\infty_c(\Omega) 
	\quad\int_\Omega \bigg(\sum_{jk} A_{jk}\,\de_j u\,\de_k\phi + (\sum_j b_j\,\de_j u  + c\,u) \,\phi \bigg) \did x = \int_\Omega f \phi \did x .
\end{gathered}
\end{equation}

\begin{exercise}
	Let $p\in [1,\infty)$ and set $p' = \frac{p}{p-1}$ the Hölder conjugate of $p$.
	Show that, if $u\in W^{1,p}(\Omega)$ satisfies $Lu=f$ as in~\eqref{eq6817a091} 
	with $f\in L^p(\Omega)$, then 
	\begin{equation}
		\forall\phi\in W^{1,p'}_\loc(\Omega) 
		\quad\int_\Omega \bigg(\langle A\grad u,\grad\phi \rangle + b\cdot\grad u \phi + cu \phi \bigg) \did x = \int_\Omega f \phi \did x .
	\end{equation}
\end{exercise}

\begin{exercise}[??]
	If $u\in W^{1,1}_\loc(\Omega)$ and $A,b,c$ are bounded and $A$ elliptic, and $Lu=f\in L^1_\loc(\Omega)$, then $\div(A\grad u) \in L^1_\loc(\Omega)$.
	Show that $u\in W^{2,1}_\loc(\Omega)$.
\end{exercise}

%%%%%%%%%%%%%%%%%%%%%%%%%%%%%%%%%%%%%%%%%%%%%%%%%%%%%%%%%%%%%%%
\subsection{The Sobolev space $H^m$}\label{subs6817a613}
Let $n\in\N$ and $\Omega\subset\R^n$ open.
If $m\in\N$, the Sobolev space $H^m(\Omega)$ is nothing else than the Sobolev space with integral exponent 2, that is,
\begin{equation}
	H^m(\Omega) = W^{m,2}(\Omega),
	\quad
	H^m_\loc(\Omega) = W^{m,2}_\loc(\Omega),
	\quad
	H^m_0(\Omega) = W^{m,2}_0(\Omega).
\end{equation}
These spaces are in fact Hilbert spaces, when endowed with a correct norm.
We will focus on $m=1$ (and $m=0$, which is $L^2$).

For $u,v\in H^1(\Omega)$, we define
\begin{equation}\label{eq682075bd}
\begin{aligned}
	\langle u,v \rangle_{H^1(\Omega)}
		&:= \llangle u,v \rrangle \\
		&:= \int_\Omega \left( \langle \grad u ,\grad\bar v \rangle + u\cdot\bar v\right) \did x  \\
		&= \langle \grad u,\grad v \rangle_{L^2(\Omega)} + \langle u,v \rangle_{L^2(\Omega)} .
\end{aligned}
\end{equation}

\begin{proposition}\label{prop6820787a}
	Let $n\ge1$ and $\Omega\subset\R^n$ open.
	Then the bilinear map $\llangle \cdot,\cdot \rrangle$ defined in~\eqref{eq682075bd}
	makes $H^1(\Omega)$ into a Hilbert space,
	with norm
	\begin{equation}
		\|u\|_{H^1(\Omega)} = \sqrt{ \langle u,u \rangle_{H^1(\Omega)} } ,
		\qquad\forall u\in H^1(\Omega) ,
	\end{equation}
	which is biLipschitz equivalent to the Sobolev norm $\|\cdot\|_{W^{1,2}(\Omega)}$.
\end{proposition}
\begin{proof}
	For every $x\in\Omega$, and every $u,v\in H^1(\Omega)$, we have 
	$|\langle \grad u(x) ,\grad\bar v(x) \rangle| 
	\le |\grad u(x)| \cdot |\grad\bar v(x)|$.
	Notice that, if $u\in H^1(\Omega)$, then
	\begin{align}
		\|u\|_{H^1(\Omega)}^2 
		&= \langle u,u \rangle_{H^1(\Omega)} \\
		&= \langle \grad u,\grad u \rangle_{L^2(\Omega)} + \langle u,u \rangle_{L^2(\Omega)} \\
		&= \|\grad u\|_{L^2(\Omega)}^2 + \|u\|_L^2(\Omega)^2 ,
	\end{align}
	that is,
	\begin{equation}
		\|u\|_{H^1(\Omega)} = \sqrt{\|\grad u\|_{L^2(\Omega)}^2 + \|u\|_L^2(\Omega)^2}
	\end{equation}
	Since there are constants $c,C\in(0,+\infty)$ such that, for every $a,b\in\R^2$,
	$c \sqrt{a^2+b^2} \le |a| + |b| \le C \sqrt{a^2+b^2}$,
	then 
	\begin{equation}\label{eq682077d7}
		c\|u\|_{H^1(\Omega)}
		\le \|u\|_{W^{1,2}(\Omega)} = \|\grad u\|_{L^2(\Omega)} + \|u\|_{L^2(\Omega)} 
		\le C \|u\|_{H^1(\Omega)} .
	\end{equation}
	See also Exercise~\ref{eq680b28e2}.
\end{proof}

\begin{remark}\label{rem6820847b}
	Proposition~\ref{prop6820787a} has already a quite interesting consequence for PDE.
	Let $f\in L^2(\Omega)$ and define $T_f:H^1(\Omega)\to\C$, $T_fu = \langle f,u \rangle_{L^2(\Omega)}$.
	The operator $T_f$ is in fact a bounded operator $L^2(\Omega)\to\C$, therefore it is bounded also on $H^1(\Omega)$;
	Explicitly, we have
	$|T_fu| \le \|f\|_{L^2(\Omega)} \|u\|_{L^2(\Omega)} \le \|f\|_{L^2(\Omega)} \|u\|_{H^1(\Omega)}$.
	Since $H^1(\Omega)$ is a Hilbert space, the Riesz Representation Theorem implies that there exists a unique $u\in H^1(\Omega)$ such that, for all $\phi\in H^1(\Omega)$, 
	$T_f\phi = \langle u,\phi \rangle_{H^1(\Omega)}$.
	In particular, 
	\begin{equation}\label{eq68207a40}
		\forall \phi\in C^\infty_c(\Omega),\qquad
		\int_\Omega f \bar\phi \did x 
		= \int_\Omega ( (\grad u\cdot \grad\bar\phi) + u\bar\phi ) \did x .
	\end{equation}
	The property~\eqref{eq68207a40} has a distributional interpretation: 
	since $\int_\Omega (\grad u\cdot \grad\bar\phi) \did x = \sum_{j=1}^n \de_ju[\de_j\bar\phi] = - \laplacian u[\phi]$,
	then~\eqref{eq68207a40} is equivalent in $\scr D'(\Omega)$ to
	\begin{equation}\label{eq68207af5}
		-\laplacian u + u = f .
	\end{equation}
	We have thus proven that, for every $f\in L^2(\Omega)$,
	there exists a unique $u\in H^1(\Omega)$ that is a distributional solution to~\eqref{eq68207af5}.
\end{remark}

\begin{exercise}
	Show that, if $\Omega\subset\R^n$ is an open and bounded set (or with finite volume),
	then, for every $\lambda\in(-\infty,0]$, there exists a unique solution to
	\begin{equation}
	\begin{cases}
		-\laplacian u = \lambda u ,\\
		u\in H^1_0(\Omega) .
	\end{cases}
	\end{equation}
	Since such solution must be $u=0$, you have shown that the half-line $(-\infty,0]$ is not in the spectrum of $-\laplacian$.
\end{exercise}

%%%%%%%%%%%%%%%%%%%%%%%%%%%%%%%%%%%%%%%%%%%%%%%%%%%%%%%%%%%%%%%
\subsection{An alternative scalar product on $H^1_0(\Omega)$}\label{subs6821ad68}
If $\Omega$ is a bounded open subset of $\R^n$,\footnote{Connected? only finite measure? \TODO}
the Poincaré inequality from Theorem~\ref{thm68207b9a}, implies that
\begin{equation}\label{eq6821ad6f}
	\llangle u,v \rrangle := \langle \grad u,\grad v \rangle_{L^2(\Omega)},
	\qquad\forall u,v\in H^1_0(\Omega),
\end{equation}
is a Hilbert scalar product on $H^1_0(\Omega)$.
Indeed, 
\begin{equation}
	\frac1{2C} \|u\|_{L^2(\Omega)} + \frac12 \|\grad u\|_{L^2(\Omega)}
	\overset{\eqref{eq68207b94}}\le \|\grad u\|_{L^2(\Omega)}
	= \sqrt{ \llangle u,u \rrangle } \le \|u\|_{H^1(\Omega)} ,
\end{equation}
which shows that the quasi-norm $u\mapsto \|\grad u\|_{L^2(\Omega)}$ is a norm bi-Lipschitz equivalent to $u\mapsto\|u\|_{H^1(\Omega)}$.

As we did in Remark~\ref{rem6820847b}, we can deduce an existence and uniqueness result for a PDE.
Indeed, if $f\in L^2(\Omega)$, then the operator $T_f\phi := \langle f,\phi \rangle_{L^2(\Omega)}$ is bounded on $H^1_0(\Omega)$.
Therefore, by the Riesz Representation Theorem, there exists a unique $u\in H^1_0(\Omega)$ such that
\begin{equation}
	\forall \phi\in C^\infty_c(\Omega),\qquad
	\llangle u,\phi \rrangle = T_f\phi .
\end{equation}
Distributionally, this reads as
\begin{equation}
	\laplacian u = f \text{ in }\Omega.
\end{equation}

We have proven the following theorem:
\begin{theorem}
	Let $\Omega\subset\R^n$ be a bounded open set.
	For every $f\in L^2(\Omega)$, the boundary problem
	\begin{equation}
		\begin{cases}
			\laplacian u = f \text{ in }\Omega, \\
			u= 0\text{ on }\de\Omega\text{, i.e., }u\in H^1_0(\Omega),
		\end{cases}
	\end{equation}
	has a unique solution $u\in H^1_0(\Omega)$.
\end{theorem}

%%%%%%%%%%%%%%%%%%%%%%%%%%%%%%%%%%%%%%%%%%%%%%%%%%%%%%%%%%%%%%%
\subsection{The dual of $H^1$}\label{subs68219d8e}
The dual space of $H^m$ is, by definition, the space
\begin{equation}\label{eq68219ddc}
	H^{-m}(\Omega) := (H^m(\Omega))',
	\qquad\text{ and }\qquad
	H^{-m}_0(\Omega) := (H^m_0(\Omega))'.
\end{equation}

Let us focus on $m=1$.
Being a Hilbert space, the dual of $H^m$ is canonically isomorphic to $H^m$ itself.
However, it is useful to keep the two spaces distinct.
The main reason, in my view, is that we want to see $L^2(\Omega)$ as a subspace of the dual of $H^1(\Omega)$.
Indeed, if $f\in L^2(\Omega)$, then, as we have seen above, $u\mapsto \int_\Omega uf \did x$ is an element of $H^m(\Omega)'$ (or $H^m_0(\Omega)'$).
However, it is clear that $L^2(\Omega)\not\subset H^1(\Omega)$: in fact, $H^1(\Omega)\into L^2(\Omega)$.

For more discussions, see\\ \url{https://math.stackexchange.com/questions/314113/dual-space-of-h1}.

%%%%%%%%%%%%%%%%%%%%%%%%%%%%%%%%%%%%%%%%%%%%%%%%%%%%%%%%%%%%%%%
\subsection{Functional Analysis: Lax--Milgram Theorem}\label{subs68208c7c}
See also \cite[Aufgabe V.6.18]{MR1787146}.

If $X$ is a Banach space, we denote by $X'$ its topological dual Banach space,
and, for $\xi\in X'$ and $x\in X$, we write the pairing $\xi[x]$ as ${}_{X'}\langle \xi | x \rangle_X$.

The following theorem will replace the role of Riesz Theorem in the previous discussion.
It is here written for Banach spaces:
notice that there are Banach spaces that are isomorphic to their duals without being ``hilbertable'';
see
\url{https://math.stackexchange.com/questions/65609/isometric-to-dual-implies-hilbertable}.

\begin{theorem}[Lax--Milgram Theorem]\label{thm68208c88}
	Let $(X,\|\cdot\|_X)$ be a Banach space and $B:X\times X\to\C$ a bilinear map.
	\begin{enumerate}
	\item	
	If $B$ is bounded (i.e., continuous), that is, there is $\beta\in\R$ such that, 
	\begin{equation}\label{eq68208d6c}
		\forall u,v\in X,\qquad
		|B[u,v]| \le \beta \|u\|_X \cdot \|v\|_X ,
	\end{equation}
	then there is a bounded linear operator $T:X\to X'$ with $\|T\|\le\beta$ and such that 
	\begin{equation}\label{eq68208e5c}
		\forall u,v\in X,\qquad
		B[u,v] = {}_{X'}{\langle} Tu|\bar v \rangle_X  .
	\end{equation}
	\item	
	If $B$ is bounded and coercive, that is, there exists $\delta>0$ 
	\begin{equation}\label{eq68208d7d}
		\forall u\in H,\qquad
		\real(B[u,u])\ge\delta\|u\|^2 ,
	\end{equation}
	then the linear operator $T:X\to X'$ is invertible and $\|T^{-1}\|\le \frac1\delta$.
	\end{enumerate}
%	In particular, for every $f\in H'$ there exists a unique $u\in H$ such that $f= B[u,\cdot]$.
\end{theorem}
\begin{proof}
	Suppose $B$ is bounded.
	If $u\in X$, then $v\mapsto B[u,v]$ is a continuous (thanks to~\eqref{eq68208d6c}) linear functional $X\to\C$ and thus there exists a unique $Tu:=w\in X'$ such that $B[u,v] = {}_{X'}{\langle} Tu|\bar v \rangle_X$ for all $v\in X$.
	It is easy to see that the so defined function $T:X\to X'$ is linear.
	We will next prove several properties of this operator $T$.
	
	We claim that $T$ is bounded, and
	\begin{equation}\label{eq68219abb}
		\|T\|_{X\to X'} \le \beta .
	\end{equation}
	Indeed, if $u\in X$, then
	\begin{align}
		\|Tu\|_{X'}
		&= \sup\{ |{}_{X'}{\langle} Tu|v \rangle_X| : v\in X,\ \|v\|_X\le 1\} \\
		&= \sup\{ |B[u,\bar v]| : v\in X,\ \|v\|_X\le 1\} \\
		&\overset{\eqref{eq68208d6c}}\le \sup\{\beta \|u\|_X \cdot \|v\|_X : v\in X,\ \|v\|_X\le 1\} \\
		&= \beta \|u\|_X .
	\end{align}
	Therefore, we have proven the claim.
	
	We claim that $T$ is coercive, i.e., 
	\begin{equation}\label{eq682090e8}
		\forall u\in X,\qquad
		\|Tu\|_{X'} \ge \delta \|x\|_X .
	\end{equation}
	Indeed, if $u\in X$, then 
	\begin{equation}
		\delta \|x\|^2 
		\overset{\eqref{eq68208d7d}}\le \real(B[x,\bar x])
		\le |B[x,\bar x]|
		= |{}_{X'}{\langle} Tu|v \rangle_X|
		\le \|Tx\|_{X'} \|x\|_X .
	\end{equation}
	Hence,~\eqref{eq682090e8} follows.

	
	We claim that $\im(T)=T[X]$, the image of $T$, is closed in $X'$.
	Indeed, if $\alpha\in X'$ and if $\{u_k\}_{k\in\N}\subset X$ is a sequence such that $\lim_{k\to\infty} Tu_k = \alpha\in X'$, then, using coercivity of $T$, we have
	$\|u_j-u_k\|_X 
	\overset{\eqref{eq682090e8}}\le \frac1\delta \|T[u_j-u_k]\|_{X'}$,
	and thus $\{u_k\}_{k\in\N}$ is a Cauchy sequence in $X$
	If $\lim_{k\to\infty} u_k=u$, then, by the boundedness of $T$, we have $\alpha=Tu\in \im(T)$.
	
	We claim that $\im(T)=X'$.
	Since $\im(T)$ is closed, we only need to show that, if $v\in X$ annihilates to $\im(T)$, then $v=0$ 
	(thanks to Hahn--Banach Theorem).
	If $v\in H$ annihilates to $\im(T)$,
	then $0 = {}_{X'}{\langle} Tu|v \rangle_X$ for every $u\in X$.
	In particular, combining this with coercivity,
	we obtain $0 = {}_{X'}{\langle} Tv|v \rangle_X \overset{\eqref{eq682090e8}}\ge \delta\|v\|_X^2$.
	Therefore, $v=0$.
	
	Finally, since $T$ is a continuous, injective and surjective linear operator, its inverse is also continuous.
\end{proof}

%%%%%%%%%%%%%%%%%%%%%%%%%%%%%%%%%%%%%%%%%%%%%%%%%%%%%%%%%%%%%%%%
% FOR HILBERT SPACES: THE ONE ABOVE IS BETTER.
%\begin{theorem}[Lax--Milgram Theorem]\label{thm68208c88}
%	Let $H$ be a Hilbert space with scalar product $\langle \cdot,\cdot \rangle$
%	and norm $\|\cdot\| = \sqrt{\langle \cdot,\cdot \rangle}$.
%	Suppose that $B:H\times H\to\C$ is a bilinear map with:
%	\begin{enumerate}
%	\item	$B$ is bounded (i.e., continuous), that is, there is $\beta\in\R$ such that, 
%		\begin{equation}\label{eq68208d6c}
%			\forall u,v\in H,\qquad
%			|B[u,v]| \le \beta \|u\|\cdot \|v\| ;
%		\end{equation}
%	\item	$B$ is coercive, that is, there exists $\delta>0$ 
%		\begin{equation}\label{eq68208d7d}
%			\forall u\in H,\qquad
%			\real(B[u,u])\ge\delta\|u\|^2 .
%		\end{equation}
%	\end{enumerate}
%	Then there exists a linear operator $T:H\to H$ such that
%	\begin{equation}\label{eq68208e5c}
%		\forall u,v\in H,\qquad
%		\langle Tu,v \rangle = B[u,v] ,
%	\end{equation}
%	and such operator is invertible.
%	In particular, for every $f\in H'$ there exists a unique $u\in H$ such that $f= B[u,\cdot]$.
%\end{theorem}
%\begin{proof}
%	The existence of a linear operator $T:H\to H$ satisfying~\eqref{eq68208e5c} is a consequence of Riesz Representation Theorem for Hilbert spaces.
%	Indeed, if $u\in H$, then $v\mapsto B[u,v]$ is a continuous (thanks to~\eqref{eq68208d6c}) linear operator $H\to\C$ and thus there exists a unique $Tu:=w\in H$ such that $B[u,v] = \langle w,v \rangle$ for all $v\in H$.
%	We will next prove several properties of this operator $T$.
%	
%	We claim that $T$ is bounded.
%	Indeed, if $x\in H$, then
%	\begin{equation}
%		\|Tx\|^2 = \langle  Tx,Tx \rangle
%		\overset{\eqref{eq68208e5c}}= B[x,Tx]
%		\overset{\eqref{eq68208d6c}}\le \beta \|x\| \|Tx\| .
%	\end{equation}
%	Hence, $\|Tx\|\le \beta\|x\|$.
%	
%	We claim that $T$ is coercive, i.e., 
%	\begin{equation}\label{eq682090e8}
%		\forall x\in H,\qquad
%		\|Tx\| \ge \delta \|x\| .
%	\end{equation}
%	Indeed, if $x\in H$, then 
%	\begin{equation}
%		\delta \|x\|^2 
%		\overset{\eqref{eq68208d7d}}\le \real(B[x,x])
%		\le |B[x,x]|
%		= |\langle Tx,x \rangle|
%		\le \|Tx\| \|x\| .
%	\end{equation}
%	Hence,~\eqref{eq682090e8} follows.
%
%	
%	We claim that $T[H]$, the image of $T$, is closed in $H$.
%	Indeed, if $y\in H$ and if $\{x_k\}_{k\in\N}\subset H$ is a sequence such that $\lim_{k\to\infty} Tx_k = y\in H$, then, using coercivity of $T$, we have
%	$\|x_j-x_k\| 
%	\overset{\eqref{eq682090e8}}\le \frac1\delta \|T[x_j-x_k]\|$,
%	and thus $\{x_k\}_{k\in\N}$ is a Cauchy sequence.
%	If $\lim_{k\to\infty} x_k=x$, then, by the boundedness of $T$, we have $y=Tx\in T[H]$.
%	
%	We claim that $T[H]=H$.
%	Since $T[H]$ is closed, we only need to show that, if $v\in H$ is orthogonal to $T[H]$, then $v=0$.
%	If $v\in H$ is orthogonal to $T[H]$,
%	then $0 = \langle Tx,v \rangle$ for every $x\in H$.
%	In particular, combining this with coercivity,
%	we obtain $0 = \langle Tv,v \rangle \overset{\eqref{eq682090e8}}\ge \delta\|v\|^2$.
%	Therefore, $v=0$.
%	
%	Finally, since $T$ is a continuous, injective and surjective linear operator, its inverse is also continuous.
%\end{proof}


%%%%%%%%%%%%%%%%%%%%%%%%%%%%%%%%%%%%%%%%%%%%%%%%%%%%%%%%%%%%%%%
\subsection{First Existence and Uniqueness result}\label{subs682095cb}

We denote by $\R^{n\times n}$ the space of all $n\times n$ matrices.
As such, if $A\in\R^{n\times n}$, then $|A|_\infty := \sup\{|Ax|:x\in\R^n,\ |x|\le 1\}$.
So, if $A\in L^\infty(\Omega;\R^{n\times n})$, then 
\begin{equation}
	\|A\|_{L^\infty(\Omega)} = \sup\{|A(x)|_\infty: x\in\Omega\} .
\end{equation}

The scalar product $\langle \cdot,\cdot \rangle$ is in fact sequilinear and defined as
\begin{equation}
	\forall x,y\in \C^n,\qquad
	\langle x,y \rangle = x\cdot\bar y = \sum_{j=1}^n x_j\bar y_j .
\end{equation}

\begin{theorem}\label{thm682098e8}
	Let $n\ge2$ and $\Omega\subset\R^n$ open and bounded.
	Let $L$ be a second order linear differential operator in divergence form, that is,
	\begin{equation}\label{eq6820a731}
		Lu = -\div(A\grad u) + b\cdot \grad u + c u ,
	\end{equation}
	where 
	\begin{equation}\label{eq6820a660}
		\text{$A\in L^\infty(\Omega;\R^{n\times n})$,
	$b\in L^\infty(\Omega;\C^n)$,
	and $c\in L^\infty(\Omega;\C)$.}
	\end{equation}
	Then $L$ is a bounded linear operator $L:H^1_0(\Omega)\to H^{-1}_0(\Omega)$.
	
	Suppose that there are $0<\theta\le \Theta<\infty$ such that
	\begin{equation}\label{eq6820970d}
		\forall x\in\Omega,\ \forall\xi\in\C^n,\qquad
		\theta |\xi|^2
		\le \langle A(x)\xi,\xi \rangle 
		\le \Theta |\xi|^2 .
	\end{equation}
	Then there is $\gamma\ge 0$ (in fact, the one for which holds~\eqref{eq6820a84f}) 
	such that for all $\lambda\in\C$ with $\real(\lambda)\ge\gamma$,
	the operator $L+\lambda :H^1_0(\Omega)\to H^{-1}_0(\Omega)$ is bounded and invertible.
	In particular, for every $f\in H^{-1}_0(\Omega)$ there exists a unique weak solution to the boundary problem
	\begin{equation}\label{eq6821b354}
		\begin{cases}
		Lu + \lambda u = f &\text{ in }\Omega, \\
		u\in H^1_0(\Omega) .
		\end{cases}
	\end{equation}
	
	Moreover,
	denoting by $\iota:H^1_0(\Omega)\into L^2(\Omega)$ the standard embedding,
	the operator $K_\lambda = \iota \circ (L-\lambda)^{-1}|_{L^2(\Omega)} : L^2(\Omega) \to L^2(\Omega)$
	is bounded, linear, and compact.
%	
%	that is (by definition of weak solution),
%	\begin{equation}
%		\forall\phi\in H^1_0(\Omega),\qquad
%		\int_\Omega \left( \langle A\grad u,\grad\phi \rangle + b\cdot\grad u \phi + (c+\mu) u \phi \right) \did x
%		= \int_\Omega f \phi \did x .
%	\end{equation}
\end{theorem}

For proving Theorem~\ref{thm682098e8}, we will study the bilinear form $\scr E_L:H^1_0(\Omega)\times H^1_0(\Omega)\to\C$, 
\begin{equation}\label{eq6820a77c}
	\scr E_L[u,v] = \int_\Omega \left( \langle A\grad u,\grad v \rangle + ( \grad u\cdot b + c u) \bar v \right) \did x .
\end{equation}
This bilinear form $\scr E_L$ is called the \emph{Dirichlet form of $L$}.

%%%%%%%%%%%%%%%%%%%%%%%%%%%%%%%%%%%%%%%%%%%%%%%%%%%%%%%%%%%%%%%
\begin{lemma}[Energy Estimates 1]\label{lem6821a526}
	Let $n\ge2$ and $\Omega\subset\R^n$ open.
	Let $L$ be a second order linear differential operator in divergence form as in~\eqref{eq6820a731}.
	Assume that $L$ has bounded coefficients, that is,~\eqref{eq6820a660}.
%	and that $L$ is elliptic, that is,~\eqref{eq6820970d}.
	Define $\scr E_L:H^1_0(\Omega)\times H^1_0(\Omega)\to\C$ as in~\eqref{eq6820a77c}.
	
	Then, there is $\alpha\ge0$ such that, for all $u,v\in H^1_0(\Omega)$,
	\begin{equation}\label{eq6820a83b}
		\forall u,v\in H^1_0(\Omega),\qquad
		|\scr E_L[u,v]| \le \alpha \|u\|_{H^1(\Omega)} \|v\|_{H^1(\Omega)} .
	\end{equation}
	In particular, the linear operator $L:H^1_0(\Omega)\to\scr D'$ is a continuous linear operator $L:H^1_0(\Omega)\to H^{-1}_0(\Omega)$ with
	\begin{equation}\label{eq6821fcc1}
		\forall u,v\in H^1_0(\Omega),\qquad
		\scr E_L[u,v] = {}_{H^{-1}_0(\Omega)} \langle Lu | \bar v \rangle_{H^1_0(\Omega)} .
	\end{equation}
\end{lemma}
\begin{proof}
	\begin{align}
		|\scr E_L[u,v]|
		&\le \int_\Omega \left| \langle A\grad u,\grad v \rangle + \langle \grad u,b \rangle \bar v + c u \bar v \right| \did x \\
		&\le \int_\Omega  |A\grad u| |\grad v| 
			+ |b| |\grad u| |v| + |c u \bar v| \did x \\
		&\le \|A\|_{L^\infty(\Omega)} \|\grad u\|_{L^2(\Omega)} \|\grad v\|_{L^2(\Omega)} 
			+ \|b\|_{L^\infty(\Omega)} \|\grad u\|_{L^2(\Omega)} \| v\|_{L^2(\Omega)} \\
		&\qquad\qquad
			+ \|c\|_{L^\infty(\Omega)} \|u\|_{L^2(\Omega)} \| v\|_{L^2(\Omega)} \\
		&\le (\|A\|_{L^\infty(\Omega)} + \|b\|_{L^\infty(\Omega)} + \|c\|_{L^\infty(\Omega)} ) \|u\|_{H^1(\Omega)} \|v\|_{H^1(\Omega)} .
	\end{align}
	So, we have~\eqref{eq6820a83b} with $\alpha=(\|A\|_{L^\infty(\Omega)} + \|b\|_{L^\infty(\Omega)} + \|c\|_{L^\infty(\Omega)} )$.
	
	By the Lax--Milgram Theorem~\ref{thm68208c88}, there is a continuous linear operator $T:H^1_0(\Omega)\to H^{-1}_0(\Omega)$ that satisfies the role of $L$ in~\eqref{eq6821fcc1}.
	We claim that $T=L$.
	Indeed, if $\phi\in\scr D(\Omega)$, then, for every $u\in H^1_0(\Omega)$ we have
	\begin{align}
		Lu[\phi] &= {}_{\scr D'(\Omega)}\langle Lu|\phi \rangle_{\scr D(\Omega)} \\
		&= (-\div(A\grad u) + b\cdot \grad u + c u )[\phi] \\
		&= - \sum_{j=1}^n \de_j(A\grad u)_j[\phi] + \int_\Omega b\cdot\grad u \phi + c u\phi \did x \\
		&= \sum_{j=1}^n (A\grad u)_j[\de_j\phi] + \int_\Omega b\cdot\grad u \phi + c u\phi \did x \\
		&= \int_\Omega ( \sum_{j=1}^n (A\grad u)_j[\de_j\phi] +  b\cdot\grad u \phi + c u\phi ) \did x \\
		&= \int_\Omega ( (A\grad u)\cdot\grad\phi +  b\cdot\grad u \phi + c u\phi ) \did x \\
		&= \langle A\grad u , \grad\bar\phi \rangle_{L^2(\Omega)} 
			+ \langle  (b\cdot\grad u + c u) , \bar \phi ) \rangle_{L^2(\Omega)} \\
		&= \scr E_L[u,\bar \phi]\\
		&= {}_{H^1_0(\Omega)'}\langle Tu,\phi \rangle_{H^1_0(\Omega)} .
	\end{align}
	Since $\scr D(\Omega)$ is dense in $H^1_0(\Omega)$, we obtain that $Lu = Tu$ not just on $\scr D(\Omega)$, but also on $H^1_0(\Omega)$.
\end{proof}

\begin{lemma}[Energy Estimates 2]\label{lem6821a53c}
	Let $n\ge2$ and $\Omega\subset\R^n$ open.
	Let $L$ be a second order linear differential operator in divergence form as in~\eqref{eq6820a731}.
	Assume that $L$ has bounded coefficients, that is,~\eqref{eq6820a660},
	and that $L$ is elliptic, that is,~\eqref{eq6820970d}.
	Define $\scr E_L:H^1_0(\Omega)\times H^1_0(\Omega)\to\C$ as in~\eqref{eq6820a77c}.
	
	Then, there are $\beta>0$ and $\gamma\ge 0$ such that, for all $u,v\in H^1_0(\Omega)$,
	\begin{equation}\label{eq6820a84f}
		\beta \|\grad u\|_{L^2(\Omega)}^2 \le \real(\scr E_L[u,u]) + \gamma \|u\|_{L^2(\Omega)}^2 .
	\end{equation}
\end{lemma}
\begin{proof}
	We have
	\begin{align}
		\theta \int_\Omega |\grad u|^2 \did x
		&\overset{\eqref{eq6820970d}}\le \int_\Omega \langle A\grad u,\grad u \rangle \did x \\
		&= \left| \scr E_L[u,u] - \int_\Omega \left( \langle \grad u,b \rangle \bar u + c |u|^2 \right) \did x \right| \\
		&\le |\scr E_L[u,u]| 
			+ \|b\|_{L^\infty(\Omega)} \|\grad u\|_{L^2(\Omega)} \|u\|_{L^2(\Omega)} 
			+ \|c\|_{L^\infty(\Omega)} \|u\|_{L^2(\Omega)}^2 \\
		&\overset{\eqref{Cauchy}}\le |\scr E_L[u,u]| 
			+ \|b\|_{L^\infty(\Omega)} \left( \epsilon \frac{\|\grad u\|_{L^2(\Omega)}^2}{2} + \frac1\epsilon \frac{\|u\|_{L^2(\Omega)}^2}{2}  \right)
			+ \|c\|_{L^\infty(\Omega)} \|u\|_{L^2(\Omega)}^2 \\
	&[\text{with }\epsilon=\frac{\theta}{\|b\|_{L^\infty(\Omega)}}]\\
		&= |\scr E_L[u,u]| 
			+ \frac{\theta}{2} \|\grad u\|_{L^2(\Omega)}^2
			+ \frac{\|b\|_{L^\infty(\Omega)}^2}{2\theta} \|u\|_{L^2(\Omega)}^2
			+ \|c\|_{L^\infty(\Omega)} \|u\|_{L^2(\Omega)}^2 \\
		&= |\scr E_L[u,u]| 
			+ \frac{\theta}{2} \|\grad u\|_{L^2(\Omega)}^2
			+ \left( \frac{\|b\|_{L^\infty(\Omega)}^2}{2\theta}
			+ \|c\|_{L^\infty(\Omega)} \right)\|u\|_{L^2(\Omega)}^2 .
	\end{align}
	So, we have~\eqref{eq6820a84f} with $\beta = \frac{\theta}{2}$ and $\gamma = \frac{\|b\|_{L^\infty(\Omega)}^2}{2\theta} + \|c\|_{L^\infty(\Omega)}$.
\end{proof}


\begin{remark}
	Notice that, if $c$ is real-valued and $c\ge c_0$ for some $c_0\in\R$, then one can take $\gamma= \frac{\|b\|_{L^\infty(\Omega)}^2}{2\theta} + c_0$ in~\eqref{eq6820a84f}. 
\end{remark}

\begin{lemma}[Cauchy inequality with $\epsilon$]
	For every $a,b\in\R$, and for every $\epsilon>0$,
	\begin{equation}\label{Cauchy}
		ab \le \epsilon \frac{a^2}{2} + \frac1\epsilon \frac{b^2}{2} .
	\end{equation}
\end{lemma}
\begin{proof}
	\begin{align}
		0
		&\le \left( \sqrt{\epsilon} a - \frac1{\sqrt{\epsilon}} b\right)^2 \\
		&= \epsilon a^2 + \frac1\epsilon b^2 - 2 ab .
	\end{align}
\end{proof}


\begin{proof}[Proof of Theorem~\ref{thm682098e8}]
	The conditions~\eqref{eq6820a660} on the coefficients of $L$ easily imply that $L$ is a continuous linear operator $H^1_0(\Omega)\to\scr D'(\Omega)$.
	Define $\scr E_L:H^1_0(\Omega)\times H^1_0(\Omega)\to\C$ as in~\eqref{eq6820a77c}.
	By Lemma~\ref{lem6821a526}, using again~\eqref{eq6820a660}, the bilinear map $\scr E_L$ is bounded.
	By Theorem~\ref{thm68208c88}, there is a continuous operator $T:H^1_0(\Omega)\to H^{-1}_0(\Omega)$ such that~\eqref{eq68208e5c} holds.
	However, see that~\eqref{eq68208e5c} is equivalent to say that $T=L$.
	Indeed, if $u\in H^1_0(\Omega)$ and $\phi\in\scr D(\Omega)\subset H^1_0(\Omega)$,
	\begin{equation}
		Lu[\phi]
		= \scr E_L[u,\phi]
		= {}_{H^{-1}_0(\Omega)}\langle Tu | \phi \rangle_{H^1_0(\Omega)} ,
	\end{equation}
	hence $Lu=Tu$ as elements of $\scr D'(\Omega)$.
	Therefore, $L$ is a bounded linear operator $L:H^1_0(\Omega)\to H^{-1}_0(\Omega)$.
	
	Next, we also assume the ellipticity condition~\eqref{eq6820970d}.
	Let $\alpha,\beta$ and $\gamma$ as in Lemmata~\ref{lem6821a526} and~\ref{lem6821a53c}, and let $\lambda\in\C$ with
	\begin{equation}\label{eq6820c2ce}
		\real(\lambda) \ge \gamma .
	\end{equation}
	Define $B_\lambda:H^1_0(\Omega)\times H^1_0(\Omega)\to\C$ by
	\begin{equation}
		\forall u,v\in H^1_0(\Omega),\qquad
		B_\lambda[u,v] = \scr E_L[u,v] + \lambda \langle u,v \rangle_{L^2(\Omega)} .
	\end{equation}
	On $H^1_0(\Omega)$, we consider the Hilbert scalar product $\llangle \cdot,\cdot \rrangle$
	defined in~\eqref{eq6821ad6f}, with norm $\|\cdot\|_{H^1_0(\Omega)} = \sqrt{\llangle \cdot,\cdot \rrangle}$.
	
	From Lemmata~\ref{lem6821a526} and~\ref{lem6821a53c}, we obtain that $B_\lambda$ is a bounded and coercive bilinear map.
	Indeed,
	on the one hand, for every $u,v\in H^1_0(\Omega)$,
	\begin{equation}
		|B_\lambda[u,v]| \overset{\eqref{eq6820a83b}}\le \alpha \|u\|_{H^1(\Omega)} \|v\|_{H^1(\Omega)} 
		\le \alpha C \|u\|_{H^1_0(\Omega)} \|v\|_{H^1_0(\Omega)} ,
	\end{equation}
	where $C$ is a constant of equivalence between the two norms on $H^1_0(\Omega)$.
	On the other hand, for every $u\in H^1_0(\Omega)$,
	\begin{align}
		\real(B_\lambda[u,u])
		&=  \real(\scr E_L[u,u]) + \real(\lambda) \|u\|_{L^2(\Omega)}^2 \\
		&\overset{\eqref{eq6820a84f}}\ge \beta \|\grad u\|_{L^2(\Omega)}^2 + (\real(\lambda)- \gamma) \|u\|_{L^2(\Omega)}^2 \\
		&\overset{\eqref{eq6820c2ce}}\ge \beta \|\grad u\|_{L^2(\Omega)}^2 
		= \beta \|u\|_{H^1_0(\Omega)} .
	\end{align}
	We apply Lax--Milgram Theorem~\ref{thm68208c88}:
	there exists a continuous, invertible linear operator $T_\lambda:H^1_0(\Omega)\to H^{-1}_0(\Omega)$ such that
	\begin{equation}
		\forall u,v\in H^1_0(\Omega)\qquad
		B_\lambda[u,v] = {}_{H^{-1}_0(\Omega)}\langle T_\lambda u | v \rangle_{H^1_0(\Omega)} .
	\end{equation}
	We claim that $T_\lambda = T+\lambda$.
	Here we mean that $T_\lambda u = Tu + \lambda \iota(u)$, where $\iota: H^1_0(\Omega)\into L^2(\Omega)$ is the standard embedding.
	Indeed, for all $u,v\in H^1_0(\Omega)$,
	\begin{align}
		{}_{H^{-1}_0(\Omega)}\langle Tu + \lambda u | v \rangle_{H^1_0(\Omega)}
		&= {}_{H^{-1}_0(\Omega)}\langle Tu  | v \rangle_{H^1_0(\Omega)}
			+ \lambda {}_{H^{-1}_0(\Omega)}\langle \iota(u)| v \rangle_{H^1_0(\Omega)} \\
		&= \scr E_L[u,v] + \lambda \langle u,v \rangle_{L^2(\Omega)} 
		= B_\lambda[u,v] \\
		&= {}_{H^{-1}_0(\Omega)}\langle T_\lambda u | v \rangle_{H^1_0(\Omega)} .
	\end{align}
	
	It follows that, for every $f\in H^{-1}_0(\Omega)$, the preimage $u=T_\lambda^{-1}[f]$ is the unique solution to~\eqref{eq6821b354}.
	
	Finally, the operator $K$ is the composition of a continuous linear operator $(T-\lambda)^{-1}$ with a continuous linear compact operator $\iota:H^1_0(\Omega)\to L^2(\Omega)$.
	The compactness of $\iota$ comes from Rellich–Kondrachov Theorem~\ref{Rellich–Kondrachov}.
%	
%	
%	If $f\in L^2(\Omega)$, then the functional $F:H^1_0(\Omega)\to\C$, $Fu = \int_\Omega fu \did x$ is continuous;
%	Hence, by Riesz Theorem, there is $u_f\in H^1_0(\Omega)$ such that $Fv = \langle u_f,v \rangle_{H^1_0(\Omega)}$ for all $v\in H^1_0(\Omega)$.
%	Then, $\hat u = T_\lambda^{-1}u_f$ satisfies, for all $v\in H^1_0(\Omega)$,
%	\begin{equation}
%		B[\hat u,v] = \langle T_\lambda \hat u,v \rangle_{H^1_0(\Omega)}
%		= \langle u_f,v \rangle_{H^1_0(\Omega)}
%		= \int_\Omega fv \did x .
%	\end{equation}
\end{proof}

%%%%%%%%%%%%%%%%%%%%%%%%%%%%%%%%%%%%%%%%%%%%%%%%%%%%%%%%%%%%%%%
\subsection{If you want kaos}\label{subs6820b196}
Here I write something that may confuse the reader quite a lot.
Read it at your own risk.

The pairing $\langle u,v \rangle_{L^2(\Omega)} = \int_\Omega u\bar v \did x$ 
is continuous on $H^1(\Omega)$, but it is NOT the Hilbert scalar product of $H^1$.
For instance, the closure of $H^1(\Omega)$ with respect to the norm $\|\cdot\|_{L^2(\Omega)} = \sqrt{\langle u,v \rangle_{L^2(\Omega)}}$, is $L^2(\Omega)$.
Continuity means that, if $f\in L^2(\Omega)$, then $u\mapsto \langle u,f \rangle_{L^2(\Omega)}$ is an element of the dual of $H^1_0(\Omega)$.

Riesz Theorem implies that there is some $v_f\in H^1(\Omega)$ such that
\begin{equation}
	\forall u\in H^1(\Omega),\qquad
	\langle u,f \rangle_{L^2(\Omega)} 
	\overset{!}= \langle \grad u, \grad v_f \rangle_{L^2(\Omega)} + \langle u, v_f \rangle_{L^2(\Omega)} 
	\overset{\text{def}}= \langle u,v_f \rangle_{H^1(\Omega)} .
\end{equation}
We have seen that, if $\Omega$ is bounded, then $\langle u,v \rangle_{H^1_0(\Omega)} = \int_\Omega \langle \grad u,\grad v \rangle \did x$ is a Hilbert scalar product on $H^1_0(\Omega)$.
Again, Riesz Theorem implies that there exists $w_f\in H^1_0(\Omega)$ such that
\begin{equation}
	\forall u\in H^1_0(\Omega),\qquad
	\langle u,f \rangle_{L^2(\Omega)} 
	\overset{!}= \langle \grad u, \grad w_f \rangle_{L^2(\Omega)} 
	\overset{\text{def}}= \langle u,w_f \rangle_{H^1_0(\Omega)} .
\end{equation}

%%%%%%%%%%%%%%%%%%%%%%%%%%%%%%%%%%%%%%%%%%%%%%%%%%%%%%%%%%%%%%%
\subsection{Functional Analysis: Fredholm Alternative Theorem}\label{subs682092ad}

Recall that, if $X$ and $Y$ are Banach spaces (e.g., Hilbert spaces),
a linear operator $K:X\to Y$ is \emph{compact operator}
if $K$ maps bounded sets to compact sets.

See also \url{https://terrytao.wordpress.com/2011/04/10/a-proof-of-the-fredholm-alternative/}
\begin{theorem}[Fredholm Alternative]\label{thm6821bc90}
	Let $H$ be a Hilbert space with scalar product $\langle \cdot,\cdot \rangle$
	and norm $\|\cdot\| = \sqrt{\langle \cdot,\cdot \rangle}$.
	Let $K:H\to H$ be a compact linear operator.
	For every $\lambda\in\C$,
	\begin{enumerate}
	\item	
	$\ker(K-\lambda)$ is finite dimensional;
	\item
	$\im(K-\lambda)$ is closed;
	\item
	$\im(K-\lambda) = \ker(K^*-\bar \lambda)^\perp$;
	\item
	$\dim(\ker(K-\lambda)) = \dim(\ker(K^*-\bar \lambda))$;
	\item
	$\ker(K-\lambda)=\{0\}$ if and only if $\im(K-\lambda) = H$.
	\end{enumerate}
\end{theorem}

%%%%%%%%%%%%%%%%%%%%%%%%%%%%%%%%%%%%%%%%%%%%%%%%%%%%%%%%%%%%%%%
\subsection{Second Existence and Uniqueness result}\label{subs6821b4db}

\begin{theorem}\label{thm6822fa7f}
	Let $n\ge 2$ and $\Omega\subset\R^n$ open and bounded.
	Let $L$ be an operator as in Theorem~\ref{thm682098e8} and define $B:H^1_0(\Omega)\times H^1_0(\Omega)\to\C$ 
	as in~\eqref{}.
	Define $\cal N,\cal M\subset H^1_0(\Omega)$ as
	\begin{equation}
	\begin{aligned}
		\cal N &= \{ u \in H^1_0(\Omega) : \scr E_L[u,v] = 0\ \forall v\in H^1_0(\Omega) \} , \\
		\cal M &= \{ v \in H^1_0(\Omega) : \scr E_L[u,v] = 0\ \forall u\in H^1_0(\Omega) \} .
	\end{aligned}
	\end{equation}
	
	Then:
	\begin{enumerate}
	\item
	Both $\cal N$ and $\cal M$ have finite dimension and $\dim(\cal N) = \dim(\cal M)$.
	\item
	$\cal N$ is the linear space of solutions to
	\begin{equation}\label{eq6821ba8f}
	\begin{cases}
		L u = 0 , \\
		u\in H^1_0(\Omega).
	\end{cases}
	\end{equation}
	\item
	for every $f\in L^2(\Omega)$, the boundary value PDE
	\begin{equation}\label{eq6821ba97}
	\begin{cases}
		L u = f , \\
		u\in H^1_0(\Omega) 
	\end{cases}
	\end{equation}
	has a solution if and only if 
	\begin{equation}\label{eq6821baaf}
		\langle f,v \rangle_{L^2(\Omega)} = 0
		\qquad\forall v\in \cal M .
	\end{equation}
	\item
	The affine space of solutions to~\eqref{eq6821ba97}, if not empty, has the same dimension of $\cal N$.
	In fact, given a solution $u_1$ to~\eqref{eq6821ba97}, then $u_1+\cal N$ is the space of solutions to~\eqref{eq6821ba97}.
	\end{enumerate}
\end{theorem}
\begin{proof}
	Let $\gamma$ as in Theorem~\ref{thm682098e8}, i.e., as in Lemma~\ref{lem6821a53c}.
	Define $K = T_\gamma^{-1} = (L+\gamma)^{-1}$, which is a compact bounded linear operator $L^2(\Omega)\to L^2(\Omega)$, as shown in Theorem~\ref{thm682098e8}.
	
	Notice that, for every $u\in H^1_0(\Omega)$ and $f\in L^2(\Omega)$
	\begin{equation}\label{eq6821bd98}
	\begin{aligned}
		u\text{ solves~\eqref{eq6821ba97}}
		&\IFF \scr E_L[u,v] = \langle f,v \rangle_{L^2(\Omega)}\qquad\forall v\in H^1_0(\Omega) \\
		&\IFF B_\gamma[u,v] = \langle f+\gamma u ,v \rangle_{L^2(\Omega)}\qquad\forall v\in H^1_0(\Omega) \\
		&\IFF T_\gamma u = f+ \gamma u \\
		&\IFF u = T_\gamma^{-1}(f+\gamma u) = Kf + \gamma Ku \\
		&\IFF u-\gamma K u = Kf .
	\end{aligned}
	\end{equation}
	So, considering the case $f=0$, we obtain that $u$ solves~\eqref{eq6821ba8f}, if and only if 
	$u\in\ker(\Id-\gamma K)$,
	that is,
	\begin{equation}
		\cal N = \ker(\Id-\gamma K) .
	\end{equation}
	Since $\gamma K$ is a compact operator $L^2(\Omega)\to L^2(\Omega)$,
	the Fredholm Alternative Theorem~\ref{thm6821bc90} implies that $\cal N = \ker(\Id-\gamma K)$ is finite dimensional.
	
%	Moreover, Theorem~\ref{thm6821bc90} tells us also that~\eqref{eq6821ba97} has a solution if and only if 
%	$Kf\in \im(\Id-\gamma K)$
%	
	Moreover, for each $f\in L^2(\Omega)$,
	\begin{align}
		\text{~\eqref{eq6821ba97} has a solution}
		&\overset{\eqref{eq6821bd98}}\IFF Kf \in \im(\Id-\gamma K) \\
		&\overset{\text{Thm~\ref{thm6821bc90}}}\IFF Kf \perp \ker(\Id-\gamma K^*) \\
		&\overset{(*)}\IFF f \perp \ker(\Id-\gamma K^*) ,
	\end{align}
	where the equivalence $(*)$ is justified as follows:
	if $v\in \ker(\Id-\gamma K^*)$, i.e., $\gamma K^*v = v$, then 
	$\gamma\langle Kf,v \rangle_{L^2(\Omega)} 
	= \gamma\langle f,K^*v \rangle_{L^2(\Omega)}
	=  \langle f,v \rangle_{L^2(\Omega)}$,
	so, $(*)$ holds, even when $\gamma=0$.
	
	We claim that 
	\begin{equation}\label{eq6822e211}
		\cal M = \ker(\Id-\gamma K^*) .
	\end{equation}
	Notice that, similarly to what we have done in~\eqref{eq6821bd98},
	\begin{equation}\label{eq6822ef43}
	\begin{aligned}
		v \in \cal M
		&\IFF \scr E_L[u,v] = 0 \quad\forall u\in H^1_0(\Omega) \\
		&\IFF B_\gamma[u,v] \overset{\text{def}}= B[u,v] + \gamma \langle u,\bar v \rangle_{L^2(\Omega)} = \langle  u , \gamma \bar v \rangle_{L^2(\Omega)}\qquad\forall u\in H^1_0(\Omega) \\
		&\IFF {}_{H^{-1}_0(\Omega)}\langle T_\gamma u | \bar v \rangle_{H^1_0(\Omega)} = \langle  u , \gamma \bar v \rangle_{L^2(\Omega)}\qquad\forall u\in H^1_0(\Omega) \\
		&\IFF {}_{H^{1}_0(\Omega)}\langle  u | T_\gamma^\top \bar v \rangle_{H^{-1}_0(\Omega)} = \langle  u , \gamma \bar v \rangle_{L^2(\Omega)}\qquad\forall u\in H^1_0(\Omega) \\
		&\IFF T_\gamma^\top \bar v = \gamma \bar v \\
		&\IFF \bar v - \gamma (T_\gamma^\top)^{-1} \bar v = 0 \\
		&\IFF v - \gamma \overline{ (T_\gamma^\top)^{-1} \bar v} = 0 
%		&\IFF \bar v\in \ker(\Id - \gamma (T_\gamma^\top)^{-1}) \\
%		&\IFF v \in \ker(\Id - \gamma \overline{(T_\gamma^\top)^{-1}}) .
	\end{aligned}
	\end{equation}
	Here we have used the notation $(\cdot)^\top$ to denote the dual map: if $T:X\to Y$ is an operator between Banach spaces, then $T^\top: Y'\to X'$ is the dual map defined by 
	\begin{equation}\label{eq6822e47f}
		{}_{Y'}\langle \zeta ,Tx \rangle_Y = {}_{X'}\langle T^\top\zeta | x \rangle_X ,
		\qquad\forall x\in X,\ \zeta\in Y' .
	\end{equation}
	To complete the proof of our claim~\eqref{eq6822e211}, we need to show that
	\begin{equation}\label{eq6822e442}
		\forall w\in L^2(\Omega),\qquad
		\overline{(T_\gamma^\top)^{-1} \bar w} = K^* w .
%		\text{ i.e., }
%		(T_\gamma^\top)^{-1}w = \overline{K^*w} .
	\end{equation}
	First, notice that $(T_\gamma^\top)^{-1} = (T_\gamma^{-1})^\top$, by an easy argument using directly~\eqref{eq6822e47f}:
	\begin{equation}
		\forall x\in X,\ \forall\xi\in X',\quad
		\langle \xi|x \rangle 
		= \langle \xi | T^{-1}Tx \rangle
		= \langle T^\top (T^{-1})^\top\xi|x \rangle 
		\quad\THEN\quad 
		T^\top (T^{-1})^\top = \Id_{X'} .
	\end{equation}
	So,~\eqref{eq6822e442} reduces to showing $\overline{ K^\top \bar w } = K^*w$ for all $w\in L^2(\Omega)$.
%	 the two conjugate are defined using different pairings, but they agree on $L^2$ functions.
	More precisely, if $\tilde K=T_\gamma^{-1}:H^{-1}_0(\Omega)\to H^1_0(\Omega)$ (which is the operator from which $K$ descends),
	then $\tilde K^\top:H^{-1}_0(\Omega)\to H^1_0(\Omega)$ and
	\begin{equation}\label{eq6822e751}
		\forall a,b\in H^{-1}_0(\Omega),\qquad
		{}_{H^{1}_0(\Omega)}\langle \tilde Ka | b \rangle_{H^{-1}_0(\Omega)}
		= {}_{H^{-1}_0(\Omega)}\langle a | \tilde K^\top b \rangle_{H^{1}_0(\Omega)}
	\end{equation}
	If $a,b\in L^2(\Omega)\subset H^{-1}_0(\Omega)$, then~\eqref{eq6822e751} says
	\begin{align}
		\langle a, K^* b \rangle_{L^2(\Omega)} 
		&= \langle K a, b \rangle_{L^2(\Omega)} \\
		&= \langle \tilde K a , b  \rangle_{L^2(\Omega)} \\
		&= \int_\Omega \tilde K a(x) \bar b(x) \did x \\
		&= {}_{H^{1}_0(\Omega)}\langle \tilde Ka | \bar b \rangle_{H^{-1}_0(\Omega)} \\
		&= {}_{H^{-1}_0(\Omega)}\langle a | \tilde K^\top \bar b \rangle_{H^{1}_0(\Omega)} \\
		&= \int_\Omega  a(x) \tilde K^\top\bar b(x) \did x \\
		&= \langle  a , \overline{\tilde K^\top \bar b}  \rangle_{L^2(\Omega)} .
	\end{align}
	We can now complete the sequence of equivalences~\eqref{eq6822ef43} with
	\begin{align}
		v \in \cal M
		&\overset{\eqref{eq6822ef43}}\IFF v - \gamma \overline{ (T_\gamma^\top)^{-1} \bar v} = 0 \\
		&\IFF v- \gamma K^*v = 0 .
	\end{align}
	We have thus proven our claim~\eqref{eq6822e211}.
	
	With Theorem~\ref{thm6821bc90} and claim~\eqref{eq6822e211},
	we have proven all the statements of Theorem~\ref{thm6822fa7f}.
\end{proof}

\begin{exercise}
	If $u\in\scr N$ means $Lu=0$, what does $v\in\scr M$ mean?
\end{exercise}

%%%%%%%%%%%%%%%%%%%%%%%%%%%%%%%%%%%%%%%%%%%%%%%%%%%%%%%%%%%%%%%
\subsection{Regularity}

We assume that $L$ has bounded coefficients, i.e.,~\eqref{eq6820970d},
and that $L$ is elliptic, i.e.,~\eqref{eq6820a660}.

\begin{lemma}[Caccioppoli Inequality\footnote{\url{https://en.wikipedia.org/wiki/Renato_Caccioppoli}}]\label{lem6822fc0f}
	Let $n\ge2$ and $\Omega\subset\R^n$ open.
	Then there exists $C\in\R$ such that the following holds.
	
	If $u\in H^1(\Omega)$ is such that $Lu=f$ in $\Omega$ for some $f\in L^2(\Omega)$,
	i.e., 
	\begin{equation}
		\forall v\in H^1_0(\Omega),\qquad
		\scr E_L(u,v) = \langle u,f \rangle_{L^2(\Omega)} ,
	\end{equation}
	then
	\begin{equation}\label{eq6822fc08}
		\|\grad u\|_{L^2(\Omega)}^2 \le C (\|u\|_{L^2(\Omega)}^2 + \|f\|_{L^2(\Omega)}^2 ) .
	\end{equation}
\end{lemma}
\begin{proof}
	Let $\Omega'\Subset\Omega$ and $\zeta\in C^\infty(\Omega;[0,1])$ be such that
	\begin{equation}
		\Omega' \subset\interior\{\zeta=1\} \subset\spt(\zeta) \Subset \Omega .
	\end{equation}
	
	\begin{align}
		\theta \|\grad u\|_{L^2(\Omega')}^2
		&= \theta \langle \grad u,\grad u \rangle_{L^2(\Omega')} \\
		&\overset{\eqref{eq6820a660}}\le \langle A\grad u,\grad u \rangle_{L^2(\Omega')} \\
		&\le \langle A\grad u,\grad (\zeta u) \rangle_{L^2(\Omega)} \\
		&\le \scr E_L(u,\zeta u) - \langle b\cdot\grad u + cu , \zeta u \rangle_{L^2(\Omega)} \\
		&= \langle f - b\cdot\grad u - cu,\zeta u \rangle_{L^2(\Omega)} \\
		&\overset{\eqref{Hölder}}\le (\|f\|_{L^2(\Omega)} + \|b\|_{L^\infty(\Omega)} \|\grad u\|_{L^2(\Omega)} + \|c\|_{L^\infty(\Omega)} \|u\|_{L^2(\Omega)} ) \|\zeta u\|_{L^2(\Omega)} \\
		&\overset{\eqref{Cauchy}}\le \frac{\|f\|_{L^2(\Omega)}^2}{2} + \frac{\|\zeta u\|_{L^2(\Omega)}^2}{2} \\
		&\qquad	+ \frac{\theta}{2} \|\grad u\|_{L^2(\Omega)}^2 + \frac{\|b\|_{L^\infty(\Omega)}}{2\theta} \|\zeta u\|_{L^2(\Omega)}^2 \\
		&\qquad + \|c\|_{L^\infty(\Omega)} \|u\|_{L^2(\Omega)}\|\zeta u\|_{L^2(\Omega)} \\
	[\text{by }0\le\zeta\le1]
		&\le \frac{\|f\|_{L^2(\Omega)}^2}{2} 
			+ \left(\frac12 + \frac{\|b\|_{L^\infty(\Omega)}}{2\theta} + \|c\|_{L^\infty(\Omega)} \right) \| u\|_{L^2(\Omega)}^2 
			+ \frac{\theta}{2} \|\grad u\|_{L^2(\Omega)}^2 .
	\end{align}
	Therefore
	\begin{equation}\label{eq6823014e}
		2 \|\grad u\|_{L^2(\Omega')}^2 - \|\grad u\|_{L^2(\Omega)}^2 
		\le \frac2\theta \frac{\|f\|_{L^2(\Omega)}^2}{2}
			+ \left( 1 + \|b\|_{L^\infty(\Omega)} + 2\theta \|c\|_{L^\infty(\Omega)} \right) \| u\|_{L^2(\Omega)}^2 . 
	\end{equation}
	Since the right-hand side of estimate~\eqref{eq6823014e} does not depend on $\Omega'$, 
	if we take the supremum over all $\Omega'\Subset\Omega$, we obtain
	\begin{equation}
		\|\grad u\|_{L^2(\Omega)}^2 
		\le \frac2\theta \frac{\|f\|_{L^2(\Omega)}^2}{2}
			+ \left( 1 + \|b\|_{L^\infty(\Omega)} + 2\theta \|c\|_{L^\infty(\Omega)} \right) \| u\|_{L^2(\Omega)}^2 .
	\end{equation}
	We have thus obtained~\eqref{eq6822fc08}.
\end{proof}

\begin{lemma}\label{lem6823059f}
	Let $n\ge2$ and $\Omega\subset\R^n$ open.
	Suppose that $u\in H^1(\Omega)$ is such that $Lu=f$ for some $f\in L^2(\Omega)$, where $L$ has $b=0$ and $c=0$.
	Equivalently, $u\in H^1(\Omega)$ and $f\in L^2(\Omega)$ are such that
	\begin{equation}\label{eq682306a7}
		\forall v\in H^1_0(\Omega),\qquad
		\langle A\grad u,\grad v \rangle_{L^2(\Omega)} = \langle f,v \rangle_{L^2(\Omega)} ,
	\end{equation}
	where $A\in L^\infty(\Omega;\C^{n\times n})$ satisfies the ellipticity condition~\eqref{eq6820a660}.
	
%	Then, for every $j\in\{1,\dots,n\}$ and $h\neq0$,
%	\begin{equation}\label{eq68230594}
%		\|\Delta_j^h\grad u\|_{L^2(\Omega)} \le C (\|u\|_{L^2(\Omega)} + \|f\|_{L^2(\Omega)}) .
%	\end{equation}
	
	For every $\Omega''\Subset\Omega$ there exist constants $C$ and $\epsilon$ depending on $\Omega''$, $\| A\|_{L^\infty(\Omega)}$, $\|\Diff A\|_{L^\infty(\Omega)}$, and $\theta$, such that,
	for every $h$ with $0<|h|<\epsilon$, and for every $j\in\{1,\dots,n\}$,
	\begin{equation}\label{eq68231837}
		\|\Delta_j^{h}\grad u\|_{L^2(\Omega'')}^2 \le C ( \|u\|_{L^2(\Omega)}^2 + \| \grad u\|_{L^2(\Omega)}^2 + \|f\|_{L^2(\Omega)}^2 ) .
	\end{equation}
\end{lemma}
\begin{proof}
	 Fix $j\in\{1,\dots,n\}$ and $h\neq0$.
	 Let $\Omega'\Subset\Omega'\Subset\Omega$ and $\zeta\in C^\infty(\Omega;[0,1])$ be such that
	\begin{equation}
		\Omega'' \subset\interior\{\zeta=1\} \subset\spt(\zeta) \Subset \Omega' \Subset\Omega .
	\end{equation}
	Then there exists $\epsilon>0$ such that, for every $h$ with $0<|h|<\epsilon$,
	\begin{equation}
		v := -\Delta_j^{-h} (\zeta^2 \Delta_j^h u) \in H^1_0(\Omega) .
	\end{equation}
	With the aim of applying~\eqref{eq682306a7} with this test function, we make the following estimates.
	We see that
	\begin{align}
		\langle A\grad u,\grad v \rangle_{L^2(\Omega)}
		&= \langle A\grad u,\grad (-\Delta_j^{-h} (\zeta^2 \Delta_j^h u)) \rangle_{L^2(\Omega)} \\
		&\overset{\eqref{eq681734d2}}= \langle \Delta_j^{h}(A\grad u),\grad (\zeta^2 \Delta_j^h u) \rangle_{L^2(\Omega)} \\
		&\overset{\eqref{eq681734d8}}= \langle (\Delta_j^{h}A)\grad u + A^{(h)} \Delta_j^{h}\grad u , 
			\Delta_j^h u 2\zeta\grad\zeta + \zeta^2 \grad(\Delta_j^h u) \rangle_{L^2(\Omega)} \\
		&= F_1 + F_2 ,
	\end{align}
	where $A^{(h)}(x) = A(x+he_j)$ and
	\begin{align}
		F_1
		&= \langle A^{(h)} \Delta_j^{h}\grad u , 
			\zeta^2 \grad(\Delta_j^h u) \rangle_{L^2(\Omega)} \\
		&= \langle A^{(h)} \zeta\Delta_j^{h}\grad u , 
			\zeta \Delta_j^h\grad u \rangle_{L^2(\Omega)} \\
		&\overset{\eqref{eq6820a660}}\ge \theta \|\zeta\Delta_j^{h}\grad u\|_{L^2(\Omega)} \\
		&\ge \theta \|\Delta_j^{h}\grad u\|_{L^2(\Omega'')} ; \\
%	\end{align}
%	\begin{align}
		F_2
		&= \langle (\Delta_j^{h}A)\grad u , \Delta_j^h u 2\zeta\grad\zeta \rangle_{L^2(\Omega)} \\
		&\qquad	+ \langle (\Delta_j^{h}A)\grad u ,\zeta^2 \grad(\Delta_j^h u) \rangle_{L^2(\Omega)} \\
		&\qquad	+ \langle A^{(h)} \Delta_j^{h}\grad u , \Delta_j^h u 2\zeta\grad\zeta \rangle_{L^2(\Omega)} \\
		%
		&\le 2 \|\Delta_j^{h}A\|_{L^\infty(\Omega')} \|\grad u \|_{L^2(\Omega')} \| \Delta_j^h u\|_{L^2(\Omega')}
			\| \grad\zeta \|_{L^\infty(\Omega')} \\
		&\qquad	+ \|\Delta_j^{h}A\|_{L^\infty(\Omega')} \|\grad u \|_{L^2(\Omega')} \| \zeta\Delta_j^h \grad u\|_{L^2(\Omega')} \\
		&\qquad	+ 2 \|A\|_{L^\infty(\Omega')} \|\zeta\Delta_j^{h}\grad u\|_{L^2(\Omega')} \| \Delta_j^h u\|_{L^2(\Omega')} 
			\| \grad\zeta \|_{L^\infty(\Omega')} \\
		%
		&\le 2 \|\Diff A\|_{L^\infty(\Omega)} \|\grad u \|_{L^2(\Omega)}^2 \| \grad\zeta \|_{L^\infty(\Omega)} \\
		&\qquad	+ \|\Diff A\|_{L^\infty(\Omega)} \|\grad u \|_{L^2(\Omega)} \|\zeta\Delta_j^{h}\grad u\|_{L^2(\Omega)} \\
		&\qquad	+ 2 \|A\|_{L^\infty(\Omega)} \|\zeta\Delta_j^{h}\grad u\|_{L^2(\Omega)} \| \grad u\|_{L^2(\Omega)} 
			\| \grad\zeta \|_{L^\infty(\Omega)} \\
		%
		&\overset{\eqref{Cauchy}}\le 2 \|\Diff A\|_{L^\infty(\Omega)} \| \grad\zeta \|_{L^\infty(\Omega)} \|\grad u \|_{L^2(\Omega)}^2\\
		&\qquad		+ 4\| \grad\zeta \|_{L^\infty(\Omega)}^2 \frac{ \|A\|_{L^\infty(\Omega)}^2 +  \|\Diff A\|_{L^\infty(\Omega)}^2 }{ \theta } \| \grad u\|_{L^2(\Omega)}^2 \\
		&\qquad	+ \frac{\theta}{2} \|\zeta \Delta_j^{h}\grad u\|_{L^2(\Omega)}^2 .
	\end{align}
	Since
	\begin{equation}
		\langle A\grad u,\grad v \rangle_{L^2(\Omega)} 
		= \langle u,f \rangle_{L^2(\Omega)}\le \|u\|_{L^2(\Omega)}\|f\|_{L^2(\Omega)} 
		\overset{\eqref{Cauchy}}\le \frac{ \|u\|_{L^2(\Omega)}^2 }{2} + \frac{ \|f\|_{L^2(\Omega)}^2 }{2} ,
	\end{equation}
	we obtain
	\begin{align}
		\frac{\theta}{2} \|\zeta \Delta_j^{h}\grad u\|_{L^2(\Omega)}^2
		&\le \frac{ \|u\|_{L^2(\Omega)}^2 }{2} + \frac{ \|f\|_{L^2(\Omega)}^2 }{2} \\
		&\qquad	+ 2 \|\Diff A\|_{L^\infty(\Omega)} \| \grad\zeta \|_{L^\infty(\Omega)} \|\grad u \|_{L^2(\Omega)}^2 \\
		&\qquad	+ 4\| \grad\zeta \|_{L^\infty(\Omega)}^2 \frac{ \|A\|_{L^\infty(\Omega)}^2 +  \|\Diff A\|_{L^\infty(\Omega)}^2 }{ \theta } \| \grad u\|_{L^2(\Omega)}^2 .
	\end{align}
	We conclude that there is a constant $C\in\R$ depending on $\| A\|_{L^\infty(\Omega)}$, $\|\Diff A\|_{L^\infty(\Omega)}$, $\| \grad\zeta \|_{L^\infty(\Omega)}$ and $\theta$ such that
	\begin{equation}
		\|\zeta \Delta_j^{h}\grad u\|_{L^2(\Omega)}^2 \le C ( \|u\|_{L^2(\Omega)}^2 + \| \grad u\|_{L^2(\Omega)}^2 + \|f\|_{L^2(\Omega)}^2 ) .
	\end{equation}
	Since $\zeta=1$ on $\Omega''$, then we get~\eqref{eq68231837}.
\end{proof}

%%%%%%%%%%%%%%%%%%%%%%%%%%%%%%%%%%%%%%%%%%%%%%%%%%%%%%%%%%%%%%%
\begin{theorem}[Regularity of solutions to Elliptic PDE]\label{thm68237c1b}
	Let $n\ge2$ and $\Omega\subset\R^n$ open.
	
	Let $L$ be a differential operator as in~\eqref{eq6820970d},
	with bounded coefficients~\eqref{eq6820a660},
	and satisfying the ellipticity condition~\eqref{eq6820a731}.
	Moreover, we assume that
	\begin{equation}
		\|\Diff A\|_{L^\infty(\Omega)} < \infty ,
	\end{equation}
	that is, that $A$ is Lipschitz.
	
	Then, for every $\Omega'\Subset\Omega$ there exists a constant $C$ such that
	\begin{equation}\label{eq68237d5d}
		\forall u\in H^1(\Omega),\qquad
		\|\Diff^2u\|_{L^2(\Omega')} \le C (\|u\|_{L^2(\Omega)} + \|Lu\|_{L^2(\Omega)} ) .
	\end{equation}
	In particular, if $u\in H^1(\Omega)$ is such that $Lu\in L^2(\Omega)$, then $u\in H^2_\loc(\Omega)$.
\end{theorem}
\begin{proof}
	Let $\Omega'\Subset\Omega$ and $u\in H^1(\Omega)$ with $f=Lu\in L^2(\Omega)$ 
	(if $\|Lu\|_{L^2(\Omega)}=\infty$, then~\eqref{eq68237d5d} is trivial.
	Define $M$ as the differential operator
	\begin{equation}
		\forall w\in H^1(\Omega),\qquad
		Mw = -\div(A\grad w) = Lw - ( b\cdot\grad w + cw ) .
	\end{equation}
	Notice that, if $w\in H^1(\Omega)$, then $b\cdot\grad w + cw\in L^2(\Omega)$.
	So,
	\begin{equation}
		 Mu = \tilde f := f - (b\cdot\grad u + cu) \in L^2(\Omega) .
	\end{equation}
	Notice that
	\begin{equation}\label{eq6823846a}
		\|\tilde f\|_{L^2(\Omega)}
		\le \|b\|_{L^\infty(\Omega)} \|\grad u\|_{L^2(\Omega)} + \|c\|_{L^\infty(\Omega)} \|u\|_{L^2(\Omega)}.
	\end{equation}
	Applying first Lemma~\ref{lem6823059f},
%	and then the Caccioppoli Inequality of Lemma~\ref{lem6822fc0f},
	we obtain a constant $C$ (independent of $u$) such that
	\begin{equation}
		\limsup_{h\to0}\|\Delta_j^{h}\grad u\|_{L^2(\Omega'')} 
		\le C (\|u\|_{L^2(\Omega)} + \|\grad u\|_{L^2(\Omega)} + \|\tilde f\|_{L^2(\Omega)} ) 
	\end{equation}
	The estimate~\eqref{eq6823846a} and the Caccioppoli Inequality of Lemma~\ref{lem6822fc0f}, then implies that there exists a constant $C$ (independent of $u$) such that
	\begin{equation}
		\limsup_{h\to0}\|\Delta_j^{h}\grad u\|_{L^2(\Omega'')} 
		\le C (\|u\|_{L^2(\Omega)} + \| f\|_{L^2(\Omega)} ) .
	\end{equation}
	By Theorem~\ref{thm681736aa} (using also the quantitative version given by Proposition~\ref{prop68176bc1}), we obtain~\eqref{eq68237d5d}.
\end{proof}

\begin{exercise}
	Re-write Theorem~\ref{thm68237c1b}, together with Lemmata~\ref{lem6822fc0f} and~\ref{lem6823059f}, and their proofs, for $L=\laplacian$ on $\R^n$.
\end{exercise}

%%%%%%%%%%%%%%%%%%%%%%%%%%%%%%%%%%%%%%%%%%%%%%%%%%%%%%%%%%%%%%%
\begin{corollary}[Regularity of solutions to Elliptic PDE]\label{cor68237f2f}
	Let $n\ge2$ and $\Omega\subset\R^n$ open.
	
	Let $L$ be a differential operator as in~\eqref{eq6820970d},
	with bounded coefficients~\eqref{eq6820a660},
	and satisfying the ellipticity condition~\eqref{eq6820a731}.
	Moreover, we assume that, for some $m\in\N$, we have
	\begin{equation}\label{eq682385e0}
		A\in W^{m+1,\infty}(\Omega;\C^{n\times n}),\qquad
		b\in W^{m,\infty}(\Omega;\C^n),\qquad
		c\in W^{m,\infty}(\Omega;\C).
	\end{equation}
		
	Then, for every $\Omega'\Subset\Omega$ there exists a constant $C$ such that
	\begin{equation}\label{eq68237d5d}
		\forall u\in H^1(\Omega),\qquad
		\|u\|_{H^{m+2}(\Omega')} \le C (\|u\|_{L^2(\Omega)} + \|Lu\|_{H^m(\Omega)} ) .
	\end{equation}
	In particular, if $u\in H^1(\Omega)$ is such that $Lu\in H^m(\Omega)$, then $u\in H^{m+2}_\loc(\Omega)$.
\end{corollary}
\begin{proof}
	We prove the statement by induction over $m$.
	Let $\ell\in\N$. If $\ell=0$, then we just apply Theorem~\ref{thm68237c1b}.
	Suppose $\ell>0$ and suppose that the statement above holds for $m<\ell$:
	we will prove it for $m=\ell$.
	Under the hypothesis above on $L$, we have that for every $\Omega'\Subset\Omega$ there exists a constant $C$ such that
	\begin{equation}\label{eq682387da}
		\forall u\in H^1(\Omega),\qquad
		\|u\|_{H^{\ell+1}(\Omega')} \le C (\|u\|_{L^2(\Omega)} + \|Lu\|_{H^{\ell-1}(\Omega)} ) .
	\end{equation}
	For each $j\in\{1,\dots,n\}$, we have
	\begin{equation}\label{eq68239171}
	\begin{aligned}
		L\de_ju
		&= -\div(A\grad\de_ju) + b\cdot\grad \de_ju + c \de_ju\\
		&= \de_j(-\div(A\grad u) + b\cdot\grad u + cu)
			- \left( -\div(\de_jA\grad u) + \de_jb\cdot\grad u + \de_jcu \right) \\
		&= \de_j Lu - L_ju .
	\end{aligned}
	\end{equation}
	Straightforward estimates give, for every $m\in\N$ for which~\eqref{eq682385e0} holds, we have $C_m'<\infty$ such that
	\begin{equation}\label{eq6823968e}
	\begin{aligned}
		\|L_ju\|_{H^m(\Omega)}
		&\le \|\div(\de_jA\grad u)\|_{H^m(\Omega)} + \|\de_jb\cdot\grad u\|_{H^m(\Omega)} + \|\de_jcu\|_{H^m(\Omega)} \\
		&\le \sum_{|\alpha|\le m} \|\Diff^\alpha(\div(\de_jA\grad u))\|_{L^2(\Omega)} 
			+ \|\Diff^\alpha\de_jb\cdot\grad u\|_{L^2(\Omega)} 
			+ \|\Diff^\alpha\de_jcu\|_{L^2(\Omega)} \\
		&\le C_m ( \|A\|_{W^{m+2,\infty}(\Omega)} 
			+ \|b\|_{W^{m+1,\infty}(\Omega)}
			+ \|c\|_{W^{m+1,\infty}(\Omega)})
			\|u\|_{H^{m+2}(\Omega)} \\
		&\overset{\eqref{eq682385e0}}\le C_m' \|u\|_{H^{m+2}(\Omega)} .
	\end{aligned}
	\end{equation}
	So
	\begin{align}
		\|u\|_{H^{\ell+2}(\Omega')}
		&\le \|u\|_{L^2(\Omega)} + \sum_{j=1}^n \|\de_ju\|_{H^{\ell+1}(\Omega')} \\
		&\overset{\eqref{eq682387da}}\le \|u\|_{L^2(\Omega)} 
			+ C \sum_{j=1}^n (\|\de_ju\|_{L^2(\Omega)} + \|L\de_ju\|_{H^{\ell-1}(\Omega')} ) \\
		&\overset{\eqref{eq68239171}}\le \|u\|_{L^2(\Omega)} 
			+ C \sum_{j=1}^n (\|\de_ju\|_{L^2(\Omega)} + \|\de_jLu\|_{H^{\ell-1}(\Omega')} + \|L_ju\|_{H^{\ell-1}(\Omega')}) \\
		&\overset{\eqref{eq6823968e}}\le C' \left( \|u\|_{H^1(\Omega)} + \|Lu\|_{H^{\ell}(\Omega)} + \|u\|_{H^{\ell+1}(\Omega)} \right) \\
		&\overset{\eqref{eq682387da}}\le C'' \left( \|Lu\|_{H^{\ell}(\Omega)} + \|u\|_{L^2(\Omega)} + \|Lu\|_{H^{\ell-1}(\Omega)} \right) \\
		&\le C''' \left( \|Lu\|_{H^{\ell}(\Omega)} + \|u\|_{L^2(\Omega)} \right) .
	\end{align}
	We eventually have~\eqref{eq68237d5d} for $m=\ell$.
\end{proof}


%%%%%%%%%%%%%%%%%%%%%%%%%%%%%%%%%%%%%%%%%%%%%%%%%%%%%%%%%%%%%%%
\newpage
\section{Exercises}

%%%%%%%%%%%%%%%%%%%%%%%%%%%%%%%%%%%%%%%%%%%%%%%%%%%%%%%%%%%%%%%
\subsection{Eigenvalues of laplacian on interval}\label{subs6890901c}
We want to study, for $\lambda\in\C$ and $\ell>0$,
\begin{equation}
\begin{cases}
	-\laplacian u + \lambda u = 0 &\text{ in }(0,\ell) ,\\
	u(0) = u(\ell) = 0 .
\end{cases}
\end{equation}
In other words,
\begin{equation}\label{eq68342701}
\begin{cases}
	u'' = \lambda u &\text{ in }(0,\ell) ,\\
	u\in H^1((0,\ell)) .
\end{cases}
\end{equation}

\begin{lemma}\label{lem6834273f}
	Let $u$ be a solution to~\eqref{eq68342701}.
	Then $u\in C^\infty((0,\ell))$.
\end{lemma}
\begin{proof}
	This is an application of Corollary~\ref{cor68237f2f}.
\end{proof}

Be aware that the regularity stated in Lemma~\ref{lem6834273f} is in the open interval $(0,\ell)$ and not on the whole $[0,\ell]$.
We will procede as follows:
\begin{enumerate}
\item
	the assumption $u\in H^1_0((0,\ell))$ alone gives us that $u$ is continuous on $[0,\ell]$
and that $u(0)=u(\ell)$;
see the following two Lemmata~\ref{lem68346d4d} and~\ref{lem68342e8e}.
\item
	Then, we will extend $u$ to $\R$ in such a way that the extension $\tilde u =E(u)$ is a bounded function on $\R$ that solves $\laplacian \tilde u=\lambda \tilde u$ on $\R$; see Lemma~\ref{lem68344adb} and Corollary~\ref{cor6834905d}.
\item
	We will find all solutions of $\laplacian v=\lambda v$ for $v\in \scr S'(\R)$ using the Fourier transform; see Proposition~\ref{prop6834897a}.
\item
	Among the solutions $v$ on $\R$, we look for those $v$ that are extensions of functions in $H^1_0((0,\ell))$;
	see Proposition~\ref{prop68349083}. 
	In this way, we will find all solutions to~\eqref{eq68342701}.
\end{enumerate}


\begin{exercise}
	Try to guess solutions to~\eqref{eq68342701}.
	Try also to show they are the only one.
\end{exercise}

%%%%%%%%%%%%%%%%%%%%%%%%%%%%%%%%%%%%%%%%%%%%%%%%%%%%%%%%%%%%%%%
\subsubsection{Fundamental properties of $u$}

\begin{lemma}\label{lem68346d4d}
	If $u\in W^{1,1}((0,\ell))$, then, up to changing $u$ on a set of measure zero, $u\in C([0,\ell])$ and 
	\begin{equation}
		\forall x\in[0,\ell],\qquad
		u(x) = u(0) + \int_0^x u'(y) \did y .
	\end{equation}
\end{lemma}
\begin{proof}
	Define $v:(0,\ell)\to\C$ by
	\begin{equation}
		v(x) = \int_{\ell/2}^x u'(y) \did y .
	\end{equation}
	We claim that
	\begin{equation}\label{eq68346c1c}
		v' = u' .
	\end{equation}
	If $\phi\in C^\infty_c((0,\ell))$, then
	\begin{align}
		\int_0^\ell v(x)\phi'(x)\did x
		&= \int_0^\ell \int_{\ell/2}^x u'(y) \phi'(x)\did y \did x \\
		&= - \int_0^{\ell/2} \int_0^y u'(y) \phi'(x) \did x\did y 
			+ \int_{\ell/2}^\ell \int_y^\ell u'(y) \phi'(x) \did x\did y \\
		&= - \int_0^\ell u'(y) \phi(y) \did y \\
		&= \int_0^\ell u(y) \phi'(y) \did y .
	\end{align}
	This shows~\eqref{eq68346c1c}.
	
	Therefore, $(u-v)'=0$ and thus, by the Constancy Theorem~\ref{thm68346c5c}, there exists $c\in \C$ such that, for almost every $x\in(0,\ell)$, $u(x)= c+v(x)$.
	We conclude that, up to changing $u$ on a set of measure zero, 
	\begin{equation}
		u(x) = u(\ell/2) + \int_{\ell/2}^x u'(y) \did y .
	\end{equation}
	By the continuity of the integral, not only $u$ is continuous, but also
	$\lim_{x\to0} u(x)$ exists and it is equal to $u(\ell/2) - \int_0^{\ell/2} u'(y) \did y$,
	and similarly for $\lim_{x\to\ell} u(x)$.
\end{proof}

\begin{lemma}\label{lem68342e8e}
	Let $v\in H^1((0,\ell)) = W^{1,2}((0,\ell))$.
	Then $v\in W^{1,1}((0,\ell))\cap C([0,\ell])$.
	Moreover, if $v\in H^1_0((0,\ell))$, then $v(0)=v(\ell)=0$.
\end{lemma}
\begin{proof}
	Since $L^2((0,\ell))\subset L^1((0,\ell))$, then $v\in W^{1,1}((0,\ell))$.
	Lemma~\ref{lem68346d4d} implies that, up to changing $v$ on a set of measure zero, $v\in C([0,\ell])$.
	
	Suppose $v\in H^1_0((0,\ell))$.
	By definition of $H^1_0((0,\ell))$, there exists a sequence $\{v_j\}_{j\in\N}\subset C^\infty_c((0,\ell))$
	such that $v_j\to v$ in $H^1((0,\ell))$,
	that is, $\lim_{j\to\infty} \|v-v_j\|_{L^2((0,\ell))} = 0$ and $\lim_{j\to\infty} \|v'-v_j'\|_{L^2((0,\ell))} = 0$.
	Since the Hölder inequality implies $\|f\|_{L^1((0,\ell))} \le \sqrt{\ell} \|f\|_{L^2((0,\ell))}$,
	we obtain $\lim_{j\to\infty} \|v-v_j\|_{L^1((0,\ell))} = 0$ and $\lim_{j\to\infty} \|v'-v_j'\|_{L^1((0,\ell))} = 0$.
	We know then that, up to passing to a subsequence, there exists a set $I\subset(0,\ell)$ of full measure so that $\lim_{j\to\infty} v_j(x) = v(x)$ for every $x\in I$.
%	If $x,y\in I$, then
%	\begin{align}
%		v(x) - v(y) 
%		&= \lim_{j\to\infty} v_j(x) - v_j(y) \\
%		&= \lim_{j\to\infty} \int_x^y v_j'(z) \did z \\
%		&= \int_x^y v'(z) \did z .
%	\end{align}
%	Since the function $x\mapsto \int_{\ell/2^x} v'(z) \did z$ is continuous, we obtain that $v$ is continuous on $I$ and thus, up to changing $v$ on a set of measure zero, and~\eqref{eq68342e86} holds.
	
%	Since $\int_0^\ell |v'(z)| \did z<\infty$, then $v(0):= \lim_{x\to0} v(x)$ and $v(\ell):=\lim_{x\to\ell} v(x)$ exist.

	We claim that in fact the convergence $v_j\to v$ is pointwise on $[0,\ell]$ (in fact, we prove that it is uniform).
	Indeed, if $\bar x\in I$ and $x\in[0,\ell]$, then, for every $j\in\N$,
	\begin{align}
		|v(x) - v_j(x)|
		&= \left| v(\bar x) - v_j(\bar x) + \int_{\bar x}^x (v'(y)-v_j'(y)) \did y \right| \\
		&\le |v(\bar x) - v_j(\bar x)| + \int_0^\ell |v'(y)-v_j'(y)| \did y .
	\end{align}
	Therefore, $\lim_{j\to\infty} |v(x) - v_j(x)| = 0$ for every $x\in[0,\ell]$.
	
	Since $v_j(0)=v_j(\ell)=0$ for all $j$, then $v(0)=v(\ell)=0$.
\end{proof}

%%%%%%%%%%%%%%%%%%%%%%%%%%%%%%%%%%%%%%%%%%%%%%%%%%%%%%%%%%%%%%%
\subsubsection{Extension of $u$ to $\R$}

\begin{lemma}\label{lem68344adb}
%	Let $v\in W^{1,1}((0,\ell))$. 
%	We should know, and in any case we assume, that, for every $x\in[0,\ell]$,
%	\begin{equation}\label{eq683435d0}
%		v(x) = v(0) + \int_0^x v'(y) \did y .
%	\end{equation}
	For $v\in L^1((0,\ell))$, define $E(v)\in L^1_\loc(\R)$ as
	\begin{equation}
		E(v)(x) = \sum_{k\in\Z} \left( u(x\mod\ell\Z) \one_{\ell[2k,2k+1]}(x) 
			- u(-x\mod\ell\Z) \one_{\ell[2k+1,2k+2]}(x) \right) .
	\end{equation}
	Then, distributionally, for every $v\in W^{1,1}((0,\ell))$,
	\begin{equation}\label{eq6834343d}
		E(v') = E(v)' .
	\end{equation}
\end{lemma}
\begin{proof}
	From Lemma~\ref{lem68346d4d}, we know that, up to modifying $v$ on a set of measure zero, $v\in C([0,\ell])$ 
	and, for all $x\in[0,\ell]$,
	\begin{equation}\label{eq683435d0}
		v(x) = v(0) + \int_0^x v'(y) \did y .
	\end{equation}
	
	The identity~\eqref{eq683435d0} implies that, for every $\phi\in\scr D(\R)$,
	\begin{equation}\label{eq6834365d}
		\int_0^\ell u(x) \phi'(x) \did x
		= u(\ell)\phi(\ell) - u(0) \phi(0) - \int_0^\ell u'(x) \phi(x) \did x .
	\end{equation}
	Let $\phi\in\scr D(\R)$.
	\begin{align}
		- \int_\R E(v)'(x) \phi(x) \did x
		&\overset{\text{def}}= \int_\R E(v)(x) \phi'(x) \did x \\
		&= \sum_{k\in\Z} \int_0^\ell \left( u(x) \phi'(2k\ell+x) - u(\ell-x) \phi'((2k+1)\ell+x) \right) \did x \\
		&\overset{\eqref{eq6834365d}}=  \sum_{k\in\Z} u(\ell) \phi'(2k\ell+\ell) - u(0) \phi'(2k\ell) \\
		&\qquad	- u(0) \phi'((2k+1)\ell+\ell) + u(\ell) \phi'((2k+1)\ell)  \\
		&\qquad - \int_0^\ell \left( u'(x) \phi(2k\ell+x) - u'(\ell-x) \phi((2k+1)\ell+x) \right) \did x \\
		&\overset{\text{def}}= - \int_\R E(v')(x) \phi(x) \did x .
	\end{align}
\end{proof}

\begin{corollary}\label{cor6834905d}
	Let $u$ be a solution to~\eqref{eq68342701}.
	Let $\tilde u = E(u)$ be the extension of $u$ as in Lemma~\ref{lem68344adb}.
	Then
	\begin{equation}\label{eq6834683a}
		\laplacian \tilde u = \lambda \tilde u \text{ in }\R .
	\end{equation}
\end{corollary}
\begin{proof}
	If $u\in H^1_0((0,\ell))$, then $u,u'\in W^{1,1}((0,\ell))\cap C([0,\ell])$ by Lemma~\ref{}.
	Therefore, by Lemma~\ref{lem68344adb}, we have the following identities of distributions on $\R$:
	\begin{equation}
		\lambda E(u) 
		= E(\lambda u)
		= E(u'')
		=  E(u')' = E(u)'' .
	\end{equation}
	This is~\eqref{eq6834683a}.
\end{proof}

%%%%%%%%%%%%%%%%%%%%%%%%%%%%%%%%%%%%%%%%%%%%%%%%%%%%%%%%%%%%%%%
\subsubsection{Solve $v''=\lambda v$ in $\R$}

\begin{proposition}\label{prop6834897a}
	Let $\lambda\in\C$.
	A Schwartz distribution $v\in\scr S'(\R)\setminus\{0\}$ solves 
	\begin{equation}\label{eq68347805}
		v''=\lambda v 
	\end{equation}
	if and only if $\lambda=0$ and $v$ is affine,
	or $\lambda\in(-\infty,0)$ and there are $a,b\in\C$ with
	\begin{equation}\label{eq68347852}
		v = a \exp(i\sqrt{-\lambda} x)  + b \exp(-i\sqrt{-\lambda} x) 
		= a \exp(i\mu x)  + b \exp(-i\mu x) ,
	\end{equation}
	wher $\mu>0$ and $\lambda=-\mu^2$.
\end{proposition}
\begin{proof}
	$\eqref{eq68347852}\THEN\eqref{eq68347805}$.
	This just a direct computation, see Exercise~\ref{ex68347888}

	$\eqref{eq68347805}\THEN\eqref{eq68347852}$.
	Apply the Fourier transform to both sides of~\eqref{eq68347805}, to obtain
	\begin{equation}
		(2\pi i\xi)^2\scr F(v)
		\overset{\eqref{eq67fe65d6}}=\scr F(v'')
		\overset{\eqref{eq68347805}}= \scr F(\lambda v)
		= \lambda \scr F(v) .
	\end{equation}
	This means that, for every $\phi\in\scr S(\R)$,
	\begin{equation}\label{eq68347c6d}
		0 = {}_{\scr S'}\langle (2\pi i\xi)^2\scr F(v) - \lambda \scr F(v) | \phi \rangle_{\scr S}
		= {}_{\scr S'}\langle \scr F(v) | (-4\pi^2 \xi^2 - \lambda)\phi \rangle_{\scr S} .
	\end{equation}
	It follows that $\scr F(v)$ is supported on $\{\xi\in\R: 4\pi^2 \xi^2 + \lambda=0\}$.
	So, if $\lambda\notin(-\infty,0]$, then $\spt(\scr F(v))=\emptyset$ and thus $v=0$.
	If $\lambda=0$, then $v''=0$ and we know that $v$ must be affine; see Exercise~\ref{ex68347bbd}.
	
	Suppose $\lambda\in(-\infty,0)$, and let $\mu>0$ such that
	\begin{equation}
		\lambda = -\mu^2 .
	\end{equation}
	Thus, we have $\sqrt{ -\frac{\lambda}{4\pi^2} } = \frac{\mu}{2\pi}$
	The above statement about the support of $\scr F(v)$ now reads
	\begin{equation}
		\spt( \scr F(v) ) \subset \left\{ \frac{\mu}{2\pi}, - \frac{\mu}{2\pi} \right\} .
	\end{equation}
	By Proposition~\ref{prop67fe6632},
	\begin{equation}
		\scr F(v) = \sum_{\alpha=0}^\infty \left( a_\alpha \de^\alpha \delta_{\frac{\mu}{2\pi}} + b_\alpha \de^\alpha \delta_{-\frac{\mu}{2\pi}} \right) .
	\end{equation}
	We use again~\eqref{eq68347c6d} to obtain that, for every $\phi\in\scr S'(\R)$,
	\begin{align}
		0 
		&= {}_{\scr S'}\langle \scr F(v) | (-4\pi^2 \xi^2 - \lambda)\phi \rangle_{\scr S} \\
		&= \sum_{\alpha=0}^\infty \left( a_\alpha \de^\alpha \delta_{\frac{\mu}{2\pi}}[(-4\pi^2 \xi^2 - \lambda)\phi] + b_\alpha \de^\alpha \delta_{-\frac{\mu}{2\pi}}[(-4\pi^2 \xi^2 - \lambda)\phi] \right) \\
		&= \sum_{\alpha=0}^\infty \sum_{\beta=0}^\alpha \binom{\alpha}{\beta} 
			\left( a_\alpha \delta_{\frac{\mu}{2\pi}}[\de^\beta(4\pi^2 \xi^2 + \lambda)\de^{\alpha-\beta}\phi] + b_\alpha \delta_{-\frac{\mu}{2\pi}}[\de^\beta(4\pi^2 \xi^2 + \lambda)\de^{\alpha-\beta}\phi] \right) \\
		%
		&\overset{(*)}= \sum_{\alpha=1}^\infty \sum_{\beta=1}^\alpha \binom{\alpha}{\beta} 
			\left( a_\alpha \delta_{\frac{\mu}{2\pi}}[\de^\beta(4\pi^2 \xi^2 + \lambda)\de^{\alpha-\beta}\phi] + b_\alpha \delta_{-\frac{\mu}{2\pi}}[\de^\beta(4\pi^2 \xi^2 + \lambda)\de^{\alpha-\beta}\phi] \right) \\
		%
		&=  a_1  8\pi^2 \frac{\mu}{2\pi} \phi(\frac{\mu}{2\pi})  
			- b_1 8\pi^2 \frac{\mu}{2\pi} \phi(-\frac{\mu}{2\pi})  \\
		&\qquad+ \sum_{\alpha=2}^\infty 
			\bigg( a_\alpha (\alpha 8\pi^2 \frac{\mu}{2\pi} \de^{\alpha-1}\phi(\frac{\mu}{2\pi}) 
				+ \alpha(\alpha-1) 8\pi^2 \de^{\alpha-2}\phi(\frac{\mu}{2\pi})  ) \\
		&\qquad\qquad	+ b_\alpha (-\alpha 8\pi^2 \frac{\mu}{2\pi} \de^{\alpha-1}\phi(-\frac{\mu}{2\pi}) 
				+ \alpha(\alpha-1) 8\pi^2 \de^{\alpha-2}\phi(-\frac{\mu}{2\pi}) )
			\bigg) ,
	\end{align}
	where in $(*)$ we observed that, if $\beta=0$, then the summand is zero.
	Since $\phi$ is an arbitrary function in $\scr S(\R)$, we obtain $a_\alpha=b_\alpha=0$ for all $\alpha\ge1$.
	Hence,
	\begin{equation}
		\scr F(v) = a_0 \delta_{\frac{\mu}{2\pi}} + b_0 \delta_{-\frac{\mu}{2\pi}} .
	\end{equation}
	From Exercise~\ref{ex6834885c}, we obtain
	\begin{align}
		v(x) 
		&\overset{\eqref{eq683488c2}}= a_0 \exp(2\pi i\frac{\mu}{2\pi} x)  + b_0 \exp(-2\pi i\frac{\mu}{2\pi} x) \\
		&= a_0 \exp(i\mu x)  + b_0 \exp(-i\mu x) .
	\end{align}
\end{proof}

\begin{exercise}\label{ex68347888}
	Show $\eqref{eq68347852}\THEN\eqref{eq68347805}$ in Proposition~\ref{prop6834897a}.
\end{exercise}

Proposition~\ref{prop6834897a} states that the spectrum of $-\de^2=-\laplacian$ on $\R$ is $[0,+\infty)$:
\begin{equation}
	\sigma(-\laplacian) = [0,+\infty) .
\end{equation}

%%%%%%%%%%%%%%%%%%%%%%%%%%%%%%%%%%%%%%%%%%%%%%%%%%%%%%%%%%%%%%%
\subsubsection{Select solutions that are extensions}

\begin{proposition}\label{prop68349083}
	A non-zero function $u\in H^1_0((0,\ell))$ solves~\eqref{eq68342701} if and only if the following two hold
	\begin{enumerate}
	\item
	there is $k\in\N\setminus\{0\}$ such that $\lambda = - \left( \frac{\pi}{\ell} k \right)^2$, and
	\item
	there is $a\in\C\setminus\{0\}$ such that
	\begin{equation}
		u(x) = a \left( \exp\left(i \frac{k\pi}{\ell} x \right) - \exp\left(-i \frac{k\pi}{\ell} x \right) \right) \\
		= 2ia \sin\left(\frac{k\pi}{\ell} x\right) .
	\end{equation}
	\end{enumerate}
\end{proposition}
\begin{proof}
	If $u$ solves~\eqref{eq68342701},
	then its extension $\tilde u = E(u)$ defined in Lemma~\ref{lem68344adb} is of the form given by Proposition~\ref{prop6834897a},
	with the additional boundary conditions $\tilde u(0)=\tilde u(\ell)=0$.
	If $\tilde u$ is affine, then $\tilde u =0$.
	Therefore, if $\tilde u\neq0$, then
	there are $a,b\in\C$ with~\eqref{eq68347852}, i.e.,
	\begin{equation}
		\tilde u(x) = a \exp(i\mu x) + b \exp(-i\mu x) ,
	\end{equation}
	where $\mu>0$ and $\lambda=-\mu^2$.
	
	The condition $\tilde u(0)=0$ implies $b=-a$.
	The condition $\tilde u(\ell)=0$, implies 
	\begin{equation}
		0 = a \exp(i\mu \ell) - a \exp(-i\mu \ell)
		= a\exp(i\mu\ell) (1-\exp(-2i\mu\ell)) .
	\end{equation}
	Therefore, either $a=v=0$, or $2\mu\ell\in 2\pi\Z$, 
	i.e., $\mu\ell\in \pi\Z$, 
	i.e., $\sqrt{-\lambda} = \mu\in \frac{\pi}{\ell}\N$.
\end{proof}

%%%%%%%%%%%%%%%%%%%%%%%%%%%%%%%%%%%%%%%%%%%%%%%%%%%%%%%%%%%%%%%
\subsection{More about the extension operator $E$}
Lemma~\ref{lem68344adb} gives a linear operator $E:L^1((0,\ell))\to L^1_\loc(\R)$.





%%%%%%%%%%%%%%%%%%%%%%%%%%%%%%%%%%%%%%%%%%%%%%%%%%%%%%%%%%%%%%%
\newpage
%\appendix
\part{Extras}

%%%%%%%%%%%%%%%%%%%%%%%%%%%%%%%%%%%%%%%%%%%%%%%%%%%%%%%%%%%%%%%
\section{Comments}

\subsection{Correct proof of Theorem~\ref{thm67bb666b}}\label{subs688c7a1a}
%%%%%%%%%%%%%%%%%%%%%%%%%%%%%%
\begin{proof}[Correct proof of Theorem~\ref{thm67bb666b}]
	Let $U\subset\R^n$ be open and $u\in C^\infty(U)$ a harmonic function.
	Fix $\hat x\in U$ and set $\hat r = \frac14 \dist(\hat x,\de U)$.
	We claim that there exists $\epsilon\in(0,1)$ such that, if
	\begin{equation}\label{eq67bb678e}
		r < %\frac{\hat r}{2^{n+2} n^3 e} = 
		\epsilon \hat r,
	\end{equation}
	then the Taylor series of $u$ centered at $\hat x$ converges on $B(\hat x,r)$ to $u$,
	that is, for every $x\in B(\hat x,r)$,
	\begin{equation}\label{eq67bb6792}
		u(x) = \lim_{N\to\infty} \sum_{k=0}^N \sum_{|\alpha|=k} \frac{\Diff^\alpha u(\hat x)}{\alpha!} (x-\hat x)^\alpha .
	\end{equation}
	
	To this aim, define the reminder function
	\[
	R_N(x) = u(x) - \sum_{k=0}^{N-1} \sum_{|\alpha|=k} \frac{\Diff^\alpha u(\hat x)}{\alpha!} (x-\hat x)^\alpha .
	\]
	For every $x\in B(\hat x,r)$ there exists $t_x\in[0,1]$ such that
	\[
	R_N(x) = \sum_{|\alpha|=N} \frac{\Diff^\alpha u(\hat x+t(x-\hat x))}{\alpha!} (x-\hat x)^\alpha .
	\]
	Using Proposition~\ref{prop67bb09b8}, we make the following estimate:
	since $\hat x+t(x-\hat x)\in B(\hat x,r) \subset B(\hat x,\hat r)$, then
	$B(\hat x+t(x-\hat x) , \hat r) \subset B(\hat x,2 \hat r) \subset U$.
	Therefore, 
	\begin{align*}
		|R_N(x)|
		&\le \sum_{|\alpha|=N} \frac{|\Diff^\alpha u(\hat x+t(x-\hat x))|}{\alpha!} |x-\hat x|^N \\
		&\overset{\eqref{eq67baf145}}\le  \frac{(2^{n+1}nN)^N \|u\|_{L^1(B(\hat x+t(x-\hat x),\hat r))}}{\omega_n \hat r^{n+N} }  \sum_{|\alpha|=N}\frac1{\alpha!} |x-\hat x|^N \\
		&\le  \frac{(2^{n+1}nN)^N \|u\|_{L^1(B(\hat x,2\hat r))}}{\omega_n \hat r^{n+N} }  r^N \sum_{|\alpha|=N}\frac1{\alpha!}\\
		&\overset{\eqref{eq67bb6ddf_bis}}= \frac{\|u\|_{L^1(B(\hat x, \hat r))}}{\omega_n \hat r^n} \epsilon^N (2^{n+1}nN)^N \frac{n^N}{N!} \\
		&= \frac{\|u\|_{L^1(B(\hat x, \hat r))} }{ \omega_n \hat r^n } \frac{\sqrt{2\pi N} (N/e)^N }{N!} 
			\frac{(\epsilon 2^{n+1}n^2N)^{N}}{\sqrt{2\pi N} (N/e)^N } \\
		&= \frac{\|u\|_{L^1(B(\hat x, \hat r))} }{ \omega_n \hat r^n } \frac{\sqrt{2\pi N} (N/e)^N }{N!} 
			\frac{(\epsilon 2^{n+1}n^2e)^{N}}{\sqrt{2\pi N} } ,\\
	\end{align*}
	where we have used the Multinomial Theorem
	\begin{equation}\label{eq67bb6ddf_bis}
		n^N = (\sum_{j=1}^n 1)^N = \sum_{|\alpha|=N} \frac{N!}{\alpha!} .
	\end{equation}
	Using Stirling's formula
	\begin{equation}\label{eq67bb6e37}
		\lim_{N\to\infty} \frac{\sqrt{2\pi N} (N/e)^N }{N!}  = 1 ,
	\end{equation}
	we conclude that, if $\epsilon\in(0,1)$ is so that $\epsilon 2^{n+1}n^2e<1$, then $\lim_{N\to\infty}|R_N(x)|=0$.
\end{proof}


%%%%%%%%%%%%%%%%%%%%%%%%%%%%%%%%%%%%%%%%%%%%%%%%%%%%%%%%%%%%%%%
\section{List of Notations}

\begin{tabular}{p{2cm}@{ : }p{9cm}}
$C^k(\Omega)$ 
	& Functions $\Omega\to\C$ that are smooth up to order $k$ \\
$C^k_c(\Omega)$ 
	& Functions in $C^k(\Omega)$ that have compact support contained in $\Omega$\\
$C^k_b(\Omega)$
	& Bounded functions belonging to $C^k(\Omega)$ \\
$C^k(\Omega;W)$, $C^k_c(\Omega;W)$ 
	& Functions $\Omega\to W$ with the required regularity\\
$\langle a,b \rangle$
	& Hermitian product: see Section~\ref{subs6889e7fa} \\
$a\cdot b$
	& Dot product: see Section~\ref{subs6889e81a} \\
$\langle a|b \rangle$
	& Pairing: see Section~\ref{subs6889e83c} \\
${}_{V^*}\langle a|b \rangle_V$
	& Pairing: see Section~\ref{subs6889e83c} \\
$\Diff f$, $\Diff^\alpha f$, $\Diff^\alpha_x f$
	& Derivatives: see Section~\ref{sec6889e87f} \\
$\grad f$
	& Gradient: see Section~\ref{sec6889e87f} \\
$\de_x f$, $\de_x^\alpha f$
	& Derivatives: see Section~\ref{sec6889e87f} \\
$\|u\|_{L^p}$, $\|u\|_{L^p(X)}$, $\|u\|_{L^p(\mu)}$
	& $L^p$ spaces: see Section~\ref{subs6889e917} \\
$\laplacian u$
	& Laplace operator: see Section~\ref{subs67b9d4b5} \\
$\abs{\alpha}$ for $\alpha\in\N^n$
	& size of multi-index: see Section~\ref{sec6889e87f} \\
$\spt(\rho)$
	& support of $\rho$, that is, the closure of $\{x:\rho(x)\neq0\}$ \\
$B(x,r)$
	& open ball with center $x$ and radius $r$, that is, $\{y: |y-x|<r\}$ \\
$\bar B(x,r)$
	& closed ball with center $x$ and radius $r$, that is, $\{y: |y-x|\le r\}$; in metric spaces, $\bar B(x,r)$ and $\overline{ B(x,r)}$ may be different \\
$\waveop u$
	& wave operator $\waveop = \de_t^2-\laplacian$; see Section~\ref{subs67cd41c5}\\
\end{tabular}





%%%%%%%%%%%%%%%%%%%%%%%%%%%%%%%%%%%%%%%%%%%%%%%%%%%%%%%%%%%%%%%
\printindex
\printbibliography
\end{document}
